% generated by GAPDoc2LaTeX from XML source (Frank Luebeck)
\documentclass[a4paper,11pt]{report}

            \usepackage[all]{xy} 
            \newcommand{\Act} {\mathrm{Act}}
            \newcommand{\Aut} {\mathrm{Aut}}
            \newcommand{\Disp}{\mathrm{Disp}}
            \newcommand{\Fix} {\mathrm{Fix}}
            \newcommand{\Inn} {\mathrm{Inn}}
            \newcommand{\Stab}{\mathrm{Stab}}
            \newcommand{\calA}{\mathcal{A}}
            \newcommand{\calB}{\mathcal{B}}
            \newcommand{\calC}{\mathcal{C}}
            \newcommand{\calD}{\mathcal{D}}
            \newcommand{\calE}{\mathcal{E}}
            \newcommand{\calF}{\mathcal{F}}
            \newcommand{\calG}{\mathcal{G}}
            \newcommand{\calL}{\mathcal{L}}
            \newcommand{\calM}{\mathcal{M}}
            \newcommand{\calN}{\mathcal{N}}
            \newcommand{\calP}{\mathcal{P}}
            \newcommand{\calQ}{\mathcal{Q}}
            \newcommand{\calS}{\mathcal{S}}
            \newcommand{\calW}{\mathcal{W}}
            \newcommand{\calX}{\mathcal{X}}
            \newcommand{\calY}{\mathcal{Y}}
        
\usepackage[top=37mm,bottom=37mm,left=27mm,right=27mm]{geometry}
\sloppy
\pagestyle{myheadings}
\usepackage{amssymb}
\usepackage[utf8]{inputenc}
\usepackage{makeidx}
\makeindex
\usepackage{color}
\definecolor{FireBrick}{rgb}{0.5812,0.0074,0.0083}
\definecolor{RoyalBlue}{rgb}{0.0236,0.0894,0.6179}
\definecolor{RoyalGreen}{rgb}{0.0236,0.6179,0.0894}
\definecolor{RoyalRed}{rgb}{0.6179,0.0236,0.0894}
\definecolor{LightBlue}{rgb}{0.8544,0.9511,1.0000}
\definecolor{Black}{rgb}{0.0,0.0,0.0}

\definecolor{linkColor}{rgb}{0.0,0.0,0.554}
\definecolor{citeColor}{rgb}{0.0,0.0,0.554}
\definecolor{fileColor}{rgb}{0.0,0.0,0.554}
\definecolor{urlColor}{rgb}{0.0,0.0,0.554}
\definecolor{promptColor}{rgb}{0.0,0.0,0.589}
\definecolor{brkpromptColor}{rgb}{0.589,0.0,0.0}
\definecolor{gapinputColor}{rgb}{0.589,0.0,0.0}
\definecolor{gapoutputColor}{rgb}{0.0,0.0,0.0}

%%  for a long time these were red and blue by default,
%%  now black, but keep variables to overwrite
\definecolor{FuncColor}{rgb}{0.0,0.0,0.0}
%% strange name because of pdflatex bug:
\definecolor{Chapter }{rgb}{0.0,0.0,0.0}
\definecolor{DarkOlive}{rgb}{0.1047,0.2412,0.0064}


\usepackage{fancyvrb}

\usepackage{mathptmx,helvet}
\usepackage[T1]{fontenc}
\usepackage{textcomp}


\usepackage[
            pdftex=true,
            bookmarks=true,        
            a4paper=true,
            pdftitle={Written with GAPDoc},
            pdfcreator={LaTeX with hyperref package / GAPDoc},
            colorlinks=true,
            backref=page,
            breaklinks=true,
            linkcolor=linkColor,
            citecolor=citeColor,
            filecolor=fileColor,
            urlcolor=urlColor,
            pdfpagemode={UseNone}, 
           ]{hyperref}

\newcommand{\maintitlesize}{\fontsize{50}{55}\selectfont}

% write page numbers to a .pnr log file for online help
\newwrite\pagenrlog
\immediate\openout\pagenrlog =\jobname.pnr
\immediate\write\pagenrlog{PAGENRS := [}
\newcommand{\logpage}[1]{\protect\write\pagenrlog{#1, \thepage,}}
%% were never documented, give conflicts with some additional packages

\newcommand{\GAP}{\textsf{GAP}}

%% nicer description environments, allows long labels
\usepackage{enumitem}
\setdescription{style=nextline}

%% depth of toc
\setcounter{tocdepth}{1}





%% command for ColorPrompt style examples
\newcommand{\gapprompt}[1]{\color{promptColor}{\bfseries #1}}
\newcommand{\gapbrkprompt}[1]{\color{brkpromptColor}{\bfseries #1}}
\newcommand{\gapinput}[1]{\color{gapinputColor}{#1}}


\begin{document}

\logpage{[ 0, 0, 0 ]}
\begin{titlepage}
\mbox{}\vfill

\begin{center}{\maintitlesize \textbf{ XMod \mbox{}}}\\
\vfill

\hypersetup{pdftitle= XMod }
\markright{\scriptsize \mbox{}\hfill  XMod  \hfill\mbox{}}
{\Huge \textbf{ Crossed Modules and Cat1\texttt{\symbol{45}}Groups \mbox{}}}\\
\vfill

{\Huge  2.98 \mbox{}}\\[1cm]
{ 12 March 2026 \mbox{}}\\[1cm]
\mbox{}\\[2cm]
{\Large \textbf{ Chris Wensley\\
   \mbox{}}}\\
{\Large \textbf{ Murat Alp\\
   \mbox{}}}\\
{\Large \textbf{ Alper Odabas\\
   \mbox{}}}\\
{\Large \textbf{ Enver Onder Uslu\\
 \mbox{}}}\\
\hypersetup{pdfauthor= Chris Wensley\\
   ;  Murat Alp\\
   ;  Alper Odabas\\
   ;  Enver Onder Uslu\\
 }
\end{center}\vfill

\mbox{}\\
{\mbox{}\\
\small \noindent \textbf{ Chris Wensley\\
   }  Email: \href{mailto://cdwensley.maths@btinternet.com} {\texttt{cdwensley.maths@btinternet.com}}\\
  Homepage: \href{https://github.com/cdwensley} {\texttt{https://github.com/cdwensley}}}\\
{\mbox{}\\
\small \noindent \textbf{ Murat Alp\\
   }  Email: \href{mailto://murat.alp@aum.edu.kw} {\texttt{murat.alp@aum.edu.kw}}\\
  Address: \begin{minipage}[t]{8cm}\noindent
 Prof. Dr. M. Alp\\
 College of Engineering and Technology\\
 American University of the Middle East\\
 Kuwait\\
 \end{minipage}
}\\
{\mbox{}\\
\small \noindent \textbf{ Alper Odabas\\
   }  Email: \href{mailto://aodabas@ogu.edu.tr} {\texttt{aodabas@ogu.edu.tr}}\\
  Address: \begin{minipage}[t]{8cm}\noindent
 Dr. A. Odabas \\
 Osmangazi University \\
 Arts and Sciences Faculty \\
 Department of Mathematics and Computer Science \\
 Eskisehir \\
 Turkey\\
 \end{minipage}
}\\
\end{titlepage}

\newpage\setcounter{page}{2}
{\small 
\section*{Abstract}
\logpage{[ 0, 0, 1 ]}
 The \textsf{XMod} package provides functions for computation with 
\begin{itemize}
\item  finite crossed modules of groups and cat1\texttt{\symbol{45}}groups, and
morphisms of these structures; 
\item  finite pre\texttt{\symbol{45}}crossed modules,
pre\texttt{\symbol{45}}cat1\texttt{\symbol{45}}groups, and their Peiffer
quotients; 
\item  isoclinism classes of groups and crossed modules; 
\item  derivations of crossed modules and sections of cat1\texttt{\symbol{45}}groups; 
\item  crossed squares and their morphisms, including the actor crossed square of a
crossed module; 
\item  crossed modules of finite groupoids (experimental version). 
\end{itemize}
 

 \textsf{XMod} was originally implemented in 1996 using the \textsf{GAP}3 language, when the second author was studying for a Ph.D. \cite{A1} in the School of Mathematics and Computer Science at Bangor University. 

 In April 2002 the first and third parts were converted to \textsf{GAP}4, the pre\texttt{\symbol{45}}structures were added, and version 2.001 was
released. The final two parts, covering derivations, sections and actors, were
included in the January 2004 release 2.002 for \textsf{GAP} 4.4. 

 In October 2015 functions for computing isoclinism classes of crossed modules,
written by Alper Odaba{\c s} and Enver Uslu, were added. These are contained
in Chapter \ref{chap-isclnc}, and are described in detail in the paper \cite{IOU1}. 

 Bug reports, suggestions and comments are, of course, welcome. To provide
these, please submit an issue at \href{https://github.com/gap-packages/xmod/issues/} {\texttt{https://github.com/gap\texttt{\symbol{45}}packages/xmod/issues/}} or send an email to the first author at \href{mailto://cdwensley@btinternet.com} {\texttt{cdwensley@btinternet.com}}. 

 \mbox{}}\\[1cm]
{\small 
\section*{Copyright}
\logpage{[ 0, 0, 2 ]}
 {\copyright} 1996\texttt{\symbol{45}}2026, Chris Wensley et al. 

 The \textsf{XMod} package is free software; you can redistribute it and/or modify it under the
terms of the GNU General Public License as published by the Free Software
Foundation; either version 2 of the License, or (at your option) any later
version. \mbox{}}\\[1cm]
{\small 
\section*{Acknowledgements}
\logpage{[ 0, 0, 3 ]}
 This documentation was prepared using the \textsf{GAPDoc} \cite{GAPDoc} and \textsf{AutoDoc} \cite{AutoDoc} packages.

 The procedure used to produce new releases uses the package \textsf{GitHubPagesForGAP} \cite{GitHubPagesForGAP} and the package \textsf{ReleaseTools}.

 \mbox{}}\\[1cm]
\newpage

\def\contentsname{Contents\logpage{[ 0, 0, 4 ]}}

\tableofcontents
\newpage

         
\chapter{\textcolor{Chapter }{Introduction}}\label{Intro}
\logpage{[ 1, 0, 0 ]}
\hyperdef{L}{X7DFB63A97E67C0A1}{}
{
  The \textsf{XMod} package provides functions for computation with 
\begin{itemize}
\item  finite crossed modules of groups and cat$^1$\texttt{\symbol{45}}groups, and morphisms of these structures; 
\item  finite pre\texttt{\symbol{45}}crossed modules, pre\texttt{\symbol{45}}cat$^1$\texttt{\symbol{45}}groups, and their Peiffer quotients; 
\item  derivations of crossed modules and sections of cat$^1$\texttt{\symbol{45}}groups; 
\item  isoclinism of groups and crossed modules; 
\item  the actor crossed square of a crossed module; 
\item  crossed squares, cat$^2$\texttt{\symbol{45}}groups, and their morphisms; 
\item  crossed modules of groupoids. 
\end{itemize}
 It is loaded with the command 
\begin{Verbatim}[commandchars=!@|,fontsize=\small,frame=single,label=Example]
  
  !gapprompt@gap>| !gapinput@LoadPackage( "xmod" ); |
  
\end{Verbatim}
 

 The term crossed module was introduced by J. H. C. Whitehead in \cite{W2}, \cite{W1}. Loday, in \cite{L1}, reformulated the notion of a crossed module as a cat$^1$\texttt{\symbol{45}}group. Norrie \cite{N1}, \cite{N2} and Gilbert \cite{G1} have studied derivations, automorphisms of crossed modules and the actor of a
crossed module, while Ellis \cite{E1} has investigated higher dimensional analogues. Properties of induced crossed
modules have been determined by Brown, Higgins and Wensley in \cite{BH1}, \cite{BW1} and \cite{BW2}. For further references see \cite{AW1}, where we discuss some of the data structures and algorithms used in this
package, and also tabulate isomorphism classes of cat$^1$\texttt{\symbol{45}}groups up to size $30$. 

 \textsf{XMod} was originally implemented in 1997 using the \textsf{GAP} 3 language. In April 2002 the first and third parts were converted to \textsf{GAP} 4, the pre\texttt{\symbol{45}}structures were added, and version 2.001 was
released. The final two parts, covering derivations, sections and actors, were
included in the January 2004 release 2.002 for \textsf{GAP} 4.4. Many of the function names have been changed during the conversion, for
example \texttt{ConjugationXMod} has become \texttt{XModByNormalSubgroup} (\ref{XModByNormalSubgroup}). For a list of name changes see the file \texttt{names.pdf} in the \texttt{doc} directory. 

 In October 2015 Alper Odaba{\c s} and Enver Uslu were added to the list of
package authors. Their functions for computing isoclinism classes of groups
and crossed modules are contained in Chapter \ref{chap-isclnc}, and are described in detail in their paper \cite{IOU1}. 

 The package may be obtained as a compressed \texttt{.tar} file \texttt{XMod\texttt{\symbol{45}}version.tar.gz} from the GitHub release site: \href{https://github.com/gap-packages/xmod/releases/tag/version} {\texttt{https://github.com/gap\texttt{\symbol{45}}packages/xmod/releases/tag/version}}. \index{GitHub repository} The package also has a GitHub repository at: \href{https://github.com/gap-packages/xmod/} {\texttt{https://github.com/gap\texttt{\symbol{45}}packages/xmod/}}. 

 Crossed modules and cat$^1$\texttt{\symbol{45}}groups are special types of \emph{2\texttt{\symbol{45}}dimensional groups} \cite{B82}, \cite{brow:hig:siv}, and are implemented as \texttt{2DimensionalDomains} and \texttt{2DimensionalGroups} having a \texttt{Source} and a \texttt{Range}.   

 The package divides into eight parts. The first part is concerned with the
standard constructions for pre\texttt{\symbol{45}}crossed modules and crossed
modules; together with direct products; normal sub\texttt{\symbol{45}}crossed
modules; and quotients. Operations for constructing pre\texttt{\symbol{45}}cat$^1$\texttt{\symbol{45}}groups and cat$^1$\texttt{\symbol{45}}groups, and for converting between cat$^1$\texttt{\symbol{45}}groups and crossed modules, are also included. 

 The second part is concerned with \emph{morphisms} of (pre\texttt{\symbol{45}})crossed modules and (pre\texttt{\symbol{45}})cat$^1$\texttt{\symbol{45}}groups, together with standard operations for morphisms,
such as composition, image and kernel. 

 The third part is the most recent part of the package, introduced in October
2015. Additional operations and properties for crossed modules are included in
Section \ref{sect-more-xmod-ops}. Then, in \ref{sect-isoclinic-groups} and \ref{sect-isoclinic-xmods} there are functions for isoclinism of groups and crossed modules. 

 The fourth part is concerned with the equivalent notions of \emph{derivation} for a crossed module and \emph{section} for a cat$^1$\texttt{\symbol{45}}group, and the monoids which they form under the Whitehead
multiplication. 

 The fifth part deals with actor crossed modules and actor cat$^1$\texttt{\symbol{45}}groups. For the actor crossed module ${\rm Act}(\calX)$ of a crossed module $\calX$ we require representations for the Whitehead group of regular derivations of $\calX$ and for the group of automorphisms of $\calX$. The construction also provides an inner morphism from $\calX$ to ${\rm Act}(\calX)$ whose kernel is the centre of $\calX$. 

 The sixth part, which remains under development, contains functions to compute
induced crossed modules. 

 Since version 2.007 there are experimental functions for \emph{crossed squares} and their morphisms, structures which arise as $3$\texttt{\symbol{45}}dimensional groups. Examples of these are inclusions of
normal sub\texttt{\symbol{45}}crossed modules, and the inner morphism from a
crossed module to its actor. 

 The eighth part has some experimental functions for crossed modules of
groupoids, interacting with the package \textsf{groupoids}. Much more work on this is needed. 

 Future plans include the implementation of \emph{group\texttt{\symbol{45}}graphs} which will provide examples of pre\texttt{\symbol{45}}crossed modules (their
implementation will require interaction with
graph\texttt{\symbol{45}}theoretic functions in \textsf{GAP} 4). There are also plans to implement cat$^2$\texttt{\symbol{45}}groups, and conversion betwen these and crossed squares. 

 The equivalent categories \texttt{XMod} (crossed modules) and \texttt{Cat1} (cat$^1$\texttt{\symbol{45}}groups) are also equivalent to \texttt{GpGpd}, the subcategory of group objects in the category \texttt{Gpd} of groupoids. Finite groupoids have been implemented in Emma Moore's package \textsf{groupoids} \cite{M1} for groupoids and crossed resolutions. 

 \index{InfoXMod@\texttt{InfoXMod}} In order that the user has some control of the verbosity of the \textsf{XMod} package's functions, an \texttt{InfoClass} \texttt{InfoXMod} is provided (see Chapter \texttt{ref:Info Functions} in the \textsf{GAP} Reference Manual for a description of the \texttt{Info} mechanism). By default, the \texttt{InfoLevel} of \texttt{InfoXMod} is \texttt{0}; progressively more information is supplied by raising the \texttt{InfoLevel} to \texttt{1}, \texttt{2} and \texttt{3}. 

 
\begin{Verbatim}[commandchars=!@|,fontsize=\small,frame=single,label=Example]
  
  !gapprompt@gap>| !gapinput@SetInfoLevel( InfoXMod, 1); #sets the InfoXMod level to 1|
  
\end{Verbatim}
 

 Once the package is loaded, the manual \texttt{doc/manual.pdf} can be found in the documentation folder. The \texttt{html} versions, with or without MathJax, should be rebuilt as follows: 

 
\begin{Verbatim}[commandchars=!@|,fontsize=\small,frame=single,label=Example]
  
  !gapprompt@gap>| !gapinput@ReadPackage( "xmod, "makedoc.g" ); |
  
\end{Verbatim}
 

 It is possible to check that the package has been installed correctly by
running the test files: 

 
\begin{Verbatim}[commandchars=!@|,fontsize=\small,frame=single,label=Example]
  
  !gapprompt@gap>| !gapinput@ReadPackage( "xmod", "tst/testall.g" );|
  #I  Testing .../pkg/xmod/tst/gp2obj.tst 
  ... 
  
\end{Verbatim}
 

 Additional information can be found on the \emph{Computational Higher\texttt{\symbol{45}}dimensional Discrete Algebra} website at: \href{https://github.com/cdwensley} {\texttt{https://github.com/cdwensley}}. }

         
\chapter{\textcolor{Chapter }{2d\texttt{\symbol{45}}groups : crossed modules and cat$^1$\texttt{\symbol{45}}groups}}\label{chap-gp2obj}
\logpage{[ 2, 0, 0 ]}
\hyperdef{L}{X7EB8288E8424F39F}{}
{
  \index{2d\texttt{\symbol{45}}domain} \index{2d\texttt{\symbol{45}}group} The term \emph{2d\texttt{\symbol{45}}group} refers to a set of equivalent categories of which the most common are the
categories of \emph{crossed modules}; \emph{cat$^1$\texttt{\symbol{45}}groups}; and \emph{group\texttt{\symbol{45}}groupoids}, all of which involve a pair of groups. 
\section{\textcolor{Chapter }{Constructions for crossed modules}}\label{sect-constructions}
\logpage{[ 2, 1, 0 ]}
\hyperdef{L}{X7BAD9A7F7AFEEC89}{}
{
  \index{crossed module} A crossed module (of groups) $\calX = (\partial : S \to R )$ consists of a group homomorphism $\partial $, called the \emph{boundary} of $\calX$, with \emph{source} $S$ and \emph{range} $R$. The group $R$ acts on itself by conjugation, and on $S$ by an \emph{action} $\alpha : R \to {\rm Aut}(S)$ such that, for all $s,s_1,s_2 \in S$ and $r \in R$, 
\[ {\bf XMod\ 1} : \partial(s^r) = r^{-1} (\partial s) r = (\partial s)^r, \qquad
{\bf XMod\ 2} : s_1^{\partial s_2} = s_2^{-1}s_1 s_2 = {s_1}^{s_2}. \]
 When only the first of these axioms is satisfied, the resulting structure is a \emph{pre\texttt{\symbol{45}}crossed module} (see section \ref{sect-precrossed-modules}). The kernel of $\partial$ is abelian. 

 (Much of the literature on crossed modules uses left actions, but we have
chosen to use right actions in this package since that is the standard choice
for group actions in \textsf{GAP}.) 

\subsection{\textcolor{Chapter }{XMod}}
\logpage{[ 2, 1, 1 ]}\nobreak
\hyperdef{L}{X7C8175AE7F76B586}{}
{\noindent\textcolor{FuncColor}{$\triangleright$\enspace\texttt{XMod({\mdseries\slshape args})\index{XMod@\texttt{XMod}}
\label{XMod}
}\hfill{\scriptsize (function)}}\\
\noindent\textcolor{FuncColor}{$\triangleright$\enspace\texttt{XModByBoundaryAndAction({\mdseries\slshape bdy, act})\index{XModByBoundaryAndAction@\texttt{XModByBoundaryAndAction}}
\label{XModByBoundaryAndAction}
}\hfill{\scriptsize (operation)}}\\


 The global function \texttt{XMod} calls one of the standard constructions described in the following
subsections. In the example the boundary is the identity mapping on \texttt{c5} and the action is trivial. }

 
\begin{Verbatim}[commandchars=!@|,fontsize=\small,frame=single,label=Example]
  
  !gapprompt@gap>| !gapinput@c5 := Group( (5,6,7,8,9) );;|
  !gapprompt@gap>| !gapinput@SetName( c5, "c5" );|
  !gapprompt@gap>| !gapinput@id5 := IdentityMapping( c5 );;|
  !gapprompt@gap>| !gapinput@ac5 := AutomorphismGroup( c5 );; |
  !gapprompt@gap>| !gapinput@act := MappingToOne( c5, ac5 );;|
  !gapprompt@gap>| !gapinput@XMod( id5, act ) = XModByBoundaryAndAction( id5, act );|
  true
  
\end{Verbatim}
 

\subsection{\textcolor{Chapter }{XModByNormalSubgroup}}
\logpage{[ 2, 1, 2 ]}\nobreak
\hyperdef{L}{X83050ED686776933}{}
{\noindent\textcolor{FuncColor}{$\triangleright$\enspace\texttt{XModByNormalSubgroup({\mdseries\slshape G, N})\index{XModByNormalSubgroup@\texttt{XModByNormalSubgroup}}
\label{XModByNormalSubgroup}
}\hfill{\scriptsize (operation)}}\\


 A \emph{conjugation crossed module} is the inclusion of a normal subgroup $S \unlhd R$, where $R$ acts on $S$ by conjugation. }

 

\subsection{\textcolor{Chapter }{XModByTrivialAction}}
\logpage{[ 2, 1, 3 ]}\nobreak
\hyperdef{L}{X867B2D53832EF05E}{}
{\noindent\textcolor{FuncColor}{$\triangleright$\enspace\texttt{XModByTrivialAction({\mdseries\slshape bdy})\index{XModByTrivialAction@\texttt{XModByTrivialAction}}
\label{XModByTrivialAction}
}\hfill{\scriptsize (operation)}}\\


 A \emph{trivial action crossed module} $(\partial : S \to R)$ has $s^r = s$ for all $s \in S, \; r \in R$, the source is abelian and the image lies in the centre of the range. }

 
\begin{Verbatim}[commandchars=!@|,fontsize=\small,frame=single,label=Example]
  
  !gapprompt@gap>| !gapinput@q8 := QuaternionGroup( IsPermGroup, 8 );|
  Group([ (1,5,3,7)(2,8,4,6), (1,2,3,4)(5,6,7,8) ])
  !gapprompt@gap>| !gapinput@SetName( q8, "q8" );|
  !gapprompt@gap>| !gapinput@c2 := Centre( q8 );                     |
  Group([ (1,3)(2,4)(5,7)(6,8) ])
  !gapprompt@gap>| !gapinput@SetName( c2, "<-1>" );|
  !gapprompt@gap>| !gapinput@bdy := InclusionMappingGroups( q8, c2 );;|
  !gapprompt@gap>| !gapinput@X8a := XModByTrivialAction( bdy );|
  [<-1>->q8]
  !gapprompt@gap>| !gapinput@c4 := Subgroup( q8, [q8.1] );;|
  !gapprompt@gap>| !gapinput@SetName( c4, "<i>" );|
  !gapprompt@gap>| !gapinput@X8b := XModByNormalSubgroup( q8, c4 );|
  [<i>->q8]
  !gapprompt@gap>| !gapinput@Display(X8b);        |
  Crossed module [<i>->q8] :- 
  : Source group has generators:
    [ (1,5,3,7)(2,8,4,6) ]
  : Range group q8 has generators:
    [ (1,5,3,7)(2,8,4,6), (1,2,3,4)(5,6,7,8) ]
  : Boundary homomorphism maps source generators to:
    [ (1,5,3,7)(2,8,4,6) ]
  : Action homomorphism maps range generators to automorphisms:
    (1,5,3,7)(2,8,4,6) --> { source gens --> [ (1,5,3,7)(2,8,4,6) ] }
    (1,2,3,4)(5,6,7,8) --> { source gens --> [ (1,7,3,5)(2,6,4,8) ] }
    These 2 automorphisms generate the group of automorphisms.
  
\end{Verbatim}
 

\subsection{\textcolor{Chapter }{XModByAutomorphismGroup}}
\logpage{[ 2, 1, 4 ]}\nobreak
\hyperdef{L}{X78B14FDA817CCEEF}{}
{\noindent\textcolor{FuncColor}{$\triangleright$\enspace\texttt{XModByAutomorphismGroup({\mdseries\slshape grp})\index{XModByAutomorphismGroup@\texttt{XModByAutomorphismGroup}}
\label{XModByAutomorphismGroup}
}\hfill{\scriptsize (attribute)}}\\
\noindent\textcolor{FuncColor}{$\triangleright$\enspace\texttt{XModByInnerAutomorphismGroup({\mdseries\slshape grp})\index{XModByInnerAutomorphismGroup@\texttt{XModByInnerAutomorphismGroup}}
\label{XModByInnerAutomorphismGroup}
}\hfill{\scriptsize (attribute)}}\\
\noindent\textcolor{FuncColor}{$\triangleright$\enspace\texttt{XModByGroupOfAutomorphisms({\mdseries\slshape G, A})\index{XModByGroupOfAutomorphisms@\texttt{XModByGroupOfAutomorphisms}}
\label{XModByGroupOfAutomorphisms}
}\hfill{\scriptsize (operation)}}\\


 An \emph{automorphism crossed module} has as range a subgroup $R$ of the automorphism group Aut$(S)$ of $S$ which contains the inner automorphism group of $S$. The boundary maps $s \in S$ to the inner automorphism of $S$ by $s$. }

 
\begin{Verbatim}[commandchars=!@|,fontsize=\small,frame=single,label=Example]
  
  !gapprompt@gap>| !gapinput@X5 := XModByAutomorphismGroup( c5 );|
  [c5 -> Aut(c5)] 
  !gapprompt@gap>| !gapinput@Display( X5 );|
  Crossed module [c5->Aut(c5)] :- 
  : Source group c5 has generators:
    [ (5,6,7,8,9) ]
  : Range group Aut(c5) has generators:
    [ GroupHomomorphismByImages( c5, c5, [ (5,6,7,8,9) ], [ (5,7,9,6,8) ] ) ]
  : Boundary homomorphism maps source generators to:
    [ IdentityMapping( c5 ) ]
  : Action homomorphism maps range generators to automorphisms:
    GroupHomomorphismByImages( c5, c5, [ (5,6,7,8,9) ], 
  [ (5,7,9,6,8) ] ) --> { source gens --> [ (5,7,9,6,8) ] }
    This automorphism generates the group of automorphisms.
  
\end{Verbatim}
 

\subsection{\textcolor{Chapter }{XModByCentralExtension}}
\logpage{[ 2, 1, 5 ]}\nobreak
\hyperdef{L}{X7D0F6FAA7AF69844}{}
{\noindent\textcolor{FuncColor}{$\triangleright$\enspace\texttt{XModByCentralExtension({\mdseries\slshape bdy})\index{XModByCentralExtension@\texttt{XModByCentralExtension}}
\label{XModByCentralExtension}
}\hfill{\scriptsize (operation)}}\\


 A \emph{central extension crossed module} has as boundary a surjection $\partial : S \to R$, with central kernel, where $r \in R$ acts on $S$ by conjugation with $\partial^{-1}r$. }

 
\begin{Verbatim}[commandchars=!@|,fontsize=\small,frame=single,label=Example]
  
  !gapprompt@gap>| !gapinput@gen12 := [ (1,2,3,4,5,6), (2,6)(3,5) ];;|
  !gapprompt@gap>| !gapinput@d12 := Group( gen12 );;                  |
  !gapprompt@gap>| !gapinput@gen6 := [ (7,8,9), (8,9) ];;|
  !gapprompt@gap>| !gapinput@s3 := Group( gen6 );;|
  !gapprompt@gap>| !gapinput@SetName( d12, "d12" );  SetName( s3, "s3" ); |
  !gapprompt@gap>| !gapinput@pr12 := GroupHomomorphismByImages( d12, s3, gen12, gen6 );;|
  !gapprompt@gap>| !gapinput@Kernel( pr12 ) = Centre( d12 );|
  true
  !gapprompt@gap>| !gapinput@X12 := XModByCentralExtension( pr12 );;|
  !gapprompt@gap>| !gapinput@Display( X12 );                         |
  Crossed module [d12->s3] :- 
  : Source group d12 has generators:
    [ (1,2,3,4,5,6), (2,6)(3,5) ]
  : Range group s3 has generators:
    [ (7,8,9), (8,9) ]
  : Boundary homomorphism maps source generators to:
    [ (7,8,9), (8,9) ]
  : Action homomorphism maps range generators to automorphisms:
    (7,8,9) --> { source gens --> [ (1,2,3,4,5,6), (1,3)(4,6) ] }
    (8,9) --> { source gens --> [ (1,6,5,4,3,2), (2,6)(3,5) ] }
    These 2 automorphisms generate the group of automorphisms.
  
\end{Verbatim}
 

\subsection{\textcolor{Chapter }{XModByPullback}}
\logpage{[ 2, 1, 6 ]}\nobreak
\hyperdef{L}{X84FA2B0A795B6997}{}
{\noindent\textcolor{FuncColor}{$\triangleright$\enspace\texttt{XModByPullback({\mdseries\slshape xmod, hom})\index{XModByPullback@\texttt{XModByPullback}}
\label{XModByPullback}
}\hfill{\scriptsize (operation)}}\\


 Let $\calX_0 = (\mu : M \to P)$ be a crossed module. If $\nu : N \to P$ is a group homomorphism with the same range as $\calX_0$, form the pullback group $L = M \times_P N$, with projection $\lambda : L \to N$ (as defined in the \textsf{Utils} package). Then $N$ acts on $L$ by $(m,n)^{n'} := (m^{\nu n'},n^{n'})$, so that $\calX_1 = (\lambda : L \to N)$ is the \emph{ pullback crossed module} determined by $\calX_0$ and $\nu$. There is also a morphism of crossed modules $(\kappa,\nu) : \calX_1 \to \calX_2$. 

 The example forms a pullback of the crossed module \texttt{X12} of the previous subsection. }

 
\begin{Verbatim}[commandchars=@|A,fontsize=\small,frame=single,label=Example]
  
  @gapprompt|gap>A @gapinput|gens4 := [ (11,12), (12,13), (13,14) ];; A
  @gapprompt|gap>A @gapinput|s4 := Group( gens4 );; A
  @gapprompt|gap>A @gapinput|theta := GroupHomomorphismByImages( s4, s3, gens4, [(7,8),(8,9),(7,8)] );;A
  @gapprompt|gap>A @gapinput|X1 := XModByPullback( X12, theta );; A
  @gapprompt|gap>A @gapinput|StructureDescription( Source( X1 ) );A
  "C2 x S4"
  @gapprompt|gap>A @gapinput|SetName( s4, "s4" );  SetName( Source( X1 ), "c2s4" ); A
  @gapprompt|gap>A @gapinput|infoX1 := PullbackInfo( Source( X1 ) );;A
  @gapprompt|gap>A @gapinput|infoX1!.directProduct;A
  Group([ (1,2,3,4,5,6), (2,6)(3,5), (7,8), (8,9), (9,10) ])
  @gapprompt|gap>A @gapinput|infoX1!.projections[1];A
  [ (7,8)(9,10), (7,9)(8,10), (2,6)(3,5)(8,9), (1,5,3)(2,6,4)(8,10,9), 
    (1,6,5,4,3,2)(8,9,10) ] -> [ (), (), (2,6)(3,5), (1,5,3)(2,6,4), 
    (1,6,5,4,3,2) ]
  @gapprompt|gap>A @gapinput|infoX1!.projections[2];A
  [ (7,8)(9,10), (7,9)(8,10), (2,6)(3,5)(8,9), (1,5,3)(2,6,4)(8,10,9), 
    (1,6,5,4,3,2)(8,9,10) ] -> [ (11,12)(13,14), (11,13)(12,14), (12,13), 
    (12,14,13), (12,13,14) ]
  
\end{Verbatim}
 

\subsection{\textcolor{Chapter }{XModByAbelianModule}}
\logpage{[ 2, 1, 7 ]}\nobreak
\hyperdef{L}{X824631577864961E}{}
{\noindent\textcolor{FuncColor}{$\triangleright$\enspace\texttt{XModByAbelianModule({\mdseries\slshape abmod})\index{XModByAbelianModule@\texttt{XModByAbelianModule}}
\label{XModByAbelianModule}
}\hfill{\scriptsize (operation)}}\\


 A \emph{crossed abelian module} has an abelian module as source and the zero map as boundary. See section \ref{sect-abelian-modules} for an example. }

 

\subsection{\textcolor{Chapter }{DirectProduct (for crossed modules)}}
\logpage{[ 2, 1, 8 ]}\nobreak
\hyperdef{L}{X81704DFB795C0D29}{}
{\noindent\textcolor{FuncColor}{$\triangleright$\enspace\texttt{DirectProduct({\mdseries\slshape X1, X2})\index{DirectProduct@\texttt{DirectProduct}!for crossed modules}
\label{DirectProduct:for crossed modules}
}\hfill{\scriptsize (operation)}}\\


 The direct product $\calX_{1} \times \calX_{2}$ of two crossed modules has source $S_1 \times S_2$, range $R_1 \times R_2$ and boundary $\partial_1 \times \partial_2$, with $R_1,\ R_2$ acting trivially on $S_2,\ S_1$ respectively. The embeddings and projections are constructed automatically,
and placed in the \texttt{DirectProductInfo} attribute, together with the two \emph{objects} $\calX_{1} \times \calX_{2}$. 

 The example constructs the product of the two crossed modules formed in
subsection \texttt{XModByTrivialAction} (\ref{XModByTrivialAction}). 

 }

 
\begin{Verbatim}[commandchars=!@|,fontsize=\small,frame=single,label=Example]
  
  !gapprompt@gap>| !gapinput@X8ab := DirectProduct( X8a, X8b );|
  [[<-1>->q8]x[<i>->q8]]
  !gapprompt@gap>| !gapinput@infoX8ab := DirectProductInfo( X8ab );|
  rec( 
    embeddings := [ [[<-1>->q8] => [<-1>x<i>->q8xq8]], 
        [[<i>->q8] => [<-1>x<i>->q8xq8]] ], objects := [ [<-1>->q8], [<i>->q8] ]
      , 
    projections := [ [[<-1>x<i>->q8xq8] => [<-1>->q8]], 
        [[<-1>x<i>->q8xq8] => [<i>->q8]] ] )
  !gapprompt@gap>| !gapinput@DirectProduct( X8a, X8b, X12 );|
  [[[<-1>->q8]x[<i>->q8]]x[d12->s3]]
  
\end{Verbatim}
 

\subsection{\textcolor{Chapter }{Source (for crossed modules)}}
\logpage{[ 2, 1, 9 ]}\nobreak
\hyperdef{L}{X790248A67CB9C33A}{}
{\noindent\textcolor{FuncColor}{$\triangleright$\enspace\texttt{Source({\mdseries\slshape X0})\index{Source@\texttt{Source}!for crossed modules}
\label{Source:for crossed modules}
}\hfill{\scriptsize (attribute)}}\\
\noindent\textcolor{FuncColor}{$\triangleright$\enspace\texttt{Range({\mdseries\slshape X0})\index{Range@\texttt{Range}!for crossed modules}
\label{Range:for crossed modules}
}\hfill{\scriptsize (attribute)}}\\
\noindent\textcolor{FuncColor}{$\triangleright$\enspace\texttt{Boundary({\mdseries\slshape X0})\index{Boundary@\texttt{Boundary}!for crossed modules}
\label{Boundary:for crossed modules}
}\hfill{\scriptsize (attribute)}}\\
\noindent\textcolor{FuncColor}{$\triangleright$\enspace\texttt{XModAction({\mdseries\slshape X0})\index{XModAction@\texttt{XModAction}!for crossed modules of groups}
\label{XModAction:for crossed modules of groups}
}\hfill{\scriptsize (attribute)}}\\


 The following attributes are used in the construction of a crossed module \texttt{X0}. 
\begin{itemize}
\item  \texttt{Source(X0)} and \texttt{Range(X0)} are the source $S$ and range $R$ of $\partial$, the boundary \texttt{Boundary(X0)}; 
\item  \texttt{XModAction(X0)} is a homomorphism from $R$ to a group of automorphisms of \texttt{X0}. 
\end{itemize}
 \index{AutoGroup} (Up until version 2.63 there was an additional attribute \texttt{AutoGroup}, the range of \texttt{XModAction(X0)}.) 

 The example uses the crossed module \texttt{X12} constructed in subsection \texttt{XModByCentralExtension} (\ref{XModByCentralExtension}). }

 
\begin{Verbatim}[commandchars=!@|,fontsize=\small,frame=single,label=Example]
  
  !gapprompt@gap>| !gapinput@[ Source( X12 ), Range( X12 ) ];    |
  [ d12, s3 ]
  !gapprompt@gap>| !gapinput@Boundary( X12 ); |
  [ (1,2,3,4,5,6), (2,6)(3,5) ] -> [ (7,8,9), (8,9) ]
  !gapprompt@gap>| !gapinput@XModAction( X12 );|
  [ (7,8,9), (8,9) ] -> 
  [ [ (1,2,3,4,5,6), (2,6)(3,5) ] -> [ (1,2,3,4,5,6), (1,3)(4,6) ], 
    [ (1,2,3,4,5,6), (2,6)(3,5) ] -> [ (1,6,5,4,3,2), (2,6)(3,5) ] ]
  
\end{Verbatim}
 

\subsection{\textcolor{Chapter }{ImageElmXModAction}}
\logpage{[ 2, 1, 10 ]}\nobreak
\hyperdef{L}{X7AF6602C87845F1D}{}
{\noindent\textcolor{FuncColor}{$\triangleright$\enspace\texttt{ImageElmXModAction({\mdseries\slshape X0, s, r})\index{ImageElmXModAction@\texttt{ImageElmXModAction}}
\label{ImageElmXModAction}
}\hfill{\scriptsize (operation)}}\\


 This function returns the element $s^r$ given by \texttt{XModAction(X0)}. 

 }

 
\begin{Verbatim}[commandchars=!@|,fontsize=\small,frame=single,label=Example]
  
  !gapprompt@gap>| !gapinput@ImageElmXModAction( X12, (1,2,3,4,5,6), (8,9) );|
  (1,6,5,4,3,2)
  
\end{Verbatim}
 

\subsection{\textcolor{Chapter }{Size2d (for crossed modules)}}
\logpage{[ 2, 1, 11 ]}\nobreak
\hyperdef{L}{X7846A7D37957B89E}{}
{\noindent\textcolor{FuncColor}{$\triangleright$\enspace\texttt{Size2d({\mdseries\slshape X0})\index{Size2d@\texttt{Size2d}!for crossed modules}
\label{Size2d:for crossed modules}
}\hfill{\scriptsize (attribute)}}\\


 The standard operation \texttt{Size} cannot be used for crossed modules because the size of a collection is
required to be a number, and we wish to return a list. \texttt{Size2d( X0 )} returns the two\texttt{\symbol{45}}element list, \texttt{[ Size( Source(X0) ), Size( Range(X0) ) ]}. 

 }

 In the simple example below, \texttt{X5} is the automorphism crossed module constructed in subsection \texttt{XModByAutomorphismGroup} (\ref{XModByAutomorphismGroup}). 

 
\begin{Verbatim}[commandchars=!@|,fontsize=\small,frame=single,label=Example]
  
  !gapprompt@gap>| !gapinput@Size2d( X5 ); |
  [ 5, 4 ]
  
\end{Verbatim}
 

\subsection{\textcolor{Chapter }{Name (for crossed modules)}}
\logpage{[ 2, 1, 12 ]}\nobreak
\hyperdef{L}{X85516B19803C01C0}{}
{\noindent\textcolor{FuncColor}{$\triangleright$\enspace\texttt{Name({\mdseries\slshape X0})\index{Name@\texttt{Name}!for crossed modules}
\label{Name:for crossed modules}
}\hfill{\scriptsize (attribute)}}\\
\noindent\textcolor{FuncColor}{$\triangleright$\enspace\texttt{IdGroup({\mdseries\slshape X0})\index{IdGroup@\texttt{IdGroup}!for crossed modules}
\label{IdGroup:for crossed modules}
}\hfill{\scriptsize (attribute)}}\\
\noindent\textcolor{FuncColor}{$\triangleright$\enspace\texttt{ExternalSetXMod({\mdseries\slshape X0})\index{ExternalSetXMod@\texttt{ExternalSetXMod}}
\label{ExternalSetXMod}
}\hfill{\scriptsize (attribute)}}\\


 More familiar attributes are \texttt{Name} and \texttt{IdGroup}. The name is formed by concatenating the names of the source and range (if
these exist). \texttt{IdGroup( X0 )} returns a two\texttt{\symbol{45}}element list \texttt{[ IdGroup( Source(X0) ), IdGroup( Range(X0) ) ]}. 

 The \texttt{ExternalSetXMod} for a crossed module is the source group considered as a
G\texttt{\symbol{45}}set of the range group using the crossed module action. 

 \index{Display for a 2d\texttt{\symbol{45}}group} The \texttt{Display} function is used to print details of 2d\texttt{\symbol{45}}groups. }

 The \texttt{Print} statements at the end of the example list the \textsf{GAP} representations and attributes of \texttt{X5}. 

 
\begin{Verbatim}[commandchars=!@|,fontsize=\small,frame=single,label=Example]
  
  !gapprompt@gap>| !gapinput@IdGroup( X5 ); |
  [ [ 5, 1 ], [ 4, 1 ] ]
  !gapprompt@gap>| !gapinput@ext := ExternalSetXMod( X5 ); |
  <xset:[ (), (5,6,7,8,9), (5,7,9,6,8), (5,8,6,9,7), (5,9,8,7,6) ]>
  !gapprompt@gap>| !gapinput@Orbits( ext );|
  [ [ () ], [ (5,6,7,8,9), (5,7,9,6,8), (5,9,8,7,6), (5,8,6,9,7) ] ]
  !gapprompt@gap>| !gapinput@a := GeneratorsOfGroup( Range( X5 ) )[1]^2; |
  [ (5,6,7,8,9) ] -> [ (5,9,8,7,6) ]
  !gapprompt@gap>| !gapinput@ImageElmXModAction( X5, (5,7,9,6,8), a );|
  (5,8,6,9,7)
  !gapprompt@gap>| !gapinput@Print( RepresentationsOfObject(X5), "\n" );|
  [ "IsComponentObjectRep", "IsAttributeStoringRep", "IsPreXModObj" ]
  !gapprompt@gap>| !gapinput@Print( KnownAttributesOfObject(X5), "\n" );|
  [ "Name", "Range", "Source", "IdGroup", "Boundary", "Size2d", "XModAction", 
    "ExternalSetXMod", "HigherDimension" ]
  
\end{Verbatim}
 }

 
\section{\textcolor{Chapter }{Properties of crossed modules}}\label{sect-properties-xmod}
\logpage{[ 2, 2, 0 ]}
\hyperdef{L}{X7CF622538749FE73}{}
{
  \index{Is2DimensionalDomain} \index{Is2DimensionalGroup} \index{IsTrivialAction2DimensionalGroup} \index{IsNormalSubgroup2DimensionalGroup} \index{IsCentralExtension2DimensionalGroup} \index{IsAutomorphismGroup2DimensionalGroup} \index{IsAbelianModule2DimensionalGroup} The underlying category structures for the objects constructed in this chapter
follow the sequence \texttt{Is2DimensionalDomain}; \texttt{Is2DimensionalMagma}; \texttt{Is2DimensionalMagmaWithOne}; \texttt{Is2DimensionalMagmaWithInverses}, mirroring the situation for (one\texttt{\symbol{45}}dimensional) groups.
From these we construct \texttt{Is2DimensionalSemigroup}, \texttt{Is2DimensionalMonoid} and \texttt{Is2DimensionalGroup}. 

 There are then a variety of properties associated with crossed modules,
starting with \texttt{IsPreXMod} and \texttt{IsXMod}. 

 

\subsection{\textcolor{Chapter }{IsXMod}}
\logpage{[ 2, 2, 1 ]}\nobreak
\hyperdef{L}{X7E77E6B881B1CE50}{}
{\noindent\textcolor{FuncColor}{$\triangleright$\enspace\texttt{IsXMod({\mdseries\slshape X0})\index{IsXMod@\texttt{IsXMod}}
\label{IsXMod}
}\hfill{\scriptsize (property)}}\\
\noindent\textcolor{FuncColor}{$\triangleright$\enspace\texttt{IsPreXMod({\mdseries\slshape X0})\index{IsPreXMod@\texttt{IsPreXMod}}
\label{IsPreXMod}
}\hfill{\scriptsize (property)}}\\
\noindent\textcolor{FuncColor}{$\triangleright$\enspace\texttt{Is2DimensionalGroup({\mdseries\slshape X0})\index{Is2DimensionalGroup@\texttt{Is2DimensionalGroup}}
\label{Is2DimensionalGroup}
}\hfill{\scriptsize (property)}}\\
\noindent\textcolor{FuncColor}{$\triangleright$\enspace\texttt{IsPerm2DimensionalGroup({\mdseries\slshape X0})\index{IsPerm2DimensionalGroup@\texttt{IsPerm2DimensionalGroup}}
\label{IsPerm2DimensionalGroup}
}\hfill{\scriptsize (property)}}\\
\noindent\textcolor{FuncColor}{$\triangleright$\enspace\texttt{IsPc2DimensionalGroup({\mdseries\slshape X0})\index{IsPc2DimensionalGroup@\texttt{IsPc2DimensionalGroup}}
\label{IsPc2DimensionalGroup}
}\hfill{\scriptsize (property)}}\\
\noindent\textcolor{FuncColor}{$\triangleright$\enspace\texttt{IsFp2DimensionalGroup({\mdseries\slshape X0})\index{IsFp2DimensionalGroup@\texttt{IsFp2DimensionalGroup}}
\label{IsFp2DimensionalGroup}
}\hfill{\scriptsize (property)}}\\


 A structure which has \texttt{Is2DimensionalGroup} is a precrossed module or a pre\texttt{\symbol{45}}cat$^1$\texttt{\symbol{45}}group (see section \ref{sect-cat1}). Such an object whose source and range are both permutation groups has \texttt{IsPerm2DimensionalGroup}. The properties \texttt{IsPc2DimensionalGroup}, \texttt{IsFp2DimensionalGroup} are defined similarly. In the example below we see that \texttt{X5} has \texttt{IsPreXMod}, \texttt{IsXMod} and \texttt{IsPerm2DimensionalGroup}. 

 There are also properties corresponding to the various construction methods
which are listed in section \ref{sect-constructions}: \texttt{IsTrivialAction2DimensionalGroup}; \texttt{IsNormalSubgroup2DimensionalGroup}; \texttt{IsCentralExtension2DimensionalGroup}; \texttt{IsAutomorphismGroup2DimensionalGroup} and \texttt{IsAbelianModule2DimensionalGroup}. }

 
\begin{Verbatim}[commandchars=!@|,fontsize=\small,frame=single,label=Example]
  
  !gapprompt@gap>| !gapinput@[ IsTrivial( X5 ),  IsNonTrivial( X5 ),  IsFinite( X5 ) ];|
  [ false, true, true ]
  !gapprompt@gap>| !gapinput@kpoX5 := KnownPropertiesOfObject(X5);;|
  !gapprompt@gap>| !gapinput@ForAll( [ "IsTrivial", "IsNonTrivial", "IsFinite", |
  !gapprompt@>| !gapinput@ "CanEasilyCompareElements", "CanEasilySortElements", "IsDuplicateFree", |
  !gapprompt@>| !gapinput@ "IsGeneratorsOfSemigroup", "IsPreXModDomain", "IsPreXMod", "IsXMod", |
  !gapprompt@>| !gapinput@ "IsAutomorphismGroup2DimensionalGroup" ], |
  !gapprompt@>| !gapinput@ s -> s in kpoX5 ); |
  true
  !gapprompt@gap>| !gapinput@Is2DimensionalGroup(X5);|
  true
  !gapprompt@gap>| !gapinput@IsAutomorphismGroup2DimensionalGroup(X5);|
  true
  
\end{Verbatim}
 

\subsection{\textcolor{Chapter }{SubXMod}}
\logpage{[ 2, 2, 2 ]}\nobreak
\hyperdef{L}{X7884284383284A87}{}
{\noindent\textcolor{FuncColor}{$\triangleright$\enspace\texttt{SubXMod({\mdseries\slshape X0, src, rng})\index{SubXMod@\texttt{SubXMod}}
\label{SubXMod}
}\hfill{\scriptsize (operation)}}\\
\noindent\textcolor{FuncColor}{$\triangleright$\enspace\texttt{TrivialSubXMod({\mdseries\slshape X0})\index{TrivialSubXMod@\texttt{TrivialSubXMod}}
\label{TrivialSubXMod}
}\hfill{\scriptsize (attribute)}}\\
\noindent\textcolor{FuncColor}{$\triangleright$\enspace\texttt{NormalSubXMods({\mdseries\slshape X0})\index{NormalSubXMods@\texttt{NormalSubXMods}}
\label{NormalSubXMods}
}\hfill{\scriptsize (attribute)}}\\
\noindent\textcolor{FuncColor}{$\triangleright$\enspace\texttt{IsNormal({\mdseries\slshape X0, S0})\index{IsNormal@\texttt{IsNormal}!for pre-crossed modules}
\label{IsNormal:for pre-crossed modules}
}\hfill{\scriptsize (operation)}}\\


 With the standard crossed module constructors listed above as building blocks,
sub\texttt{\symbol{45}}crossed modules, normal sub\texttt{\symbol{45}}crossed
modules $\calN \lhd \calX$, and also quotients $\calX/\calN$ may be constructed. A sub\texttt{\symbol{45}}crossed module $\calS = (\delta : N \to M)$ is \emph{normal} in $\calX = (\partial : S \to R)$ if 
\begin{itemize}
\item  $N,M$ are normal subgroups of $S,R$ respectively, 
\item  $\delta$ is the restriction of $\partial$, 
\item  $n^r \in N$ for all $n \in N,~r \in R$, 
\item  $(s^{-1})^ms \in N$ for all $m \in M,~s \in S$. 
\end{itemize}
 These conditions ensure that $M \ltimes N$ is normal in the semidirect product $R \ltimes S$. (Note that $\langle s,m \rangle = (s^{-1})^ms$ is a displacement: see \texttt{Displacement} (\ref{Displacement}).) 

 A method for \texttt{IsNormal} for precrossed modules is provided. See section \ref{sect-more-xmod-ops} for factor crossed modules and their natural morphisms. 

 The five normal subcrossed modules of \texttt{X4} found in the following example are \texttt{[id,id], [k4,k4], [k4,a4], [a4,a4]} and \texttt{X4} itself. }

 

 
\begin{Verbatim}[commandchars=!@|,fontsize=\small,frame=single,label=Example]
  
  !gapprompt@gap>| !gapinput@s4 := Group( (1,2), (2,3), (3,4) );; |
  !gapprompt@gap>| !gapinput@a4 := Subgroup( s4, [ (1,2,3), (2,3,4) ] );; |
  !gapprompt@gap>| !gapinput@k4 := Subgroup( a4, [ (1,2)(3,4), (1,3)(2,4) ] );; |
  !gapprompt@gap>| !gapinput@SetName(s4,"s4");  SetName(a4,"a4");  SetName(k4,"k4"); |
  !gapprompt@gap>| !gapinput@X4 := XModByNormalSubgroup( s4, a4 );|
  [a4->s4]
  !gapprompt@gap>| !gapinput@Y4 := SubXMod( X4, k4, a4 ); |
  [k4->a4]
  !gapprompt@gap>| !gapinput@IsNormal( X4, Y4 );|
  true
  !gapprompt@gap>| !gapinput@NX4 := NormalSubXMods( X4 );;|
  !gapprompt@gap>| !gapinput@Length( NX4 ); |
  5
  
\end{Verbatim}
 

\subsection{\textcolor{Chapter }{KernelCokernelXMod}}
\logpage{[ 2, 2, 3 ]}\nobreak
\hyperdef{L}{X7D8165F77B23BCF6}{}
{\noindent\textcolor{FuncColor}{$\triangleright$\enspace\texttt{KernelCokernelXMod({\mdseries\slshape X0})\index{KernelCokernelXMod@\texttt{KernelCokernelXMod}}
\label{KernelCokernelXMod}
}\hfill{\scriptsize (attribute)}}\\


 Let $\calX = (\partial : S \to R)$. If $K \leqslant S$ is the kernel of $\partial$, and $J \leqslant R$ is the image of $\partial$, form $C = R/J$. Then $(\nu\partial|_K : K \to C)$ is a crossed module where $\nu : R \to C, r \mapsto Jr$ is the natural map, and the action of $C$ on $K$ is given by $k^{Jr} = k^r$. }

 

 
\begin{Verbatim}[commandchars=!@|,fontsize=\small,frame=single,label=Example]
  
  !gapprompt@gap>| !gapinput@d8d8 := Group( (1,2,3,4), (1,3), (5,6,7,8), (5,7) );;|
  !gapprompt@gap>| !gapinput@X88 := XModByAutomorphismGroup( d8d8 );;|
  !gapprompt@gap>| !gapinput@Size2d( X88 );|
  [ 64, 2048 ]
  !gapprompt@gap>| !gapinput@Y88 := KernelCokernelXMod( X88 );;|
  !gapprompt@gap>| !gapinput@IdGroup(Y88);|
  [ [ 4, 2 ], [ 128, 928 ] ]
  !gapprompt@gap>| !gapinput@StructureDescription( Y88 );|
  [ "C2 x C2", "(D8 x D8) : C2" ] 
  
\end{Verbatim}
 }

 
\section{\textcolor{Chapter }{Pre\texttt{\symbol{45}}crossed modules}}\label{sect-precrossed-modules}
\logpage{[ 2, 3, 0 ]}
\hyperdef{L}{X7D435B6279032D4D}{}
{
  \index{pre\texttt{\symbol{45}}crossed module} 

\subsection{\textcolor{Chapter }{PreXModByBoundaryAndAction}}
\logpage{[ 2, 3, 1 ]}\nobreak
\hyperdef{L}{X8487BE427858C5C9}{}
{\noindent\textcolor{FuncColor}{$\triangleright$\enspace\texttt{PreXModByBoundaryAndAction({\mdseries\slshape bdy, act})\index{PreXModByBoundaryAndAction@\texttt{PreXModByBoundaryAndAction}}
\label{PreXModByBoundaryAndAction}
}\hfill{\scriptsize (operation)}}\\
\noindent\textcolor{FuncColor}{$\triangleright$\enspace\texttt{PreXModWithTrivialRange({\mdseries\slshape src, rng})\index{PreXModWithTrivialRange@\texttt{PreXModWithTrivialRange}}
\label{PreXModWithTrivialRange}
}\hfill{\scriptsize (operation)}}\\
\noindent\textcolor{FuncColor}{$\triangleright$\enspace\texttt{SubPreXMod({\mdseries\slshape X0, src, rng})\index{SubPreXMod@\texttt{SubPreXMod}}
\label{SubPreXMod}
}\hfill{\scriptsize (operation)}}\\


 If axiom ${\bf XMod\ 2}$ is \emph{not} satisfied, the corresponding structure is known as a \emph{pre\texttt{\symbol{45}}crossed module}. 

 A special case of this operation is when the range is a trivial group (not
necessarily a subgroup of the source), and so the action is trivial. This case
will be used when constructing a special type of double groupoid in Chapter \ref{chap-double}. }

 

 
\begin{Verbatim}[commandchars=!@|,fontsize=\small,frame=single,label=Example]
  
  !gapprompt@gap>| !gapinput@b1 := (11,12,13,14,15,16,17,18);;  b2 := (12,18)(13,17)(14,16);;|
  !gapprompt@gap>| !gapinput@d16 := Group( b1, b2 );;|
  !gapprompt@gap>| !gapinput@sk4 := Subgroup( d16, [ b1^4, b2 ] );;|
  !gapprompt@gap>| !gapinput@SetName( d16, "d16" );  SetName( sk4, "sk4" );|
  !gapprompt@gap>| !gapinput@bdy16 := GroupHomomorphismByImages( d16, sk4, [b1,b2], [b1^4,b2] );;|
  !gapprompt@gap>| !gapinput@aut1 := GroupHomomorphismByImages( d16, d16, [b1,b2], [b1^5,b2] );;|
  !gapprompt@gap>| !gapinput@aut2 := GroupHomomorphismByImages( d16, d16, [b1,b2], [b1,b2^4*b2] );;|
  !gapprompt@gap>| !gapinput@aut16 := Group( [ aut1, aut2 ] );;|
  !gapprompt@gap>| !gapinput@act16 := GroupHomomorphismByImages( sk4, aut16, [b1^4,b2], [aut1,aut2] );;|
  !gapprompt@gap>| !gapinput@P16 := PreXModByBoundaryAndAction( bdy16, act16 );|
  [d16->sk4]
  !gapprompt@gap>| !gapinput@IsXMod( P16 );|
  false
  !gapprompt@gap>| !gapinput@Q16 := PreXModWithTrivialRange( d16, d16 ); |
  [d16->Group( [ () ] )]
  !gapprompt@gap>| !gapinput@SQ16 := SubPreXMod( Q16, sk4, Group( [()] ) );; |
  !gapprompt@gap>| !gapinput@Display(SQ16);|
  Crossed module :- 
  : Source group has generators:
    [ (11,15)(12,16)(13,17)(14,18), (12,18)(13,17)(14,16) ]
  : Range group has generators:
    [ () ]
  : Boundary homomorphism maps source generators to:
    [ (), () ]
    The automorphism group is trivial
  
\end{Verbatim}
 \index{Peiffer subgroup} 

\subsection{\textcolor{Chapter }{PeifferSubgroup}}
\logpage{[ 2, 3, 2 ]}\nobreak
\hyperdef{L}{X8527F4C07A8F359E}{}
{\noindent\textcolor{FuncColor}{$\triangleright$\enspace\texttt{PeifferSubgroup({\mdseries\slshape X0})\index{PeifferSubgroup@\texttt{PeifferSubgroup}}
\label{PeifferSubgroup}
}\hfill{\scriptsize (attribute)}}\\
\noindent\textcolor{FuncColor}{$\triangleright$\enspace\texttt{XModByPeifferQuotient({\mdseries\slshape prexmod})\index{XModByPeifferQuotient@\texttt{XModByPeifferQuotient}}
\label{XModByPeifferQuotient}
}\hfill{\scriptsize (attribute)}}\\


 The \emph{Peiffer subgroup} $P$ of a pre\texttt{\symbol{45}}crossed module $\calX$ is the subgroup of ${\rm ker}(\partial)$ generated by \emph{Peiffer commutators} 
\[ \lfloor s_1,s_2 \rfloor ~=~ (s_1^{-1})^{\partial s_2}~s_2^{-1}~s_1~s_2 ~=~
\langle \partial s_2, s_1 \rangle\ [s_1,s_2]~. \]
 Then $\calP = (0 : P \to \{1_R\})$ is a normal sub\texttt{\symbol{45}}pre\texttt{\symbol{45}}crossed module of $\calX$ and $\calX/\calP = (\partial : S/P \to R)$ is a crossed module. 

 In the following example the Peiffer subgroup is cyclic of size $4$. }

 

 
\begin{Verbatim}[commandchars=!@|,fontsize=\small,frame=single,label=Example]
  
  !gapprompt@gap>| !gapinput@P := PeifferSubgroup( P16 );|
  Group( [ (11,15)(12,16)(13,17)(14,18), (11,17,15,13)(12,18,16,14) ] )
  !gapprompt@gap>| !gapinput@X16 := XModByPeifferQuotient( P16 );|
  Peiffer([d16->sk4])
  !gapprompt@gap>| !gapinput@Display( X16 );|
  Crossed module Peiffer([d16->sk4]) :-
  : Source group has generators:
    [ f1, f2 ]
  : Range group has generators:
    [ (11,15)(12,16)(13,17)(14,18), (12,18)(13,17)(14,16) ]
  : Boundary homomorphism maps source generators to:
    [ (12,18)(13,17)(14,16), (11,15)(12,16)(13,17)(14,18) ]
    The automorphism group is trivial
  !gapprompt@gap>| !gapinput@iso16 := IsomorphismPermGroup( Source( X16 ) );;|
  !gapprompt@gap>| !gapinput@S16 := Image( iso16 );|
  Group([ (1,2), (3,4) ])   
  
\end{Verbatim}
 }

 
\section{\textcolor{Chapter }{Cat$^1$\texttt{\symbol{45}}groups and pre\texttt{\symbol{45}}cat$^1$\texttt{\symbol{45}}groups}}\label{sect-cat1}
\logpage{[ 2, 4, 0 ]}
\hyperdef{L}{X7AAABC1D7E110988}{}
{
  \index{cat$^1$\texttt{\symbol{45}}group} In \cite{L1}, Loday reformulated the notion of a crossed module as a cat$^1$\texttt{\symbol{45}}group, namely a group $G$ with a pair of endomorphisms $t,h : G \to G$ having a common image $R$ and satisfying certain axioms. We find it computationally convenient to define
a cat$^1$\texttt{\symbol{45}}group $\calC = (e;t,h : G \to R )$ as having source group $G$, range group $R$, and three homomorphisms: two surjections $t,h : G \to R$ and an embedding $e : R \to G$ satisfying: 
\[ {\bf Cat\ 1} : ~t \circ e ~=~ h \circ e = {\rm id}_R, \qquad {\bf Cat\ 2} :
~[\ker t, \ker h] ~=~ \{ 1_G \}. \]
 It follows that $\;t \circ e \circ h = h,~ h \circ e \circ t = t,~ t \circ e \circ t = t~$ and $~h \circ e \circ h = h$. 

 See section \texttt{IsPreCat1GroupWithIdentityEmbedding} (\ref{IsPreCat1GroupWithIdentityEmbedding}) for the case when $t$ and $h$ are endomorphisms and $e$ an inclusion. 

\subsection{\textcolor{Chapter }{Cat1Group}}
\logpage{[ 2, 4, 1 ]}\nobreak
\label{mansect-cat1}
\hyperdef{L}{X7CF4C37F87D27EBA}{}
{\noindent\textcolor{FuncColor}{$\triangleright$\enspace\texttt{Cat1Group({\mdseries\slshape args})\index{Cat1Group@\texttt{Cat1Group}}
\label{Cat1Group}
}\hfill{\scriptsize (function)}}\\
\noindent\textcolor{FuncColor}{$\triangleright$\enspace\texttt{PreCat1Group({\mdseries\slshape args})\index{PreCat1Group@\texttt{PreCat1Group}}
\label{PreCat1Group}
}\hfill{\scriptsize (function)}}\\
\noindent\textcolor{FuncColor}{$\triangleright$\enspace\texttt{PreCat1GroupByTailHeadEmbedding({\mdseries\slshape t, h, e})\index{PreCat1GroupByTailHeadEmbedding@\texttt{PreCat1GroupByTailHeadEmbedding}}
\label{PreCat1GroupByTailHeadEmbedding}
}\hfill{\scriptsize (operation)}}\\
\noindent\textcolor{FuncColor}{$\triangleright$\enspace\texttt{PreCat1GroupWithIdentityEmbedding({\mdseries\slshape t, h})\index{PreCat1GroupWithIdentityEmbedding@\texttt{PreCat1GroupWithIdentityEmbedding}}
\label{PreCat1GroupWithIdentityEmbedding}
}\hfill{\scriptsize (operation)}}\\


 The global functions \texttt{Cat1Group} and \texttt{PreCat1Group} can be called in various ways. 
\begin{itemize}
\item  as \texttt{Cat1Group(t,h,e);} when $t,h,e$ are three homomorphisms, which is equivalent to \texttt{PreCat1GroupByTailHeadEmbedding(t,h,e);} 
\item  as \texttt{Cat1Group(t,h);} when $t,h$ are two endomorphisms, which is equivalent to \texttt{PreCat1GroupWithIdentityEmbedding(t,h);} 
\item  as \texttt{Cat1Group(t);} when $t=h$ is an endomorphism, which is equivalent to \texttt{PreCat1GroupWithIdentityEmbedding(t,t);} 
\item  as \texttt{Cat1Group(t,e);} when $t=h$ and $e$ are homomorphisms, which is equivalent to \texttt{PreCat1GroupByTailHeadEmbedding(t,t,e);} 
\item  as \texttt{Cat1Group(i,j,k);} when $i,j,k$ are integers, which is equivalent to \texttt{Cat1Select(i,j,k);} as described in section \ref{sect-cat1select}. 
\end{itemize}
 }

 

 
\begin{Verbatim}[commandchars=!@|,fontsize=\small,frame=single,label=Example]
  
  !gapprompt@gap>| !gapinput@g18gens := [ (1,2,3), (4,5,6), (2,3)(5,6) ];;     |
  !gapprompt@gap>| !gapinput@s3agens := [ (7,8,9), (8,9) ];;                |
  !gapprompt@gap>| !gapinput@g18 := Group( g18gens );;  SetName( g18, "g18" ); |
  !gapprompt@gap>| !gapinput@s3a := Group( s3agens );;  SetName( s3a, "s3a" );|
  !gapprompt@gap>| !gapinput@t1 := GroupHomomorphismByImages(g18,s3a,g18gens,[(7,8,9),(),(8,9)]);     |
  [ (1,2,3), (4,5,6), (2,3)(5,6) ] -> [ (7,8,9), (), (8,9) ]
  !gapprompt@gap>| !gapinput@h1 := GroupHomomorphismByImages(g18,s3a,g18gens,[(7,8,9),(7,8,9),(8,9)]);|
  [ (1,2,3), (4,5,6), (2,3)(5,6) ] -> [ (7,8,9), (7,8,9), (8,9) ]
  !gapprompt@gap>| !gapinput@e1 := GroupHomomorphismByImages(s3a,g18,s3agens,[(1,2,3),(2,3)(5,6)]);   |
  [ (7,8,9), (8,9) ] -> [ (1,2,3), (2,3)(5,6) ]
  !gapprompt@gap>| !gapinput@C18 := Cat1Group( t1, h1, e1 );|
  [g18=>s3a]
  
\end{Verbatim}
 

\subsection{\textcolor{Chapter }{Source (for cat1-groups)}}
\logpage{[ 2, 4, 2 ]}\nobreak
\hyperdef{L}{X7C4FFC4086531157}{}
{\noindent\textcolor{FuncColor}{$\triangleright$\enspace\texttt{Source({\mdseries\slshape C})\index{Source@\texttt{Source}!for cat1-groups}
\label{Source:for cat1-groups}
}\hfill{\scriptsize (attribute)}}\\
\noindent\textcolor{FuncColor}{$\triangleright$\enspace\texttt{Range({\mdseries\slshape C})\index{Range@\texttt{Range}!for cat1-groups}
\label{Range:for cat1-groups}
}\hfill{\scriptsize (attribute)}}\\
\noindent\textcolor{FuncColor}{$\triangleright$\enspace\texttt{TailMap({\mdseries\slshape C})\index{TailMap@\texttt{TailMap}}
\label{TailMap}
}\hfill{\scriptsize (attribute)}}\\
\noindent\textcolor{FuncColor}{$\triangleright$\enspace\texttt{HeadMap({\mdseries\slshape C})\index{HeadMap@\texttt{HeadMap}}
\label{HeadMap}
}\hfill{\scriptsize (attribute)}}\\
\noindent\textcolor{FuncColor}{$\triangleright$\enspace\texttt{RangeEmbedding({\mdseries\slshape C})\index{RangeEmbedding@\texttt{RangeEmbedding}}
\label{RangeEmbedding}
}\hfill{\scriptsize (attribute)}}\\
\noindent\textcolor{FuncColor}{$\triangleright$\enspace\texttt{KernelEmbedding({\mdseries\slshape C})\index{KernelEmbedding@\texttt{KernelEmbedding}}
\label{KernelEmbedding}
}\hfill{\scriptsize (attribute)}}\\
\noindent\textcolor{FuncColor}{$\triangleright$\enspace\texttt{Boundary({\mdseries\slshape C})\index{Boundary@\texttt{Boundary}!for cat1-groups}
\label{Boundary:for cat1-groups}
}\hfill{\scriptsize (attribute)}}\\
\noindent\textcolor{FuncColor}{$\triangleright$\enspace\texttt{Name({\mdseries\slshape C})\index{Name@\texttt{Name}!for cat1-groups}
\label{Name:for cat1-groups}
}\hfill{\scriptsize (attribute)}}\\
\noindent\textcolor{FuncColor}{$\triangleright$\enspace\texttt{Size2d({\mdseries\slshape C})\index{Size2d@\texttt{Size2d}!for cat1-groups}
\label{Size2d:for cat1-groups}
}\hfill{\scriptsize (attribute)}}\\


 These are the attributes of a cat$^1$\texttt{\symbol{45}}group $\calC$ in this implementation. 

 The maps $t,h$ are often referred to as the \emph{source} and \emph{target}, but we choose to call them the \emph{tail} and \emph{head} of $\calC$, because \emph{source} is the \textsf{GAP} term for the domain of a function. The \texttt{RangeEmbedding} is the embedding of \texttt{R} in \texttt{G}, the \texttt{KernelEmbedding} is the inclusion of the kernel of \texttt{t} in \texttt{G}, and the \texttt{Boundary} is the restriction of \texttt{h} to the kernel of \texttt{t}. It is frequently the case that $t=h$, but not in the example \texttt{C18} above. }

 
\begin{Verbatim}[commandchars=!@|,fontsize=\small,frame=single,label=Example]
  
  !gapprompt@gap>| !gapinput@[ Source( C18 ), Range( C18 ) ];|
  [ g18, s3a ]
  !gapprompt@gap>| !gapinput@TailMap( C18 );|
  [ (1,2,3), (4,5,6), (2,3)(5,6) ] -> [ (7,8,9), (), (8,9) ]
  !gapprompt@gap>| !gapinput@HeadMap( C18 );|
  [ (1,2,3), (4,5,6), (2,3)(5,6) ] -> [ (7,8,9), (7,8,9), (8,9) ]
  !gapprompt@gap>| !gapinput@RangeEmbedding( C18 );|
  [ (7,8,9), (8,9) ] -> [ (1,2,3), (2,3)(5,6) ]
  !gapprompt@gap>| !gapinput@Kernel( C18 );|
  Group([ (4,5,6) ])
  !gapprompt@gap>| !gapinput@KernelEmbedding( C18 );|
  [ (4,5,6) ] -> [ (4,5,6) ]
  !gapprompt@gap>| !gapinput@Name( C18 );|
  "[g18=>s3a]"
  !gapprompt@gap>| !gapinput@Size2d( C18 );|
  [ 18, 6 ]
  !gapprompt@gap>| !gapinput@StructureDescription( C18 );|
  [ "(C3 x C3) : C2", "S3" ]
  
\end{Verbatim}
 The next four subsections contain some more constructors for cat$^1$\texttt{\symbol{45}}groups. 

\subsection{\textcolor{Chapter }{DiagonalCat1Group}}
\logpage{[ 2, 4, 3 ]}\nobreak
\hyperdef{L}{X79944C7B87F767FD}{}
{\noindent\textcolor{FuncColor}{$\triangleright$\enspace\texttt{DiagonalCat1Group({\mdseries\slshape genG})\index{DiagonalCat1Group@\texttt{DiagonalCat1Group}}
\label{DiagonalCat1Group}
}\hfill{\scriptsize (operation)}}\\


 This operation constructs examples of cat$^1$\texttt{\symbol{45}}groups of the form $G \times G \Rightarrow G$. The tail map is the identity on the first factor and kills of the second,
while the head map does the reverse. The range embedding maps $G$ to the diagonal in $G \times G$. The corresponding crossed module is isomorphic to the identity crossed
module on $G$. }

 

 
\begin{Verbatim}[commandchars=!@|,fontsize=\small,frame=single,label=Example]
  
  !gapprompt@gap>| !gapinput@C4 := DiagonalCat1Group( [ (1,2,3), (2,3,4) ] );;|
  !gapprompt@gap>| !gapinput@SetName( Source(C4), "a4a4" );  SetName( Range(C4_, "a4d" );|
  !gapprompt@gap>| !gapinput@Display( C4 );|
  Cat1-group [a4a4=>a4d] :- 
  : Source group a4a4 has generators:
    [ (1,2,3), (2,3,4), (5,6,7), (6,7,8) ]
  : Range group a4d has generators:
    [ ( 9,10,11), (10,11,12) ]
  : tail homomorphism maps source generators to:
    [ ( 9,10,11), (10,11,12), (), () ]
  : head homomorphism maps source generators to:
    [ (), (), ( 9,10,11), (10,11,12) ]
  : range embedding maps range generators to:
    [ (1,2,3)(5,6,7), (2,3,4)(6,7,8) ]
  : kernel has generators:
    [ (5,6,7), (6,7,8) ]
  : boundary homomorphism maps generators of kernel to:
    [ ( 9,10,11), (10,11,12) ]
  : kernel embedding maps generators of kernel to:
    [ (5,6,7), (6,7,8) ]
  
\end{Verbatim}
 

\subsection{\textcolor{Chapter }{TransposeCat1Group}}
\logpage{[ 2, 4, 4 ]}\nobreak
\hyperdef{L}{X79385660821E54A3}{}
{\noindent\textcolor{FuncColor}{$\triangleright$\enspace\texttt{TransposeCat1Group({\mdseries\slshape C0})\index{TransposeCat1Group@\texttt{TransposeCat1Group}}
\label{TransposeCat1Group}
}\hfill{\scriptsize (attribute)}}\\
\noindent\textcolor{FuncColor}{$\triangleright$\enspace\texttt{TransposeIsomorphism({\mdseries\slshape C0})\index{TransposeIsomorphism@\texttt{TransposeIsomorphism}}
\label{TransposeIsomorphism}
}\hfill{\scriptsize (attribute)}}\\


 The \emph{transpose} of a cat$^1$\texttt{\symbol{45}}group $C$ has the same source, range and embedding, but has the tail and head maps
interchanged. The \texttt{TransposeIsomorphism} gives the isomorphism between the two. }

 

 
\begin{Verbatim}[commandchars=!@|,fontsize=\small,frame=single,label=Example]
  
  !gapprompt@gap>| !gapinput@R4 := TransposeCat1Group( C4 );|
  [a4a4=>a4d]
  !gapprompt@gap>| !gapinput@Boundary( R4 );|
  [ (2,3,4), (1,2,3) ] -> [ (10,11,12), (9,10,11) ]
  !gapprompt@gap>| !gapinput@TailMap( R4 ) = HeadMap( R4 ); |
  false
  !gapprompt@gap>| !gapinput@TailMap( R4 ) = HeadMap( C4 ); |
  true
  !gapprompt@gap>| !gapinput@MappingGeneratorsImages( TransposeIsomorphism(C4) );|
  [ [ [ (1,2,3), (2,3,4), (5,6,7), (6,7,8) ], 
        [ (5,6,7), (6,7,8), (1,2,3), (2,3,4) ] ], 
    [ [ (9,10,11), (10,11,12) ], [ (9,10,11), (10,11,12) ] ] ]
  
\end{Verbatim}
 

\subsection{\textcolor{Chapter }{Cat1GroupByPeifferQuotient}}
\logpage{[ 2, 4, 5 ]}\nobreak
\hyperdef{L}{X87544FAD873672E1}{}
{\noindent\textcolor{FuncColor}{$\triangleright$\enspace\texttt{Cat1GroupByPeifferQuotient({\mdseries\slshape P})\index{Cat1GroupByPeifferQuotient@\texttt{Cat1GroupByPeifferQuotient}}
\label{Cat1GroupByPeifferQuotient}
}\hfill{\scriptsize (operation)}}\\


 If $C = (e;t,h : G \to R)$ is a pre\texttt{\symbol{45}}cat$^1$\texttt{\symbol{45}}group, its Peiffer subgroup is $P = [\ker t,\ker h]$ and the associated cat$^1$\texttt{\symbol{45}}group $C_2$ has source $G/P$. In the example, $t=h : s4 \to c2$ with $\ker t = \ker h = a4$ and $P = [a4,a4]=k4$, so that $G/P = s4/k4 \cong s3$. }

 

 
\begin{Verbatim}[commandchars=!@E,fontsize=\small,frame=single,label=Example]
  
  !gapprompt@gap>E !gapinput@s4 := Group( (1,2,3), (3,4) );;  SetName( s4, "s4" ); E
  !gapprompt@gap>E !gapinput@h := GroupHomomorphismByImages( s4, s4, [(1,2,3),(3,4)], [(),(3,4)] );;E
  !gapprompt@gap>E !gapinput@c2 := Image( h );;  SetName( c2, "c2" );E
  !gapprompt@gap>E !gapinput@C := PreCat1Group( h, h );E
  [s4=>c2]
  !gapprompt@gap>E !gapinput@P := PeifferSubgroupPreCat1Group( C );E
  Group([ (1,3)(2,4), (1,2)(3,4) ])
  !gapprompt@gap>E !gapinput@C2 := Cat1GroupByPeifferQuotient( C );E
  [Group( [ f1, f2 ] )=>c2]
  !gapprompt@gap>E !gapinput@StructureDescription( C2 );E
  [ "S3", "C2" ]
  !gapprompt@gap>E !gapinput@rec2 := PreXModRecordOfPreCat1Group( C );;E
  !gapprompt@gap>E !gapinput@XC := rec2.prexmod;;E
  !gapprompt@gap>E !gapinput@StructureDescription( XC );  E
  [ "A4", "C2" ]
  !gapprompt@gap>E !gapinput@XC2 := XModByPeifferQuotient( XC );;E
  !gapprompt@gap>E !gapinput@StructureDescription( XC2 );E
  [ "C3", "C2" ]
  !gapprompt@gap>E !gapinput@CXC2 := Cat1GroupOfXMod( XC2 );;E
  !gapprompt@gap>E !gapinput@StructureDescription( CXC2 );E
  [ "S3", "C2" ]
  !gapprompt@gap>E !gapinput@IsomorphismCat1Groups( C2, CXC2 );E
  [[Group( [ f1, f2 ] ) => c2] => [(..|X..) => c2]]
  
\end{Verbatim}
 

\subsection{\textcolor{Chapter }{SubCat1Group}}
\logpage{[ 2, 4, 6 ]}\nobreak
\hyperdef{L}{X85D9C5F881DBA9FC}{}
{\noindent\textcolor{FuncColor}{$\triangleright$\enspace\texttt{SubCat1Group({\mdseries\slshape C1, S1})\index{SubCat1Group@\texttt{SubCat1Group}}
\label{SubCat1Group}
}\hfill{\scriptsize (operation)}}\\
\noindent\textcolor{FuncColor}{$\triangleright$\enspace\texttt{SubPreCat1Group({\mdseries\slshape C1, S1})\index{SubPreCat1Group@\texttt{SubPreCat1Group}}
\label{SubPreCat1Group}
}\hfill{\scriptsize (operation)}}\\


 $S_1$ is a sub\texttt{\symbol{45}}cat$^1$\texttt{\symbol{45}}group of $C_1$ provided the source and range of $S_1$ are subgroups of the source and range of $C_1$ and that the tail, head and embedding of $S_1$ are the appropriate restrictions of those of $C_1$. 

 }

 
\begin{Verbatim}[commandchars=!@|,fontsize=\small,frame=single,label=Example]
  
  !gapprompt@gap>| !gapinput@s3 := Subgroup( s4, [(2,3),(3,4)] );; |
  !gapprompt@gap>| !gapinput@res := DoublyRestrictedMapping( h, s3, s3 );; |
  !gapprompt@gap>| !gapinput@S := PreCat1Group( res, res );|
  [Group( [ (2,3), (3,4) ] )=>Group( [ (3,4), (3,4) ] )]
  
\end{Verbatim}
 

\subsection{\textcolor{Chapter }{DirectProduct (for cat1-groups)}}
\logpage{[ 2, 4, 7 ]}\nobreak
\hyperdef{L}{X7CE4F14585F6D473}{}
{\noindent\textcolor{FuncColor}{$\triangleright$\enspace\texttt{DirectProduct({\mdseries\slshape C1, C2})\index{DirectProduct@\texttt{DirectProduct}!for cat1-groups}
\label{DirectProduct:for cat1-groups}
}\hfill{\scriptsize (operation)}}\\


 The direct product $\calC_{1} \times \calC_{2}$ of two cat$^1$\texttt{\symbol{45}}groups has source $G_1 \times G_2$ and range $R_1 \times R_2$. The tail, head and embedding maps are $t_1 \times t_2$, $h_1 \times h_2$ and $e_1 \times e_2$. The embeddings and projections are constructed automatically, and placed in
the \texttt{DirectProductInfo} attribute, together with the two \emph{objects} $\calC_{1}$ and $\calC_{2}$. 

 The example constructs the product of two of the cat$^1$\texttt{\symbol{45}}groups constructed above. 

 }

 
\begin{Verbatim}[commandchars=!@|,fontsize=\small,frame=single,label=Example]
  
  !gapprompt@gap>| !gapinput@C418 := DirectProduct( C4, C18 );|
  [(a4a4xg18)=>(a4d x s3a)]
  !gapprompt@gap>| !gapinput@infoC418 := DirectProductInfo( C418 );|
  rec( 
    embeddings := [ [[a4a4=>a4d] => [(a4a4xg18)=>(a4d x s3a)]], 
        [[g18=>s3a] => [(a4a4xg18)=>(a4d x s3a)]] ], 
    objects := [ [a4a4=>a4d], [g18=>s3a] ], 
    projections := [ [[(a4a4xg18)=>(a4d x s3a)] => [a4a4=>a4d]], 
        [[(a4a4xg18)=>(a4d x s3a)] => [g18=>s3a]] ] )
  !gapprompt@gap>| !gapinput@t418 := TailMap( C418 );|
  [ (1,2,3), (2,3,4), (5,6,7), (6,7,8), (9,10,11), (12,13,14), (10,11)(13,14) 
   ] -> [ (1,2,3), (2,3,4), (), (), (5,6,7), (), (6,7) ]
  !gapprompt@gap>| !gapinput@h418 := HeadMap( C418 );|
  [ (1,2,3), (2,3,4), (5,6,7), (6,7,8), (9,10,11), (12,13,14), (10,11)(13,14) 
   ] -> [ (), (), (1,2,3), (2,3,4), (5,6,7), (5,6,7), (6,7) ]
  !gapprompt@gap>| !gapinput@e418 := RangeEmbedding( C418 );|
  [ (1,2,3), (2,3,4), (5,6,7), (6,7) ] -> [ (1,2,3)(5,6,7), (2,3,4)(6,7,8), 
    (9,10,11), (10,11)(13,14) ]
  
\end{Verbatim}
 }

 
\section{\textcolor{Chapter }{ Properties of cat$^1$\texttt{\symbol{45}}groups and pre\texttt{\symbol{45}}cat$^1$\texttt{\symbol{45}}groups }}\label{sect-properties-cat1}
\logpage{[ 2, 5, 0 ]}
\hyperdef{L}{X8317816A8361F88C}{}
{
  Many of the properties listed in section \ref{sect-properties-xmod} apply to pre\texttt{\symbol{45}}cat$^1$\texttt{\symbol{45}}groups and to cat$^1$\texttt{\symbol{45}}groups since these are also 2d\texttt{\symbol{45}}groups.
There are also more specific properties. 

\subsection{\textcolor{Chapter }{IsCat1Group}}
\logpage{[ 2, 5, 1 ]}\nobreak
\hyperdef{L}{X78E03FAB84A57D03}{}
{\noindent\textcolor{FuncColor}{$\triangleright$\enspace\texttt{IsCat1Group({\mdseries\slshape C0})\index{IsCat1Group@\texttt{IsCat1Group}}
\label{IsCat1Group}
}\hfill{\scriptsize (property)}}\\
\noindent\textcolor{FuncColor}{$\triangleright$\enspace\texttt{IsPreXCat1Group({\mdseries\slshape C0})\index{IsPreXCat1Group@\texttt{IsPreXCat1Group}}
\label{IsPreXCat1Group}
}\hfill{\scriptsize (property)}}\\
\noindent\textcolor{FuncColor}{$\triangleright$\enspace\texttt{IsIdentityCat1Group({\mdseries\slshape C0})\index{IsIdentityCat1Group@\texttt{IsIdentityCat1Group}}
\label{IsIdentityCat1Group}
}\hfill{\scriptsize (property)}}\\


 \texttt{IsIdentityCat1Group(C0)} is true when the head and tail maps of \texttt{C0} are identity mappings. }

 

 
\begin{Verbatim}[commandchars=!@|,fontsize=\small,frame=single,label=Example]
  
  !gapprompt@gap>| !gapinput@G8 := SmallGroup( 288, 956 );  SetName( G8, "G8" );|
  <pc group of size 288 with 7 generators>
  !gapprompt@gap>| !gapinput@d12 := DihedralGroup( 12 );  SetName( d12, "d12" );|
  <pc group of size 12 with 3 generators>
  !gapprompt@gap>| !gapinput@a1 := d12.1;;  a2 := d12.2;;  a3 := d12.3;;  a0 := One( d12 );;|
  !gapprompt@gap>| !gapinput@gensG8 := GeneratorsOfGroup( G8 );;|
  !gapprompt@gap>| !gapinput@t8 := GroupHomomorphismByImages( G8, d12, gensG8,|
  !gapprompt@>| !gapinput@          [ a0, a1*a3, a2*a3, a0, a0, a3, a0 ] );;|
  !gapprompt@gap>| !gapinput@h8 := GroupHomomorphismByImages( G8, d12, gensG8,|
  !gapprompt@>| !gapinput@          [ a1*a2*a3, a0, a0, a2*a3, a0, a0, a3^2 ] );;                   |
  !gapprompt@gap>| !gapinput@e8 := GroupHomomorphismByImages( d12, G8, [a1,a2,a3],|
  !gapprompt@>| !gapinput@       [ G8.1*G8.2*G8.4*G8.6^2, G8.3*G8.4*G8.6^2*G8.7, G8.6*G8.7^2 ] );|
  [ f1, f2, f3 ] -> [ f1*f2*f4*f6^2, f3*f4*f6^2*f7, f6*f7^2 ]
  !gapprompt@gap>| !gapinput@C8 := PreCat1GroupByTailHeadEmbedding( t8, h8, e8 );|
  [G8=>d12]
  !gapprompt@gap>| !gapinput@IsCat1Group( C8 );|
  true
  !gapprompt@gap>| !gapinput@KnownPropertiesOfObject( C8 );|
  [ "CanEasilyCompareElements", "CanEasilySortElements", "IsDuplicateFree", 
    "IsGeneratorsOfSemigroup", "IsPreCat1Domain", "IsPc2DimensionalGroup", 
    "IsPreXMod", "IsPreCat1Group", "IsCat1Group", "IsIdentityPreCat1Group", 
    "IsPreCat1GroupWithIdentityEmbedding" ]
  
\end{Verbatim}
 

\subsection{\textcolor{Chapter }{IsPreCat1GroupWithIdentityEmbedding}}
\logpage{[ 2, 5, 2 ]}\nobreak
\hyperdef{L}{X7B7CF88F83B0129D}{}
{\noindent\textcolor{FuncColor}{$\triangleright$\enspace\texttt{IsPreCat1GroupWithIdentityEmbedding({\mdseries\slshape C0})\index{IsPreCat1GroupWithIdentityEmbedding@\texttt{IsPreCat1GroupWithIdentityEmbedding}}
\label{IsPreCat1GroupWithIdentityEmbedding}
}\hfill{\scriptsize (property)}}\\
\noindent\textcolor{FuncColor}{$\triangleright$\enspace\texttt{IsomorphicPreCat1GroupWithIdentityEmbedding({\mdseries\slshape C0})\index{IsomorphicPreCat1GroupWithIdentityEmbedding@\texttt{Isomorphic}\-\texttt{Pre}\-\texttt{Cat1}\-\texttt{Group}\-\texttt{With}\-\texttt{Identity}\-\texttt{Embedding}}
\label{IsomorphicPreCat1GroupWithIdentityEmbedding}
}\hfill{\scriptsize (attribute)}}\\
\noindent\textcolor{FuncColor}{$\triangleright$\enspace\texttt{IsomorphismToPreCat1GroupWithIdentityEmbedding({\mdseries\slshape C0})\index{IsomorphismToPreCat1GroupWithIdentityEmbedding@\texttt{Isomorphism}\-\texttt{To}\-\texttt{Pre}\-\texttt{Cat1}\-\texttt{Group}\-\texttt{With}\-\texttt{Identity}\-\texttt{Embedding}}
\label{IsomorphismToPreCat1GroupWithIdentityEmbedding}
}\hfill{\scriptsize (attribute)}}\\


 \texttt{IsPreCat1GroupWithIdentityEmbedding(C0)} is true when the range embedding of \texttt{C0} is an inclusion mapping. (This property used to be called \texttt{IsPreCat1GroupByEndomorphisms} but, as the example below shows, when the tail and head maps are endomorphisms
the range embedding need not be a subgroup inclusion.) When this is not the
case, we may replace $t$ and $h$ by $t*e$ and $h*e$, and $e$ by the inclusion mapping of the image of $e$. This gives an isomorphic cat$^1$\texttt{\symbol{45}}group for which \texttt{IsPreCat1GroupWithIdentityEmbedding} is true. This is the \texttt{IsomorphicPreCat1GroupWithIdentityEmbedding} of \texttt{C0}, and \texttt{IsomorphismToPreCat1GroupWithIdentityEmbedding} is the isomorphism between them. (See the next chapter for mappings of cat$^1$\texttt{\symbol{45}}groups.) }

 

 
\begin{Verbatim}[commandchars=!@|,fontsize=\small,frame=single,label=Example]
  
  !gapprompt@gap>| !gapinput@G5 := Group( (1,2,3,4,5) );;                                             |
  !gapprompt@gap>| !gapinput@t := GroupHomomorphismByImages( G5, G5, [(1,2,3,4,5)], [(1,5,4,3,2)] );;|
  !gapprompt@gap>| !gapinput@PC5 := PreCat1GroupByTailHeadEmbedding( t, t, t );|
  [Group( [ (1,2,3,4,5) ] )=>Group( [ (1,2,3,4,5) ] )]
  !gapprompt@gap>| !gapinput@IsPreCat1GroupWithIdentityEmbedding( PC5 );|
  false
  !gapprompt@gap>| !gapinput@IPC5 := IsomorphicPreCat1GroupWithIdentityEmbedding( PC5 );|
  [Group( [ (1,2,3,4,5) ] )=>Group( [ (1,2,3,4,5) ] )]
  !gapprompt@gap>| !gapinput@TailMap( IPC5 ); RangeEmbedding( IPC5 );|
  [ (1,2,3,4,5) ] -> [ (1,2,3,4,5) ]
  [ (1,2,3,4,5) ] -> [ (1,2,3,4,5) ]
  
\end{Verbatim}
 

\subsection{\textcolor{Chapter }{Cat1GroupOfXMod}}
\logpage{[ 2, 5, 3 ]}\nobreak
\hyperdef{L}{X82F10A59867C765D}{}
{\noindent\textcolor{FuncColor}{$\triangleright$\enspace\texttt{Cat1GroupOfXMod({\mdseries\slshape X0})\index{Cat1GroupOfXMod@\texttt{Cat1GroupOfXMod}}
\label{Cat1GroupOfXMod}
}\hfill{\scriptsize (attribute)}}\\
\noindent\textcolor{FuncColor}{$\triangleright$\enspace\texttt{XModOfCat1Group({\mdseries\slshape C0})\index{XModOfCat1Group@\texttt{XModOfCat1Group}}
\label{XModOfCat1Group}
}\hfill{\scriptsize (attribute)}}\\
\noindent\textcolor{FuncColor}{$\triangleright$\enspace\texttt{PreCat1GroupRecordOfPreXMod({\mdseries\slshape P0})\index{PreCat1GroupRecordOfPreXMod@\texttt{PreCat1GroupRecordOfPreXMod}}
\label{PreCat1GroupRecordOfPreXMod}
}\hfill{\scriptsize (attribute)}}\\
\noindent\textcolor{FuncColor}{$\triangleright$\enspace\texttt{PreXModRecordOfPreCat1Group({\mdseries\slshape P0})\index{PreXModRecordOfPreCat1Group@\texttt{PreXModRecordOfPreCat1Group}}
\label{PreXModRecordOfPreCat1Group}
}\hfill{\scriptsize (attribute)}}\\


 The category of crossed modules is equivalent to the category of cat$^1$\texttt{\symbol{45}}groups, and the functors between these two categories may
be described as follows. Starting with the crossed module $\calX = (\partial : S \to R)$ the group $G$ is defined as the semidirect product $G = R \ltimes S$ using the action from $\calX$, with multiplication rule 
\[ (r_1,s_1)(r_2,s_2) ~=~ (r_1r_2,{s_1}^{r_2}s_2). \]
 The structural morphisms are given by 
\[ t(r,s) = r, \quad h(r,s) = r (\partial s), \quad er = (r,1). \]
 On the other hand, starting with a cat$^1$\texttt{\symbol{45}}group $ \calC = (e;t,h : G \to R)$, we define $ S = \ker t$, the range $R$ is unchanged, and $ \partial = h\!\mid_S $. The action of $R$ on $S$ is conjugation in $G$ via the embedding of $R$ in $G$. 

 As from version 2.74, the attribute \texttt{PreCat1GroupRecordOfPreXMod} of a pre\texttt{\symbol{45}}crossed modute $X = (\partial : S \to R)$ returns a record with fields 
\begin{itemize}
\item  \texttt{.precat1}, the pre\texttt{\symbol{45}}cat$^1$\texttt{\symbol{45}}group $C = (e;t,h: G \to R)$ of $X$, where $G = R \ltimes S$; 
\item  \texttt{.iscat1}, true if $C$ is a cat$^1$\texttt{\symbol{45}}group; 
\item  \texttt{.xmodSourceEmbedding}, the image $S'$ of $S$ in $G$; 
\item  \texttt{.xmodSourceEmbeddingIsomorphism}, the isomorphism $S \to S'$; 
\item  \texttt{.xmodRangeEmbedding}, the image $R'$ of $R$ in $G$; 
\item  \texttt{.xmodRangeEmbeddingIsomorphism}, the isomorphism $R \to R'$; 
\end{itemize}
 

 }

 

 
\begin{Verbatim}[commandchars=!@|,fontsize=\small,frame=single,label=Example]
  
  !gapprompt@gap>| !gapinput@X8 := XModOfCat1Group( C8 );;|
  !gapprompt@gap>| !gapinput@Display( X8 );|
  
  Crossed module X([G8=>d12]) :- 
  : Source group has generators:
    [ f1, f4, f5, f7 ]
  : Range group d12 has generators:
    [ f1, f2, f3 ]
  : Boundary homomorphism maps source generators to:
    [ f1*f2*f3, f2*f3, <identity> of ..., f3^2 ]
  : Action homomorphism maps range generators to automorphisms:
    f1 --> { source gens --> [ f1*f5, f4*f5, f5, f7^2 ] }
    f2 --> { source gens --> [ f1*f5*f7^2, f4, f5, f7 ] }
    f3 --> { source gens --> [ f1*f7, f4, f5, f7 ] }
    These 3 automorphisms generate the group of automorphisms.
  : associated cat1-group is [G8=>d12]
  
  !gapprompt@gap>| !gapinput@StructureDescription(X8);|
  [ "D24", "D12" ]
  
  
\end{Verbatim}
 }

 
\section{\textcolor{Chapter }{Enumerating cat$^1$\texttt{\symbol{45}}groups with a given source}}\label{sect-allcat1}
\logpage{[ 2, 6, 0 ]}
\hyperdef{L}{X80D6CB4080417BFA}{}
{
  As the size of a group $G$ increases, the number of cat$^1$\texttt{\symbol{45}}groups with source $G$ increases rapidly. However, one is usually only interested in the isomorphism
classes of cat$^1$\texttt{\symbol{45}}groups with source $G$. An iterator \texttt{AllCat1GroupsIterator} is provided, which runs through the various cat$^1$\texttt{\symbol{45}}groups. This iterator finds, for each subgroup $R$ of $G$, the cat$^1$\texttt{\symbol{45}}groups with range $R$. It does this by running through the \texttt{AllSubgroupsIterator(G)} provided by the \textsf{Utils} package, and then using the iterator \texttt{AllCat1GroupsWithImageIterator(G,R)}. 

\subsection{\textcolor{Chapter }{AllCat1GroupsWithImage}}
\logpage{[ 2, 6, 1 ]}\nobreak
\hyperdef{L}{X7BDEBBF17CE6A6D4}{}
{\noindent\textcolor{FuncColor}{$\triangleright$\enspace\texttt{AllCat1GroupsWithImage({\mdseries\slshape G, R})\index{AllCat1GroupsWithImage@\texttt{AllCat1GroupsWithImage}}
\label{AllCat1GroupsWithImage}
}\hfill{\scriptsize (operation)}}\\
\noindent\textcolor{FuncColor}{$\triangleright$\enspace\texttt{AllCat1GroupsWithImageIterator({\mdseries\slshape G, R})\index{AllCat1GroupsWithImageIterator@\texttt{AllCat1GroupsWithImageIterator}}
\label{AllCat1GroupsWithImageIterator}
}\hfill{\scriptsize (operation)}}\\
\noindent\textcolor{FuncColor}{$\triangleright$\enspace\texttt{AllCat1GroupsWithImageNumber({\mdseries\slshape G, R})\index{AllCat1GroupsWithImageNumber@\texttt{AllCat1GroupsWithImageNumber}}
\label{AllCat1GroupsWithImageNumber}
}\hfill{\scriptsize (attribute)}}\\
\noindent\textcolor{FuncColor}{$\triangleright$\enspace\texttt{AllCat1GroupsWithImageUpToIsomorphism({\mdseries\slshape G, R})\index{AllCat1GroupsWithImageUpToIsomorphism@\texttt{All}\-\texttt{Cat1}\-\texttt{Groups}\-\texttt{With}\-\texttt{Image}\-\texttt{Up}\-\texttt{To}\-\texttt{Isomorphism}}
\label{AllCat1GroupsWithImageUpToIsomorphism}
}\hfill{\scriptsize (operation)}}\\


 The iterator \texttt{AllCat1GroupsWithImageIterator(G,R)} iterates through all the cat$^1$\texttt{\symbol{45}}groups with source \texttt{G} and range $R$. The attribute \texttt{AllCat1GroupsWithImageNumber(G)} runs through this iterator to count the number $n_R$ of these cat$^1$\texttt{\symbol{45}}groups. The operation \texttt{AllCat1GroupsWithImage(G)} returns a list containing these $n_R$ cat$^1$\texttt{\symbol{45}}groups. Since these lists can get very long, this
operation should only be used for simple cases. The operation \texttt{AllCat1GroupsWithImageUpToIsomorphism(G)} returns representatives of the isomorphism classes of these cat$^1$\texttt{\symbol{45}}groups. }

 

 
\begin{Verbatim}[commandchars=!@|,fontsize=\small,frame=single,label=Example]
  
  !gapprompt@gap>| !gapinput@d12 := DihedralGroup( IsPermGroup, 12 );  SetName( d12, "d12" );|
  Group([ (1,2,3,4,5,6), (2,6)(3,5) ])
  !gapprompt@gap>| !gapinput@c2 := Subgroup( d12, [ (1,6)(2,5)(3,4) ] );; |
  !gapprompt@gap>| !gapinput@AllCat1GroupsWithImageNumber( d12, c2 );|
  1
  !gapprompt@gap>| !gapinput@L12 := AllCat1GroupsWithImage( d12, c2 );|
  [ [d12=>Group( [ (), (1,6)(2,5)(3,4) ] )] ]
  
\end{Verbatim}
 

\subsection{\textcolor{Chapter }{AllCat1GroupsMatrix}}
\logpage{[ 2, 6, 2 ]}\nobreak
\hyperdef{L}{X7FBFC8C87FC1AC5A}{}
{\noindent\textcolor{FuncColor}{$\triangleright$\enspace\texttt{AllCat1GroupsMatrix({\mdseries\slshape G})\index{AllCat1GroupsMatrix@\texttt{AllCat1GroupsMatrix}}
\label{AllCat1GroupsMatrix}
}\hfill{\scriptsize (attribute)}}\\


 The operation \texttt{AllCat1GroupsMatrix(G)} constructs a symmetric matrix $M$ with rows and columns labelled by the idempotent endomorphisms $e_i$ on $G$, where $M_{ij} = 2$ if $e_i,e_j$ combine to form a cat$^1$\texttt{\symbol{45}}group; $M_{ij} = 1$ if they only form a pre\texttt{\symbol{45}}cat$^1$\texttt{\symbol{45}}group; and $M_{ij} = 0$ otherwise. The matrix is automatically printed out with dots in place of
zeroes. 

 In the example we see that the group $QD_{16}$ has $10$ idempotent endomorphisms and $5$ cat$^1$\texttt{\symbol{45}}groups, all of which are symmetric ($t=h$), and a further $9$ pre\texttt{\symbol{45}}cat$^1$\texttt{\symbol{45}}groups, $5$ of which are symmetric. (A cat$^1$\texttt{\symbol{45}}group and its transpose are not counted twice.) This
operation is intended to be used to illustrate how cat$^1$\texttt{\symbol{45}}groups are formed, and should only be used with groups of
low order. 

 The attribute \texttt{AllCat1GroupsNumber(G)} returns the number $n$ of these cat$^1$\texttt{\symbol{45}}groups. }

 

 
\begin{Verbatim}[commandchars=!@|,fontsize=\small,frame=single,label=Example]
  
  !gapprompt@gap>| !gapinput@qd16 := SmallGroup( 16, 8 );; |
  !gapprompt@gap>| !gapinput@AllCat1GroupsMatrix( qd16 );;                 |
  number of idempotent endomorphisms found = 10
  number of cat1-groups found = 5
  number of additional pre-cat1-groups found = 9
  1.........
  .21.......
  .11.......
  ...21.....
  ...11.....
  .....21...
  .....11...
  .......21.
  .......11.
  .........2
  
\end{Verbatim}
 

\subsection{\textcolor{Chapter }{AllCat1GroupsIterator}}
\logpage{[ 2, 6, 3 ]}\nobreak
\hyperdef{L}{X7FEB2FCE7D9ADA85}{}
{\noindent\textcolor{FuncColor}{$\triangleright$\enspace\texttt{AllCat1GroupsIterator({\mdseries\slshape G})\index{AllCat1GroupsIterator@\texttt{AllCat1GroupsIterator}}
\label{AllCat1GroupsIterator}
}\hfill{\scriptsize (operation)}}\\
\noindent\textcolor{FuncColor}{$\triangleright$\enspace\texttt{AllCat1GroupsUpToIsomorphism({\mdseries\slshape G})\index{AllCat1GroupsUpToIsomorphism@\texttt{AllCat1GroupsUpToIsomorphism}}
\label{AllCat1GroupsUpToIsomorphism}
}\hfill{\scriptsize (operation)}}\\
\noindent\textcolor{FuncColor}{$\triangleright$\enspace\texttt{AllCat1Groups({\mdseries\slshape G})\index{AllCat1Groups@\texttt{AllCat1Groups}}
\label{AllCat1Groups}
}\hfill{\scriptsize (operation)}}\\


 The iterator \texttt{AllCat1GroupsIterator(G)} iterates through all the cat$^1$\texttt{\symbol{45}}groups with source \texttt{G}. The operation \texttt{AllCat1Groups(G)} returns a list containing these $n$ cat$^1$\texttt{\symbol{45}}groups. Since these lists can get very long, this
operation should only be used for simple cases. The operation \texttt{AllCat1GroupsUpToIsomorphism(G)} returns representatives of the isomorphism classes of these subgroups. }

 

 
\begin{Verbatim}[commandchars=!@|,fontsize=\small,frame=single,label=Example]
  
  !gapprompt@gap>| !gapinput@iter := AllCat1GroupsIterator( d12 );;|
  !gapprompt@gap>| !gapinput@AllCat1GroupsNumber( d12 );|
  12
  !gapprompt@gap>| !gapinput@iso12 := AllCat1GroupsUpToIsomorphism( d12 );|
  [ [d12=>Group( [ (), (2,6)(3,5) ] )], 
    [d12=>Group( [ (1,4)(2,5)(3,6), (2,6)(3,5) ] )], 
    [d12=>Group( [ (1,5,3)(2,6,4), (2,6)(3,5) ] )], 
    [d12=>Group( [ (1,2,3,4,5,6), (2,6)(3,5) ] )] ]
  
\end{Verbatim}
 

\subsection{\textcolor{Chapter }{CatnGroupNumbers (for cat1-groups)}}
\logpage{[ 2, 6, 4 ]}\nobreak
\hyperdef{L}{X7C6346A17FEEDFA1}{}
{\noindent\textcolor{FuncColor}{$\triangleright$\enspace\texttt{CatnGroupNumbers({\mdseries\slshape G})\index{CatnGroupNumbers@\texttt{CatnGroupNumbers}!for cat1-groups}
\label{CatnGroupNumbers:for cat1-groups}
}\hfill{\scriptsize (attribute)}}\\
\noindent\textcolor{FuncColor}{$\triangleright$\enspace\texttt{CatnGroupLists({\mdseries\slshape G})\index{CatnGroupLists@\texttt{CatnGroupLists}!for cat1-groups}
\label{CatnGroupLists:for cat1-groups}
}\hfill{\scriptsize (attribute)}}\\
\noindent\textcolor{FuncColor}{$\triangleright$\enspace\texttt{InitCatnGroupRecords({\mdseries\slshape G})\index{InitCatnGroupRecords@\texttt{InitCatnGroupRecords}}
\label{InitCatnGroupRecords}
}\hfill{\scriptsize (operation)}}\\


 The attribute \texttt{CatnGroupNumbers} for a group $G$ is a mutable record which stores numbers of cat$^1$\texttt{\symbol{45}}groups, cat$^2$\texttt{\symbol{45}}groups, etc. as they are calculated. The field \texttt{CatnGroupNumbers(G).idem} is the number of idempotent endomorphisms of $G$. Similarly, \texttt{CatnGroupNumbers(G).cat1} is the number of cat$^1$\texttt{\symbol{45}}groups on $G$, while \texttt{CatnGroupNumbers(G).iso1} is the number of isomorphism classes of these cat$^1$\texttt{\symbol{45}}groups. Also \texttt{CatnGroupNumbers(G).symm} is the number of cat$^1$\texttt{\symbol{45}}groups whose \texttt{TailMap} is the same as the \texttt{HeadMap}, while \texttt{CatnGroupNumbers(G).siso} is the number of isomorphism classes of these symmetric cat$^1$\texttt{\symbol{45}}groups. Symmetric cat$^1$\texttt{\symbol{45}}groups are in one\texttt{\symbol{45}}one correspondence
with symmetric cat$^2$\texttt{\symbol{45}}groups. The attribute \texttt{CatnGroupLists} is used for storing results of cat$^2$\texttt{\symbol{45}}group calculations. }

 

 
\begin{Verbatim}[commandchars=!@|,fontsize=\small,frame=single,label=Example]
  
  !gapprompt@gap>| !gapinput@CatnGroupNumbers( d12 );|
  rec( cat1 := 12, idem := 21, iso1 := 4, siso := 4, symm := 12 )
  
\end{Verbatim}
 }

 
\section{\textcolor{Chapter }{Selection of a small cat$^1$\texttt{\symbol{45}}group}}\label{sect-cat1select}
\logpage{[ 2, 7, 0 ]}
\hyperdef{L}{X7A6A70BD86DE458D}{}
{
  \index{selection of a small cat$^1$\texttt{\symbol{45}}group} The \texttt{Cat1Group} function may also be used to select a cat$^1$\texttt{\symbol{45}}group from a data file. All cat$^1$\texttt{\symbol{45}}structures on groups of size up to $60$ (ordered according to the \textsf{GAP} 4 numbering of small groups) are stored in a list in file \texttt{cat1data.g}. Global variables \texttt{CAT1{\textunderscore}LIST{\textunderscore}MAX{\textunderscore}SIZE := 60}, \index{CAT1{\textunderscore}LIST{\textunderscore}MAX{\textunderscore}SIZE} \texttt{CAT1{\textunderscore}LIST{\textunderscore}CLASS{\textunderscore}SIZES} and \texttt{CAT1{\textunderscore}LIST{\textunderscore}NUMBERS} are also stored. \index{CAT1{\textunderscore}LIST{\textunderscore}NUMBERS} The second of these just stores the number of isomorphism classes of groups of
size \texttt{size}. The third stores the numbers of isomorphism classes of cat$^1$\texttt{\symbol{45}}groups for each of these groups. The data is read into the
list \texttt{CAT1{\textunderscore}LIST} only when this function is called. 

 This data was available in early versions of \textsf{XMod} with groups up to order $70$ covered. More recently a larger range of groups has become available in the
package \textsf{HAP}. The authors are indebted to Van Luyen Le in Galway for pointing out a number
of errors in the version of this list distributed up to version \textsc{2.24} of this package. 

\subsection{\textcolor{Chapter }{Cat1Select}}
\logpage{[ 2, 7, 1 ]}\nobreak
\hyperdef{L}{X7B8E67D880E380C8}{}
{\noindent\textcolor{FuncColor}{$\triangleright$\enspace\texttt{Cat1Select({\mdseries\slshape size, gpnum, num})\index{Cat1Select@\texttt{Cat1Select}}
\label{Cat1Select}
}\hfill{\scriptsize (operation)}}\\


 The function \texttt{Cat1Select} returns the cat$^1$\texttt{\symbol{45}}group numbered \texttt{num} whose source is the group \texttt{G := SmallGroup(size,gpnum)}. When $|G| \leqslant 60$ the data file in this package is used. For larger groups \texttt{SmallCat1Group} (see \ref{sect-hap}) is called, accessing the datafile in package \textsf{HAP}. 

 

 The example below is the first case in which $t \neq h$ and the associated conjugation crossed module is given by the normal subgroup \texttt{c3} of \texttt{s3}. }

 

 
\begin{Verbatim}[commandchars=!@|,fontsize=\small,frame=single,label=Example]
  
  !gapprompt@gap>| !gapinput@L18 := Cat1Select( 18 ); |
  Usage:  Cat1Select( size, gpnum, num );  where gpnum <= 5
  fail
  !gapprompt@gap>| !gapinput@## check the number of cat1-structures on the fourth group of order 18 |
  !gapprompt@gap>| !gapinput@Cat1Select( 18, 4 );|
  Usage:  Cat1Select( size, gpnum, num );  where num <= 4
  fail
  !gapprompt@gap>| !gapinput@## select the second of these cat1-structures |
  !gapprompt@gap>| !gapinput@B18 := Cat1Select( 18, 4, 2 );|
  [(C3 x C3) : C2=>Group( [ f1, <identity> of ..., f3 ] )]
  !gapprompt@gap>| !gapinput@## convert from a pc-cat1-group to a permutation cat1-group|
  !gapprompt@gap>| !gapinput@iso18 := IsomorphismPermObject( B18 );;|
  !gapprompt@gap>| !gapinput@PB18 := Image( iso18 );;|
  !gapprompt@gap>| !gapinput@Display( PB18 );|
  Cat1-group :- 
  : Source group has generators:
    [ (4,5,6), (1,2,3), (2,3)(5,6) ]
  : Range group has generators:
    [ (4,5,6), (2,3)(5,6) ]
  : tail homomorphism maps source generators to:
    [ (4,5,6), (), (2,3)(5,6) ]
  : head homomorphism maps source generators to:
    [ (4,5,6), (), (2,3)(5,6) ]
  : range embedding maps range generators to:
    [ (4,5,6), (2,3)(5,6) ]
  : kernel has generators:
    [ (1,2,3) ]
  : boundary homomorphism maps generators of kernel to:
    [ () ]
  : kernel embedding maps generators of kernel to:
    [ (1,2,3) ]
  
  !gapprompt@gap>| !gapinput@convert the result to the associated permutation crossed module |
  !gapprompt@gap>| !gapinput@Y18 := XModOfCat1Group( PB18 );; |
  !gapprompt@gap>| !gapinput@Display( Y18 ); |
  Crossed module :- 
  : Source group has generators:
    [ (1,2,3) ]
  : Range group has generators:
    [ (4,5,6), (2,3)(5,6) ]
  : Boundary homomorphism maps source generators to:
    [ () ]
  : Action homomorphism maps range generators to automorphisms:
    (4,5,6) --> { source gens --> [ (1,2,3) ] }
    (2,3)(5,6) --> { source gens --> [ (1,3,2) ] }
    These 2 automorphisms generate the group of automorphisms.
  : associated cat1-group is [Group( [ (4,5,6), (1,2,3), (2,3)(5,6) 
   ] ) => Group( [ (4,5,6), (2,3)(5,6) ] )]
  
\end{Verbatim}
 }

 
\section{\textcolor{Chapter }{More functions for crossed modules and cat$^1$\texttt{\symbol{45}}groups}}\label{sect-extra-fns}
\logpage{[ 2, 8, 0 ]}
\hyperdef{L}{X8614CDCF8063117F}{}
{
  Chapter \ref{chap-isclnc} contains functions for quotient crossed modules; centre of a crossed module;
commutator and derived subcrossed modules; etc. 

 Here we mention two functions for groups which have been extended to the
two\texttt{\symbol{45}}dimensional case. 

 

\subsection{\textcolor{Chapter }{IdGroup (for 2d-groups)}}
\logpage{[ 2, 8, 1 ]}\nobreak
\hyperdef{L}{X7831DB527CF9DD57}{}
{\noindent\textcolor{FuncColor}{$\triangleright$\enspace\texttt{IdGroup({\mdseries\slshape 2DimensionalGroup})\index{IdGroup@\texttt{IdGroup}!for 2d-groups}
\label{IdGroup:for 2d-groups}
}\hfill{\scriptsize (operation)}}\\
\noindent\textcolor{FuncColor}{$\triangleright$\enspace\texttt{StructureDescription({\mdseries\slshape 2DimensionalGroup})\index{StructureDescription@\texttt{StructureDescription}!for 2d-groups}
\label{StructureDescription:for 2d-groups}
}\hfill{\scriptsize (operation)}}\\


 These functions return two\texttt{\symbol{45}}element lists formed by applying
the function to the source and range of the 2d\texttt{\symbol{45}}group. }

 

 
\begin{Verbatim}[commandchars=!@|,fontsize=\small,frame=single,label=Example]
  
  !gapprompt@gap>| !gapinput@IdGroup( X8 );|
  [ [ 24, 6 ], [ 12, 4 ] ]
  !gapprompt@gap>| !gapinput@StructureDescription( C8 );|
  [ "(S3 x D24) : C2", "D12" ]
  
\end{Verbatim}
 There are also a number of functions which test for
sub\texttt{\symbol{45}}structures. 

 

\subsection{\textcolor{Chapter }{IsSubXMod}}
\logpage{[ 2, 8, 2 ]}\nobreak
\hyperdef{L}{X83BBC6818168C282}{}
{\noindent\textcolor{FuncColor}{$\triangleright$\enspace\texttt{IsSubXMod({\mdseries\slshape X0, S0})\index{IsSubXMod@\texttt{IsSubXMod}}
\label{IsSubXMod}
}\hfill{\scriptsize (operation)}}\\
\noindent\textcolor{FuncColor}{$\triangleright$\enspace\texttt{IsSubPreXMod({\mdseries\slshape X0, S0})\index{IsSubPreXMod@\texttt{IsSubPreXMod}}
\label{IsSubPreXMod}
}\hfill{\scriptsize (operation)}}\\
\noindent\textcolor{FuncColor}{$\triangleright$\enspace\texttt{IsSubCat1Group({\mdseries\slshape G0, R0})\index{IsSubCat1Group@\texttt{IsSubCat1Group}}
\label{IsSubCat1Group}
}\hfill{\scriptsize (operation)}}\\
\noindent\textcolor{FuncColor}{$\triangleright$\enspace\texttt{IsSubPreCat1Group({\mdseries\slshape G0, R0})\index{IsSubPreCat1Group@\texttt{IsSubPreCat1Group}}
\label{IsSubPreCat1Group}
}\hfill{\scriptsize (operation)}}\\
\noindent\textcolor{FuncColor}{$\triangleright$\enspace\texttt{IsSub2DimensionalGroup({\mdseries\slshape G0, R0})\index{IsSub2DimensionalGroup@\texttt{IsSub2DimensionalGroup}}
\label{IsSub2DimensionalGroup}
}\hfill{\scriptsize (operation)}}\\


 These functions test whether the second argument is a
sub\texttt{\symbol{45}}2d\texttt{\symbol{45}}group of the first argument. The
examples refer back to sub\texttt{\symbol{45}}2d\texttt{\symbol{45}}groups
created in sections \ref{sect-properties-xmod} and \ref{sect-cat1}. }

 

 
\begin{Verbatim}[commandchars=!@|,fontsize=\small,frame=single,label=Example]
  
  !gapprompt@gap>| !gapinput@IsSubXMod( X4, Y4 );|
  true
  !gapprompt@gap>| !gapinput@IsSubPreCat1Group( C, S );|
  true
  
\end{Verbatim}
 }

 
\section{\textcolor{Chapter }{The group groupoid associated to a cat$^1$\texttt{\symbol{45}}group}}\label{sect-gpgpd}
\logpage{[ 2, 9, 0 ]}
\hyperdef{L}{X7CFAB044817E5E91}{}
{
  A \emph{group groupoid} is an algebraic object which is both a groupoid and a group. The category of
group groupoids is equivalent to the categories of precrossed modules and
precat$^1$\texttt{\symbol{45}}groups. Starting with a (pre)cat$^1$\texttt{\symbol{45}}group $\calC = (e;t,h : G \to R)$, we form the groupoid $\calG$ having the elements of $R$ as objects and the elements of $G$ as arrows. The arrow $g$ has tail $tg$ and head $hg$. $\calG$ has one connected component for each coset of $tG$ in $R$. 

 The groupoid (partial) multiplication $*$ on these arrows is defined by: 
\[ (g_1 : r_1 \to r_2) * (g_2 : r_2 \to r_3) ~=~ (g_1(er_2^{-1})g_2 : r_1 \to
r_3). \]
 

 

\subsection{\textcolor{Chapter }{GroupGroupoid}}
\logpage{[ 2, 9, 1 ]}\nobreak
\hyperdef{L}{X7AF5AF668331321E}{}
{\noindent\textcolor{FuncColor}{$\triangleright$\enspace\texttt{GroupGroupoid({\mdseries\slshape precat1})\index{GroupGroupoid@\texttt{GroupGroupoid}}
\label{GroupGroupoid}
}\hfill{\scriptsize (attribute)}}\\


 The operation \texttt{GroupGroupoid} implements this construction. In the example we start with a crossed module $(C_3^2 \to S_3)$, form the associated cat$^1$\texttt{\symbol{45}}group $(S_3 \ltimes C_3^2 \Rightarrow S_3)$, and then form the group groupoid \texttt{gpd33}. Since the image of the boundary of the crossed module is $C_3$, with index $2$ in the range, the groupoid has two connected components, and the root objects
are $\{(),(12,13)\}$. The size of the vertex groups is $|\ker t \cap \ker h| = 3$, and the generators at the root objects are $() \to ( 4, 5, 6)( 7, 9, 8) \to ()$ and $(12,13) \to ( 2, 3)( 4, 6)( 7, 8) \to (12,13)$. 

 }

 

 
\begin{Verbatim}[commandchars=!@|,fontsize=\small,frame=single,label=Example]
  
  !gapprompt@gap>| !gapinput@s3 := Group( (11,12), (12,13) );; |
  !gapprompt@gap>| !gapinput@c3c3 := Group( [ (14,15,16), (17,18,19) ] );; |
  !gapprompt@gap>| !gapinput@bdy3 := GroupHomomorphismByImages( c3c3, s3, |
  !gapprompt@>| !gapinput@       [(14,15,16),(17,18,19)], [(11,12,13),(11,12,13)] );;|
  !gapprompt@gap>| !gapinput@a := GroupHomomorphismByImages( c3c3, c3c3, |
  !gapprompt@>| !gapinput@       [(14,15,16),(17,18,19)], [(14,16,15),(17,19,18)] );; |
  !gapprompt@gap>| !gapinput@aut := Group( [a] );; |
  !gapprompt@gap>| !gapinput@act := GroupHomomorphismByImages( s3, aut, [(11,12),(12,13)], [a,a] );;|
  !gapprompt@gap>| !gapinput@X33 := XModByBoundaryAndAction( bdy3, act );; |
  !gapprompt@gap>| !gapinput@C33 := Cat1GroupOfXMod( X33 );; |
  !gapprompt@gap>| !gapinput@G33 := Source( C33 );; |
  !gapprompt@gap>| !gapinput@gpd33 := GroupGroupoid( C33 ); |
  groupoid with 2 pieces:
  1:  single piece groupoid with rays: < Group( [ ()>-(4,5,6)(7,9,8)->() ] ), 
  [ (), (11,12,13), (11,13,12) ], [ ()>-()->(), ()>-(7,8,9)->(11,12,13), 
    ()>-(7,9,8)->(11,13,12) ] >
  2:  single piece groupoid with rays: < Group( 
  [ (12,13)>-(2,3)(4,6)(7,8)->(12,13) ] ), [ (12,13), (11,12), (11,13) ], 
  [ (12,13)>-(2,3)(5,6)(8,9)->(12,13), (12,13)>-(2,3)(5,6)(7,9)->(11,13), 
    (12,13)>-(2,3)(5,6)(7,8)->(11,12) ] >
  
\end{Verbatim}
 

\subsection{\textcolor{Chapter }{GroupGroupoidElement}}
\logpage{[ 2, 9, 2 ]}\nobreak
\hyperdef{L}{X8578AB6D7C1FC4F3}{}
{\noindent\textcolor{FuncColor}{$\triangleright$\enspace\texttt{GroupGroupoidElement({\mdseries\slshape precat1, root, g})\index{GroupGroupoidElement@\texttt{GroupGroupoidElement}}
\label{GroupGroupoidElement}
}\hfill{\scriptsize (operation)}}\\


 Since we need to define a second multiplication on the elements of $G$, we have to convert $g \in G$ into a new type of object, \texttt{GroupGroupoidElementType}, a record $e$ with fields: 
\begin{itemize}
\item  \texttt{e!.precat1}, the precat$^1$\texttt{\symbol{45}}group from which $\calG$ was formed; 
\item  \texttt{e!.root}, the root object of the component containing $e$; 
\item  \texttt{e!.element}, the element $g \in G$; 
\item  \texttt{e!.tail}, the tail object of the element $e$; 
\item  \texttt{e!.head}, the head object of the element $e$; 
\item  \texttt{e!.tailid}, the identity element at the tail object; 
\item  \texttt{e!.headid}, the identity element at the head object; 
\end{itemize}
 In the example we pick a particular pair of elements $g_1,g_2 \in G$, construct group groupoid elements $e_1,e_2$ from them, and show that $g_1*g_2$ and $e_1*e_2$ give very different results. (Warning: at present iterators for object groups
and homsets do not work.) }

 

 
\begin{Verbatim}[commandchars=@|A,fontsize=\small,frame=single,label=Example]
  
  @gapprompt|gap>A @gapinput|piece2 := Pieces( gpd33 )[2];;A
  @gapprompt|gap>A @gapinput|obs2 := piece2!.objects; A
  [ (12,13), (11,12), (11,13) ]
  @gapprompt|gap>A @gapinput|RaysOfGroupoid( piece2 );A
  [ (12,13)>-(2,3)(5,6)(8,9)->(12,13), (12,13)>-(2,3)(5,6)(7,9)->(11,13), 
    (12,13)>-(2,3)(5,6)(7,8)->(11,12) ]
  @gapprompt|gap>A @gapinput|g1 := (1,2)(5,6)(7,9);; A
  @gapprompt|gap>A @gapinput|g2 := (2,3)(4,5)(7,8);;                         A
  @gapprompt|gap>A @gapinput|g1 * g2;A
  (1,3,2)(4,5,6)(7,9,8)
  @gapprompt|gap>A @gapinput|e1 := GroupGroupoidElement( C33, (12,13), g1 ); A
  (11,12)>-(1,2)(5,6)(7,9)->(12,13)
  @gapprompt|gap>A @gapinput|e2 := GroupGroupoidElement( C33, (12,13), g2 );A
  (12,13)>-(2,3)(4,5)(7,8)->(11,13)
  @gapprompt|gap>A @gapinput|e1*e2;A
  (11,12)>-(1,2)(4,5)(8,9)->(11,13)
  @gapprompt|gap>A @gapinput|e2^-1;A
  (11,13)>-(1,3)(4,6)(7,9)->(12,13)
  @gapprompt|gap>A @gapinput|obgp := ObjectGroup( gpd33, (11,12) );;A
  @gapprompt|gap>A @gapinput|GeneratorsOfGroup( obgp )[1];A
  (11,13)>-( 1, 3)( 4, 6)( 7, 8)->(11,13)
  @gapprompt|gap>A @gapinput|Homset( gpd33, (11,12), (11,13) );A
  <homset (11,12) -> (11,13) with head group Group( 
  [ (11,12)>-( 1, 2)( 4, 6)( 7, 8)->(11,12) ] )>
  
\end{Verbatim}
 }

 }

         
\chapter{\textcolor{Chapter }{2d\texttt{\symbol{45}}mappings}}\label{chap-gpmap2}
\logpage{[ 3, 0, 0 ]}
\hyperdef{L}{X815144D67C1D1AE3}{}
{
  \index{2d\texttt{\symbol{45}}mapping} 
\section{\textcolor{Chapter }{Morphisms of 2\texttt{\symbol{45}}dimensional groups}}\logpage{[ 3, 1, 0 ]}
\hyperdef{L}{X7BBEA95E7AE1F317}{}
{
 \index{morphism of 2d\texttt{\symbol{45}}group} \index{crossed module morphism} This chapter describes morphisms of (pre\texttt{\symbol{45}})crossed modules
and (pre\texttt{\symbol{45}})cat$^1$\texttt{\symbol{45}}groups. 

\subsection{\textcolor{Chapter }{Source (for 2d-group mappings)}}
\logpage{[ 3, 1, 1 ]}\nobreak
\hyperdef{L}{X7FFD094F7FFB1F17}{}
{\noindent\textcolor{FuncColor}{$\triangleright$\enspace\texttt{Source({\mdseries\slshape map})\index{Source@\texttt{Source}!for 2d-group mappings}
\label{Source:for 2d-group mappings}
}\hfill{\scriptsize (attribute)}}\\
\noindent\textcolor{FuncColor}{$\triangleright$\enspace\texttt{Range({\mdseries\slshape map})\index{Range@\texttt{Range}!for 2d-group mappings}
\label{Range:for 2d-group mappings}
}\hfill{\scriptsize (attribute)}}\\
\noindent\textcolor{FuncColor}{$\triangleright$\enspace\texttt{SourceHom({\mdseries\slshape map})\index{SourceHom@\texttt{SourceHom}}
\label{SourceHom}
}\hfill{\scriptsize (attribute)}}\\
\noindent\textcolor{FuncColor}{$\triangleright$\enspace\texttt{RangeHom({\mdseries\slshape map})\index{RangeHom@\texttt{RangeHom}}
\label{RangeHom}
}\hfill{\scriptsize (attribute)}}\\


 Morphisms of \emph{2\texttt{\symbol{45}}dimensional groups} are implemented as \emph{2\texttt{\symbol{45}}dimensional mappings}. These have a pair of 2\texttt{\symbol{45}}dimensional groups as source and
range, together with two group homomorphisms mapping between corresponding
source and range groups. These functions return \texttt{fail} when invalid data is supplied. }

 }

 
\section{\textcolor{Chapter }{Morphisms of pre\texttt{\symbol{45}}crossed modules}}\logpage{[ 3, 2, 0 ]}
\hyperdef{L}{X78CADE4D7EB1EA44}{}
{
 \index{morphism} 

\subsection{\textcolor{Chapter }{IsXModMorphism}}
\logpage{[ 3, 2, 1 ]}\nobreak
\hyperdef{L}{X82B912B18127A42A}{}
{\noindent\textcolor{FuncColor}{$\triangleright$\enspace\texttt{IsXModMorphism({\mdseries\slshape map})\index{IsXModMorphism@\texttt{IsXModMorphism}}
\label{IsXModMorphism}
}\hfill{\scriptsize (property)}}\\
\noindent\textcolor{FuncColor}{$\triangleright$\enspace\texttt{IsPreXModMorphism({\mdseries\slshape map})\index{IsPreXModMorphism@\texttt{IsPreXModMorphism}}
\label{IsPreXModMorphism}
}\hfill{\scriptsize (property)}}\\


 A morphism between two pre\texttt{\symbol{45}}crossed modules $\calX_{1} = (\partial_1 : S_1 \to R_1)$ and $\calX_{2} = (\partial_2 : S_2 \to R_2)$ is a pair $(\sigma, \rho)$, where $\sigma : S_1 \to S_2$ and $\rho : R_1 \to R_2$ commute with the two boundary maps and are morphisms for the two actions: 
\[ \partial_2 \circ \sigma ~=~ \rho \circ \partial_1, \qquad \sigma(s^r) ~=~
(\sigma s)^{\rho r}. \]
 Here $\sigma$ is the \texttt{SourceHom} (\ref{SourceHom}) and $\rho$ is the \texttt{RangeHom} (\ref{RangeHom}) of the morphism. When $\calX_{1} = \calX_{2}$ and $\sigma, \rho$ are automorphisms then $(\sigma, \rho)$ is an automorphism of $\calX_1$. The group of automorphisms is denoted by ${\rm Aut}(\calX_1 )$. }

 

\subsection{\textcolor{Chapter }{IsInjective (for pre-xmod morphisms)}}
\logpage{[ 3, 2, 2 ]}\nobreak
\hyperdef{L}{X7E078F497F4EFA9F}{}
{\noindent\textcolor{FuncColor}{$\triangleright$\enspace\texttt{IsInjective({\mdseries\slshape map})\index{IsInjective@\texttt{IsInjective}!for pre-xmod morphisms}
\label{IsInjective:for pre-xmod morphisms}
}\hfill{\scriptsize (method)}}\\
\noindent\textcolor{FuncColor}{$\triangleright$\enspace\texttt{IsSurjective({\mdseries\slshape map})\index{IsSurjective@\texttt{IsSurjective}!for pre-xmod morphisms}
\label{IsSurjective:for pre-xmod morphisms}
}\hfill{\scriptsize (method)}}\\
\noindent\textcolor{FuncColor}{$\triangleright$\enspace\texttt{IsSingleValued({\mdseries\slshape map})\index{IsSingleValued@\texttt{IsSingleValued}!for pre-xmod morphisms}
\label{IsSingleValued:for pre-xmod morphisms}
}\hfill{\scriptsize (method)}}\\
\noindent\textcolor{FuncColor}{$\triangleright$\enspace\texttt{IsTotal({\mdseries\slshape map})\index{IsTotal@\texttt{IsTotal}!for pre-xmod morphisms}
\label{IsTotal:for pre-xmod morphisms}
}\hfill{\scriptsize (method)}}\\
\noindent\textcolor{FuncColor}{$\triangleright$\enspace\texttt{IsBijective({\mdseries\slshape map})\index{IsBijective@\texttt{IsBijective}!for pre-xmod morphisms}
\label{IsBijective:for pre-xmod morphisms}
}\hfill{\scriptsize (method)}}\\
\noindent\textcolor{FuncColor}{$\triangleright$\enspace\texttt{IsEndo2DimensionalMapping({\mdseries\slshape map})\index{IsEndo2DimensionalMapping@\texttt{IsEndo2DimensionalMapping}}
\label{IsEndo2DimensionalMapping}
}\hfill{\scriptsize (property)}}\\


 The usual properties of mappings are easily checked. It is usually sufficient
to verify that both the \texttt{SourceHom} (\ref{SourceHom}) and the \texttt{RangeHom} (\ref{RangeHom}) have the required property. }

 

\subsection{\textcolor{Chapter }{XModMorphism}}
\logpage{[ 3, 2, 3 ]}\nobreak
\hyperdef{L}{X7CEABD6487CF2A38}{}
{\noindent\textcolor{FuncColor}{$\triangleright$\enspace\texttt{XModMorphism({\mdseries\slshape args})\index{XModMorphism@\texttt{XModMorphism}}
\label{XModMorphism}
}\hfill{\scriptsize (function)}}\\
\noindent\textcolor{FuncColor}{$\triangleright$\enspace\texttt{XModMorphismByGroupHomomorphisms({\mdseries\slshape X1, X2, sigma, rho})\index{XModMorphismByGroupHomomorphisms@\texttt{XModMorphismByGroupHomomorphisms}}
\label{XModMorphismByGroupHomomorphisms}
}\hfill{\scriptsize (operation)}}\\
\noindent\textcolor{FuncColor}{$\triangleright$\enspace\texttt{PreXModMorphism({\mdseries\slshape args})\index{PreXModMorphism@\texttt{PreXModMorphism}}
\label{PreXModMorphism}
}\hfill{\scriptsize (function)}}\\
\noindent\textcolor{FuncColor}{$\triangleright$\enspace\texttt{PreXModMorphismByGroupHomomorphisms({\mdseries\slshape P1, P2, sigma, rho})\index{PreXModMorphismByGroupHomomorphisms@\texttt{PreXModMorphismByGroupHomomorphisms}}
\label{PreXModMorphismByGroupHomomorphisms}
}\hfill{\scriptsize (operation)}}\\
\noindent\textcolor{FuncColor}{$\triangleright$\enspace\texttt{InclusionMorphism2DimensionalDomains({\mdseries\slshape X1, S1})\index{InclusionMorphism2DimensionalDomains@\texttt{Inclusion}\-\texttt{Morphism2}\-\texttt{Dimensional}\-\texttt{Domains}!for crossed modules}
\label{InclusionMorphism2DimensionalDomains:for crossed modules}
}\hfill{\scriptsize (operation)}}\\
\noindent\textcolor{FuncColor}{$\triangleright$\enspace\texttt{IdentityMapping({\mdseries\slshape X1})\index{IdentityMapping@\texttt{IdentityMapping}!for pre-xmods}
\label{IdentityMapping:for pre-xmods}
}\hfill{\scriptsize (attribute)}}\\


 These are the constructors for morphisms of pre\texttt{\symbol{45}}crossed and
crossed modules. 

 In the following example we construct a simple automorphism of the crossed
module \texttt{X5} constructed in the previous chapter. }

 

 \index{display a 2d\texttt{\symbol{45}}mapping} \index{order of a 2d\texttt{\symbol{45}}automorphism} 
\begin{Verbatim}[commandchars=!@|,fontsize=\small,frame=single,label=Example]
  
  !gapprompt@gap>| !gapinput@sigma5 := GroupHomomorphismByImages( c5, c5, [ (5,6,7,8,9) ]|
          [ (5,9,8,7,6) ] );;
  !gapprompt@gap>| !gapinput@rho5 := IdentityMapping( Range( X1 ) );|
  IdentityMapping( PAut(c5) )
  !gapprompt@gap>| !gapinput@mor5 := XModMorphism( X5, X5, sigma5, rho5 );|
  [[c5->Aut(c5))] => [c5->Aut(c5))]] 
  !gapprompt@gap>| !gapinput@Display( mor5 );|
  Morphism of crossed modules :- 
  : Source = [c5->Aut(c5)] with generating sets:
    [ (5,6,7,8,9) ]
    [ GroupHomomorphismByImages( c5, c5, [ (5,6,7,8,9) ], [ (5,7,9,6,8) ] ) ]
  : Range = Source
  : Source Homomorphism maps source generators to:
    [ (5,9,8,7,6) ]
  : Range Homomorphism maps range generators to:
    [ GroupHomomorphismByImages( c5, c5, [ (5,6,7,8,9) ], [ (5,7,9,6,8) ] ) ]
  !gapprompt@gap>| !gapinput@IsAutomorphism2DimensionalDomain( mor5 );|
  true 
  !gapprompt@gap>| !gapinput@Order( mor5 );|
  2
  !gapprompt@gap>| !gapinput@RepresentationsOfObject( mor5 );|
  [ "IsComponentObjectRep", "IsAttributeStoringRep", "Is2DimensionalMappingRep" ]
  !gapprompt@gap>| !gapinput@KnownPropertiesOfObject( mor5 );|
  [ "CanEasilyCompareElements", "CanEasilySortElements", "IsTotal", 
    "IsSingleValued", "IsInjective", "IsSurjective", "RespectsMultiplication", 
    "IsPreXModMorphism", "IsXModMorphism", "IsEndomorphism2DimensionalDomain", 
    "IsAutomorphism2DimensionalDomain" ]
  !gapprompt@gap>| !gapinput@KnownAttributesOfObject( mor5 );|
  [ "Name", "Order", "Range", "Source", "SourceHom", "RangeHom" ]
  
\end{Verbatim}
 

\subsection{\textcolor{Chapter }{InnerAutomorphismXMod}}
\logpage{[ 3, 2, 4 ]}\nobreak
\hyperdef{L}{X7B0381EC85628CF6}{}
{\noindent\textcolor{FuncColor}{$\triangleright$\enspace\texttt{InnerAutomorphismXMod({\mdseries\slshape X1, r})\index{InnerAutomorphismXMod@\texttt{InnerAutomorphismXMod}}
\label{InnerAutomorphismXMod}
}\hfill{\scriptsize (operation)}}\\


 For each $r \in R$ the \emph{inner automorphism} of $\calX = (\partial : S \to R)$ is the automorphism $\alpha_r = (\sigma_r,\rho_r)$ of $\calX$ where $\sigma_r s = s^r$ and $\rho_r r' = r^{-1}r'r$. 

 Since $\sigma_r \ast \sigma_{r'} = \sigma_{rr'}$ and $\rho_r \ast \rho_{r'} = \rho_{rr'}$ it follows that $\alpha : R \to \Aut(\calX), r \mapsto \alpha_r$ is a group homomorphism. We shall see in Chapter 6 (see \texttt{NorrieXMod} (\ref{NorrieXMod})) that $\alpha$ is the boundary of the Norrie crossed module $\calN(\calX)$. }

 

 
\begin{Verbatim}[commandchars=!@|,fontsize=\small,frame=single,label=Example]
  
  !gapprompt@gap>| !gapinput@a := GeneratorsOfGroup( Range( X5 ) )[1];|
  [ (5,6,7,8,9) ] -> [ (5,7,9,6,8) ]
  !gapprompt@gap>| !gapinput@inna := InnerAutomorphismXMod( X5, a );;|
  !gapprompt@gap>| !gapinput@Display( inna );|
  Morphism of crossed modules :- 
  : Source = [c5->Aut(c5)] with generating sets:
    [ (5,6,7,8,9) ]
    [ GroupHomomorphismByImages( c5, c5, [ (5,6,7,8,9) ], [ (5,7,9,6,8) ] ) ]
  : Range = Source
  : Source Homomorphism maps source generators to:
    [ (5,7,9,6,8) ]
  : Range Homomorphism maps range generators to:
    [ GroupHomomorphismByImages( c5, c5, [ (5,6,7,8,9) ], [ (5,7,9,6,8) ] ) ]
  
\end{Verbatim}
 

\subsection{\textcolor{Chapter }{IsomorphismPerm2DimensionalGroup (for pre-xmod morphisms)}}
\logpage{[ 3, 2, 5 ]}\nobreak
\hyperdef{L}{X854FC0C781AD62EC}{}
{\noindent\textcolor{FuncColor}{$\triangleright$\enspace\texttt{IsomorphismPerm2DimensionalGroup({\mdseries\slshape obj})\index{IsomorphismPerm2DimensionalGroup@\texttt{IsomorphismPerm2DimensionalGroup}!for pre-xmod morphisms}
\label{IsomorphismPerm2DimensionalGroup:for pre-xmod morphisms}
}\hfill{\scriptsize (attribute)}}\\
\noindent\textcolor{FuncColor}{$\triangleright$\enspace\texttt{IsomorphismPc2DimensionalGroup({\mdseries\slshape obj})\index{IsomorphismPc2DimensionalGroup@\texttt{IsomorphismPc2DimensionalGroup}!for pre-xmod morphisms}
\label{IsomorphismPc2DimensionalGroup:for pre-xmod morphisms}
}\hfill{\scriptsize (attribute)}}\\
\noindent\textcolor{FuncColor}{$\triangleright$\enspace\texttt{IsomorphismByIsomorphisms({\mdseries\slshape D, list})\index{IsomorphismByIsomorphisms@\texttt{IsomorphismByIsomorphisms}}
\label{IsomorphismByIsomorphisms}
}\hfill{\scriptsize (operation)}}\\


 When $\calD$ is a $2$\texttt{\symbol{45}}dimensional domain with source $S$ and range $R$ and $\sigma : S \to S',~ \rho : R \to R'$ are isomorphisms, then \texttt{IsomorphismByIsomorphisms(D,[sigma,rho])} returns an isomorphism $(\sigma,\rho) : \calD \to \calD'$ where $\calD'$ has source $S'$ and range $R'$. Be sure to test \texttt{IsBijective} for the two functions $\sigma,\rho$ before applying this operation. 

 Using \texttt{IsomorphismByIsomorphisms} with a pair of isomorphisms obtained using \texttt{IsomorphismPermGroup} or \texttt{IsomorphismPcGroup}, we may construct a crossed module or a cat$^1$\texttt{\symbol{45}}group of permutation groups or
pc\texttt{\symbol{45}}groups. }

 

 
\begin{Verbatim}[commandchars=!@|,fontsize=\small,frame=single,label=Example]
  
  !gapprompt@gap>| !gapinput@q8 := SmallGroup(8,4);;   ## quaternion group |
  !gapprompt@gap>| !gapinput@XAq8 := XModByAutomorphismGroup( q8 );|
  [Group( [ f1, f2, f3 ] )->Group( [ Pcgs([ f1, f2, f3 ]) -> [ f1*f2, f2, f3 ], 
    Pcgs([ f1, f2, f3 ]) -> [ f2, f1*f2, f3 ], 
    Pcgs([ f1, f2, f3 ]) -> [ f1*f3, f2, f3 ], 
    Pcgs([ f1, f2, f3 ]) -> [ f1, f2*f3, f3 ] ] )]
  !gapprompt@gap>| !gapinput@iso := IsomorphismPerm2DimensionalGroup( XAq8 );;|
  !gapprompt@gap>| !gapinput@YAq8 := Image( iso );|
  [Group( [ (1,2,4,6)(3,8,7,5), (1,3,4,7)(2,5,6,8), (1,4)(2,6)(3,7)(5,8) 
   ] )->Group( [ (1,3,4,6), (1,2,3)(4,5,6), (1,4)(3,6), (2,5)(3,6) ] )]
  !gapprompt@gap>| !gapinput@s4 := SymmetricGroup(4);; |
  !gapprompt@gap>| !gapinput@isos4 := IsomorphismGroups( Range(YAq8), s4 );;|
  !gapprompt@gap>| !gapinput@id := IdentityMapping( Source( YAq8 ) );; |
  !gapprompt@gap>| !gapinput@IsBijective( id );;  IsBijective( isos4 );;|
  !gapprompt@gap>| !gapinput@mor := IsomorphismByIsomorphisms( YAq8, [id,isos4] );;|
  !gapprompt@gap>| !gapinput@ZAq8 := Image( mor );|
  [Group( [ (1,2,4,6)(3,8,7,5), (1,3,4,7)(2,5,6,8), (1,4)(2,6)(3,7)(5,8) 
   ] )->SymmetricGroup( [ 1 .. 4 ] )]
  
\end{Verbatim}
 

\subsection{\textcolor{Chapter }{MorphismOfPullback (for a crossed module by pullback)}}
\logpage{[ 3, 2, 6 ]}\nobreak
\hyperdef{L}{X87BCAAF787A7FF69}{}
{\noindent\textcolor{FuncColor}{$\triangleright$\enspace\texttt{MorphismOfPullback({\mdseries\slshape xmod})\index{MorphismOfPullback@\texttt{MorphismOfPullback}!for a crossed module by pullback}
\label{MorphismOfPullback:for a crossed module by pullback}
}\hfill{\scriptsize (attribute)}}\\


 Let $\calX_1 = (\lambda : L \to N)$ be the pullback crossed module obtained from a crossed module $\calX_0 = (\mu : M \to P)$ and a group homomorphism $\nu : N \to P$. Then the associated crossed module morphism is $(\kappa,\nu) : \calX_1 \to \calX_0$ where $\kappa$ is the projection from $L$ to $M$. }

 }

 
\section{\textcolor{Chapter }{Morphisms of pre\texttt{\symbol{45}}cat$^1$\texttt{\symbol{45}}groups}}\label{sect-mor-pre-cat1}
\logpage{[ 3, 3, 0 ]}
\hyperdef{L}{X7B9D3C1F7A395FF2}{}
{
  A morphism of pre\texttt{\symbol{45}}cat$^1$\texttt{\symbol{45}}groups from $\calC_1 = (e_1;t_1,h_1 : G_1 \to R_1)$ to $\calC_2 = (e_2;t_2,h_2 : G_2 \to R_2)$ is a pair $(\gamma, \rho)$ where $\gamma : G_1 \to G_2$ and $\rho : R_1 \to R_2$ are homomorphisms satisfying 
\[ h_2 \circ \gamma ~=~ \rho \circ h_1, \qquad t_2 \circ \gamma ~=~ \rho \circ
t_1, \qquad e_2 \circ \rho ~=~ \gamma \circ e_1. \]
 

\subsection{\textcolor{Chapter }{IsCat1GroupMorphism}}
\logpage{[ 3, 3, 1 ]}\nobreak
\hyperdef{L}{X7C47D0EC782D4C40}{}
{\noindent\textcolor{FuncColor}{$\triangleright$\enspace\texttt{IsCat1GroupMorphism({\mdseries\slshape map})\index{IsCat1GroupMorphism@\texttt{IsCat1GroupMorphism}}
\label{IsCat1GroupMorphism}
}\hfill{\scriptsize (property)}}\\
\noindent\textcolor{FuncColor}{$\triangleright$\enspace\texttt{IsPreCat1GroupMorphism({\mdseries\slshape map})\index{IsPreCat1GroupMorphism@\texttt{IsPreCat1GroupMorphism}}
\label{IsPreCat1GroupMorphism}
}\hfill{\scriptsize (property)}}\\
\noindent\textcolor{FuncColor}{$\triangleright$\enspace\texttt{Cat1GroupMorphism({\mdseries\slshape args})\index{Cat1GroupMorphism@\texttt{Cat1GroupMorphism}}
\label{Cat1GroupMorphism}
}\hfill{\scriptsize (function)}}\\
\noindent\textcolor{FuncColor}{$\triangleright$\enspace\texttt{Cat1GroupMorphismByGroupHomomorphisms({\mdseries\slshape C1, C2, gamma, rho})\index{Cat1GroupMorphismByGroupHomomorphisms@\texttt{Cat1}\-\texttt{Group}\-\texttt{Morphism}\-\texttt{By}\-\texttt{Group}\-\texttt{Homomorphisms}}
\label{Cat1GroupMorphismByGroupHomomorphisms}
}\hfill{\scriptsize (operation)}}\\
\noindent\textcolor{FuncColor}{$\triangleright$\enspace\texttt{PreCat1GroupMorphism({\mdseries\slshape args})\index{PreCat1GroupMorphism@\texttt{PreCat1GroupMorphism}}
\label{PreCat1GroupMorphism}
}\hfill{\scriptsize (function)}}\\
\noindent\textcolor{FuncColor}{$\triangleright$\enspace\texttt{PreCat1GroupMorphismByGroupHomomorphisms({\mdseries\slshape P1, P2, gamma, rho})\index{PreCat1GroupMorphismByGroupHomomorphisms@\texttt{Pre}\-\texttt{Cat1}\-\texttt{Group}\-\texttt{Morphism}\-\texttt{By}\-\texttt{Group}\-\texttt{Homomorphisms}}
\label{PreCat1GroupMorphismByGroupHomomorphisms}
}\hfill{\scriptsize (operation)}}\\
\noindent\textcolor{FuncColor}{$\triangleright$\enspace\texttt{InclusionMorphism2DimensionalDomains({\mdseries\slshape C1, S1})\index{InclusionMorphism2DimensionalDomains@\texttt{Inclusion}\-\texttt{Morphism2}\-\texttt{Dimensional}\-\texttt{Domains}!for cat1-groups}
\label{InclusionMorphism2DimensionalDomains:for cat1-groups}
}\hfill{\scriptsize (operation)}}\\
\noindent\textcolor{FuncColor}{$\triangleright$\enspace\texttt{InnerAutomorphismCat1({\mdseries\slshape C1, r})\index{InnerAutomorphismCat1@\texttt{InnerAutomorphismCat1}}
\label{InnerAutomorphismCat1}
}\hfill{\scriptsize (operation)}}\\
\noindent\textcolor{FuncColor}{$\triangleright$\enspace\texttt{IdentityMapping({\mdseries\slshape C1})\index{IdentityMapping@\texttt{IdentityMapping}!for precat1-morphisms}
\label{IdentityMapping:for precat1-morphisms}
}\hfill{\scriptsize (attribute)}}\\


 For an example we form a second cat$^1$\texttt{\symbol{45}}group \texttt{C2=[g18={\textgreater}s3a]}, similar to \texttt{C1} in \ref{mansect-cat1}, then construct an isomorphism $(\gamma,\rho)$ between them. }

 

 
\begin{Verbatim}[commandchars=!@|,fontsize=\small,frame=single,label=Example]
  
  !gapprompt@gap>| !gapinput@t3 := GroupHomomorphismByImages(g18,s3a,g18gens,[(),(7,8,9),(8,9)]);;     |
  !gapprompt@gap>| !gapinput@e3 := GroupHomomorphismByImages(s3a,g18,s3agens,[(4,5,6),(2,3)(5,6)]);;   |
  !gapprompt@gap>| !gapinput@C3 := Cat1Group( t3, h1, e3 );; |
  !gapprompt@gap>| !gapinput@imgamma := [ (4,5,6), (1,2,3), (2,3)(5,6) ];; |
  !gapprompt@gap>| !gapinput@gamma := GroupHomomorphismByImages( g18, g18, g18gens, imgamma );;|
  !gapprompt@gap>| !gapinput@rho := IdentityMapping( s3a );; |
  !gapprompt@gap>| !gapinput@phi3 := Cat1GroupMorphism( C18, C3, gamma, rho );;|
  !gapprompt@gap>| !gapinput@Display( phi3 );;|
  Morphism of cat1-groups :- 
  : Source = [g18=>s3a] with generating sets:
    [ (1,2,3), (4,5,6), (2,3)(5,6) ]
    [ (7,8,9), (8,9) ]
  :  Range = [g18=>s3a] with generating sets:
    [ (1,2,3), (4,5,6), (2,3)(5,6) ]
    [ (7,8,9), (8,9) ]
  : Source Homomorphism maps source generators to:
    [ (4,5,6), (1,2,3), (2,3)(5,6) ]
  : Range Homomorphism maps range generators to:
    [ (7,8,9), (8,9) ]
  
\end{Verbatim}
 

\subsection{\textcolor{Chapter }{Cat1GroupMorphismOfXModMorphism}}
\logpage{[ 3, 3, 2 ]}\nobreak
\hyperdef{L}{X7D7459E67B568B44}{}
{\noindent\textcolor{FuncColor}{$\triangleright$\enspace\texttt{Cat1GroupMorphismOfXModMorphism({\mdseries\slshape IsXModMorphism})\index{Cat1GroupMorphismOfXModMorphism@\texttt{Cat1GroupMorphismOfXModMorphism}}
\label{Cat1GroupMorphismOfXModMorphism}
}\hfill{\scriptsize (attribute)}}\\
\noindent\textcolor{FuncColor}{$\triangleright$\enspace\texttt{XModMorphismOfCat1GroupMorphism({\mdseries\slshape IsCat1GroupMorphism})\index{XModMorphismOfCat1GroupMorphism@\texttt{XModMorphismOfCat1GroupMorphism}}
\label{XModMorphismOfCat1GroupMorphism}
}\hfill{\scriptsize (attribute)}}\\


 If $(\sigma,\rho) : \calX_1 \to \calX_2$ and $\calC_1,\calC_2$ are the cat$^1$\texttt{\symbol{45}}groups accociated to $\calX_1, \calX_2$, then the associated morphism of cat$^1$\texttt{\symbol{45}}groups is $(\gamma,\rho)$ where $\gamma(r_1,s_1) = (\rho r_1, \sigma s_1)$. 

 Similarly, given a morphism $(\gamma,\rho) : \calC_1 \to \calC_2$ of cat$^1$\texttt{\symbol{45}}groups, the associated morphism of crossed modules is $(\sigma,\rho) : \calX_1 \to \calX_2$ where $\sigma s_1 = \gamma(1,s_1)$. . }

 

 
\begin{Verbatim}[commandchars=!@B,fontsize=\small,frame=single,label=Example]
  
  !gapprompt@gap>B !gapinput@phi5 := Cat1GroupMorphismOfXModMorphism( mor5 );B
  [[(Aut(c5) |X c5)=>Aut(c5)] => [(Aut(c5) |X c5)=>Aut(c5)]]
  !gapprompt@gap>B !gapinput@mor3 := XModMorphismOfCat1GroupMorphism( phi3 );;B
  !gapprompt@gap>B !gapinput@Display( mor3 );B
  Morphism of crossed modules :- 
  : Source = xmod([g18=>s3a]) with generating sets:
    [ (4,5,6) ]
    [ (7,8,9), (8,9) ]
  :  Range = xmod([g18=>s3a]) with generating sets:
    [ (1,2,3) ]
    [ (7,8,9), (8,9) ]
  : Source Homomorphism maps source generators to:
    [ (1,2,3) ]
  : Range Homomorphism maps range generators to:
    [ (7,8,9), (8,9) ]
  
\end{Verbatim}
 

\subsection{\textcolor{Chapter }{IsomorphismPermObject}}
\logpage{[ 3, 3, 3 ]}\nobreak
\hyperdef{L}{X7C6AF7C285D546B2}{}
{\noindent\textcolor{FuncColor}{$\triangleright$\enspace\texttt{IsomorphismPermObject({\mdseries\slshape obj})\index{IsomorphismPermObject@\texttt{IsomorphismPermObject}}
\label{IsomorphismPermObject}
}\hfill{\scriptsize (function)}}\\
\noindent\textcolor{FuncColor}{$\triangleright$\enspace\texttt{IsomorphismPerm2DimensionalGroup({\mdseries\slshape 2DimensionalGroup})\index{IsomorphismPerm2DimensionalGroup@\texttt{IsomorphismPerm2DimensionalGroup}!for pre-cat1 morphisms}
\label{IsomorphismPerm2DimensionalGroup:for pre-cat1 morphisms}
}\hfill{\scriptsize (attribute)}}\\
\noindent\textcolor{FuncColor}{$\triangleright$\enspace\texttt{IsomorphismFp2DimensionalGroup({\mdseries\slshape 2DimensionalGroup})\index{IsomorphismFp2DimensionalGroup@\texttt{IsomorphismFp2DimensionalGroup}!for pre-cat1 morphisms}
\label{IsomorphismFp2DimensionalGroup:for pre-cat1 morphisms}
}\hfill{\scriptsize (attribute)}}\\
\noindent\textcolor{FuncColor}{$\triangleright$\enspace\texttt{IsomorphismPc2DimensionalGroup({\mdseries\slshape 2DimensionalGroup})\index{IsomorphismPc2DimensionalGroup@\texttt{IsomorphismPc2DimensionalGroup}!for pre-cat1 morphisms}
\label{IsomorphismPc2DimensionalGroup:for pre-cat1 morphisms}
}\hfill{\scriptsize (attribute)}}\\
\noindent\textcolor{FuncColor}{$\triangleright$\enspace\texttt{RegularActionHomomorphism2DimensionalGroup({\mdseries\slshape 2DimensionalGroup})\index{RegularActionHomomorphism2DimensionalGroup@\texttt{Regular}\-\texttt{Action}\-\texttt{Homomorphism2}\-\texttt{Dimensional}\-\texttt{Group}!for pre-cat1 morphisms}
\label{RegularActionHomomorphism2DimensionalGroup:for pre-cat1 morphisms}
}\hfill{\scriptsize (attribute)}}\\


 The global function \texttt{IsomorphismPermObject} calls \texttt{IsomorphismPerm2DimensionalGroup}, which constructs a morphism whose \texttt{SourceHom} (\ref{SourceHom}) and \texttt{RangeHom} (\ref{RangeHom}) are calculated using \texttt{IsomorphismPermGroup} on the source and range. 

 The global function \texttt{RegularActionHomomorphism2DimensionalGroup} is similar, but uses \texttt{RegularActionHomomorphism} in place of \texttt{IsomorphismPermGroup}. }

 

 
\begin{Verbatim}[commandchars=!@|,fontsize=\small,frame=single,label=Example]
  
  !gapprompt@gap>| !gapinput@iso8 := IsomorphismPerm2DimensionalGroup( C8 );|
  [[G8=>d12] => [..]]
  
\end{Verbatim}
 

\subsection{\textcolor{Chapter }{InnerAutomorphismCat1Group}}
\logpage{[ 3, 3, 4 ]}\nobreak
\hyperdef{L}{X7C31C47C80D95D5E}{}
{\noindent\textcolor{FuncColor}{$\triangleright$\enspace\texttt{InnerAutomorphismCat1Group({\mdseries\slshape C1, r})\index{InnerAutomorphismCat1Group@\texttt{InnerAutomorphismCat1Group}}
\label{InnerAutomorphismCat1Group}
}\hfill{\scriptsize (operation)}}\\


 For $r \in R$, the \emph{inner automorphism} of $\calC = (e;t,h : G \to R )$ is the automorphism $\alpha_r = (\gamma_r,\rho_r)$ of $\calC$ where $\gamma_r s = s^{er}$ and $\rho_r r' = r^{-1}r'r$. }

 

 
\begin{Verbatim}[commandchars=!@|,fontsize=\small,frame=single,label=Example]
  
  !gapprompt@gap>| !gapinput@mor4 := InnerAutomorphismCat1Group( C3, (7,9) );;  |
  !gapprompt@gap>| !gapinput@Display( mor4 );                                 |
  Morphism of cat1-groups :- 
  : Source = [g18 => s3a] with generating sets:
    [ (1,2,3), (4,5,6), (2,3)(5,6) ]
    [ (7,8,9), (8,9) ]
  : Range = Source
  : Source Homomorphism maps source generators to:
    [ (1,3,2), (4,6,5), (2,3)(4,5) ]
  : Range Homomorphism maps range generators to:
    [ (7,9,8), (7,8) ]
  
\end{Verbatim}
 

\subsection{\textcolor{Chapter }{SmallerDegreePermutationRepresentation2DimensionalGroup (for perm 2d-groups)}}
\logpage{[ 3, 3, 5 ]}\nobreak
\hyperdef{L}{X837B0299846C2391}{}
{\noindent\textcolor{FuncColor}{$\triangleright$\enspace\texttt{SmallerDegreePermutationRepresentation2DimensionalGroup({\mdseries\slshape Perm2DimensionalGroup})\index{SmallerDegreePermutationRepresentation2DimensionalGroup@\texttt{Smaller}\-\texttt{Degree}\-\texttt{Permutation}\-\texttt{Representation2}\-\texttt{Dimensional}\-\texttt{Group}!for perm 2d-groups}
\label{SmallerDegreePermutationRepresentation2DimensionalGroup:for perm 2d-groups}
}\hfill{\scriptsize (attribute)}}\\


 The attribute \texttt{SmallerDegreePermutationRepresentation2DimensionalGroup} is obtained by calling \texttt{SmallerDegreePermutationRepresentation} on the source and range to obtain the an isomorphism for the
pre\texttt{\symbol{45}}xmod or pre\texttt{\symbol{45}}cat$^1$\texttt{\symbol{45}}group. }

 

 
\begin{Verbatim}[commandchars=!@|,fontsize=\small,frame=single,label=Example]
  
  !gapprompt@gap>| !gapinput@G := Group( (1,2,3,4)(5,6,7,8) );; |
  !gapprompt@gap>| !gapinput@H := Subgroup( G, [ (1,3)(2,4)(5,7)(6,8) ] );;|
  !gapprompt@gap>| !gapinput@XG := XModByNormalSubgroup( G, H );|
  [Group( [ (1,3)(2,4)(5,7)(6,8) ] )->Group( [ (1,2,3,4)(5,6,7,8) ] )]
  !gapprompt@gap>| !gapinput@sdpr := SmallerDegreePermutationRepresentation2DimensionalGroup( XG );; |
  !gapprompt@gap>| !gapinput@Range( sdpr );|
  [Group( [ (1,2) ] )->Group( [ (1,2,3,4) ] )]
  
\end{Verbatim}
 }

 
\section{\textcolor{Chapter }{Operations on morphisms}}\label{sect-oper-mor}
\logpage{[ 3, 4, 0 ]}
\hyperdef{L}{X7B09A28579707CAF}{}
{
  \index{operations on morphisms} 

\subsection{\textcolor{Chapter }{CompositionMorphism}}
\logpage{[ 3, 4, 1 ]}\nobreak
\hyperdef{L}{X811F886081AAB95F}{}
{\noindent\textcolor{FuncColor}{$\triangleright$\enspace\texttt{CompositionMorphism({\mdseries\slshape map2, map1})\index{CompositionMorphism@\texttt{CompositionMorphism}}
\label{CompositionMorphism}
}\hfill{\scriptsize (operation)}}\\


 Composition of morphisms (written \texttt{({\textless}map1{\textgreater} * {\textless}map2{\textgreater})} when maps act on the right) calls the \texttt{CompositionMorphism} function for maps (acting on the left), applied to the appropriate type of
2d\texttt{\symbol{45}}mapping. }

 

 
\begin{Verbatim}[commandchars=!@|,fontsize=\small,frame=single,label=Example]
  
  !gapprompt@gap>| !gapinput@H8 := Subgroup(G8,[G8.3,G8.4,G8.6,G8.7]);  SetName( H8, "H8" );|
  Group([ f3, f4, f6, f7 ])
  !gapprompt@gap>| !gapinput@c6 := Subgroup( d12, [b,c] );  SetName( c6, "c6" );|
  Group([ f2, f3 ])
  !gapprompt@gap>| !gapinput@SC8 := Sub2DimensionalGroup( C8, H8, c6 );|
  [H8=>c6]
  !gapprompt@gap>| !gapinput@IsCat1Group( SC8 );|
  true
  !gapprompt@gap>| !gapinput@inc8 := InclusionMorphism2DimensionalDomains( C8, SC8 );|
  [[H8=>c6] => [G8=>d12]]
  !gapprompt@gap>| !gapinput@CompositionMorphism( iso8, inc );                  |
  [[H8=>c6] => P[G8=>d12]]
  
\end{Verbatim}
 

\subsection{\textcolor{Chapter }{Kernel (for 2d-mappings)}}
\logpage{[ 3, 4, 2 ]}\nobreak
\hyperdef{L}{X83C3A2478159DE76}{}
{\noindent\textcolor{FuncColor}{$\triangleright$\enspace\texttt{Kernel({\mdseries\slshape map})\index{Kernel@\texttt{Kernel}!for 2d-mappings}
\label{Kernel:for 2d-mappings}
}\hfill{\scriptsize (operation)}}\\
\noindent\textcolor{FuncColor}{$\triangleright$\enspace\texttt{Kernel2DimensionalMapping({\mdseries\slshape map})\index{Kernel2DimensionalMapping@\texttt{Kernel2DimensionalMapping}}
\label{Kernel2DimensionalMapping}
}\hfill{\scriptsize (attribute)}}\\


 The kernel of a morphism of crossed modules is a normal subcrossed module
whose groups are the kernels of the source and target homomorphisms. The
inclusion of the kernel is a standard example of a crossed square, but these
have not yet been implemented. }

 

 
\begin{Verbatim}[commandchars=!@|,fontsize=\small,frame=single,label=Example]
  
  !gapprompt@gap>| !gapinput@c2 := Group( (19,20) );                                    |
  Group([ (19,20) ])
  !gapprompt@gap>| !gapinput@X0 := XModByNormalSubgroup( c2, c2 );  SetName( X0, "X0" );|
  [Group( [ (19,20) ] )->Group( [ (19,20) ] )]
  !gapprompt@gap>| !gapinput@SX8 := Source( X8 );;|
  !gapprompt@gap>| !gapinput@genSX8 := GeneratorsOfGroup( SX8 ); |
  [ f1, f4, f5, f7 ]
  !gapprompt@gap>| !gapinput@sigma0 := GroupHomomorphismByImages(SX8,c2,genSX8,[(19,20),(),(),()]);|
  [ f1, f4, f5, f7 ] -> [ (19,20), (), (), () ]
  !gapprompt@gap>| !gapinput@rho0 := GroupHomomorphismByImages(d12,c2,[a1,a2,a3],[(19,20),(),()]);|
  [ f1, f2, f3 ] -> [ (19,20), (), () ]
  !gapprompt@gap>| !gapinput@mor0 := XModMorphism( X8, X0, sigma0, rho0 );;           |
  !gapprompt@gap>| !gapinput@K0 := Kernel( mor0 );;|
  !gapprompt@gap>| !gapinput@StructureDescription( K0 );|
  [ "C12", "C6" ]
\end{Verbatim}
 }

 
\section{\textcolor{Chapter }{Quasi\texttt{\symbol{45}}isomorphisms}}\label{sect-quasi-iso}
\logpage{[ 3, 5, 0 ]}
\hyperdef{L}{X79C47E3D7855A117}{}
{
  \index{quasi isomorphisms} A morphism of crossed modules $ \phi : \calX = (\partial : S \to R) \to \calX' = (\partial' : S' \to R') $ induces homomorphisms $\pi_1(\phi) : \pi_1(\partial) \to \pi_1(\partial')$ and $\pi_2(\phi) : \pi_2(\partial) \to \pi_2(\partial')$. A morphism $\phi$ is a \emph{quasi\texttt{\symbol{45}}isomorphism} if both $\pi_1(\phi)$ and $\pi_2(\phi)$ are isomorphisms. Two crossed modules $\calX,\calX'$ are \emph{quasi\texttt{\symbol{45}}isomorphic} is there exists a sequence of quasi\texttt{\symbol{45}}isomorphisms 
\[ \calX ~=~ \calX_1 ~\leftrightarrow~ \calX_2 ~\leftrightarrow~ \calX_3
~\leftrightarrow~ \cdots ~\longleftrightarrow~ \calX_{\ell} ~=~ \calX' \]
 of length $\ell-1$. Here $\calX_i \leftrightarrow \calX_j$ means that \emph{either} $\calX_i \to \calX_j$ \emph{or} $\calX_j \to \calX_i$. When $\calX,\calX'$ are quasi\texttt{\symbol{45}}isomorphic we write $\calX \simeq \calX'$. Clearly $\simeq$ is an equivalence relation. Mac\texttt{\symbol{92}} Lane and Whitehead in \cite{MW} showed that there is a one\texttt{\symbol{45}}to\texttt{\symbol{45}}one
correspondence between homotopy $2$\texttt{\symbol{45}}types and quasi\texttt{\symbol{45}}isomorphism classes. We
say that $\calX$ represents a \emph{trivial} quasi\texttt{\symbol{45}}isomorphism class if $\partial=0$. 

 Two cat$^1$\texttt{\symbol{45}}groups are quasi\texttt{\symbol{45}}isomorphic if their
corresponding crossed modules are. The procedure for constructing a
representative for the quasi\texttt{\symbol{45}}isomorphism class of a cat$^1$\texttt{\symbol{45}}group $\calC$, as described by Ellis and Le in \cite{EL}, is as follows. The \emph{quotient process} consists of finding all normal sub\texttt{\symbol{45}}crossed modules $\calN$ of the crossed module $\calX$ associated to $\calC$; constructing the quotient crossed module morphisms $\nu : \calX \to \calX/\calN$; and converting these $\nu$ to morphisms from $\calC$. 

 The \emph{sub\texttt{\symbol{45}}crossed module process} consists of finding all sub\texttt{\symbol{45}}crossed modules $\calS$ of $\calX$ such that the inclusion $\iota : \calS \to \calX$ is a quasi\texttt{\symbol{45}}isomorphism; then converting $\iota$ to a morphism to $\calC$. 

 The procedure for finding all quasi\texttt{\symbol{45}}isomorphism reductions
consists of repeating the quotient process, followed by the
sub\texttt{\symbol{45}}crossed module process, until no further reductions are
possible. 

 It may happen that $\calC_1 \simeq \calC_2$ without either having a quasi\texttt{\symbol{45}}isomorphism reduction. In
this case it is necessary to find a suitable $\calC_3$ with reductions $\calC_3 \to \calC_1$ and $\calC_3 \to \calC_2$. No such automated process is available in \textsf{XMod}. 

 Functions for these computations were first implemented in the package \textsf{HAP} and are available as \texttt{QuotientQuasiIsomorph}, \texttt{SubQuasiIsomorph} and \texttt{QuasiIsomorph}. 

\subsection{\textcolor{Chapter }{QuotientQuasiIsomorphism}}
\logpage{[ 3, 5, 1 ]}\nobreak
\hyperdef{L}{X86F08B1981618400}{}
{\noindent\textcolor{FuncColor}{$\triangleright$\enspace\texttt{QuotientQuasiIsomorphism({\mdseries\slshape cat1, bool})\index{QuotientQuasiIsomorphism@\texttt{QuotientQuasiIsomorphism}}
\label{QuotientQuasiIsomorphism}
}\hfill{\scriptsize (operation)}}\\


 This function implements the quotient process. The second parameter is a
boolean which, when true, causes the results of some intermediate calculations
to be printed. The output shows the identity of the reduced cat$^1$\texttt{\symbol{45}}group, if there is one. }

 

 
\begin{Verbatim}[commandchars=!@|,fontsize=\small,frame=single,label=Example]
  
  !gapprompt@gap>| !gapinput@C18a := Cat1Select( 18, 4, 4 );;          |
  !gapprompt@gap>| !gapinput@StructureDescription( C18a );             |
  [ "(C3 x C3) : C2", "S3" ]
  !gapprompt@gap>| !gapinput@QuotientQuasiIsomorphism( C18a, true );   |
  quo: [ f2 ][ f3 ], [ "1", "C2" ]
  [ [ 2, 1 ], [ 2, 1 ] ], [ 2, 1, 1 ]
  [ [ 2, 1, 1 ] ]
  
\end{Verbatim}
 

\subsection{\textcolor{Chapter }{SubQuasiIsomorphism}}
\logpage{[ 3, 5, 2 ]}\nobreak
\hyperdef{L}{X7C4C1C587B8932A7}{}
{\noindent\textcolor{FuncColor}{$\triangleright$\enspace\texttt{SubQuasiIsomorphism({\mdseries\slshape cat1, bool})\index{SubQuasiIsomorphism@\texttt{SubQuasiIsomorphism}}
\label{SubQuasiIsomorphism}
}\hfill{\scriptsize (operation)}}\\


 This function implements the sub\texttt{\symbol{45}}crossed module process. }

 

 
\begin{Verbatim}[commandchars=!@|,fontsize=\small,frame=single,label=Example]
  
  !gapprompt@gap>| !gapinput@SubQuasiIsomorphism( C18a, false );|
  [ [ 2, 1, 1 ], [ 2, 1, 1 ], [ 2, 1, 1 ] ]
  
\end{Verbatim}
 

\subsection{\textcolor{Chapter }{QuasiIsomorphism}}
\logpage{[ 3, 5, 3 ]}\nobreak
\hyperdef{L}{X83365AF2812E5C04}{}
{\noindent\textcolor{FuncColor}{$\triangleright$\enspace\texttt{QuasiIsomorphism({\mdseries\slshape cat1, list, bool})\index{QuasiIsomorphism@\texttt{QuasiIsomorphism}}
\label{QuasiIsomorphism}
}\hfill{\scriptsize (operation)}}\\


 This function implements the general process. }

 

 
\begin{Verbatim}[commandchars=!@|,fontsize=\small,frame=single,label=Example]
  
  !gapprompt@gap>| !gapinput@L18a := QuasiIsomorphism( C18a, [18,4,4], false );|
  [ [ 2, 1, 1 ], [ 18, 4, 4 ] ]
  
\end{Verbatim}
 The logs above show that \texttt{C18a} has just one normal sub\texttt{\symbol{45}}crossed module $\calN$ leading to a reduction, and that there are three
sub\texttt{\symbol{45}}crossed modules $\calS$ all giving the same reduction. The conclusion is that \texttt{C18a} is quasi\texttt{\symbol{45}}isomorphic to the identity cat$^1$\texttt{\symbol{45}}group on the cyclic group of order $2$. }

 }

         
\chapter{\textcolor{Chapter }{Isoclinism of groups and crossed modules}}\label{chap-isclnc}
\logpage{[ 4, 0, 0 ]}
\hyperdef{L}{X802AFE8E7EDB435E}{}
{
  This chapter describes some functions written by Alper Odaba{\c s} and Enver
Uslu, and reported in their paper \cite{IOU1}. Section \ref{sect-more-xmod-ops} contains some additional basic functions for crossed modules, constructing
quotients, centres, centralizers and normalizers. In Sections \ref{sect-isoclinic-groups} and \ref{sect-isoclinic-xmods} there are functions dealing specifically with isoclinism for groups and for
crossed modules. Since these functions represent a recent addition to the
package (as of November 2015), the function names are liable to change in
future versions. The notion of isoclinism has been crucial to the enumeration
of groups of prime power order, see for example James, Newman and O'Brien, \cite{JNO}. 
\section{\textcolor{Chapter }{More operations for crossed modules}}\label{sect-more-xmod-ops}
\logpage{[ 4, 1, 0 ]}
\hyperdef{L}{X7E373BF3836B3A9C}{}
{
  

\subsection{\textcolor{Chapter }{FactorPreXMod}}
\logpage{[ 4, 1, 1 ]}\nobreak
\hyperdef{L}{X873ED97185D9176E}{}
{\noindent\textcolor{FuncColor}{$\triangleright$\enspace\texttt{FactorPreXMod({\mdseries\slshape X1, X2})\index{FactorPreXMod@\texttt{FactorPreXMod}}
\label{FactorPreXMod}
}\hfill{\scriptsize (operation)}}\\
\noindent\textcolor{FuncColor}{$\triangleright$\enspace\texttt{NaturalMorphismByNormalSubPreXMod({\mdseries\slshape X1, X2})\index{NaturalMorphismByNormalSubPreXMod@\texttt{NaturalMorphismByNormalSubPreXMod}}
\label{NaturalMorphismByNormalSubPreXMod}
}\hfill{\scriptsize (operation)}}\\


 When $\calX_2 = (\partial_2 : S_2 \to R_2)$ is a normal sub\texttt{\symbol{45}}precrossed module of $\calX_1 = (\partial_1 : S_1 \to R_1)$, then the quotient precrossed module is $(\partial : S_2/S_1 \to R_2/R_1)$ with the induced boundary and action maps. Quotienting a precrossed module by
it's Peiffer subgroup is a special case of this construction. (Permutation
representations vary between different versions of \textsf{GAP}.) }

 

 
\begin{Verbatim}[commandchars=!@|,fontsize=\small,frame=single,label=Example]
  
  !gapprompt@gap>| !gapinput@d24 := DihedralGroup( IsPermGroup, 24 );; |
  !gapprompt@gap>| !gapinput@SetName( d24, "d24" );|
  !gapprompt@gap>| !gapinput@Y24 := XModByAutomorphismGroup( d24 );; |
  !gapprompt@gap>| !gapinput@Size2d( Y24 );|
  [ 24, 48 ]
  !gapprompt@gap>| !gapinput@X24 := Image( IsomorphismPerm2DimensionalGroup( Y24 ) );|
  [d24->Group([ (2,4), (1,2,3,4), (6,7), (5,6,7) ])]
  !gapprompt@gap>| !gapinput@nsx := NormalSubXMods( X24 );; |
  !gapprompt@gap>| !gapinput@Length( nsx );|
  40
  !gapprompt@gap>| !gapinput@ids := List( nsx, n -> IdGroup(n) );; |
  !gapprompt@gap>| !gapinput@pos1 := Position( ids, [ [4,1], [8,3] ] );;|
  !gapprompt@gap>| !gapinput@Xn1 := nsx[pos1]; |
  [Group( [ f2*f4^2, f3*f4 ] )->Group( [ f3, f4, f5 ] )]
  !gapprompt@gap>| !gapinput@nat1 := NaturalMorphismByNormalSubPreXMod( X24, Xn1 );; |
  !gapprompt@gap>| !gapinput@Qn1 := FactorPreXMod( X24, Xn1 );; |
  !gapprompt@gap>| !gapinput@[ Size2d( Xn1 ), Size2d( Qn1 ) ];|
  [ [ 4, 8 ], [ 6, 6 ] ]
  
\end{Verbatim}
 

\subsection{\textcolor{Chapter }{IntersectionSubXMods}}
\logpage{[ 4, 1, 2 ]}\nobreak
\hyperdef{L}{X8591E25680C5C575}{}
{\noindent\textcolor{FuncColor}{$\triangleright$\enspace\texttt{IntersectionSubXMods({\mdseries\slshape X0, X1, X2})\index{IntersectionSubXMods@\texttt{IntersectionSubXMods}}
\label{IntersectionSubXMods}
}\hfill{\scriptsize (operation)}}\\


 When \texttt{X1,X2} are subcrossed modules of \texttt{X0}, then the source and range of their intersection are the intersections of the
sources and ranges of \texttt{X1} and \texttt{X2} respectively. }

 

 
\begin{Verbatim}[commandchars=!@|,fontsize=\small,frame=single,label=Example]
  
  !gapprompt@gap>| !gapinput@pos2 := Position( ids, [ [24,6], [12,4] ] );;|
  !gapprompt@gap>| !gapinput@Xn2 := nsx[pos2];; |
  !gapprompt@gap>| !gapinput@IdGroup( Xn2 );         |
  [ [ 24, 6 ], [ 12, 4 ] ]
  !gapprompt@gap>| !gapinput@pos3 := Position( ids, [ [12,2], [24,5] ] );;|
  !gapprompt@gap>| !gapinput@Xn3 := nsx[pos3];; |
  !gapprompt@gap>| !gapinput@IdGroup( Xn3 );|
  [ [ 12, 2 ], [ 24, 5 ] ]
  !gapprompt@gap>| !gapinput@Xn23 := IntersectionSubXMods( X24, Xn2, Xn3 );;|
  !gapprompt@gap>| !gapinput@IdGroup( Xn23 );|
  [ [ 12, 2 ], [ 6, 2 ] ]
  
\end{Verbatim}
 

\subsection{\textcolor{Chapter }{Displacement}}
\logpage{[ 4, 1, 3 ]}\nobreak
\hyperdef{L}{X7E20208279038BB8}{}
{\noindent\textcolor{FuncColor}{$\triangleright$\enspace\texttt{Displacement({\mdseries\slshape alpha, r, s})\index{Displacement@\texttt{Displacement}}
\label{Displacement}
}\hfill{\scriptsize (operation)}}\\
\noindent\textcolor{FuncColor}{$\triangleright$\enspace\texttt{DisplacementGroup({\mdseries\slshape X0, Q, T})\index{DisplacementGroup@\texttt{DisplacementGroup}}
\label{DisplacementGroup}
}\hfill{\scriptsize (operation)}}\\
\noindent\textcolor{FuncColor}{$\triangleright$\enspace\texttt{DisplacementSubgroup({\mdseries\slshape X0})\index{DisplacementSubgroup@\texttt{DisplacementSubgroup}}
\label{DisplacementSubgroup}
}\hfill{\scriptsize (attribute)}}\\


 Commutators may be written $[r,q] = r^{-1}q^{-1}rq = (q^{-1})^rq = r^{-1}r^q$, and satisfy identities 
\[ [r,q]^p = [r^p,q^p], \qquad [pr,q] = [p,q]^r[r,q], \qquad [r,pq] =
[r,q][r,p]^q, \qquad [r,q]^{-1} = [q,r]. \]
 In a similar way, when a group $R$ acts on a group $S$, the \emph{displacement} of $s \in S$ by $r \in R$ is defined to be $\langle r,s \rangle := (s^{-1})^rs \in S$. When $\calX = (\partial : S \to R)$ is a pre\texttt{\symbol{45}}crossed module, the first crossed module axiom
requires $\partial\langle r,s \rangle = [r,\partial s]$. When $\alpha$ is the action of $\calX$, the \texttt{Displacement} function may be used to calculate $\langle r,s \rangle$. Displacements satisfy the following identities, where $s,t \in S,~ p,q,r \in R$: 
\[ \langle r,s \rangle^p = \langle r^p,s^p \rangle, \qquad \langle qr,s \rangle =
\langle q,s \rangle^r \langle r,s \rangle, \qquad \langle r,st \rangle =
\langle r,t \rangle \langle r,s \rangle^t, \qquad \langle r,s \rangle^{-1} =
\langle r^{-1},s^r \rangle. \]
 The operation \texttt{DisplacementGroup} applied to \texttt{X0,Q,T} is the subgroup of \texttt{S} consisting of all the displacements $\langle r,s \rangle, r \in Q \leqslant R, s \in T \leqslant S$. The \texttt{DisplacementSubgroup} of $\calX$ is the subgroup $\Disp(\calX)$ of $S$ given by \texttt{DisplacementGroup(X0,R,S)}. The identities imply $\langle r,s \rangle^t = \langle r,st^{r^{-1}} \rangle \langle r^{-1},t \rangle$, so $\Disp(\calX)$ is normal in $S$. }

 

 
\begin{Verbatim}[commandchars=!@|,fontsize=\small,frame=single,label=Example]
  
  !gapprompt@gap>| !gapinput@pos4 := Position( ids, [ [6,2], [24,14] ] );;|
  !gapprompt@gap>| !gapinput@Xn4 := nsx[pos4];; |
  !gapprompt@gap>| !gapinput@bn4 := Boundary( Xn4 );;|
  !gapprompt@gap>| !gapinput@Sn4 := Source(Xn4);; |
  !gapprompt@gap>| !gapinput@Rn4 := Range(Xn4);; |
  !gapprompt@gap>| !gapinput@genRn4 := GeneratorsOfGroup( Rn4 );;|
  !gapprompt@gap>| !gapinput@L := List( genRn4, g -> ( Order(g) = 2 ) and |
  !gapprompt@>| !gapinput@                not ( IsNormal( Rn4, Subgroup( Rn4, [g] ) ) ) );;|
  !gapprompt@gap>| !gapinput@pos := Position( L, true );;|
  !gapprompt@gap>| !gapinput@s := Sn4.1;  r := genRn4[pos]; |
  (1,3,5,7,9,11)(2,4,6,8,10,12)
  (6,7)
  !gapprompt@gap>| !gapinput@act := XModAction( Xn4 );; |
  !gapprompt@gap>| !gapinput@d := Displacement( act, r, s );|
  (1,5,9)(2,6,10)(3,7,11)(4,8,12)
  !gapprompt@gap>| !gapinput@Image( bn4, d ) = Comm( r, Image( bn4, s ) );  |
  true
  !gapprompt@gap>| !gapinput@Qn4 := Subgroup( Rn4, [ (6,7), (1,3), (2,4) ] );;   |
  !gapprompt@gap>| !gapinput@Tn4 := Subgroup( Sn4, [ (1,3,5,7,9,11)(2,4,6,8,10,12) ] );;|
  !gapprompt@gap>| !gapinput@DisplacementGroup( Xn4, Qn4, Tn4 );                        |
  Group([ (1,5,9)(2,6,10)(3,7,11)(4,8,12) ])
  !gapprompt@gap>| !gapinput@DisplacementSubgroup( Xn4 );|
  Group([ (1,5,9)(2,6,10)(3,7,11)(4,8,12) ])
  
\end{Verbatim}
 

\subsection{\textcolor{Chapter }{CommutatorSubXMod}}
\logpage{[ 4, 1, 4 ]}\nobreak
\hyperdef{L}{X86ACB83E7D70C625}{}
{\noindent\textcolor{FuncColor}{$\triangleright$\enspace\texttt{CommutatorSubXMod({\mdseries\slshape X, X1, X2})\index{CommutatorSubXMod@\texttt{CommutatorSubXMod}}
\label{CommutatorSubXMod}
}\hfill{\scriptsize (operation)}}\\
\noindent\textcolor{FuncColor}{$\triangleright$\enspace\texttt{CrossActionSubgroup({\mdseries\slshape X, X1, X2})\index{CrossActionSubgroup@\texttt{CrossActionSubgroup}}
\label{CrossActionSubgroup}
}\hfill{\scriptsize (operation)}}\\


 When $\calX_1 = (N \to Q), \calX_2 = (M \to P)$ are two normal subcrossed modules of $\calX = (\partial : S \to R)$, the displacements $\langle p,n \rangle$ and $\langle q,m \rangle$ all map by $\partial$ into $[Q,P]$. These displacements form a normal subgroup of $S$, called the \texttt{CrossActionSubgroup}. The \texttt{CommutatorSubXMod} $[\calX_1,\calX_2]$ has this subgroup as source and $[P,Q]$ as range, and is normal in $\calX$. }

 

 
\begin{Verbatim}[commandchars=!@|,fontsize=\small,frame=single,label=Example]
  
  !gapprompt@gap>| !gapinput@CAn23 := CrossActionSubgroup( X24, Xn2, Xn3 );;|
  !gapprompt@gap>| !gapinput@IdGroup( CAn23 );|
  [ 12, 2 ]
  !gapprompt@gap>| !gapinput@Cn23 := CommutatorSubXMod( X24, Xn2, Xn3 );;|
  !gapprompt@gap>| !gapinput@IdGroup( Cn23 );|
  [ [ 12, 2 ], [ 6, 2 ] ]
  !gapprompt@gap>| !gapinput@Xn23 = Cn23;|
  true
  
\end{Verbatim}
 

\subsection{\textcolor{Chapter }{DerivedSubXMod}}
\logpage{[ 4, 1, 5 ]}\nobreak
\hyperdef{L}{X86E0804B780A7FD6}{}
{\noindent\textcolor{FuncColor}{$\triangleright$\enspace\texttt{DerivedSubXMod({\mdseries\slshape X0})\index{DerivedSubXMod@\texttt{DerivedSubXMod}}
\label{DerivedSubXMod}
}\hfill{\scriptsize (attribute)}}\\


 The \texttt{DerivedSubXMod} of $\calX$ is the normal subcrossed module $[\calX,\calX] = (\partial' : \Disp(\calX) \to [R,R])$ where $\partial'$ is the restriction of $\partial$ (see page 66 of Norrie's thesis \cite{N2}). }

 

 
\begin{Verbatim}[commandchars=!@|,fontsize=\small,frame=single,label=Example]
  
  !gapprompt@gap>| !gapinput@DXn4 := DerivedSubXMod( Xn4 );;|
  !gapprompt@gap>| !gapinput@IdGroup( DXn4 );|
  [ [ 3, 1 ], [ 3, 1 ] ]
  
\end{Verbatim}
 

\subsection{\textcolor{Chapter }{FixedPointSubgroupXMod}}
\logpage{[ 4, 1, 6 ]}\nobreak
\hyperdef{L}{X85640DD17F5A2949}{}
{\noindent\textcolor{FuncColor}{$\triangleright$\enspace\texttt{FixedPointSubgroupXMod({\mdseries\slshape X0, T, Q})\index{FixedPointSubgroupXMod@\texttt{FixedPointSubgroupXMod}}
\label{FixedPointSubgroupXMod}
}\hfill{\scriptsize (operation)}}\\
\noindent\textcolor{FuncColor}{$\triangleright$\enspace\texttt{StabilizerSubgroupXMod({\mdseries\slshape X0, T, Q})\index{StabilizerSubgroupXMod@\texttt{StabilizerSubgroupXMod}}
\label{StabilizerSubgroupXMod}
}\hfill{\scriptsize (operation)}}\\


 The \texttt{FixedPointSubgroupXMod(X,T,Q)} for $\calX=(\partial : S \to R)$ is the subgroup $\Fix(\calX,T,Q)$ of elements $t \in T \leqslant S$ individually fixed under the action of $Q \leqslant R$. 

 The \texttt{StabilizerSubgroupXMod(X,T,Q)} for $\calX$ is the subgroup $\Stab(\calX,T,Q)$ of $Q \leqslant R$ whose elements act trivially on the whole of $T \leqslant S$ (see page 19 of Norrie's thesis \cite{N2}). }

 

 
\begin{Verbatim}[commandchars=!@|,fontsize=\small,frame=single,label=Example]
  
  !gapprompt@gap>| !gapinput@fix := FixedPointSubgroupXMod( Xn4, Sn4, Rn4 );|
  Group([ (1,7)(2,8)(3,9)(4,10)(5,11)(6,12) ])
  !gapprompt@gap>| !gapinput@stab := StabilizerSubgroupXMod( Xn4, Sn4, Rn4 );;|
  !gapprompt@gap>| !gapinput@IdGroup( stab );|
  [ 12, 5 ]
  
\end{Verbatim}
 

\subsection{\textcolor{Chapter }{CentreXMod}}
\logpage{[ 4, 1, 7 ]}\nobreak
\hyperdef{L}{X7B57446086BA1BF0}{}
{\noindent\textcolor{FuncColor}{$\triangleright$\enspace\texttt{CentreXMod({\mdseries\slshape X0})\index{CentreXMod@\texttt{CentreXMod}}
\label{CentreXMod}
}\hfill{\scriptsize (attribute)}}\\
\noindent\textcolor{FuncColor}{$\triangleright$\enspace\texttt{Centralizer({\mdseries\slshape X, Y})\index{Centralizer@\texttt{Centralizer}}
\label{Centralizer}
}\hfill{\scriptsize (operation)}}\\
\noindent\textcolor{FuncColor}{$\triangleright$\enspace\texttt{Normalizer({\mdseries\slshape X, Y})\index{Normalizer@\texttt{Normalizer}}
\label{Normalizer}
}\hfill{\scriptsize (operation)}}\\


 The \emph{centre} $Z(\calX)$ of $\calX = (\partial : S \to R)$ has as source the fixed point subgroup $\Fix(\calX,S,R)$. The range is the intersection of the centre $Z(R)$ with the stabilizer subgroup. $Z(\calX)$ is constructed in a different way as \texttt{XModCentre} (\ref{XModCentre}). 

 When $\calY = (T \to Q)$ is a subcrossed module of $\calX = (\partial : S \to R)$, the \emph{centralizer} $C_{\calX}(\calY)$ of $\calY$ has as source the fixed point subgroup $\Fix(\calX,S,Q)$. The range is the intersection of the centralizer $C_R(Q)$ with $\Stab(\calX,T,R)$. 

 The \emph{normalizer} $N_{\calX}(\calY)$ of $\calY$ has as source the subgroup of $S$ consisting of the displacements $\langle s,q \rangle$ which lie in $S$. }

 

 
\begin{Verbatim}[commandchars=!@|,fontsize=\small,frame=single,label=Example]
  
  !gapprompt@gap>| !gapinput@ZXn4 := CentreXMod( Xn4 );      |
  [Group( [ f3*f4 ] )->Group( [ f3, f5 ] )]
  !gapprompt@gap>| !gapinput@IdGroup( ZXn4 );|
  [ [ 2, 1 ], [ 4, 2 ] ]
  !gapprompt@gap>| !gapinput@CDXn4 := Centralizer( Xn4, DXn4 );|
  [Group( [ f3*f4 ] )->Group( [ f2 ] )]
  !gapprompt@gap>| !gapinput@IdGroup( CDXn4 );    |
  [ [ 2, 1 ], [ 3, 1 ] ]
  !gapprompt@gap>| !gapinput@NDXn4 := Normalizer( Xn4, DXn4 ); |
  [Group( <identity> of ... )->Group( [ f5, f2*f3 ] )]
  !gapprompt@gap>| !gapinput@IdGroup( NDXn4 );|
  [ [ 1, 1 ], [ 12, 5 ] ]
  
\end{Verbatim}
 

\subsection{\textcolor{Chapter }{CentralQuotient}}
\logpage{[ 4, 1, 8 ]}\nobreak
\hyperdef{L}{X814D9E1E78EEE665}{}
{\noindent\textcolor{FuncColor}{$\triangleright$\enspace\texttt{CentralQuotient({\mdseries\slshape G})\index{CentralQuotient@\texttt{CentralQuotient}}
\label{CentralQuotient}
}\hfill{\scriptsize (attribute)}}\\


 The \texttt{CentralQuotient} of a group $G$ is the crossed module $(G \to G/Z(G))$ with the natural homomorphism as the boundary map. This is a special case of \texttt{XModByCentralExtension} (\ref{XModByCentralExtension}). 

 Similarly, the central quotient of a crossed module $\calX$ is the crossed square $(\calX \Rightarrow \calX/Z(\calX))$ (see section \ref{sect-xsq-constructions}). }

 

 
\begin{Verbatim}[commandchars=!@|,fontsize=\small,frame=single,label=Example]
  
  !gapprompt@gap>| !gapinput@Q24 := CentralQuotient( d24);  IdGroup( Q24 );                     |
  [d24->d24/Z(d24)]
  [ [ 24, 6 ], [ 12, 4 ] ]
  
\end{Verbatim}
 

\subsection{\textcolor{Chapter }{IsAbelian2DimensionalGroup}}
\logpage{[ 4, 1, 9 ]}\nobreak
\hyperdef{L}{X7F4222757B0E08B6}{}
{\noindent\textcolor{FuncColor}{$\triangleright$\enspace\texttt{IsAbelian2DimensionalGroup({\mdseries\slshape X0})\index{IsAbelian2DimensionalGroup@\texttt{IsAbelian2DimensionalGroup}}
\label{IsAbelian2DimensionalGroup}
}\hfill{\scriptsize (property)}}\\
\noindent\textcolor{FuncColor}{$\triangleright$\enspace\texttt{IsAspherical2DimensionalGroup({\mdseries\slshape X0})\index{IsAspherical2DimensionalGroup@\texttt{IsAspherical2DimensionalGroup}}
\label{IsAspherical2DimensionalGroup}
}\hfill{\scriptsize (property)}}\\
\noindent\textcolor{FuncColor}{$\triangleright$\enspace\texttt{IsSimplyConnected2DimensionalGroup({\mdseries\slshape X0})\index{IsSimplyConnected2DimensionalGroup@\texttt{IsSimplyConnected2DimensionalGroup}}
\label{IsSimplyConnected2DimensionalGroup}
}\hfill{\scriptsize (property)}}\\
\noindent\textcolor{FuncColor}{$\triangleright$\enspace\texttt{IsFaithful2DimensionalGroup({\mdseries\slshape X0})\index{IsFaithful2DimensionalGroup@\texttt{IsFaithful2DimensionalGroup}}
\label{IsFaithful2DimensionalGroup}
}\hfill{\scriptsize (property)}}\\


 A crossed module is \emph{abelian} if it equal to its centre. This is the case when the range group is abelian
and the action is trivial. 

 A crossed module is \emph{aspherical} if the boundary has trivial kernel. 

 A crossed module is \emph{simply connected} if the boundary has trivial cokernel. 

 A crossed module is \emph{faithful} if the action is faithful. }

 

 
\begin{Verbatim}[commandchars=!@|,fontsize=\small,frame=single,label=Example]
  
  !gapprompt@gap>| !gapinput@[ IsAbelian2DimensionalGroup(Xn4), IsAbelian2DimensionalGroup(X24) ];
|
  [ false, false ]
  !gapprompt@gap>| !gapinput@pos7 := Position( ids, [ [3,1], [6,1] ] );;
|
  !gapprompt@gap>| !gapinput@[ IsAspherical2DimensionalGroup(nsx[pos7]), IsAspherical2DimensionalGroup(X24) ];
|
  [ true, false ]
  !gapprompt@gap>| !gapinput@[ IsSimplyConnected2DimensionalGroup(Xn4), IsSimplyConnected2DimensionalGroup(X24) ];|
  [ true, true ]
  !gapprompt@gap>| !gapinput@[ IsFaithful2DimensionalGroup(Xn4), IsFaithful2DimensionalGroup(X24) ];              |
  [ false, true ] 
  
\end{Verbatim}
 

\subsection{\textcolor{Chapter }{LowerCentralSeriesOfXMod}}
\logpage{[ 4, 1, 10 ]}\nobreak
\hyperdef{L}{X87C524C08588AAC0}{}
{\noindent\textcolor{FuncColor}{$\triangleright$\enspace\texttt{LowerCentralSeriesOfXMod({\mdseries\slshape X0})\index{LowerCentralSeriesOfXMod@\texttt{LowerCentralSeriesOfXMod}}
\label{LowerCentralSeriesOfXMod}
}\hfill{\scriptsize (attribute)}}\\
\noindent\textcolor{FuncColor}{$\triangleright$\enspace\texttt{IsNilpotent2DimensionalGroup({\mdseries\slshape X0})\index{IsNilpotent2DimensionalGroup@\texttt{IsNilpotent2DimensionalGroup}}
\label{IsNilpotent2DimensionalGroup}
}\hfill{\scriptsize (property)}}\\
\noindent\textcolor{FuncColor}{$\triangleright$\enspace\texttt{NilpotencyClass2DimensionalGroup({\mdseries\slshape X0})\index{NilpotencyClass2DimensionalGroup@\texttt{NilpotencyClass2DimensionalGroup}}
\label{NilpotencyClass2DimensionalGroup}
}\hfill{\scriptsize (attribute)}}\\


 Let $\calY$ be a subcrossed module of $\calX$. A \emph{ series of length} $n$ from $\calX$ to $\calY$ has the form 
\[ \calX ~=~ \calX_0 ~\unrhd~ \calX_1 ~\unrhd~ \cdots ~\unrhd~ \calX_i ~\unrhd~
\cdots ~\unrhd~ \calX_n ~=~ \calY \quad (1 \leqslant i \leqslant n). \]
 The quotients $\calF_i = \calX_i / \calX_{i-1}$ are the \emph{factors} of the series. 

 A factor $\calF_i$ is \emph{central} if $\calX_{i-1} \unlhd \calX$ and $\calF_i$ is a subcrossed module of the centre of $\calX / \calX_{i-1}$. 

 A series is \emph{central} if all its factors are central. 

 $\calX$ is \emph{soluble} if it has a series all of whose factors are abelian. 

 $\calX$ is \emph{nilpotent} is it has a series all of whose factors are central factors of $\calX$. 

 The \emph{lower central series} of $\calX$ is the sequence (see \cite{N2}, p.77): 
\[ \calX ~=~ \Gamma_1(\calX) ~\unrhd~ \Gamma_2(\calX) ~\unrhd~ \cdots \qquad
\mbox{where} \qquad \Gamma_j(\calX) ~=~ [ \Gamma_{j-1}(\calX), \calX]. \]
 If $\calX$ is nilpotent, then its lower central series is its most rapidly descending
central series. 

 The least integer $c$ such that $\Gamma_{c+1}(\calX)$ is the trivial crossed module is the \emph{nilpotency class} of $\calX$. }

 

 
\begin{Verbatim}[commandchars=!@|,fontsize=\small,frame=single,label=Example]
  
  !gapprompt@gap>| !gapinput@lcs := LowerCentralSeries( X24 );;      |
  !gapprompt@gap>| !gapinput@List( lcs, g -> IdGroup(g) );|
  [ [ [ 24, 6 ], [ 48, 38 ] ], [ [ 12, 2 ], [ 6, 2 ] ], [ [ 6, 2 ], [ 3, 1 ] ], 
    [ [ 3, 1 ], [ 3, 1 ] ] ]
  !gapprompt@gap>| !gapinput@IsNilpotent2DimensionalGroup( X24 );      |
  false
  !gapprompt@gap>| !gapinput@NilpotencyClassOf2DimensionalGroup( X24 );|
  0
  
\end{Verbatim}
 

\subsection{\textcolor{Chapter }{IsomorphismXMods}}
\logpage{[ 4, 1, 11 ]}\nobreak
\hyperdef{L}{X7C67623F797A0301}{}
{\noindent\textcolor{FuncColor}{$\triangleright$\enspace\texttt{IsomorphismXMods({\mdseries\slshape X1, X2})\index{IsomorphismXMods@\texttt{IsomorphismXMods}}
\label{IsomorphismXMods}
}\hfill{\scriptsize (operation)}}\\


 The function \texttt{IsomorphismXMods} computes an isomorphism $\mu : \calX_1 \to \calX_2$, provided one exists, or else returns \texttt{fail}. }

 

 
\begin{Verbatim}[commandchars=!@|,fontsize=\small,frame=single,label=Example]
  
  !gapprompt@gap>| !gapinput@gend24 := GeneratorsOfGroup( d24 );;             |
  !gapprompt@gap>| !gapinput@a := gend24[1];; b:= gend24[2];; |
  !gapprompt@gap>| !gapinput@J := Subgroup( d24, [a^2,b] ); |
  Group([ (1,3,5,7,9,11)(2,4,6,8,10,12), (2,12)(3,11)(4,10)(5,9)(6,8) ])
  !gapprompt@gap>| !gapinput@K := Subgroup( d24, [a^2,a*b] );|
  Group([ (1,3,5,7,9,11)(2,4,6,8,10,12), (1,12)(2,11)(3,10)(4,9)(5,8)(6,7) ])
  !gapprompt@gap>| !gapinput@XJ := XModByNormalSubgroup( d24, J );;|
  !gapprompt@gap>| !gapinput@XK := XModByNormalSubgroup( d24, K );;|
  !gapprompt@gap>| !gapinput@iso := IsomorphismXMods( XJ, XK );;|
  !gapprompt@gap>| !gapinput@SourceHom( iso );|
  [ (1,3,5,7,9,11)(2,4,6,8,10,12), (2,12)(3,11)(4,10)(5,9)(6,8) ] -> 
  [ (1,3,5,7,9,11)(2,4,6,8,10,12), (1,12)(2,11)(3,10)(4,9)(5,8)(6,7) ]
  !gapprompt@gap>| !gapinput@RangeHom( iso ); |
  [ (1,2,3,4,5,6,7,8,9,10,11,12), (2,12)(3,11)(4,10)(5,9)(6,8) ] -> 
  [ (1,2,3,4,5,6,7,8,9,10,11,12), (1,12)(2,11)(3,10)(4,9)(5,8)(6,7) ]
  
\end{Verbatim}
 

\subsection{\textcolor{Chapter }{AllXMods}}
\logpage{[ 4, 1, 12 ]}\nobreak
\hyperdef{L}{X81EE2188863E6E85}{}
{\noindent\textcolor{FuncColor}{$\triangleright$\enspace\texttt{AllXMods({\mdseries\slshape args})\index{AllXMods@\texttt{AllXMods}}
\label{AllXMods}
}\hfill{\scriptsize (function)}}\\
\noindent\textcolor{FuncColor}{$\triangleright$\enspace\texttt{AllXModsWithGroups({\mdseries\slshape src, rng})\index{AllXModsWithGroups@\texttt{AllXModsWithGroups}}
\label{AllXModsWithGroups}
}\hfill{\scriptsize (operation)}}\\
\noindent\textcolor{FuncColor}{$\triangleright$\enspace\texttt{AllXModsUpToIsomorphism({\mdseries\slshape src, rng})\index{AllXModsUpToIsomorphism@\texttt{AllXModsUpToIsomorphism}}
\label{AllXModsUpToIsomorphism}
}\hfill{\scriptsize (operation)}}\\
\noindent\textcolor{FuncColor}{$\triangleright$\enspace\texttt{IsomorphismClassRepresentatives2dGroups({\mdseries\slshape L})\index{IsomorphismClassRepresentatives2dGroups@\texttt{Isomorphism}\-\texttt{Class}\-\texttt{Representatives2d}\-\texttt{Groups}}
\label{IsomorphismClassRepresentatives2dGroups}
}\hfill{\scriptsize (operation)}}\\


 The global function \texttt{AllXMods} may be called in three ways. Firstly, as \texttt{AllXMods(S,R)} to compute all crossed modules with chosen source and range groups: this calls \texttt{AllXModsWithGroups(S,R)}. Secondly, \texttt{AllXMods([n,m])} computes all crossed modules with a given size \texttt{[n,m]}. Thirdly \texttt{AllXMods(ord)} to compute all crossed modules whose associated cat1\texttt{\symbol{45}}groups
have a given size \texttt{ord}. 

 The function \texttt{AllXModsUpToIsomorphism(S,R)} returns a list of representatives of the isomorphism classes of crossed
modules with source $S$ and range $R$. 

 If \texttt{L} is a list returned by, for example, \texttt{AllXModsWithGroups(S,R)}, then the isomorphism class representatives for this list is returned by \texttt{IsomorphismClassRepresentatives2dGroups(L)}. This result is the same as that given by \texttt{AllXModsUpToIsomorphism(S,R)}. 

 In the example we see that there are $4$ crossed modules, in $3$ isomorphism classes, $(C_6 \to S_3)$; forming a subset of the $17$ crossed modules with size \texttt{[6,6]}; and that these form a subset of the $205$ crossed modules whose cat1\texttt{\symbol{45}}group has size $36$. There are $40$ precrossed modules with size \texttt{[6,6]}. }

 

 
\begin{Verbatim}[commandchars=!@|,fontsize=\small,frame=single,label=Example]
  
  !gapprompt@gap>| !gapinput@c6 := SmallGroup( 6, 2 );; |
  !gapprompt@gap>| !gapinput@s3 := SmallGroup( 6, 1 );; |
  !gapprompt@gap>| !gapinput@Ac6s3 := AllXMods( c6, s3 );;   |
  !gapprompt@gap>| !gapinput@Length( Ac6s3 );           |
  4
  !gapprompt@gap>| !gapinput@Ic6s3 := AllXModsUpToIsomorphism( c6, s3 );; |
  !gapprompt@gap>| !gapinput@List( Ic6s3, obj -> IsTrivialAction2DimensionalGroup( obj ) ); |
  [ true, false, false ]
  !gapprompt@gap>| !gapinput@Kc6s3 := List( Ic6s3, obj -> KernelCokernelXMod( obj ) );; |
  !gapprompt@gap>| !gapinput@List( Kc6s3, obj -> IdGroup( obj ) );|
  [ [ [ 6, 2 ], [ 6, 1 ] ], [ [ 6, 2 ], [ 6, 1 ] ], [ [ 2, 1 ], [ 2, 1 ] ] ]
  [ ]
  !gapprompt@gap>| !gapinput@A66 := AllXMods( [6,6] );;   |
  !gapprompt@gap>| !gapinput@Length( A66 );|
  17
  !gapprompt@gap>| !gapinput@IA66 := IsomorphismClassRepresentatives2dGroups( A66 );;|
  !gapprompt@gap>| !gapinput@Length( IA66 );|
  9
  !gapprompt@gap>| !gapinput@x36 := AllXMods( 36 );; |
  !gapprompt@gap>| !gapinput@Length( x36 ); |
  205
  !gapprompt@gap>| !gapinput@size36 := List( x36, x -> Size2d( x ) );;|
  !gapprompt@gap>| !gapinput@Collected( size36 );|
  [ [ [ 1, 36 ], 14 ], [ [ 2, 18 ], 7 ], [ [ 3, 12 ], 21 ], [ [ 4, 9 ], 14 ], 
    [ [ 6, 6 ], 17 ], [ [ 9, 4 ], 102 ], [ [ 12, 3 ], 8 ], [ [ 18, 2 ], 18 ], 
    [ [ 36, 1 ], 4 ] ]
  
\end{Verbatim}
 }

 
\section{\textcolor{Chapter }{Isoclinism for groups}}\label{sect-isoclinic-groups}
\logpage{[ 4, 2, 0 ]}
\hyperdef{L}{X7B0D511A82FD945E}{}
{
  

\subsection{\textcolor{Chapter }{Isoclinism (for groups)}}
\logpage{[ 4, 2, 1 ]}\nobreak
\hyperdef{L}{X7B0D511A82FD945E}{}
{\noindent\textcolor{FuncColor}{$\triangleright$\enspace\texttt{Isoclinism({\mdseries\slshape G, H})\index{Isoclinism@\texttt{Isoclinism}!for groups}
\label{Isoclinism:for groups}
}\hfill{\scriptsize (operation)}}\\
\noindent\textcolor{FuncColor}{$\triangleright$\enspace\texttt{AreIsoclinicDomains({\mdseries\slshape G, H})\index{AreIsoclinicDomains@\texttt{AreIsoclinicDomains}!for groups}
\label{AreIsoclinicDomains:for groups}
}\hfill{\scriptsize (operation)}}\\


 Let $G,H$ be groups with central quotients $Q(G)$ and $Q(H)$ and derived subgroups $[G,G]$ and $[H,H]$ respectively. Let $c_G : G/Z(G) \times G/Z(G) \to [G,G]$ and $c_H : H/Z(H) \times H/Z(H) \to [H,H]$ be the two commutator maps. An \emph{isoclinism} $G \sim H$ is a pair of isomorphisms $(\eta,\xi)$ where $\eta : G/Z(G) \to H/Z(H)$ and $\xi : [G,G] \to [H,H]$ such that $c_G * \xi = (\eta \times \eta) * c_H$. Isoclinism is an equivalence relation, and all abelian groups are isoclinic
to the trivial group. 

 }

 

 
\begin{Verbatim}[commandchars=!@|,fontsize=\small,frame=single,label=Example]
  
  !gapprompt@gap>| !gapinput@G := SmallGroup( 64, 6 );;  StructureDescription( G ); |
  "(C8 x C4) : C2"
  !gapprompt@gap>| !gapinput@QG := CentralQuotient( G );;  IdGroup( QG );|
  [ [ 64, 6 ], [ 8, 3 ] ]
  !gapprompt@gap>| !gapinput@H := SmallGroup( 32, 41 );;  StructureDescription( H );|
  "C2 x Q16"
  !gapprompt@gap>| !gapinput@QH := CentralQuotient( H );;  IdGroup( QH );|
  [ [ 32, 41 ], [ 8, 3 ] ]
  !gapprompt@gap>| !gapinput@Isoclinism( G, H );|
  [ [ f1, f2, f3 ] -> [ f1, f2*f3, f3 ], [ f3, f5 ] -> [ f4*f5, f5 ] ]
  !gapprompt@gap>| !gapinput@K := SmallGroup( 32, 43 );;  StructureDescription( K );|
  "(C2 x D8) : C2"
  !gapprompt@gap>| !gapinput@QK := CentralQuotient( K );;  IdGroup( QK );                       |
  [ [ 32, 43 ], [ 16, 11 ] ]
  !gapprompt@gap>| !gapinput@AreIsoclinicDomains( G, K );|
  false
  
\end{Verbatim}
 

\subsection{\textcolor{Chapter }{IsStemDomain (for groups)}}
\logpage{[ 4, 2, 2 ]}\nobreak
\hyperdef{L}{X7C72991985B58DB8}{}
{\noindent\textcolor{FuncColor}{$\triangleright$\enspace\texttt{IsStemDomain({\mdseries\slshape G})\index{IsStemDomain@\texttt{IsStemDomain}!for groups}
\label{IsStemDomain:for groups}
}\hfill{\scriptsize (property)}}\\
\noindent\textcolor{FuncColor}{$\triangleright$\enspace\texttt{IsoclinicStemDomain({\mdseries\slshape G})\index{IsoclinicStemDomain@\texttt{IsoclinicStemDomain}!for groups}
\label{IsoclinicStemDomain:for groups}
}\hfill{\scriptsize (attribute)}}\\
\noindent\textcolor{FuncColor}{$\triangleright$\enspace\texttt{AllStemGroupIds({\mdseries\slshape n})\index{AllStemGroupIds@\texttt{AllStemGroupIds}}
\label{AllStemGroupIds}
}\hfill{\scriptsize (operation)}}\\
\noindent\textcolor{FuncColor}{$\triangleright$\enspace\texttt{AllStemGroupFamilies({\mdseries\slshape n})\index{AllStemGroupFamilies@\texttt{AllStemGroupFamilies}}
\label{AllStemGroupFamilies}
}\hfill{\scriptsize (operation)}}\\


 A group $G$ is a \emph{stem group} if $Z(G) \leq [G,G]$. Every group is isoclinic to a stem group, but distinct stem groups may be
isoclinic. For example, groups $D_8, Q_8$ are two isoclinic stem groups. 

 The function \texttt{IsoclinicStemDomain } returns a stem group isoclinic to $G$. 

 The function \texttt{AllStemGroupIds} returns the \texttt{IdGroup} list of the stem groups of a specified size, while \texttt{AllStemGroupFamilies} splits this list into isoclinism classes. }

 

 
\begin{Verbatim}[commandchars=!@|,fontsize=\small,frame=single,label=Example]
  
  !gapprompt@gap>| !gapinput@DerivedSubgroup(G);     |
  Group([ f3, f5 ])
  !gapprompt@gap>| !gapinput@IsStemDomain( G );|
  false
  !gapprompt@gap>| !gapinput@IsoclinicStemDomain( G );|
  <pc group of size 16 with 4 generators>
  !gapprompt@gap>| !gapinput@AllStemGroupIds( 32 );|
  [ [ 32, 6 ], [ 32, 7 ], [ 32, 8 ], [ 32, 18 ], [ 32, 19 ], [ 32, 20 ], 
    [ 32, 27 ], [ 32, 28 ], [ 32, 29 ], [ 32, 30 ], [ 32, 31 ], [ 32, 32 ], 
    [ 32, 33 ], [ 32, 34 ], [ 32, 35 ], [ 32, 43 ], [ 32, 44 ], [ 32, 49 ], 
    [ 32, 50 ] ]
  !gapprompt@gap>| !gapinput@AllStemGroupFamilies( 32 );|
  [ [ [ 32, 6 ], [ 32, 7 ], [ 32, 8 ] ], [ [ 32, 18 ], [ 32, 19 ], [ 32, 20 ] ],
    [ [ 32, 27 ], [ 32, 28 ], [ 32, 29 ], [ 32, 30 ], [ 32, 31 ], [ 32, 32 ], 
        [ 32, 33 ], [ 32, 34 ], [ 32, 35 ] ], [ [ 32, 43 ], [ 32, 44 ] ], 
    [ [ 32, 49 ], [ 32, 50 ] ] ]
  
\end{Verbatim}
 

\subsection{\textcolor{Chapter }{IsoclinicRank (for groups)}}
\logpage{[ 4, 2, 3 ]}\nobreak
\label{group-rank}
\hyperdef{L}{X82DD52587F81C95C}{}
{\noindent\textcolor{FuncColor}{$\triangleright$\enspace\texttt{IsoclinicRank({\mdseries\slshape G})\index{IsoclinicRank@\texttt{IsoclinicRank}!for groups}
\label{IsoclinicRank:for groups}
}\hfill{\scriptsize (attribute)}}\\
\noindent\textcolor{FuncColor}{$\triangleright$\enspace\texttt{IsoclinicMiddleLength({\mdseries\slshape G})\index{IsoclinicMiddleLength@\texttt{IsoclinicMiddleLength}!for groups}
\label{IsoclinicMiddleLength:for groups}
}\hfill{\scriptsize (attribute)}}\\


 Let $G$ be a finite $p$\texttt{\symbol{45}}group. Then $\log_p |[G,G] / (Z(G) \cap [G,G])|$ is called the \emph{middle length} of $G$. Also $\log_p |Z(G) \cap [G,G]| + \log_p |G/Z(G)|$ is called the \emph{rank} of $G$. These invariants appear in the tables of isoclinism families of groups of
order 128 in \cite{JNO}. }

 

 
\begin{Verbatim}[commandchars=!@|,fontsize=\small,frame=single,label=Example]
  
  !gapprompt@gap>| !gapinput@IsoclinicMiddleLength( G );|
  1
  !gapprompt@gap>| !gapinput@IsoclinicRank( G );|
  4
  
\end{Verbatim}
 }

 
\section{\textcolor{Chapter }{Isoclinism for crossed modules}}\label{sect-isoclinic-xmods}
\logpage{[ 4, 3, 0 ]}
\hyperdef{L}{X81338C977972AD83}{}
{
  

\subsection{\textcolor{Chapter }{Isoclinism (for crossed modules)}}
\logpage{[ 4, 3, 1 ]}\nobreak
\hyperdef{L}{X81338C977972AD83}{}
{\noindent\textcolor{FuncColor}{$\triangleright$\enspace\texttt{Isoclinism({\mdseries\slshape X0, Y0})\index{Isoclinism@\texttt{Isoclinism}!for crossed modules}
\label{Isoclinism:for crossed modules}
}\hfill{\scriptsize (operation)}}\\
\noindent\textcolor{FuncColor}{$\triangleright$\enspace\texttt{AreIsoclinicDomains({\mdseries\slshape X0, Y0})\index{AreIsoclinicDomains@\texttt{AreIsoclinicDomains}!for crossed modules of groups}
\label{AreIsoclinicDomains:for crossed modules of groups}
}\hfill{\scriptsize (operation)}}\\


 Let $\calX,\calY$ be crossed modules with central quotients $Q(\calX)$ and $ Q(\calY)$, and derived subcrossed modules $[\calX,\calX]$ and $[\calY,\calY]$ respectively. Let $c_\calX : Q(\calX) \times Q(\calX) \to [\calX,\calX]$ and $c_\calY : Q(\calY) \times Q(\calY) \to [\calY,\calY]$ be the two commutator maps. An \emph{isoclinism} $\calX \sim \calY$ is a pair of bijective morphisms $(\eta,\xi)$ where $\eta : Q(\calX) \to Q(\calY)$ and $\xi : [\calX,\calX] \to [\calY,\calY]$ such that $c_\calX * \xi = (\eta \times \eta) * c_\calY$. Isoclinism is an equivalence relation, and all abelian crossed modules are
isoclinic to the trivial crossed module. 

 }

 

 
\begin{Verbatim}[commandchars=!@|,fontsize=\small,frame=single,label=Example]
  
  !gapprompt@gap>| !gapinput@C8 := Cat1Group( 16, 8, 2 );;|
  !gapprompt@gap>| !gapinput@X8 := XMod(C8);  IdGroup( X8 );|
  [Group( [ f1*f2*f3, f3, f4 ] )->Group( [ f2, f2 ] )]
  [ [ 8, 1 ], [ 2, 1 ] ]
  !gapprompt@gap>| !gapinput@C9 := Cat1Group( 32, 9, 2 );|
  [(C8 x C2) : C2 => Group( [ f2, f2 ] )]
  !gapprompt@gap>| !gapinput@X9 := XMod( C9 );  IdGroup( X9 );|
  [Group( [ f1*f2*f3, f3, f4, f5 ] )->Group( [ f2, f2 ] )]
  [ [ 16, 5 ], [ 2, 1 ] ]
  !gapprompt@gap>| !gapinput@AreIsoclinicDomains( X8, X9 );|
  true
  !gapprompt@gap>| !gapinput@ism89 := Isoclinism( X8, X9 );;|
  !gapprompt@gap>| !gapinput@Display( ism89 );|
  [ [[Group( [ f1, f2, <identity> of ... ] ) -> Group( [ f2, f2 ] )] => [Group( 
      [ f1, f2, <identity> of ..., <identity> of ... ] ) -> Group( 
      [ f2, f2 ] )]], 
    [[Group( [ f3 ] ) -> Group( <identity> of ... )] => [Group( 
      [ f3 ] ) -> Group( <identity> of ... )]] ]
  
\end{Verbatim}
 

\subsection{\textcolor{Chapter }{IsStemDomain (for crossed modules of groups)}}
\logpage{[ 4, 3, 2 ]}\nobreak
\hyperdef{L}{X7E86DCB083CA5915}{}
{\noindent\textcolor{FuncColor}{$\triangleright$\enspace\texttt{IsStemDomain({\mdseries\slshape X0})\index{IsStemDomain@\texttt{IsStemDomain}!for crossed modules of groups}
\label{IsStemDomain:for crossed modules of groups}
}\hfill{\scriptsize (property)}}\\
\noindent\textcolor{FuncColor}{$\triangleright$\enspace\texttt{IsoclinicStemDomain({\mdseries\slshape X0})\index{IsoclinicStemDomain@\texttt{IsoclinicStemDomain}!for crossed modules of groups}
\label{IsoclinicStemDomain:for crossed modules of groups}
}\hfill{\scriptsize (attribute)}}\\


 A crossed module $\calX$ is a \emph{stem crossed module} if $Z(\calX) \leq [\calX,\calX]$. Every crossed module is isoclinic to a stem crossed module, but distinct
stem crossed modules may be isoclinic. 

 A method for \texttt{IsoclinicStemDomain} has yet to be implemented. }

 

 
\begin{Verbatim}[commandchars=!@|,fontsize=\small,frame=single,label=Example]
  
  !gapprompt@gap>| !gapinput@IsStemDomain(X8);|
  true
  !gapprompt@gap>| !gapinput@IsStemDomain(X9);|
  false
  
\end{Verbatim}
 

\subsection{\textcolor{Chapter }{IsoclinicRank (for crossed modules of groups)}}
\logpage{[ 4, 3, 3 ]}\nobreak
\hyperdef{L}{X820C412679910975}{}
{\noindent\textcolor{FuncColor}{$\triangleright$\enspace\texttt{IsoclinicRank({\mdseries\slshape X0})\index{IsoclinicRank@\texttt{IsoclinicRank}!for crossed modules of groups}
\label{IsoclinicRank:for crossed modules of groups}
}\hfill{\scriptsize (attribute)}}\\
\noindent\textcolor{FuncColor}{$\triangleright$\enspace\texttt{IsoclinicMiddleLength({\mdseries\slshape X0})\index{IsoclinicMiddleLength@\texttt{IsoclinicMiddleLength}!for crossed modules of groups}
\label{IsoclinicMiddleLength:for crossed modules of groups}
}\hfill{\scriptsize (attribute)}}\\


 The formulae in subsection \ref{group-rank} are applied to the crossed module. }

 

 
\begin{Verbatim}[commandchars=!@|,fontsize=\small,frame=single,label=Example]
  
  !gapprompt@gap>| !gapinput@IsoclinicMiddleLength(X8);|
  [ 1, 0 ]
  !gapprompt@gap>| !gapinput@IsoclinicRank(X8);        |
  [ 3, 1 ]
  
\end{Verbatim}
 }

 }

         
\chapter{\textcolor{Chapter }{Whitehead group of a crossed module}}\label{chap-gp2up}
\logpage{[ 5, 0, 0 ]}
\hyperdef{L}{X85CD9A43847AE1B8}{}
{
  \index{up 2d\texttt{\symbol{45}}mapping of 2d\texttt{\symbol{45}}group} 
\section{\textcolor{Chapter }{Derivations and Sections}}\label{sect-der-and-sect}
\logpage{[ 5, 1, 0 ]}
\hyperdef{L}{X7C01AE7783898705}{}
{
  \index{derivation, of crossed module} \index{Whitehead monoid} The Whitehead monoid ${\rm Der}(\calX)$ of $\calX$ was defined in \cite{W2} to be the monoid of all \emph{derivations} from $R$ to $S$, that is the set of all maps $\chi : R \to S$, with \emph{Whitehead product} $\star$ (on the \emph{right}) satisfying: 
\[ {\bf Der\ 1}: \chi(qr) ~=~ (\chi q)^{r} \; (\chi r), \qquad {\bf Der\ 2}:
(\chi_1 \star \chi_2)(r) ~=~ (\chi_2 r)(\chi_1 r)(\chi_2 \partial \chi_1 r). \]
 It easily follows that $\chi 1 = 1$ and $\chi(r^{-1}) = ((\chi r)^{-1})^{r^{-1}}$, and that the zero map is the identity for this composition. \index{regular derivation} Invertible elements in the monoid are called \emph{regular}. \index{Whitehead group} The Whitehead group of $\calX$ is the group of regular derivations in ${\rm Der}(\calX )$. In the next chapter the \emph{actor} of $\calX$ is defined as a crossed module whose source and range are permutation
representations of the Whitehead group and the automorphism group of $\calX$. 

 \index{section, of cat1\texttt{\symbol{45}}group} The construction for cat$^1$\texttt{\symbol{45}}groups equivalent to the derivation of a crossed module is
the \emph{section}. \index{Whitehead multiplication} The monoid of sections of $\calC = (e;t,h : G \to R)$ is the set of group homomorphisms $\xi : R \to G$, with Whitehead multiplication $\star$ (on the \emph{right}) satisfying: 
\[ {\bf Sect\ 1}: t \circ \xi ~=~ {\rm id}_R, \quad {\bf Sect\ 2}: (\xi_1 \star
\xi_2)(r) ~=~ (\xi_1 r)(e h \xi_1 r)^{-1}(\xi_2 h \xi_1 r) ~=~ (\xi_2 h \xi_1
r)(e h \xi_1 r)^{-1}(\xi_1 r). \]
 The embedding $e$ is the identity for this composition, and $h(\xi_1 \star \xi_2) = (h \xi_1)(h \xi_2)$. A section is \emph{regular} when $h \xi$ is an automorphism, and the group of regular sections is isomorphic to the
Whitehead group. 

 If $\epsilon$ denotes the inclusion of $S = \ker\ t$ in $G$ then $\partial = h \epsilon : S \to R$ and 
\[ \xi r ~=~ (e r)(\epsilon \chi r), \quad\mbox{which equals}\quad (r, \chi r)
~\in~ R \ltimes S, \]
 determines a section $\xi$ of $\calC$ in terms of the corresponding derivation $\chi$ of $\calX$, and conversely. 

\subsection{\textcolor{Chapter }{DerivationByImages}}
\logpage{[ 5, 1, 1 ]}\nobreak
\hyperdef{L}{X83EC6F7780F5636E}{}
{\noindent\textcolor{FuncColor}{$\triangleright$\enspace\texttt{DerivationByImages({\mdseries\slshape X0, ims})\index{DerivationByImages@\texttt{DerivationByImages}}
\label{DerivationByImages}
}\hfill{\scriptsize (operation)}}\\
\noindent\textcolor{FuncColor}{$\triangleright$\enspace\texttt{IsDerivation({\mdseries\slshape map})\index{IsDerivation@\texttt{IsDerivation}}
\label{IsDerivation}
}\hfill{\scriptsize (property)}}\\
\noindent\textcolor{FuncColor}{$\triangleright$\enspace\texttt{IsUp2DimensionalMapping({\mdseries\slshape chi})\index{IsUp2DimensionalMapping@\texttt{IsUp2DimensionalMapping}}
\label{IsUp2DimensionalMapping}
}\hfill{\scriptsize (property)}}\\
\noindent\textcolor{FuncColor}{$\triangleright$\enspace\texttt{UpGeneratorImages({\mdseries\slshape chi})\index{UpGeneratorImages@\texttt{UpGeneratorImages}}
\label{UpGeneratorImages}
}\hfill{\scriptsize (attribute)}}\\
\noindent\textcolor{FuncColor}{$\triangleright$\enspace\texttt{UpImagePositions({\mdseries\slshape chi})\index{UpImagePositions@\texttt{UpImagePositions}}
\label{UpImagePositions}
}\hfill{\scriptsize (attribute)}}\\
\noindent\textcolor{FuncColor}{$\triangleright$\enspace\texttt{Object2d({\mdseries\slshape chi})\index{Object2d@\texttt{Object2d}}
\label{Object2d}
}\hfill{\scriptsize (attribute)}}\\
\noindent\textcolor{FuncColor}{$\triangleright$\enspace\texttt{DerivationImage({\mdseries\slshape chi, r})\index{DerivationImage@\texttt{DerivationImage}}
\label{DerivationImage}
}\hfill{\scriptsize (operation)}}\\


 A derivation $\chi$ is stored like a group homomorphisms by specifying the images of the
generating set \texttt{StrongGeneratorsStabChain( StabChain(R) )} of the range $R$. This set of images is stored as the attribute \texttt{UpGeneratorImages} of $\chi$. The function \texttt{IsDerivation} is automatically called to check that this procedure is
well\texttt{\symbol{45}}defined. The attribute \texttt{Object2d}$(\chi)$ returns the underlying crossed module. 

 Images of the remaining elements may be obtained using axiom ${\bf Der\ 1}$. \texttt{UpImagePositions(chi)} is the list of the images under $\chi$ of \texttt{Elements(R)} and \texttt{DerivationImage(chi,r)} returns $\chi r$. 

 In the following example a cat$^1$\texttt{\symbol{45}}group \texttt{C3} and the associated crossed module \texttt{X3} are constructed, where \texttt{X3} is isomorphic to the inclusion of the normal cyclic group \texttt{c3} in the symmetric group \texttt{s3}. The derivation $\chi_1$ maps \texttt{c3} to the identity and the other $3$ elements to $(1,2,3)(4,6,5)$. }

 

 
\begin{Verbatim}[commandchars=!@|,fontsize=\small,frame=single,label=Example]
  
  !gapprompt@gap>| !gapinput@g18 := Group( (1,2,3), (4,5,6), (2,3)(5,6) );;|
  !gapprompt@gap>| !gapinput@SetName( g18, "g18" );|
  !gapprompt@gap>| !gapinput@gen18 := GeneratorsOfGroup( g18 );;|
  !gapprompt@gap>| !gapinput@g1 := gen18[1];;  g2 := gen18[2];;  g3 := gen18[3];;|
  !gapprompt@gap>| !gapinput@s3 := Subgroup( g18, gen18{[2..3]} );;|
  !gapprompt@gap>| !gapinput@SetName( s3, "s3" );;|
  !gapprompt@gap>| !gapinput@t := GroupHomomorphismByImages( g18, s3, gen18, [g2,g2,g3] );;|
  !gapprompt@gap>| !gapinput@h := GroupHomomorphismByImages( g18, s3, gen18, [(),g2,g3] );;|
  !gapprompt@gap>| !gapinput@e := GroupHomomorphismByImages( s3, g18, [g2,g3], [g2,g3] );;|
  !gapprompt@gap>| !gapinput@C3 := Cat1Group( t, h, e );|
  [g18=>s3]
  !gapprompt@gap>| !gapinput@SetName( Kernel(t), "c3" );;|
  !gapprompt@gap>| !gapinput@X3 := XModOfCat1Group( C3 );|
  [c3->s3]
  !gapprompt@gap>| !gapinput@R3 := Range( X3 );;|
  !gapprompt@gap>| !gapinput@StrongGeneratorsStabChain( StabChain( R3 ) );|
  [ (4,5,6), (2,3)(5,6) ]
  !gapprompt@gap>| !gapinput@chi1 := DerivationByImages( X3, [ (), (1,2,3)(4,6,5) ] );|
  DerivationByImages( s3, c3, [ (4,5,6), (2,3)(5,6) ], 
  [ (), (1,2,3)(4,6,5) ] )
  !gapprompt@gap>| !gapinput@[ IsUp2DimensionalMapping( chi1 ), IsDerivation( chi1 ) ];|
  [ true, true ]
  !gapprompt@gap>| !gapinput@Object2d( chi1 );|
  [c3->s3]
  !gapprompt@gap>| !gapinput@UpGeneratorImages( chi1 ); |
  [ (), (1,2,3)(4,6,5) ]
  !gapprompt@gap>| !gapinput@UpImagePositions( chi1 );|
  [ 1, 1, 1, 2, 2, 2 ]
  !gapprompt@gap>| !gapinput@DerivationImage( chi1, (2,3)(4,5) );|
  (1,2,3)(4,6,5)
  
\end{Verbatim}
 

\subsection{\textcolor{Chapter }{PrincipalDerivation}}
\logpage{[ 5, 1, 2 ]}\nobreak
\hyperdef{L}{X7F119F9580C150B1}{}
{\noindent\textcolor{FuncColor}{$\triangleright$\enspace\texttt{PrincipalDerivation({\mdseries\slshape X0, s})\index{PrincipalDerivation@\texttt{PrincipalDerivation}}
\label{PrincipalDerivation}
}\hfill{\scriptsize (operation)}}\\


 The \emph{principal derivation} determined by $s \in S$ is the derivation $\eta_s : R \to S,\; r \mapsto (s^{-1})^rs$. }

 

 
\begin{Verbatim}[commandchars=!@|,fontsize=\small,frame=single,label=Example]
  
  !gapprompt@gap>| !gapinput@eta := PrincipalDerivation( X3, (1,2,3)(4,6,5) );|
  DerivationByImages( s3, c3, [ (4,5,6), (2,3)(5,6) ], [ (), (1,3,2)(4,5,6) ] )
  
\end{Verbatim}
 

\subsection{\textcolor{Chapter }{SectionByHomomorphism}}
\logpage{[ 5, 1, 3 ]}\nobreak
\hyperdef{L}{X7E1C72897CD31D66}{}
{\noindent\textcolor{FuncColor}{$\triangleright$\enspace\texttt{SectionByHomomorphism({\mdseries\slshape C, hom})\index{SectionByHomomorphism@\texttt{SectionByHomomorphism}}
\label{SectionByHomomorphism}
}\hfill{\scriptsize (operation)}}\\
\noindent\textcolor{FuncColor}{$\triangleright$\enspace\texttt{IsSection({\mdseries\slshape xi})\index{IsSection@\texttt{IsSection}}
\label{IsSection}
}\hfill{\scriptsize (property)}}\\
\noindent\textcolor{FuncColor}{$\triangleright$\enspace\texttt{UpHomomorphism({\mdseries\slshape xi})\index{UpHomomorphism@\texttt{UpHomomorphism}}
\label{UpHomomorphism}
}\hfill{\scriptsize (attribute)}}\\
\noindent\textcolor{FuncColor}{$\triangleright$\enspace\texttt{SectionByDerivation({\mdseries\slshape chi})\index{SectionByDerivation@\texttt{SectionByDerivation}}
\label{SectionByDerivation}
}\hfill{\scriptsize (operation)}}\\
\noindent\textcolor{FuncColor}{$\triangleright$\enspace\texttt{DerivationBySection({\mdseries\slshape xi})\index{DerivationBySection@\texttt{DerivationBySection}}
\label{DerivationBySection}
}\hfill{\scriptsize (operation)}}\\


 Sections \emph{are} group homomorphisms but, although they do not require a special
representation, one is provided in the category \texttt{IsUp2DimensionalMapping} having attributes \texttt{Object2d}, \texttt{UpHomomorphism} and \texttt{UpGeneratorsImages}. Operations \texttt{SectionByDerivation} and \texttt{DerivationBySection} convert derivations to sections, and vice\texttt{\symbol{45}}versa, calling \texttt{Cat1GroupOfXMod} (\ref{Cat1GroupOfXMod}) and \texttt{XModOfCat1Group} (\ref{XModOfCat1Group}) automatically. 

 Two strategies for calculating derivations and sections are implemented, see \cite{AW1}. The default method for \texttt{AllDerivations} (\ref{AllDerivations}) is to search for all possible sets of images using a backtracking procedure,
and when all the derivations are found it is not known which are regular. In
early versions of this package, the default method for \texttt{AllSections( {\textless}C{\textgreater} )} was to compute all endomorphisms on the range group \texttt{R} of \texttt{C} as possibilities for the composite $h \xi$. A backtrack method then found possible images for such a section. In the
current version the derivations of the associated crossed module are
calculated, and these are all converted to sections using \texttt{SectionByDerivation}. }

 

 
\begin{Verbatim}[commandchars=!@|,fontsize=\small,frame=single,label=Example]
  
  !gapprompt@gap>| !gapinput@hom2 := GroupHomomorphismByImages( s3, g18, [ (4,5,6), (2,3)(5,6) ], |
  !gapprompt@>| !gapinput@[ (1,3,2)(4,6,5), (1,2)(4,6) ] );;|
  !gapprompt@gap>| !gapinput@xi2 := SectionByHomomorphism( C3, hom2 );                                 |
  SectionByHomomorphism( s3, g18, [ (4,5,6), (2,3)(5,6) ], 
  [ (1,3,2)(4,6,5), (1,2)(4,6) ] )
  !gapprompt@gap>| !gapinput@[ IsUp2DimensionalMapping( xi2 ), IsSection( xi2 ) ];|
  [ true, true ]
  !gapprompt@gap>| !gapinput@Object2d( xi2 );|
  [g18 => s3]
  !gapprompt@gap>| !gapinput@UpHomomorphism( xi2 );         |
  [ (4,5,6), (2,3)(5,6) ] -> [ (1,3,2)(4,6,5), (1,2)(4,6) ]
  !gapprompt@gap>| !gapinput@UpGeneratorImages( xi2 );|
  [ (1,3,2)(4,6,5), (1,2)(4,6) ]
  !gapprompt@gap>| !gapinput@chi2 := DerivationBySection( xi2 );|
  DerivationByImages( s3, c3, [ (4,5,6), (2,3)(5,6) ], 
  [ (1,3,2)(4,5,6), (1,2,3)(4,6,5) ] )
  !gapprompt@gap>| !gapinput@xi1 := SectionByDerivation( chi1 );|
  SectionByHomomorphism( s3, g18, [ (4,5,6), (2,3)(5,6) ], 
  [ (1,2,3), (1,2)(4,6) ] )
  
\end{Verbatim}
 

\subsection{\textcolor{Chapter }{IdentityDerivation}}
\logpage{[ 5, 1, 4 ]}\nobreak
\hyperdef{L}{X87D9F7257DFF0236}{}
{\noindent\textcolor{FuncColor}{$\triangleright$\enspace\texttt{IdentityDerivation({\mdseries\slshape X0})\index{IdentityDerivation@\texttt{IdentityDerivation}}
\label{IdentityDerivation}
}\hfill{\scriptsize (attribute)}}\\
\noindent\textcolor{FuncColor}{$\triangleright$\enspace\texttt{IdentitySection({\mdseries\slshape C0})\index{IdentitySection@\texttt{IdentitySection}}
\label{IdentitySection}
}\hfill{\scriptsize (attribute)}}\\


 The identity derivation maps the range group to the identity subgroup of the
source, while the identity section is just the range embedding considered as a
section. }

 

 
\begin{Verbatim}[commandchars=!@|,fontsize=\small,frame=single,label=Example]
  
  !gapprompt@gap>| !gapinput@IdentityDerivation( X3 ); |
  DerivationByImages( s3, c3, [ (4,5,6), (2,3)(5,6) ], [ (), () ] )
  !gapprompt@gap>| !gapinput@IdentitySection( C3 );     |
  SectionByHomomorphism( s3, g18, [ (4,5,6), (2,3)(5,6) ], 
  [ (4,5,6), (2,3)(5,6) ] )
  
\end{Verbatim}
 

\subsection{\textcolor{Chapter }{WhiteheadProduct}}
\logpage{[ 5, 1, 5 ]}\nobreak
\hyperdef{L}{X7AD6E23F8254F400}{}
{\noindent\textcolor{FuncColor}{$\triangleright$\enspace\texttt{WhiteheadProduct({\mdseries\slshape upi, upj})\index{WhiteheadProduct@\texttt{WhiteheadProduct}}
\label{WhiteheadProduct}
}\hfill{\scriptsize (operation)}}\\
\noindent\textcolor{FuncColor}{$\triangleright$\enspace\texttt{WhiteheadOrder({\mdseries\slshape up})\index{WhiteheadOrder@\texttt{WhiteheadOrder}}
\label{WhiteheadOrder}
}\hfill{\scriptsize (operation)}}\\


 The \texttt{WhiteheadProduct}, defined in section \ref{sect-der-and-sect}, may be applied to two derivations to form $\chi_i \star \chi_j$, or to two sections to form $\xi_i \star \xi_j$. The \texttt{WhiteheadOrder} of a regular derivation $\chi$ is the smallest power of $\chi$, using this product, equal to the \texttt{IdentityDerivation} (\ref{IdentityDerivation}). }

 

 
\begin{Verbatim}[commandchars=!@|,fontsize=\small,frame=single,label=Example]
  
  !gapprompt@gap>| !gapinput@chi12 := WhiteheadProduct( chi1, chi2 );|
  DerivationByImages( s3, c3, [ (4,5,6), (2,3)(5,6) ], [ (1,3,2)(4,5,6), () ] )
  !gapprompt@gap>| !gapinput@xi12 := WhiteheadProduct( xi1, xi2 );|
  SectionByHomomorphism( s3, g18, [ (4,5,6), (2,3)(5,6) ], 
  [ (1,3,2)(4,6,5), (2,3)(5,6) ] )
  !gapprompt@gap>| !gapinput@xi12 = SectionByDerivation( chi12 ); |
  true 
  !gapprompt@gap>| !gapinput@[ WhiteheadOrder( chi2 ), WhiteheadOrder( xi2 ) ];|
  [ 2, 2 ]
  
\end{Verbatim}
 }

 
\section{\textcolor{Chapter }{Whitehead Monoids and Groups}}\label{sect-whitehead-monoids}
\logpage{[ 5, 2, 0 ]}
\hyperdef{L}{X861A52407D3C627D}{}
{
  As mentioned at the beginning of this chapter, the Whitehead monoid ${\rm Der}(\calX)$ of $\calX$ is the monoid of all derivations from $R$ to $S$. Monoids of derivations have representation \texttt{IsMonoidOfUp2DimensionalMappingsObj}. \index{IsMonoidOfUp2DimensionalMappingsObj} Multiplication tables for Whitehead monoids enable the construction of
transformation representations. 

\subsection{\textcolor{Chapter }{AllDerivations}}
\logpage{[ 5, 2, 1 ]}\nobreak
\hyperdef{L}{X788884E48534F7CB}{}
{\noindent\textcolor{FuncColor}{$\triangleright$\enspace\texttt{AllDerivations({\mdseries\slshape X0})\index{AllDerivations@\texttt{AllDerivations}}
\label{AllDerivations}
}\hfill{\scriptsize (attribute)}}\\
\noindent\textcolor{FuncColor}{$\triangleright$\enspace\texttt{ImagesList({\mdseries\slshape obj})\index{ImagesList@\texttt{ImagesList}}
\label{ImagesList}
}\hfill{\scriptsize (attribute)}}\\
\noindent\textcolor{FuncColor}{$\triangleright$\enspace\texttt{DerivationClass({\mdseries\slshape mon})\index{DerivationClass@\texttt{DerivationClass}}
\label{DerivationClass}
}\hfill{\scriptsize (attribute)}}\\
\noindent\textcolor{FuncColor}{$\triangleright$\enspace\texttt{ImagesTable({\mdseries\slshape obj})\index{ImagesTable@\texttt{ImagesTable}}
\label{ImagesTable}
}\hfill{\scriptsize (attribute)}}\\


 Using our example \texttt{X3} we find that there are just nine derivations. Here \texttt{AllDerivations} returns a collection of up mappings with attributes: 
\begin{itemize}
\item  \texttt{Object2d} \texttt{\symbol{45}} the crossed modules $\calX$; 
\item  \texttt{ImagesList} \texttt{\symbol{45}} a list, for each derivation, of the images of the
generators of the range group; 
\item  \texttt{DerivationClass} \texttt{\symbol{45}} the string "all"; other classes include "regular" and
"principal"; 
\item  \texttt{ImagesTable} \texttt{\symbol{45}} this is a table whose $[i,j]$\texttt{\symbol{45}}th entry is the position in the list of elements of $S$ of the image under the $i$\texttt{\symbol{45}}th derivation of the $j$\texttt{\symbol{45}}th element of $R$. 
\end{itemize}
 }

 

 
\begin{Verbatim}[commandchars=!@|,fontsize=\small,frame=single,label=Example]
  
  !gapprompt@gap>| !gapinput@all3 := AllDerivations( X3 );|
  monoid of derivations with images list:
  [ (), () ]
  [ (), (1,3,2)(4,5,6) ]
  [ (), (1,2,3)(4,6,5) ]
  [ (1,3,2)(4,5,6), () ]
  [ (1,3,2)(4,5,6), (1,3,2)(4,5,6) ]
  [ (1,3,2)(4,5,6), (1,2,3)(4,6,5) ]
  [ (1,2,3)(4,6,5), () ]
  [ (1,2,3)(4,6,5), (1,3,2)(4,5,6) ]
  [ (1,2,3)(4,6,5), (1,2,3)(4,6,5) ]
  !gapprompt@gap>| !gapinput@DerivationClass( all3 );|
  "all"
  !gapprompt@gap>| !gapinput@Perform( ImagesTable( all3 ), Display );|
  [ 1, 1, 1, 1, 1, 1 ]
  [ 1, 1, 1, 3, 3, 3 ]
  [ 1, 1, 1, 2, 2, 2 ]
  [ 1, 3, 2, 1, 3, 2 ]
  [ 1, 3, 2, 3, 2, 1 ]
  [ 1, 3, 2, 2, 1, 3 ]
  [ 1, 2, 3, 1, 2, 3 ]
  [ 1, 2, 3, 3, 1, 2 ]
  [ 1, 2, 3, 2, 3, 1 ]
  
\end{Verbatim}
 

\subsection{\textcolor{Chapter }{WhiteheadMonoidTable}}
\logpage{[ 5, 2, 2 ]}\nobreak
\hyperdef{L}{X7CB1614E7EC58A84}{}
{\noindent\textcolor{FuncColor}{$\triangleright$\enspace\texttt{WhiteheadMonoidTable({\mdseries\slshape X0})\index{WhiteheadMonoidTable@\texttt{WhiteheadMonoidTable}}
\label{WhiteheadMonoidTable}
}\hfill{\scriptsize (attribute)}}\\
\noindent\textcolor{FuncColor}{$\triangleright$\enspace\texttt{WhiteheadTransformationMonoid({\mdseries\slshape X0})\index{WhiteheadTransformationMonoid@\texttt{WhiteheadTransformationMonoid}}
\label{WhiteheadTransformationMonoid}
}\hfill{\scriptsize (attribute)}}\\


 The \texttt{WhiteheadMonoidTable} of $\calX$ is the multiplication table whose $[i,j]$\texttt{\symbol{45}}th entry is the position $k$ in the list of derivations of the Whitehead product $\chi_i*\chi_j = \chi_k$. 

 Using the rows of the table as transformations, we may construct the \texttt{WhiteheadTransformationMonoid} of $\calX$. }

 

 
\begin{Verbatim}[commandchars=!@|,fontsize=\small,frame=single,label=Example]
  
  !gapprompt@gap>| !gapinput@wmt3 := WhiteheadMonoidTable( X3 );; |
  !gapprompt@gap>| !gapinput@Perform( wmt3, Display );|
  [ 1, 2, 3, 4, 5, 6, 7, 8, 9 ]
  [ 2, 3, 1, 5, 6, 4, 8, 9, 7 ]
  [ 3, 1, 2, 6, 4, 5, 9, 7, 8 ]
  [ 4, 6, 5, 1, 3, 2, 7, 9, 8 ]
  [ 5, 4, 6, 2, 1, 3, 8, 7, 9 ]
  [ 6, 5, 4, 3, 2, 1, 9, 8, 7 ]
  [ 7, 7, 7, 7, 7, 7, 7, 7, 7 ]
  [ 8, 8, 8, 8, 8, 8, 8, 8, 8 ]
  [ 9, 9, 9, 9, 9, 9, 9, 9, 9 ]
  !gapprompt@gap>| !gapinput@wtm3 := WhiteheadTransformationMonoid( X3 );|
  <transformation monoid of degree 9 with 3 generators>
  !gapprompt@gap>| !gapinput@GeneratorsOfMonoid( wtm3 ); |
  [ Transformation( [ 2, 3, 1, 5, 6, 4, 8, 9, 7 ] ), 
    Transformation( [ 4, 6, 5, 1, 3, 2, 7, 9, 8 ] ), 
    Transformation( [ 7, 7, 7, 7, 7, 7, 7, 7, 7 ] ) ]
  
\end{Verbatim}
 

\subsection{\textcolor{Chapter }{WhiteheadPermGroup}}
\logpage{[ 5, 2, 3 ]}\nobreak
\hyperdef{L}{X78680F197B48BFA9}{}
{\noindent\textcolor{FuncColor}{$\triangleright$\enspace\texttt{WhiteheadPermGroup({\mdseries\slshape X0})\index{WhiteheadPermGroup@\texttt{WhiteheadPermGroup}}
\label{WhiteheadPermGroup}
}\hfill{\scriptsize (attribute)}}\\
\noindent\textcolor{FuncColor}{$\triangleright$\enspace\texttt{RegularDerivations({\mdseries\slshape X0})\index{RegularDerivations@\texttt{RegularDerivations}}
\label{RegularDerivations}
}\hfill{\scriptsize (attribute)}}\\
\noindent\textcolor{FuncColor}{$\triangleright$\enspace\texttt{WhiteheadGroupTable({\mdseries\slshape X0})\index{WhiteheadGroupTable@\texttt{WhiteheadGroupTable}}
\label{WhiteheadGroupTable}
}\hfill{\scriptsize (attribute)}}\\
\noindent\textcolor{FuncColor}{$\triangleright$\enspace\texttt{IsWhiteheadPermGroup({\mdseries\slshape G})\index{IsWhiteheadPermGroup@\texttt{IsWhiteheadPermGroup}}
\label{IsWhiteheadPermGroup}
}\hfill{\scriptsize (property)}}\\
\noindent\textcolor{FuncColor}{$\triangleright$\enspace\texttt{WhiteheadRegularGroup({\mdseries\slshape X0})\index{WhiteheadRegularGroup@\texttt{WhiteheadRegularGroup}}
\label{WhiteheadRegularGroup}
}\hfill{\scriptsize (attribute)}}\\
\noindent\textcolor{FuncColor}{$\triangleright$\enspace\texttt{WhiteheadGroupIsomorphism({\mdseries\slshape X0})\index{WhiteheadGroupIsomorphism@\texttt{WhiteheadGroupIsomorphism}}
\label{WhiteheadGroupIsomorphism}
}\hfill{\scriptsize (attribute)}}\\


 \texttt{RegularDerivations} are those derivations which are invertible in the monoid. The multiplication
table for a Whitehead group \texttt{\symbol{45}} a subtable of the Whitehead
monoid table \texttt{\symbol{45}} enables the construction of a permutation
representation \texttt{WhiteheadPermGroup} $W\calX$ of $\calX$. This group satisfies the property \texttt{IsWhiteheadPermGroup}. and $\calX$ is its attribute \texttt{Object2d}. 

 Of the nine derivations of \texttt{X3} just six are regular. The associated group is isomorphic to the symmetric
group \texttt{s3}. 

 Taking the rows of the \texttt{WhiteheadGroupTable} as permutations, we may construct the \texttt{WhiteheadRegularGroup} of $\calX$. Then, seeking a \texttt{SmallerDegreePermutationRepresentation}, we obtain the \texttt{WhiteheadGroupIsomorphism} whose image is the \texttt{WhiteheadPermGroup} of $\calX$. }

 

 
\begin{Verbatim}[commandchars=!@|,fontsize=\small,frame=single,label=Example]
  
  !gapprompt@gap>| !gapinput@reg3 := RegularDerivations( X3 );|
  monoid of derivations with images list:
  [ (), () ]
  [ (), (1,3,2)(4,5,6) ]
  [ (), (1,2,3)(4,6,5) ]
  [ (1,3,2)(4,5,6), () ]
  [ (1,3,2)(4,5,6), (1,3,2)(4,5,6) ]
  [ (1,3,2)(4,5,6), (1,2,3)(4,6,5) ]
  !gapprompt@gap>| !gapinput@DerivationClass( reg3 );|
  "regular"
  !gapprompt@gap>| !gapinput@wgt3 := WhiteheadGroupTable( X3 );; |
  !gapprompt@gap>| !gapinput@Perform( wgt3, Display );|
  [ [ 1, 2, 3, 4, 5, 6 ],
    [ 2, 3, 1, 5, 6, 4 ],
    [ 3, 1, 2, 6, 4, 5 ],
    [ 4, 6, 5, 1, 3, 2 ],
    [ 5, 4, 6, 2, 1, 3 ],
    [ 6, 5, 4, 3, 2, 1 ] ]
  !gapprompt@gap>| !gapinput@wpg3 := WhiteheadPermGroup( X3 );|
  Group([ (1,2,3), (1,2) ])
  !gapprompt@gap>| !gapinput@IsWhiteheadPermGroup( wpg3 );|
  true
  !gapprompt@gap>| !gapinput@Object2d( wpg3 );|
  [c3->s3]
  !gapprompt@gap>| !gapinput@WhiteheadRegularGroup( X3 );|
  Group([ (1,2,3)(4,5,6), (1,4)(2,6)(3,5) ])
  !gapprompt@gap>| !gapinput@MappingGeneratorsImages( WhiteheadGroupIsomorphism( X3 ) );|
  [ [ (1,2,3)(4,5,6), (1,4)(2,6)(3,5) ], [ (1,2,3), (1,2) ] ]
  
\end{Verbatim}
 

\subsection{\textcolor{Chapter }{WhiteheadGroupElement}}
\logpage{[ 5, 2, 4 ]}\nobreak
\hyperdef{L}{X85C192717BDE886A}{}
{\noindent\textcolor{FuncColor}{$\triangleright$\enspace\texttt{WhiteheadGroupElement({\mdseries\slshape chi})\index{WhiteheadGroupElement@\texttt{WhiteheadGroupElement}}
\label{WhiteheadGroupElement}
}\hfill{\scriptsize (operation)}}\\


 This operation returns the image of a derivation $\chi$ in the \texttt{WhiteheadPermGroup}. 

 }

 

 
\begin{Verbatim}[commandchars=!@|,fontsize=\small,frame=single,label=Example]
  
  !gapprompt@gap>| !gapinput@chi3 := DerivationByImages( X3, [ (1,3,2)(4,5,6), (1,2,3)(4,6,5) ] );;|
  !gapprompt@gap>| !gapinput@WhiteheadGroupElement( chi3 );|
  (2,3)
  
\end{Verbatim}
 

\subsection{\textcolor{Chapter }{PrincipalDerivations}}
\logpage{[ 5, 2, 5 ]}\nobreak
\hyperdef{L}{X8452E45878691CD3}{}
{\noindent\textcolor{FuncColor}{$\triangleright$\enspace\texttt{PrincipalDerivations({\mdseries\slshape X0})\index{PrincipalDerivations@\texttt{PrincipalDerivations}}
\label{PrincipalDerivations}
}\hfill{\scriptsize (attribute)}}\\
\noindent\textcolor{FuncColor}{$\triangleright$\enspace\texttt{PrincipalDerivationSubgroup({\mdseries\slshape X0})\index{PrincipalDerivationSubgroup@\texttt{PrincipalDerivationSubgroup}}
\label{PrincipalDerivationSubgroup}
}\hfill{\scriptsize (attribute)}}\\
\noindent\textcolor{FuncColor}{$\triangleright$\enspace\texttt{WhiteheadHomomorphism({\mdseries\slshape X0})\index{WhiteheadHomomorphism@\texttt{WhiteheadHomomorphism}}
\label{WhiteheadHomomorphism}
}\hfill{\scriptsize (attribute)}}\\


 The principal derivations form a subgroup of the Whitehead group. The \texttt{PrincipalDerivationSubgroup} is the corresponding subgroup of the \texttt{WhiteheadPermGroup}. 

 The Whitehead homomorphism $\eta : S \to W\calX,~ s \mapsto \eta_s$ for $\calX$ maps the source group of $\calX$ to the Whitehead group of $\calX$. }

 

 
\begin{Verbatim}[commandchars=!@|,fontsize=\small,frame=single,label=Example]
  
  !gapprompt@gap>| !gapinput@PDX3 := PrincipalDerivations( X3 );|
  monoid of derivations with images list:
  [ (), () ]
  [ (), (1,3,2)(4,5,6) ]
  [ (), (1,2,3)(4,6,5) ]
  !gapprompt@gap>| !gapinput@PDSX3 := PrincipalDerivationSubgroup( X3 );|
  Group([ (1,2,3) ])
  !gapprompt@gap>| !gapinput@Whom3 := WhiteheadHomomorphism( X3 );|
  [ (1,2,3)(4,6,5) ] -> [ (1,2,3) ]
  
\end{Verbatim}
 \textsc{Exercise:}$~$ Use the two crossed module axioms to show that $\eta_{s_1} \star \eta_{s_2} = \eta_{s_1s_2}$. }

 
\section{\textcolor{Chapter }{Endomorphisms determined by a derivation}}\label{sect-derivation-endo}
\logpage{[ 5, 3, 0 ]}
\hyperdef{L}{X7E53AF1884B2D03D}{}
{
  

\subsection{\textcolor{Chapter }{SourceEndomorphism}}
\logpage{[ 5, 3, 1 ]}\nobreak
\hyperdef{L}{X86AF32CA84217C46}{}
{\noindent\textcolor{FuncColor}{$\triangleright$\enspace\texttt{SourceEndomorphism({\mdseries\slshape chi})\index{SourceEndomorphism@\texttt{SourceEndomorphism}}
\label{SourceEndomorphism}
}\hfill{\scriptsize (operation)}}\\


 A derivation $\chi$ of $\calX = (\partial : S \to R)$ determines an endomorphism $\sigma_{\chi} : S \to S,~ s \mapsto s(\chi \partial s)$. We may verify that $\sigma_{\chi}$ is a homomorphism by: 
\[ \sigma_{\chi}(s_1s_2) ~=~ s_1s_2\chi((\partial s_1)(\partial s_2)) ~=~
s_1s_2(\chi\partial s_1)^{\partial s_2}(\chi\partial s_2) ~=~ s_1(\chi\partial
s_1)s_2(\chi\partial s_2) ~=~ (\sigma_{\chi}s_1)(\sigma_{\chi}s_2). \]
 }

 

 
\begin{Verbatim}[commandchars=!@|,fontsize=\small,frame=single,label=Example]
  
  !gapprompt@gap>| !gapinput@sigma2 := SourceEndomorphism( chi2 );|
  [ (1,2,3)(4,6,5) ] -> [ (1,3,2)(4,5,6) ]
  
\end{Verbatim}
 

\subsection{\textcolor{Chapter }{RangeEndomorphism}}
\logpage{[ 5, 3, 2 ]}\nobreak
\hyperdef{L}{X84463DE2872AA709}{}
{\noindent\textcolor{FuncColor}{$\triangleright$\enspace\texttt{RangeEndomorphism({\mdseries\slshape chi})\index{RangeEndomorphism@\texttt{RangeEndomorphism}}
\label{RangeEndomorphism}
}\hfill{\scriptsize (operation)}}\\


 A derivation $\chi$ of $\calX = (\partial : S \to R)$ determines an endomorphism $\rho_{\chi} : R \to R,~ r \mapsto r(\partial \chi r)$. }

 

 
\begin{Verbatim}[commandchars=!@|,fontsize=\small,frame=single,label=Example]
  
  !gapprompt@gap>| !gapinput@rho2 := RangeEndomorphism( chi2 );|
  [ (4,5,6), (2,3)(5,6) ] -> [ (4,6,5), (2,3)(4,6) ]
  
\end{Verbatim}
 

\subsection{\textcolor{Chapter }{Object2dEndomorphism}}
\logpage{[ 5, 3, 3 ]}\nobreak
\hyperdef{L}{X7E27AE6478566A94}{}
{\noindent\textcolor{FuncColor}{$\triangleright$\enspace\texttt{Object2dEndomorphism({\mdseries\slshape chi})\index{Object2dEndomorphism@\texttt{Object2dEndomorphism}}
\label{Object2dEndomorphism}
}\hfill{\scriptsize (operation)}}\\


 A derivation $\chi$ of $\calX = (\partial : S \to R)$ determines an endomorphism $\alpha_{\chi} : \calX \to \calX$ whose source and range endomorphisms are given by the previous two operations. }

 

 
\begin{Verbatim}[commandchars=!@|,fontsize=\small,frame=single,label=Example]
  
  !gapprompt@gap>| !gapinput@alpha2 := Object2dEndomorphism( chi2 );;|
  !gapprompt@gap>| !gapinput@Display( alpha2 );|
  Morphism of crossed modules :- 
  : Source = [c3->s3] with generating sets:
    [ (1,2,3)(4,6,5) ]
    [ (4,5,6), (2,3)(5,6) ]
  : Range = Source
  : Source Homomorphism maps source generators to:
    [ (1,3,2)(4,5,6) ]
  : Range Homomorphism maps range generators to:
    [ (4,6,5), (2,3)(4,6) ]
  
\end{Verbatim}
 }

 
\section{\textcolor{Chapter }{Whitehead groups for cat$^1$\texttt{\symbol{45}}groups}}\label{sect-whitehead-cat1}
\logpage{[ 5, 4, 0 ]}
\hyperdef{L}{X820501BF83D1D6D7}{}
{
  

\subsection{\textcolor{Chapter }{AllSections}}
\logpage{[ 5, 4, 1 ]}\nobreak
\hyperdef{L}{X813CAA17855172E4}{}
{\noindent\textcolor{FuncColor}{$\triangleright$\enspace\texttt{AllSections({\mdseries\slshape C0})\index{AllSections@\texttt{AllSections}}
\label{AllSections}
}\hfill{\scriptsize (attribute)}}\\
\noindent\textcolor{FuncColor}{$\triangleright$\enspace\texttt{RegularSections({\mdseries\slshape C0})\index{RegularSections@\texttt{RegularSections}}
\label{RegularSections}
}\hfill{\scriptsize (attribute)}}\\


 These operations are currently obtained by running the equivalent operation
for derivations and then converting the result to sections. }

 

 
\begin{Verbatim}[commandchars=!@|,fontsize=\small,frame=single,label=Example]
  
  !gapprompt@gap>| !gapinput@AllSections( C3 );|
  monoid of sections with images list:
  [ (4,5,6), (2,3)(5,6) ]
  [ (4,5,6), (1,3)(4,5) ]
  [ (4,5,6), (1,2)(4,6) ]
  [ (1,3,2)(4,6,5), (2,3)(5,6) ]
  [ (1,3,2)(4,6,5), (1,3)(4,5) ]
  [ (1,3,2)(4,6,5), (1,2)(4,6) ]
  [ (1,2,3), (2,3)(5,6) ]
  [ (1,2,3), (1,3)(4,5) ]
  [ (1,2,3), (1,2)(4,6) ]
  !gapprompt@gap>| !gapinput@RegularSections( C3 );         |
  monoid of sections with images list:
  [ (4,5,6), (2,3)(5,6) ]
  [ (4,5,6), (1,3)(4,5) ]
  [ (4,5,6), (1,2)(4,6) ]
  [ (1,3,2)(4,6,5), (2,3)(5,6) ]
  [ (1,3,2)(4,6,5), (1,3)(4,5) ]
  [ (1,3,2)(4,6,5), (1,2)(4,6) ]
  
\end{Verbatim}
 }

 }

         
\chapter{\textcolor{Chapter }{Actors of 2d\texttt{\symbol{45}}groups}}\label{chap-gp2act}
\logpage{[ 6, 0, 0 ]}
\hyperdef{L}{X84C872BB7F1E5F25}{}
{
  
\section{\textcolor{Chapter }{Actor of a crossed module}}\logpage{[ 6, 1, 0 ]}
\hyperdef{L}{X7B853602873FC7AB}{}
{
 \index{actor} The \emph{actor} of $\calX$ is a crossed module $\Act(\calX) = (\Delta : \calW(\calX) \to \Aut(\calX))$ which was shown by Lue and Norrie, in \cite{N2} and \cite{N1} to give the automorphism object of a crossed module $\calX$. In this implementation, the source of the actor is a permutation
representation $W$ of the Whitehead group of regular derivations, and the range of the actor is a
permutation representation $A$ of the automorphism group $\Aut(\calX)$ of $\calX$. 

\subsection{\textcolor{Chapter }{AutomorphismPermGroup}}
\logpage{[ 6, 1, 1 ]}\nobreak
\hyperdef{L}{X80F121357F06E72D}{}
{\noindent\textcolor{FuncColor}{$\triangleright$\enspace\texttt{AutomorphismPermGroup({\mdseries\slshape 2d\texttt{\symbol{45}}gp})\index{AutomorphismPermGroup@\texttt{AutomorphismPermGroup}}
\label{AutomorphismPermGroup}
}\hfill{\scriptsize (attribute)}}\\
\noindent\textcolor{FuncColor}{$\triangleright$\enspace\texttt{GeneratingAutomorphisms({\mdseries\slshape 2d\texttt{\symbol{45}}gp})\index{GeneratingAutomorphisms@\texttt{GeneratingAutomorphisms}}
\label{GeneratingAutomorphisms}
}\hfill{\scriptsize (attribute)}}\\
\noindent\textcolor{FuncColor}{$\triangleright$\enspace\texttt{PermAutomorphismAs2dGroupMorphism({\mdseries\slshape 2d\texttt{\symbol{45}}gp, perm})\index{PermAutomorphismAs2dGroupMorphism@\texttt{PermAutomorphismAs2dGroupMorphism}}
\label{PermAutomorphismAs2dGroupMorphism}
}\hfill{\scriptsize (operation)}}\\


 The automorphisms $( \sigma, \rho )$ of $\calX$ form a group $\Aut(\calX)$ of crossed module isomorphisms. The function \texttt{AutomorphismPermGroup} finds a set of \texttt{GeneratingAutomorphisms} for $\Aut(\calX)$, and then constructs a permutation representation of this group, which is
used as the range of the actor crossed module of $\calX$. The individual automorphisms can be constructed from the permutation group
using the function \texttt{PermAutomorphismAs2dGroupMorphism}. The example below uses the crossed module \texttt{X3=[c3\texttt{\symbol{45}}{\textgreater}s3]} constructed in section \ref{DerivationByImages}. }

 

 
\begin{Verbatim}[commandchars=!@|,fontsize=\small,frame=single,label=Example]
  
  !gapprompt@gap>| !gapinput@APX3 := AutomorphismPermGroup( X3 );|
  Group([ (5,7,6), (1,2)(3,4)(6,7) ])
  !gapprompt@gap>| !gapinput@Size( APX3 );|
  6
  !gapprompt@gap>| !gapinput@genX3 := GeneratingAutomorphisms( X3 );    |
  [ [[c3->s3] => [c3->s3]], [[c3->s3] => [c3->s3]] ]
  !gapprompt@gap>| !gapinput@e6 := Elements( APX3 )[6];|
  (1,2)(3,4)(5,7)
  !gapprompt@gap>| !gapinput@m6 := PermAutomorphismAs2dGroupMorphism( X3, e6 );;|
  !gapprompt@gap>| !gapinput@Display( m6 );|
  Morphism of crossed modules :- 
  : Source = [c3->s3] with generating sets:
    [ (1,2,3)(4,6,5) ]
    [ (4,5,6), (2,3)(5,6) ]
  : Range = Source
  : Source Homomorphism maps source generators to:
    [ (1,3,2)(4,5,6) ]
  : Range Homomorphism maps range generators to:
    [ (4,6,5), (2,3)(4,5) ]
  
\end{Verbatim}
 The automorphisms $( \gamma, \rho )$ of a cat$^1$\texttt{\symbol{45}}group $\calC$ form a group $\Aut(\calC)$ of cat$^1$\texttt{\symbol{45}}group isomorphisms. The function \texttt{AutomorphismPermGroup} (\ref{AutomorphismPermGroup}) constructs a permutation representation of this group, which is used as the
range of the actor crossed module of $\calC$. The individual automorphisms can be constructed from the permutation group
using the function \texttt{PermAutomorphismAs2dGroupMorphism} (\ref{PermAutomorphismAs2dGroupMorphism}). The example below uses the cat$^1$\texttt{\symbol{45}}group \texttt{C3} constructed in section \texttt{DerivationByImages} (\ref{DerivationByImages}). 

 
\begin{Verbatim}[commandchars=!@|,fontsize=\small,frame=single,label=Example]
  
  !gapprompt@gap>| !gapinput@APC3 := AutomorphismPermGroup( C3 );|
  Group([ (1,3,2)(4,6,5)(7,9,8)(12,13,14), (2,3)(4,7)(5,9)(6,8)(10,11)(13,14) ])
  !gapprompt@gap>| !gapinput@IdGroup( APC3 );|
  [ 6, 1 ]
  !gapprompt@gap>| !gapinput@a := GeneratorsOfGroup( APC3 )[1];;|
  !gapprompt@gap>| !gapinput@m := PermAutomorphismAs2dGroupMorphism( C3, a );|
  [[g18 => s3] => [g18 => s3]]
  
\end{Verbatim}
 

\subsection{\textcolor{Chapter }{InnerAutomorphismPermGroup}}
\logpage{[ 6, 1, 2 ]}\nobreak
\hyperdef{L}{X7CA87CA382CF5967}{}
{\noindent\textcolor{FuncColor}{$\triangleright$\enspace\texttt{InnerAutomorphismPermGroup({\mdseries\slshape 2d\texttt{\symbol{45}}gp})\index{InnerAutomorphismPermGroup@\texttt{InnerAutomorphismPermGroup}}
\label{InnerAutomorphismPermGroup}
}\hfill{\scriptsize (attribute)}}\\


 The crosssed module inner automorphisms $\alpha_r = (\sigma_r,\rho_r)$ for $r \in R$ described in \texttt{InnerAutomorphismXMod} (\ref{InnerAutomorphismXMod}) form a subgroup of the automorphism group, and so the \texttt{InnerAutomorphismPermGroup} of $\calX$ is a subgroup of the \texttt{AutomorphismPermGroup}, the image of $\alpha : R \to A,~ r \mapsto \alpha_r$. 

 All the automorphisms of \texttt{X3} are inner, so we take as an alternative example the crossed module \texttt{X12} constructed in \texttt{XModByCentralExtension} (\ref{XModByCentralExtension}). 

 A similar argument applies to cat$^1$\texttt{\symbol{45}}groups. }

 

 
\begin{Verbatim}[commandchars=!@|,fontsize=\small,frame=single,label=Example]
  
  !gapprompt@gap>| !gapinput@InnerAutomorphismPermGroup( X3 ) = APX3;  |
  true
  !gapprompt@gap>| !gapinput@Display( X12 );|
  Crossed module X12 :- 
  : Source group d12 has generators:
    [ (1,2,3,4,5,6), (2,6)(3,5) ]
  : Range group s3 has generators:
    [ (7,8,9), (8,9) ]
  : Boundary homomorphism maps source generators to:
    [ (7,8,9), (8,9) ]
  : Action homomorphism maps range generators to automorphisms:
    (7,8,9) --> { source gens --> [ (1,2,3,4,5,6), (1,3)(4,6) ] }
    (8,9) --> { source gens --> [ (1,6,5,4,3,2), (2,6)(3,5) ] }
    These 2 automorphisms generate the group of automorphisms.
  !gapprompt@gap>| !gapinput@APX12 := AutomorphismPermGroup( X12 );     |
  Group([ (1,4,6)(2,5,8)(9,10,11), (1,5)(2,6)(4,8), (1,6)(2,5)(3,7)(9,11) ])
  !gapprompt@gap>| !gapinput@IAPX12 := InnerAutomorphismPermGroup( X12 );        |
  Group([ (1,4,6)(2,5,8)(9,10,11), (2,8)(3,7)(4,6)(10,11) ])
  !gapprompt@gap>| !gapinput@[ StructureDescription( APX12 ), StructureDescription( IAPX12 ) ];|
  [ "D12", "S3" ]
  !gapprompt@gap>| !gapinput@InnerAutomorphismPermGroup( C3 );|
  Group([ (1,3,2)(4,6,5)(7,9,8)(12,13,14), (2,3)(4,7)(5,9)(6,8)(10,11)(13,14) ])
  
  
\end{Verbatim}
 

\subsection{\textcolor{Chapter }{WhiteheadXMod}}
\logpage{[ 6, 1, 3 ]}\nobreak
\hyperdef{L}{X790EBC7C7D320C03}{}
{\noindent\textcolor{FuncColor}{$\triangleright$\enspace\texttt{WhiteheadXMod({\mdseries\slshape xmod})\index{WhiteheadXMod@\texttt{WhiteheadXMod}}
\label{WhiteheadXMod}
}\hfill{\scriptsize (attribute)}}\\
\noindent\textcolor{FuncColor}{$\triangleright$\enspace\texttt{LueXMod({\mdseries\slshape xmod})\index{LueXMod@\texttt{LueXMod}}
\label{LueXMod}
}\hfill{\scriptsize (attribute)}}\\
\noindent\textcolor{FuncColor}{$\triangleright$\enspace\texttt{NorrieXMod({\mdseries\slshape xmod})\index{NorrieXMod@\texttt{NorrieXMod}}
\label{NorrieXMod}
}\hfill{\scriptsize (attribute)}}\\
\noindent\textcolor{FuncColor}{$\triangleright$\enspace\texttt{ActorXMod({\mdseries\slshape xmod})\index{ActorXMod@\texttt{ActorXMod}}
\label{ActorXMod}
}\hfill{\scriptsize (attribute)}}\\


 An automorphism $( \sigma, \rho )$ of $\calX$ acts on the Whitehead monoid by $\chi^{(\sigma,\rho)} = \sigma \circ \chi \circ \rho^{-1}$, and this determines the action for the actor. In fact the four groups $S, W, R, A$, the homomorphisms between them, and the various actions, give five crossed
modules forming a \emph{crossed square} (see \texttt{ActorCrossedSquare} (\ref{ActorCrossedSquare})). \index{crossed square} 
\begin{itemize}
\item  $\calW(\calX) = (\eta : S \to W),~$ the Whitehead crossed module of $\calX$, at the top, 
\item  $\calX = (\partial : S \to R),~$ the initial crossed module, on the left, 
\item  $\Act(\calX) = ( \Delta : W \to A),~$ the actor crossed module of $\calX$, on the right, 
\item  $\calN(X) = (\alpha : R \to A),~$ the Norrie crossed module of $\calX$, on the bottom, and 
\item  $\calL(\calX) = (\Delta\circ\eta = \alpha\circ\partial : S \to A),~$ the Lue crossed module of $\calX$, along the top\texttt{\symbol{45}}left to bottom\texttt{\symbol{45}}right
diagonal. 
\end{itemize}
 }

 

 
\begin{Verbatim}[commandchars=!@|,fontsize=\small,frame=single,label=Example]
  
  !gapprompt@gap>| !gapinput@WGX3 := WhiteheadPermGroup( X3 );|
  Group([ (1,2,3), (1,2) ])
  !gapprompt@gap>| !gapinput@APX3 := AutomorphismPermGroup( X3 );|
  Group([ (5,7,6), (1,2)(3,4)(6,7) ])
  !gapprompt@gap>| !gapinput@WX3 := WhiteheadXMod( X3 );; |
  !gapprompt@gap>| !gapinput@Display( WX3 );|
  Crossed module Whitehead[c3->s3] :- 
  : Source group has generators:
    [ (1,2,3)(4,6,5) ]
  : Range group has generators:
    [ (1,2,3), (1,2) ]
  : Boundary homomorphism maps source generators to:
    [ (1,2,3) ]
  : Action homomorphism maps range generators to automorphisms:
    (1,2,3) --> { source gens --> [ (1,2,3)(4,6,5) ] }
    (1,2) --> { source gens --> [ (1,3,2)(4,5,6) ] }
    These 2 automorphisms generate the group of automorphisms.
  !gapprompt@gap>| !gapinput@LX3 := LueXMod( X3 );;|
  !gapprompt@gap>| !gapinput@Display( LX3 );|
  Crossed module Lue[c3->s3] :- 
  : Source group has generators:
    [ (1,2,3)(4,6,5) ]
  : Range group has generators:
    [ (5,7,6), (1,2)(3,4)(6,7) ]
  : Boundary homomorphism maps source generators to:
    [ (5,7,6) ]
  : Action homomorphism maps range generators to automorphisms:
    (5,7,6) --> { source gens --> [ (1,2,3)(4,6,5) ] }
    (1,2)(3,4)(6,7) --> { source gens --> [ (1,3,2)(4,5,6) ] }
    These 2 automorphisms generate the group of automorphisms.
  !gapprompt@gap>| !gapinput@NX3 := NorrieXMod( X3 );; |
  !gapprompt@gap>| !gapinput@Display( NX3 );|
  Crossed module Norrie[c3->s3] :- 
  : Source group has generators:
    [ (4,5,6), (2,3)(5,6) ]
  : Range group has generators:
    [ (5,7,6), (1,2)(3,4)(6,7) ]
  : Boundary homomorphism maps source generators to:
    [ (5,6,7), (1,2)(3,4)(6,7) ]
  : Action homomorphism maps range generators to automorphisms:
    (5,7,6) --> { source gens --> [ (4,5,6), (2,3)(4,5) ] }
    (1,2)(3,4)(6,7) --> { source gens --> [ (4,6,5), (2,3)(5,6) ] }
    These 2 automorphisms generate the group of automorphisms.
  !gapprompt@gap>| !gapinput@AX3 := ActorXMod( X3 );; |
  !gapprompt@gap>| !gapinput@Display( AX3);|
  Crossed module Actor[c3->s3] :- 
  : Source group has generators:
    [ (1,2,3), (1,2) ]
  : Range group has generators:
    [ (5,7,6), (1,2)(3,4)(6,7) ]
  : Boundary homomorphism maps source generators to:
    [ (5,7,6), (1,2)(3,4)(6,7) ]
  : Action homomorphism maps range generators to automorphisms:
    (5,7,6) --> { source gens --> [ (1,2,3), (2,3) ] }
    (1,2)(3,4)(6,7) --> { source gens --> [ (1,3,2), (1,2) ] }
    These 2 automorphisms generate the group of automorphisms.
  
\end{Verbatim}
 The main methods for these operations are written for permutation crossed
modules. For other crossed modules an isomorphism to a permutation crossed
module is found first, and then the main method is applied to the image. In
the example the crossed module \texttt{XAq8} is the automorphism crossed module of the quaternion group. 
\begin{Verbatim}[commandchars=!@|,fontsize=\small,frame=single,label=Example]
  
  !gapprompt@gap>| !gapinput@q8 := Group( (1,2,3,4)(5,8,7,6), (1,5,3,7)(2,6,4,8) );;|
  !gapprompt@gap>| !gapinput@SetName( q8, "q8" );|
  !gapprompt@gap>| !gapinput@XAq8 := XModByAutomorphismGroup( q8 );; |
  !gapprompt@gap>| !gapinput@StructureDescription( WhiteheadXMod( XAq8 ) ); |
  [ "Q8", "C2 x C2 x C2" ]
  !gapprompt@gap>| !gapinput@StructureDescription( LueXMod( XAq8 ) );      |
  [ "Q8", "S4" ]
  !gapprompt@gap>| !gapinput@StructureDescription( NorrieXMod( XAq8 ) );|
  [ "S4", "S4" ]
  !gapprompt@gap>| !gapinput@StructureDescription( ActorXMod( XAq8 ) ); |
  [ "C2 x C2 x C2", "S4" ]
  
\end{Verbatim}
 

\subsection{\textcolor{Chapter }{XModCentre}}
\logpage{[ 6, 1, 4 ]}\nobreak
\hyperdef{L}{X85CF21F57F0F1329}{}
{\noindent\textcolor{FuncColor}{$\triangleright$\enspace\texttt{XModCentre({\mdseries\slshape xmod})\index{XModCentre@\texttt{XModCentre}}
\label{XModCentre}
}\hfill{\scriptsize (attribute)}}\\
\noindent\textcolor{FuncColor}{$\triangleright$\enspace\texttt{InnerActorXMod({\mdseries\slshape xmod})\index{InnerActorXMod@\texttt{InnerActorXMod}}
\label{InnerActorXMod}
}\hfill{\scriptsize (attribute)}}\\
\noindent\textcolor{FuncColor}{$\triangleright$\enspace\texttt{InnerMorphism({\mdseries\slshape xmod})\index{InnerMorphism@\texttt{InnerMorphism}}
\label{InnerMorphism}
}\hfill{\scriptsize (attribute)}}\\


 Pairs of boundaries or identity mappings provide six morphisms of crossed
modules. In particular, the boundaries $(\eta,\alpha)$ of $\calW(\calX)$ and $\calN(\calX)$ form the \emph{inner morphism} of $\calX$, mapping source elements to principal derivations and range elements to inner
automorphisms. The image of $\calX$ under this morphism is the \emph{inner actor} of $\calX$, while the kernel is the \emph{centre} of $\calX$. In the example which follows, the inner morphism of \texttt{X3=(c3\texttt{\symbol{45}}{\textgreater}s3)}, from Chapter \ref{chap-gp2up}, is an inclusion of crossed modules. 

 Note that we appear to have defined \emph{two} sorts of \emph{centre} for a crossed module $\calX$: \texttt{XModCentre} here, and \texttt{CentreXMod} (\ref{CentreXMod}) in the chapter on isoclinism. It is straightforward to show that these are the
same. Derivation $\eta_s$ is in the kernel of $\eta$ provided $(s^{-1})^rs = 1$ for all $r \in R$, so that $r \in \Fix(\calX,S,R)$. Similarly, inner automorphism $\alpha_r$ is in the kernel of $\alpha$ provided $s^r = s$ for all $s \in S$ and $r^{-1}r'r = r'$ for all $r' \in R$, so that $r$ is also in the centre of $R$. }

 

 
\begin{Verbatim}[commandchars=!@|,fontsize=\small,frame=single,label=Example]
  
  !gapprompt@gap>| !gapinput@IMX3 := InnerMorphism( X3 );; |
  !gapprompt@gap>| !gapinput@Display( IMX3 );|
  Morphism of crossed modules :- 
  : Source = [c3->s3] with generating sets:
    [ (1,2,3)(4,6,5) ]
    [ (4,5,6), (2,3)(5,6) ]
  :  Range = Actor[c3->s3] with generating sets:
    [ (1,2,3), (1,2) ]
    [ (5,7,6), (1,2)(3,4)(6,7) ]
  : Source Homomorphism maps source generators to:
    [ (1,2,3) ]
  : Range Homomorphism maps range generators to:
    [ (5,6,7), (1,2)(3,4)(6,7) ]
  !gapprompt@gap>| !gapinput@IsInjective( IMX3 );|
  true
  !gapprompt@gap>| !gapinput@ZX3 := XModCentre( X3 ); |
  [Group( () )->Group( () )]
  !gapprompt@gap>| !gapinput@IAX3 := InnerActorXMod( X3 );;  |
  !gapprompt@gap>| !gapinput@Display( IAX3 );|
  Crossed module InnerActor[c3->s3] :- 
  : Source group has generators:
    [ (1,2,3) ]
  : Range group has generators:
    [ (5,6,7), (1,2)(3,4)(6,7) ]
  : Boundary homomorphism maps source generators to:
    [ (5,7,6) ]
  : Action homomorphism maps range generators to automorphisms:
    (5,6,7) --> { source gens --> [ (1,2,3) ] }
    (1,2)(3,4)(6,7) --> { source gens --> [ (1,3,2) ] }
    These 2 automorphisms generate the group of automorphisms.
  
\end{Verbatim}
 }

 
\section{\textcolor{Chapter }{Actor of a cat$^1$\texttt{\symbol{45}}group}}\logpage{[ 6, 2, 0 ]}
\hyperdef{L}{X81BFAD86831097E3}{}
{
 

\subsection{\textcolor{Chapter }{ActorCat1Group}}
\logpage{[ 6, 2, 1 ]}\nobreak
\hyperdef{L}{X82097EE1866D0C2B}{}
{\noindent\textcolor{FuncColor}{$\triangleright$\enspace\texttt{ActorCat1Group({\mdseries\slshape cat1})\index{ActorCat1Group@\texttt{ActorCat1Group}}
\label{ActorCat1Group}
}\hfill{\scriptsize (attribute)}}\\
\noindent\textcolor{FuncColor}{$\triangleright$\enspace\texttt{InnerActorCat1Group({\mdseries\slshape cat1})\index{InnerActorCat1Group@\texttt{InnerActorCat1Group}}
\label{InnerActorCat1Group}
}\hfill{\scriptsize (function)}}\\


 The actor of a cat$^1$\texttt{\symbol{45}}group $C$ is obtained by converting $C$ to a crossed module; forming the actor of that crossed module; and then
converting that actor into a cat$^1$\texttt{\symbol{45}}group. 

 A similar procedure is followed for the inner actor. }

 

 
\begin{Verbatim}[commandchars=!@|,fontsize=\small,frame=single,label=Example]
  
  !gapprompt@gap>| !gapinput@C3;|
  [g18 => s3]
  !gapprompt@gap>| !gapinput@AC3 := ActorCat1Group( C3 );|
  cat1(Actor[c3->s3])
  !gapprompt@gap>| !gapinput@Display( AC3 );             |
  Cat1-group cat1(Actor[c3->s3]) :- 
  : Source group has generators:
    [ ( 9,10), ( 8, 9,10), ( 5, 7, 6)( 8, 9,10), (1,2)(3,4)(6,7)(8,9) ]
  : Range group has generators:
    [ (5,7,6), (1,2)(3,4)(6,7) ]
  : tail homomorphism maps source generators to:
    [ (), (), (5,7,6), (1,2)(3,4)(6,7) ]
  : head homomorphism maps source generators to:
    [ (1,2)(3,4)(5,6), (5,7,6), (5,7,6), (1,2)(3,4)(6,7) ]
  : range embedding maps range generators to:
    [ ( 5, 7, 6)( 8, 9,10), (1,2)(3,4)(6,7)(8,9) ]
  : kernel has generators:
    [ ( 9,10), ( 8, 9,10) ]
  : boundary homomorphism maps generators of kernel to:
    [ (1,2)(3,4)(5,6), (5,7,6) ]
  : kernel embedding maps generators of kernel to:
    [ ( 9,10), ( 8, 9,10) ]
  : associated crossed module is Actor[c3->s3]
  !gapprompt@gap>| !gapinput@StructureDescription( AC3 );|
  [ "S3 x S3", "S3" ]
  !gapprompt@gap>| !gapinput@IAC3 := InnerActorCat1Group( C3 );|
  cat1(InnerActor[c3->s3])
  !gapprompt@gap>| !gapinput@StructureDescription( IAC3 );|
  [ "(C3 x C3) : C2", "S3" ]
  
\end{Verbatim}
 

\subsection{\textcolor{Chapter }{Actor}}
\logpage{[ 6, 2, 2 ]}\nobreak
\hyperdef{L}{X7FE056707BB983B3}{}
{\noindent\textcolor{FuncColor}{$\triangleright$\enspace\texttt{Actor({\mdseries\slshape args})\index{Actor@\texttt{Actor}}
\label{Actor}
}\hfill{\scriptsize (function)}}\\
\noindent\textcolor{FuncColor}{$\triangleright$\enspace\texttt{InnerActor({\mdseries\slshape args})\index{InnerActor@\texttt{InnerActor}}
\label{InnerActor}
}\hfill{\scriptsize (function)}}\\


 The global functions \texttt{Actor} and \texttt{InnerActor} will call operations \texttt{ActorXMod} and \texttt{InnerActorXMod} or \texttt{ActorCat1Group} and \texttt{InnerActorCat1Group} as appropriate. }

 

 
\begin{Verbatim}[commandchars=!@|,fontsize=\small,frame=single,label=Example]
  
  !gapprompt@gap>| !gapinput@c4q := Subgroup( q8, [ (1,2,3,4)(5,8,7,6) ] );;|
  !gapprompt@gap>| !gapinput@SetName( c4q, "c4q" );                         |
  !gapprompt@gap>| !gapinput@Xc4q := XModByNormalSubgroup( q8, c4q );;      |
  !gapprompt@gap>| !gapinput@AXc4q := Actor( Xc4q );|
  Actor[c4q->q8]
  !gapprompt@gap>| !gapinput@StructureDescription( AXc4q );|
  [ "D8", "D8" ]
  !gapprompt@gap>| !gapinput@IAXc4q := InnerActor( Xc4q );|
  InnerActor[c4q->q8]
  !gapprompt@gap>| !gapinput@StructureDescription( IAXc4q );|
  [ "C2", "C2 x C2" ]
  
\end{Verbatim}
 }

 }

         
\chapter{\textcolor{Chapter }{Induced constructions}}\label{chap-gp2ind}
\logpage{[ 7, 0, 0 ]}
\hyperdef{L}{X8339DF98872D2E1C}{}
{
  Before describing general functions for computing induced structures, we
consider coproducts of crossed modules which provide induced crossed modules
in certain cases. 
\section{\textcolor{Chapter }{Coproducts of crossed modules}}\logpage{[ 7, 1, 0 ]}
\hyperdef{L}{X80D3C8F97A10D5E5}{}
{
 Need to add here a reference (or two) for coproducts. 

\subsection{\textcolor{Chapter }{CoproductXMod}}
\logpage{[ 7, 1, 1 ]}\nobreak
\hyperdef{L}{X7C01F5D98046E44B}{}
{\noindent\textcolor{FuncColor}{$\triangleright$\enspace\texttt{CoproductXMod({\mdseries\slshape X1, X2})\index{CoproductXMod@\texttt{CoproductXMod}}
\label{CoproductXMod}
}\hfill{\scriptsize (operation)}}\\
\noindent\textcolor{FuncColor}{$\triangleright$\enspace\texttt{CoproductInfo({\mdseries\slshape X0})\index{CoproductInfo@\texttt{CoproductInfo}}
\label{CoproductInfo}
}\hfill{\scriptsize (attribute)}}\\


 This function calculates the coproduct crossed module of two or more crossed
modules which have a common range $R$. The standard method applies to $\calX_1 = (\partial_1 : S_1 \to R)$ and $\calX_2 = (\partial_2 : S_2 \to R)$. See below for the case of three or more crossed modules. 

 The source $S_2$ of $\calX_2$ acts on $S_1$ via $\partial_2$ and the action of $\calX_1$, so we can form a precrossed module $(\partial' : S_1 \ltimes S_2 \to R)$ where $\partial'(s_1,s_2) = (\partial_1 s_1)(\partial_2 s_2)$. The action of this precrossed module is the diagonal action $(s_1,s_2)^r = (s_1^r,s_2^r)$. Factoring out by the Peiffer subgroup, we obtain the coproduct crossed
module $\calX_1 \circ \calX_2$. 

 In the example the structure descriptions of the precrossed module, the
Peiffer subgroup, and the resulting coproduct are printed out when \texttt{InfoLevel(InfoXMod)} is at least $1$. The coproduct comes supplied with attribute \texttt{CoproductInfo}, which includes the embedding morphisms of the two factors. }

 

 
\begin{Verbatim}[commandchars=!@|,fontsize=\small,frame=single,label=Example]
  
  !gapprompt@gap>| !gapinput@q8 := Group( (1,2,3,4)(5,8,7,6), (1,5,3,7)(2,6,4,8) );;|
  !gapprompt@gap>| !gapinput@XAq8 := XModByAutomorphismGroup( q8 );;|
  !gapprompt@gap>| !gapinput@s4b := Range( XAq8 );; |
  !gapprompt@gap>| !gapinput@SetName( q8, "q8" );  SetName( s4b, "s4b" ); |
  !gapprompt@gap>| !gapinput@a := q8.1;;  b := q8.2;; |
  !gapprompt@gap>| !gapinput@alpha := GroupHomomorphismByImages( q8, q8, [a,b], [a^-1,b] );;|
  !gapprompt@gap>| !gapinput@beta := GroupHomomorphismByImages( q8, q8, [a,b], [a,b^-1] );;|
  !gapprompt@gap>| !gapinput@k4b := Subgroup( s4b, [ alpha, beta ] );;  SetName( k4b, "k4b" );|
  !gapprompt@gap>| !gapinput@Z8 := XModByNormalSubgroup( s4b, k4b );;|
  !gapprompt@gap>| !gapinput@SetName( XAq8, "XAq8" );  SetName( Z8, "Z8" );  |
  !gapprompt@gap>| !gapinput@SetInfoLevel( InfoXMod, 1 ); |
  !gapprompt@gap>| !gapinput@XZ8 := CoproductXMod( XAq8, Z8 );|
  #I  prexmod is [ [ 32, 47 ], [ 24, 12 ] ]
  #I  peiffer subgroup is C2, [ 2, 1 ]
  #I  the coproduct is [ "C2 x C2 x C2 x C2", "S4" ], [ [ 16, 14 ], [ 24, 12 ] ]
  [Group( [ f1, f2, f3, f4 ] )->s4b]
  !gapprompt@gap>| !gapinput@SetName( XZ8, "XZ8" ); |
  !gapprompt@gap>| !gapinput@info := CoproductInfo( XZ8 );|
  rec( embeddings := [ [XAq8 => XZ8], [Z8 => XZ8] ], xmods := [ XAq8, Z8 ] )
  !gapprompt@gap>| !gapinput@SetInfoLevel( InfoXMod, 0 ); |
  
\end{Verbatim}
 Given a list of more than two crossed modules with a common range $R$, then an iterated coproduct is formed: 
\[ \bigcirc~\left[ \calX_1,\calX_2,\ldots,\calX_n\right] ~=~ \calX_1 \circ
(\calX_2 \circ ( \ldots (\calX_{n-1} \circ \calX_n) \ldots ) ). \]
 The \texttt{embeddings} field of the \texttt{CoproductInfo} of the resulting crossed module $\calY$ contains the $n$ morphisms $\epsilon_i : \calX_i \to \calY (1 \leqslant i \leqslant n)$. 

 
\begin{Verbatim}[commandchars=!@|,fontsize=\small,frame=single,label=Example]
  
  !gapprompt@gap>| !gapinput@Y := CoproductXMod( [ XAq8, XAq8, Z8, Z8 ] );|
  [Group( [ f1, f2, f3, f4, f5, f6, f7, f8 ] )->s4b]
  !gapprompt@gap>| !gapinput@StructureDescription( Y );          |
  [ "C2 x C2 x C2 x C2 x C2 x C2 x C2 x C2", "S4" ]
  !gapprompt@gap>| !gapinput@CoproductInfo( Y );|
  rec( 
    embeddings := 
      [ [XAq8 => [Group( [ f1, f2, f3, f4, f5, f6, f7, f8 ] ) -> s4b]], 
        [XAq8 => [Group( [ f1, f2, f3, f4, f5, f6, f7, f8 ] ) -> s4b]], 
        [Z8 => [Group( [ f1, f2, f3, f4, f5, f6, f7, f8 ] ) -> s4b]], 
        [Z8 => [Group( [ f1, f2, f3, f4, f5, f6, f7, f8 ] ) -> s4b]] ], 
    xmods := [ XAq8, XAq8, Z8, Z8 ] )
  
\end{Verbatim}
 }

 
\section{\textcolor{Chapter }{Induced crossed modules}}\logpage{[ 7, 2, 0 ]}
\hyperdef{L}{X7966FF497C36C465}{}
{
 \index{induced crossed module} 

\subsection{\textcolor{Chapter }{InducedXMod}}
\logpage{[ 7, 2, 1 ]}\nobreak
\hyperdef{L}{X874CB2A278AADE3A}{}
{\noindent\textcolor{FuncColor}{$\triangleright$\enspace\texttt{InducedXMod({\mdseries\slshape args})\index{InducedXMod@\texttt{InducedXMod}}
\label{InducedXMod}
}\hfill{\scriptsize (function)}}\\
\noindent\textcolor{FuncColor}{$\triangleright$\enspace\texttt{IsInducedXMod({\mdseries\slshape xmod})\index{IsInducedXMod@\texttt{IsInducedXMod}}
\label{IsInducedXMod}
}\hfill{\scriptsize (property)}}\\
\noindent\textcolor{FuncColor}{$\triangleright$\enspace\texttt{InducedXModBySurjection({\mdseries\slshape xmod, hom})\index{InducedXModBySurjection@\texttt{InducedXModBySurjection}}
\label{InducedXModBySurjection}
}\hfill{\scriptsize (operation)}}\\
\noindent\textcolor{FuncColor}{$\triangleright$\enspace\texttt{InducedXModByCopower({\mdseries\slshape xmod, hom, list})\index{InducedXModByCopower@\texttt{InducedXModByCopower}}
\label{InducedXModByCopower}
}\hfill{\scriptsize (operation)}}\\
\noindent\textcolor{FuncColor}{$\triangleright$\enspace\texttt{MorphismOfInducedXMod({\mdseries\slshape xmod})\index{MorphismOfInducedXMod@\texttt{MorphismOfInducedXMod}}
\label{MorphismOfInducedXMod}
}\hfill{\scriptsize (attribute)}}\\


 A morphism of crossed modules $(\sigma, \rho) : \calX_1 \to \calX_2$ factors uniquely through an induced crossed module $\rho_{\ast} \calX_1 = (\delta : \rho_{\ast} S_1 \to R_2)$. Similarly, a morphism of cat$^1$\texttt{\symbol{45}}groups factors through an induced cat$^1$\texttt{\symbol{45}}group. Calculation of induced crossed modules of $\calX$ also provides an algebraic means of determining the homotopy $2$\texttt{\symbol{45}}type of homotopy pushouts of the classifying space of $\calX$. For more background from algebraic topology see references in \cite{BH1}, \cite{BW1}, \cite{BW2}. Induced crossed modules and induced cat$^1$\texttt{\symbol{45}}groups also provide the building blocks for constructing
pushouts in the categories \emph{XMod} and \emph{Cat1}. 

 Data for the cases of algebraic interest is provided by a crossed module $\calX = (\partial : S \to R)$ and a homomorphism $\iota : R \to Q$. The output from the calculation is a crossed module $\iota_{\ast}\calX = (\delta : \iota_{\ast}S \to Q)$ together with a morphism of crossed modules $\calX \to \iota_{\ast}\calX$. When $\iota$ is a surjection with kernel $K$ then $\iota_{\ast}S = S/[K,S]$ where $[K,S]$ is the subgroup of $S$ generated by elements of the form $s^{-1}s^k, s \in S, k \in K$ (see \cite{BH1}, Prop.9). (For many years, up until June 2018, this manual has stated the
result to be $[K,S]$, though the correct quotient had been calculated.) When $\iota$ is an inclusion the induced crossed module may be calculated using a copower
construction \cite{BW1} or, in the case when $R$ is normal in $Q$, as a coproduct of crossed modules (\cite{BW2}, but not yet implemented). When $\iota$ is neither a surjection nor an inclusion, $\iota$ is factored as the composite of the surjection onto the image and the
inclusion of the image in $Q$, and then the composite induced crossed module is constructed. These
constructions use Tietze transformation routines in the library file \texttt{tietze.gi}. 

 As a first, surjective example, we take for $\calX$ a central extension crossed module of dihedral groups, $(d_{24} \to d_{12})$, and for $\iota$ a surjection $d_{12} \to s_3$ with kernel $c_2$. The induced crossed module is isomorphic to $(d_{12} \to s_3)$. }

 

 
\begin{Verbatim}[commandchars=!@|,fontsize=\small,frame=single,label=Example]
  
  !gapprompt@gap>| !gapinput@a := (6,7,8,9)(10,11,12);;  b := (7,9)(11,12);;|
  !gapprompt@gap>| !gapinput@d24 := Group( [ a, b ] );;|
  !gapprompt@gap>| !gapinput@SetName( d24, "d24" );|
  !gapprompt@gap>| !gapinput@c := (1,2)(3,4,5);;  d := (4,5);;|
  !gapprompt@gap>| !gapinput@d12 := Group( [ c, d ] );;|
  !gapprompt@gap>| !gapinput@SetName( d12, "d12" );|
  !gapprompt@gap>| !gapinput@bdy := GroupHomomorphismByImages( d24, d12, [a,b], [c,d] );;|
  !gapprompt@gap>| !gapinput@X24 := XModByCentralExtension( bdy );|
  [d24->d12]
  !gapprompt@gap>| !gapinput@e := (13,14,15);;  f := (14,15);;|
  !gapprompt@gap>| !gapinput@s3 := Group( [ e, f ] );;|
  !gapprompt@gap>| !gapinput@SetName( s3, "s3" );;|
  !gapprompt@gap>| !gapinput@epi := GroupHomomorphismByImages( d12, s3, [c,d], [e,f] );;|
  !gapprompt@gap>| !gapinput@iX24 := InducedXModBySurjection( X24, epi );|
  [d24/ker->s3]
  !gapprompt@gap>| !gapinput@Display( iX24 );               |
  Crossed module [d24/ker->s3] :- 
  : Source group d24/ker has generators:
    [ ( 1,11, 5, 4,10, 8)( 2,12, 6, 3, 9, 7), 
    ( 1, 2)( 3, 4)( 5, 9)( 6,10)( 7,11)( 8,12) ]
  : Range group s3 has generators:
    [ (13,14,15), (14,15) ]
  : Boundary homomorphism maps source generators to:
    [ (13,14,15), (14,15) ]
  : Action homomorphism maps range generators to automorphisms:
    (13,14,15) --> { source gens --> [ ( 1,11, 5, 4,10, 8)( 2,12, 6, 3, 9, 7), 
    ( 1, 6)( 2, 5)( 3, 8)( 4, 7)( 9,10)(11,12) ] }
    (14,15) --> { source gens --> [ ( 1, 8,10, 4, 5,11)( 2, 7, 9, 3, 6,12), 
    ( 1, 2)( 3, 4)( 5, 9)( 6,10)( 7,11)( 8,12) ] }
    These 2 automorphisms generate the group of automorphisms.
  !gapprompt@gap>| !gapinput@morX24 := MorphismOfInducedXMod( iX24 );|
  [[d24->d12] => [d24/ker->s3]]
  
\end{Verbatim}
 For a second, injective example we take for $\calX$ the result \texttt{iX24} of the previous example and for $\iota$ an inclusion of $s_3$ in $s_4$. The resulting source group has size $96$. 

 
\begin{Verbatim}[commandchars=!@|,fontsize=\small,frame=single,label=Example]
  
  !gapprompt@gap>| !gapinput@g := (16,17,18);;  h := (16,17,18,19);;|
  !gapprompt@gap>| !gapinput@s4 := Group( [ g, h ] );;|
  !gapprompt@gap>| !gapinput@SetName( s4, "s4" );;|
  !gapprompt@gap>| !gapinput@iota := GroupHomomorphismByImages( s3, s4, [e,f], [g^2*h^2,g*h^-1] );|
  [ (13,14,15), (14,15) ] -> [ (17,18,19), (18,19) ]
  !gapprompt@gap>| !gapinput@iiX24 := InducedXModByCopower( iX24, iota, [ ] );|
  i*([d24/ker->s3])
  !gapprompt@gap>| !gapinput@Size2d( iiX24 );               |
  [ 96, 24 ]
  !gapprompt@gap>| !gapinput@StructureDescription( iiX24 );|
  [ "C2 x GL(2,3)", "S4" ]
  
\end{Verbatim}
 For a third example we combine the previous two examples by taking for $\iota$ the more general case \texttt{alpha = theta*iota}. The resulting \texttt{jX24} is isomorphic to, but not identical to, \texttt{iiX24}. 

 
\begin{Verbatim}[commandchars=!@|,fontsize=\small,frame=single,label=Example]
  
  !gapprompt@gap>| !gapinput@alpha := CompositionMapping( iota, epi );|
  [ (1,2)(3,4,5), (4,5) ] -> [ (17,18,19), (18,19) ]
  !gapprompt@gap>| !gapinput@jX24 := InducedXMod( X24, alpha );;|
  !gapprompt@gap>| !gapinput@StructureDescription( jX24 );|
  [ "C2 x GL(2,3)", "S4" ]
  
\end{Verbatim}
 For a fourth example we use the version \texttt{InducedXMod(Q,R,S)} of this global function, with a normal inclusion crossed module $(S \to R)$ and an inclusion mapping $R \to Q$. We take $(c_6 \to d_{12})$ as $\calX$ and the inclusion of $d_{12}$ in $d_{24}$ as $\iota$. 

 
\begin{Verbatim}[commandchars=!@|,fontsize=\small,frame=single,label=Example]
  
  ## Section 7.2.1 : Example 4 
  !gapprompt@gap>| !gapinput@d12b := Subgroup( d24, [ a^2, b ] );;|
  !gapprompt@gap>| !gapinput@SetName( d12b, "d12b" ); |
  !gapprompt@gap>| !gapinput@c6b := Subgroup( d12b, [ a^2 ] );; |
  !gapprompt@gap>| !gapinput@SetName( c6b, "c6b" );  |
  !gapprompt@gap>| !gapinput@X12 := InducedXMod( d24, d12b, c6b );|
  i*([c6b->d12b])
  !gapprompt@gap>| !gapinput@StructureDescription( X12 );|
  [ "C6 x C6", "D24" ]
  !gapprompt@gap>| !gapinput@Display( MorphismOfInducedXMod( X12 ) );|
  Morphism of crossed modules :- 
  : Source = [c6b->d12b] with generating sets:
    [ ( 6, 8)( 7, 9)(10,12,11) ]
    [ ( 6, 8)( 7, 9)(10,12,11), ( 7, 9)(11,12) ]
  :  Range = i*([c6b->d12b]) with generating sets:
    [ ( 4, 5)( 6, 7)( 8, 9)(10,11)(12,13)(14,15), 
    ( 4, 6, 8)( 5, 7, 9)(10,12,14)(11,13,15), 
    ( 4,10)( 5,11)( 6,12)( 7,13)( 8,14)( 9,15), (1,2,3) ]
    [ ( 6, 7, 8, 9)(10,11,12), ( 7, 9)(11,12) ]
  : Source Homomorphism maps source generators to:
    [ ( 4, 9, 6, 5, 8, 7)(10,15,12,11,14,13) ]
  : Range Homomorphism maps range generators to:
    [ ( 6, 8)( 7, 9)(10,12,11), ( 7, 9)(11,12) ]
  #
\end{Verbatim}
 

\subsection{\textcolor{Chapter }{AllInducedXMods}}
\logpage{[ 7, 2, 2 ]}\nobreak
\hyperdef{L}{X7B24D47F8078540F}{}
{\noindent\textcolor{FuncColor}{$\triangleright$\enspace\texttt{AllInducedXMods({\mdseries\slshape Q})\index{AllInducedXMods@\texttt{AllInducedXMods}}
\label{AllInducedXMods}
}\hfill{\scriptsize (operation)}}\\


 This function calculates all the induced crossed modules \texttt{InducedXMod(Q,R,S)}, where \texttt{R} runs over all conjugacy classes of subgroups of \texttt{Q} and \texttt{S} runs over all non\texttt{\symbol{45}}trivial normal subgroups of \texttt{R}. }

 
\begin{Verbatim}[commandchars=!@|,fontsize=\small,frame=single,label=Example]
  
  !gapprompt@gap>| !gapinput@all := AllInducedXMods( q8 );;|
  !gapprompt@gap>| !gapinput@L := List( all, x -> Source( x ) );;|
  !gapprompt@gap>| !gapinput@Sort( L, function(g,h) return Size(g) < Size(h); end );;|
  !gapprompt@gap>| !gapinput@List( L, x -> StructureDescription( x ) );|
  [ "1", "1", "1", "1", "C2 x C2", "C2 x C2", "C2 x C2", "C4 x C4", "C4 x C4", 
    "C4 x C4", "C2 x C2 x C2 x C2" ]
  
\end{Verbatim}
 }

 
\section{\textcolor{Chapter }{Induced cat$^1$\texttt{\symbol{45}}groups}}\logpage{[ 7, 3, 0 ]}
\hyperdef{L}{X814A695779706E22}{}
{
 \index{induced cat1\texttt{\symbol{45}}groups} 

\subsection{\textcolor{Chapter }{InducedCat1Group}}
\logpage{[ 7, 3, 1 ]}\nobreak
\hyperdef{L}{X7BCE57BE7F6E6B08}{}
{\noindent\textcolor{FuncColor}{$\triangleright$\enspace\texttt{InducedCat1Group({\mdseries\slshape args})\index{InducedCat1Group@\texttt{InducedCat1Group}}
\label{InducedCat1Group}
}\hfill{\scriptsize (function)}}\\
\noindent\textcolor{FuncColor}{$\triangleright$\enspace\texttt{InducedCat1GroupByFreeProduct({\mdseries\slshape grp, hom})\index{InducedCat1GroupByFreeProduct@\texttt{InducedCat1GroupByFreeProduct}}
\label{InducedCat1GroupByFreeProduct}
}\hfill{\scriptsize (property)}}\\


 This area awaits development. }

 }

 }

         
\chapter{\textcolor{Chapter }{Crossed squares and Cat$^2$\texttt{\symbol{45}}groups}}\label{chap-obj3}
\logpage{[ 8, 0, 0 ]}
\hyperdef{L}{X780368C083C76EDC}{}
{
  \index{3d\texttt{\symbol{45}}group} \index{3d\texttt{\symbol{45}}domain} The term \emph{3d\texttt{\symbol{45}}group} refers to a set of equivalent categories of which the most common are the
categories of \emph{crossed squares} and \emph{cat$^2$\texttt{\symbol{45}}groups}. A \emph{3d\texttt{\symbol{45}}mapping} is a function between two 3d\texttt{\symbol{45}}groups which preserves all the
structure. 

 The material in this chapter should be considered experimental. A major
overhaul took place in time for \textsf{XMod} version 2.73, with the names of a number of operations being changed. 
\section{\textcolor{Chapter }{Definition of a crossed square and a crossed $n$\texttt{\symbol{45}}cube of groups}}\label{sect-xsq-definition}
\logpage{[ 8, 1, 0 ]}
\hyperdef{L}{X7C4AFE8D85848C8F}{}
{
  \index{crossed square} Crossed squares were introduced by Guin\texttt{\symbol{45}}Wal{\a'e}ry and
Loday (see, for example, \cite{brow:lod}) as fundamental crossed squares of commutative squares of spaces, but are
also of purely algebraic interest. We denote by $[n]$ the set $\{1,2,\ldots,n\}$. We use the $n=2$ version of the definition of crossed $n$\texttt{\symbol{45}}cube as given by Ellis and Steiner \cite{ell:st}. 

 A \emph{crossed square} $\calS$ consists of the following: 
\begin{itemize}
\item  groups $S_J$ for each of the four subsets $J \subseteq [2]$ (we often find it convenient to write $L = S_{[2]},~ M = S_{\{1\}},~ N = S_{\{2\}}$ and $P = S_{\emptyset}$); 
\item  a commutative diagram of group homomorphisms: 
\[ \ddot{\partial}_1 : S_{[2]} \to S_{\{2\}}, \quad \ddot{\partial}_2 : S_{[2]}
\to S_{\{1\}}, \quad \dot{\partial}_2 : S_{\{2\}} \to S_{\emptyset}, \quad
\dot{\partial}_1 : S_{\{1\}} \to S_{\emptyset} \]
 (again we often write $\kappa = \ddot{\partial}_1,~ \lambda = \ddot{\partial}_2,~ \mu =
\dot{\partial}_2$ and $\nu = \dot{\partial}_1$); 
\item  actions of $S_{\emptyset}$ on $S_{\{1\}}, S_{\{2\}}$ and $S_{[2]}$ which determine actions of $S_{\{1\}}$ on $S_{\{2\}}$ and $S_{[2]}$ via $\dot{\partial}_1$ and actions of $S_{\{2\}}$ on $S_{\{1\}}$ and $S_{[2]}$ via $\dot{\partial}_2\;$; 
\item  a function $\boxtimes : S_{\{1\}} \times S_{\{2\}} \to S_{[2]}$. 
\end{itemize}
 Here is a picture of the situation: 
\[ \vcenter{\xymatrix{ & & S_{[2]} \ar[rr]^{\ddot{\partial}_1}
\ar[dd]_{\ddot{\partial}_2} && S_{\{2\}} \ar[dd]^{\dot{\partial}_2} && L
\ar[rr]^{\kappa} \ar[dd]_{\lambda} && M \ar[dd]^{\mu} & \\ \mathbb{S} & = & &&
& = && \\ & & S_{\{1\}} \ar[rr]_{\dot{\partial}_1} && S_{\emptyset} && N
\ar[rr]_{\nu} && P }} \]
 The following axioms must be satisfied for all $l \in L,\; m,m_1,m_2 \in M,\; n,n_1,n_2 \in N,\; p \in P$. 
\begin{itemize}
\item  The homomorphisms $\kappa, \lambda$ preserve the action of $P\;$. 
\item  Each of the upper, left\texttt{\symbol{45}}hand, right\texttt{\symbol{45}}hand
and lower sides of the square, 
\[ \ddot{\calS}_1 = (\kappa : L \to M), \quad \ddot{\calS}_2 = (\lambda : L \to
N), \quad \dot{\calS}_2 = (\mu : M \to P), \quad \dot{\calS}_1 = (\nu : N \to
P), \]
 and the diagonal 
\[ \calS_{12} = (\partial_{12} := \mu \circ \kappa = \nu \circ \lambda : L \to P) \]
 are crossed modules (with actions via $P$). 

 These will be called the \emph{up, left, right, down} and \emph{diagonal} crossed modules of $\calS$. 
\item  \index{crossed pairing} $\boxtimes$ is a \emph{crossed pairing}: 
\begin{itemize}
\item  $(n_1n_2 \boxtimes m)\;=\; {(n_1 \boxtimes m)}^{n_2}\;(n_2 \boxtimes m)$, 
\item  $(n \boxtimes m_1m_2) \;=\; (n \boxtimes m_2)\;{(n \boxtimes m_1)}^{m_2}$, 
\item  $(n \boxtimes m)^{p} \;=\; (n^p \boxtimes m^p)$. 
\end{itemize}
 
\item  $\kappa(n \boxtimes m) \;=\; (m^{-1})^{n}\;m \quad \mbox{and} \quad \lambda(n
\boxtimes m) \;=\; n^{-1}\;n^{m}$. 
\item  $(n \boxtimes \kappa l) \;=\; (l^{-1})^{n}\;l \quad \mbox{and} \quad (\lambda l
\boxtimes m) \;=\; l^{-1}\;l^m$. 
\end{itemize}
 Note that the actions of $M$ on $N$ and $N$ on $M$ via $P$ are compatible since 
\[ {n_1}^{(m^n)} \;=\; {n_1}^{\mu(m^n)} \;=\; {n_1}^{n^{-1}(\mu m)n} \;=\;
(({n_1}^{n^{-1}})^m)^n. \]
 

 A \emph{precrossed square} is a similar structure which satisfies some subset of these axioms. (This
theoretical notion needs to be clarified.) In this implementation \texttt{IsPreCrossedSquare} only checks that the five pre\texttt{\symbol{45}}crossed modules fit together
correctly and that $(\kappa,\nu)$ and $(\lambda,\mu)$ are both morphisms of pre\texttt{\symbol{45}}crossed modules. 

 Crossed squares are the $n=2$ case of a crossed $n$\texttt{\symbol{45}}cube of groups, defined as follows. (This is an attempt to
translate Definition 2.1 in Ronnie Brown's \emph{Computing homotopy types using crossed n\texttt{\symbol{45}}cubes of groups} into right actions \texttt{\symbol{45}}\texttt{\symbol{45}} but this
definition is not yet completely understood!) 

 \index{crossed n\texttt{\symbol{45}}cube} A \emph{crossed} $n$\emph{\texttt{\symbol{45}}cube of groups} consists of the following: 
\begin{itemize}
\item  groups $S_A$ for every subset $A \subseteq [n]$; 
\item  a commutative diagram of group homomorphisms $\partial_i : S_A \to S_{A \setminus \{i\}},\; i \in [n]$; with composites $\partial_B : S_A \to S_{A \setminus B},\; B \subseteq [n]$; 
\item  actions of $S_{\emptyset}$ on each $S_A$; and hence actions of $S_B$ on $S_A$ via $\partial_B$ for each $B \subseteq [n]$; 
\item  functions $\boxtimes_{A,B} : S_A \times S_B \to S_{A \cup B}, (A,B \subseteq [n])$. 
\end{itemize}
 There is then a long list of axioms which must be satisfied. }

 
\section{\textcolor{Chapter }{Constructions for crossed squares}}\label{sect-xsq-constructions}
\logpage{[ 8, 2, 0 ]}
\hyperdef{L}{X820A7D30847BC828}{}
{
  Analogously to the data structure used for crossed modules, crossed squares
are implemented as \texttt{3d\texttt{\symbol{45}}groups}. There are also experimental implementations of cat$^2$\texttt{\symbol{45}}groups, with conversion between the two types of
structure. Some standard constructions of crossed squares are listed below. At
present, a limited number of constructions is implemented. Morphisms of
crossed squares have also been implemented, though there is still a great deal
to be done. 

\subsection{\textcolor{Chapter }{CrossedSquareByXMods}}
\logpage{[ 8, 2, 1 ]}\nobreak
\hyperdef{L}{X866A7FAC7FCB62C2}{}
{\noindent\textcolor{FuncColor}{$\triangleright$\enspace\texttt{CrossedSquareByXMods({\mdseries\slshape up, left, right, down, diag, pairing})\index{CrossedSquareByXMods@\texttt{CrossedSquareByXMods}}
\label{CrossedSquareByXMods}
}\hfill{\scriptsize (operation)}}\\
\noindent\textcolor{FuncColor}{$\triangleright$\enspace\texttt{PreCrossedSquareByPreXMods({\mdseries\slshape up, left, right, down, diag, pairing})\index{PreCrossedSquareByPreXMods@\texttt{PreCrossedSquareByPreXMods}}
\label{PreCrossedSquareByPreXMods}
}\hfill{\scriptsize (operation)}}\\


 If \emph{up,left,right,down,diag} are five (pre\texttt{\symbol{45}})crossed modules whose sources and ranges
agree, as above, then we just have to add a crossed pairing to complete the
data for a (pre\texttt{\symbol{45}})crossed square. 

 \index{Display for a 3d\texttt{\symbol{45}}group} The \texttt{Display} function is used to print details of 3d\texttt{\symbol{45}}groups. 

 We take as our example a simple, but significant case. We start with five
crossed modules formed from subgroups of $D_8$ with generators $[(1,2,3,4),(3,4)$. The result is a pre\texttt{\symbol{45}}crossed square which is \emph{not} a crossed square. }

 

 
\begin{Verbatim}[commandchars=!@|,fontsize=\small,frame=single,label=Example]
  
  !gapprompt@gap>| !gapinput@b := (2,4);; c := (1,2)(3,4);; p := (1,2,3,4);; |
  !gapprompt@gap>| !gapinput@d8 := Group( b, c );; |
  !gapprompt@gap>| !gapinput@SetName( d8, "d8" );; |
  !gapprompt@gap>| !gapinput@L := Subgroup( d8, [p^2] );; |
  !gapprompt@gap>| !gapinput@M := Subgroup( d8, [b] );; |
  !gapprompt@gap>| !gapinput@N := Subgroup( d8, [c] );; |
  !gapprompt@gap>| !gapinput@P := TrivialSubgroup( d8 );; |
  !gapprompt@gap>| !gapinput@kappa := GroupHomomorphismByImages( L, M, [p^2], [b] );; |
  !gapprompt@gap>| !gapinput@lambda := GroupHomomorphismByImages( L, N, [p^2], [c] );; |
  !gapprompt@gap>| !gapinput@delta := GroupHomomorphismByImages( L, P, [p^2], [()] );; |
  !gapprompt@gap>| !gapinput@mu := GroupHomomorphismByImages( M, P, [b], [()] );; |
  !gapprompt@gap>| !gapinput@nu := GroupHomomorphismByImages( N, P, [c], [()] );; |
  !gapprompt@gap>| !gapinput@up := XModByTrivialAction( kappa );; |
  !gapprompt@gap>| !gapinput@left := XModByTrivialAction( lambda );; |
  !gapprompt@gap>| !gapinput@diag := XModByTrivialAction( delta );; |
  !gapprompt@gap>| !gapinput@right := XModByTrivialAction( mu );; |
  !gapprompt@gap>| !gapinput@down := XModByTrivialAction( nu );; |
  !gapprompt@gap>| !gapinput@xp := CrossedPairingByCommutators( N, M, L );; |
  !gapprompt@gap>| !gapinput@Print( "xp([c,b]) = ", xp( c, b ), "\n" ); |
  xp([c,b]) = (1,3)(2,4)
  !gapprompt@gap>| !gapinput@PXS := PreCrossedSquareByPreXMods( up, left, right, down, diag, xp );;|
  !gapprompt@gap>| !gapinput@Display( PXS ); |
  (pre-)crossed square with (pre-)crossed modules:
        up = [Group( [ (1,3)(2,4) ] ) -> Group( [ (2,4) ] )]
      left = [Group( [ (1,3)(2,4) ] ) -> Group( [ (1,2)(3,4) ] )]
     right = [Group( [ (2,4) ] ) -> Group( () )]
      down = [Group( [ (1,2)(3,4) ] ) -> Group( () )]
  !gapprompt@gap>| !gapinput@IsCrossedSquare( PXS ); |
  false
  
\end{Verbatim}
 

\subsection{\textcolor{Chapter }{Size3d (for 3d-objects)}}
\logpage{[ 8, 2, 2 ]}\nobreak
\hyperdef{L}{X7FA98BD47FF9B044}{}
{\noindent\textcolor{FuncColor}{$\triangleright$\enspace\texttt{Size3d({\mdseries\slshape XS})\index{Size3d@\texttt{Size3d}!for 3d-objects}
\label{Size3d:for 3d-objects}
}\hfill{\scriptsize (attribute)}}\\


 Just as \texttt{Size2d} was used in place of \texttt{Size} for crossed modules, so \texttt{Size3d} is used for crossed squares: \texttt{Size3d( XS )} returns a four\texttt{\symbol{45}}element list containing the sizes of the
four groups at the corners of the square. 

 }

 
\begin{Verbatim}[commandchars=!@|,fontsize=\small,frame=single,label=Example]
  
  !gapprompt@gap>| !gapinput@Size3d( PXS ); |
  [ 2, 2, 2, 1 ]
  
\end{Verbatim}
 

\subsection{\textcolor{Chapter }{CrossedSquareByNormalSubgroups}}
\logpage{[ 8, 2, 3 ]}\nobreak
\hyperdef{L}{X7896DAF786F46234}{}
{\noindent\textcolor{FuncColor}{$\triangleright$\enspace\texttt{CrossedSquareByNormalSubgroups({\mdseries\slshape L, M, N, P})\index{CrossedSquareByNormalSubgroups@\texttt{CrossedSquareByNormalSubgroups}}
\label{CrossedSquareByNormalSubgroups}
}\hfill{\scriptsize (operation)}}\\
\noindent\textcolor{FuncColor}{$\triangleright$\enspace\texttt{CrossedPairingByCommutators({\mdseries\slshape N, M, L})\index{CrossedPairingByCommutators@\texttt{CrossedPairingByCommutators}}
\label{CrossedPairingByCommutators}
}\hfill{\scriptsize (operation)}}\\


 If $L, M, N$ are normal subgroups of a group $P$, and $[M,N] \leqslant L \leqslant M \cap N$, then the four inclusions $L \to M,~ L \to N,~ M \to P,~ N \to P$, together with the actions of $P$ on $M, N$ and $L$ given by conjugation, form a crossed square with crossed pairing 
\[ \boxtimes \;:\; N \times M \to L, \quad (n,m) \mapsto [n,m] \,=\,
n^{-1}m^{-1}nm \,=\,(m^{-1})^nm \,=\, n^{-1}n^m\,. \]
 This construction is implemented as \texttt{CrossedSquareByNormalSubgroups(L,M,N,P)} (note that the parent group comes last). }

 

 
\begin{Verbatim}[commandchars=!@A,fontsize=\small,frame=single,label=Example]
  
  !gapprompt@gap>A !gapinput@d20 := DihedralGroup( IsPermGroup, 20 );;A
  !gapprompt@gap>A !gapinput@gend20 := GeneratorsOfGroup( d20 ); A
  [ (1,2,3,4,5,6,7,8,9,10), (2,10)(3,9)(4,8)(5,7) ]
  !gapprompt@gap>A !gapinput@p1 := gend20[1];;  p2 := gend20[2];;  p12 := p1*p2; A
  (1,10)(2,9)(3,8)(4,7)(5,6)
  !gapprompt@gap>A !gapinput@d10a := Subgroup( d20, [ p1^2, p2 ] );;A
  !gapprompt@gap>A !gapinput@d10b := Subgroup( d20, [ p1^2, p12 ] );;A
  !gapprompt@gap>A !gapinput@c5d := Subgroup( d20, [ p1^2 ] );;A
  !gapprompt@gap>A !gapinput@SetName( d20, "d20" );  SetName( d10a, "d10a" ); A
  !gapprompt@gap>A !gapinput@SetName( d10b, "d10b" );  SetName( c5d, "c5d" ); A
  !gapprompt@gap>A !gapinput@XSconj := CrossedSquareByNormalSubgroups( c5d, d10a, d10b, d20 );A
  [  c5d -> d10a ]
  [   |      |   ]
  [ d10b -> d20  ]
  !gapprompt@gap>A !gapinput@xpc := CrossedPairing( XSconj );;A
  !gapprompt@gap>A !gapinput@xpc( p12, p2 );A
  (1,3,5,7,9)(2,4,6,8,10)
  
\end{Verbatim}
 

\subsection{\textcolor{Chapter }{CrossedSquareByNormalSubXMod}}
\logpage{[ 8, 2, 4 ]}\nobreak
\hyperdef{L}{X7FA367977B1895A7}{}
{\noindent\textcolor{FuncColor}{$\triangleright$\enspace\texttt{CrossedSquareByNormalSubXMod({\mdseries\slshape X0, X1})\index{CrossedSquareByNormalSubXMod@\texttt{CrossedSquareByNormalSubXMod}}
\label{CrossedSquareByNormalSubXMod}
}\hfill{\scriptsize (operation)}}\\
\noindent\textcolor{FuncColor}{$\triangleright$\enspace\texttt{CrossedPairingBySingleXModAction({\mdseries\slshape X0, X1})\index{CrossedPairingBySingleXModAction@\texttt{CrossedPairingBySingleXModAction}}
\label{CrossedPairingBySingleXModAction}
}\hfill{\scriptsize (operation)}}\\


 If $\calX_1 = (\partial_1 : S_1 \to R_1)$ is a normal sub\texttt{\symbol{45}}crossed module of $\calX_0 = (\partial_0 : S_0 \to R_0)$ then the inclusion morphism gives a crossed square with crossed pairing 
\[ \boxtimes \;:\; R_1 \times S_0 \to S_1, \quad (r_1,s_0) \mapsto
(s_0^{-1})^{r_1} s_0. \]
 

 The example constructs the same crossed square as in the previous subsection. }

 

 
\begin{Verbatim}[commandchars=!@A,fontsize=\small,frame=single,label=Example]
  
  !gapprompt@gap>A !gapinput@X20 := XModByNormalSubgroup( d20, d10a );; A
  !gapprompt@gap>A !gapinput@X10 := XModByNormalSubgroup( d10b, c5d );; A
  !gapprompt@gap>A !gapinput@ok := IsNormalSub2DimensionalDomain( X20, X10 ); A
  true 
  !gapprompt@gap>A !gapinput@XS20 := CrossedSquareByNormalSubXMod( X20, X10 ); A
  [  c5d -> d10a ]
  [   |      |   ]
  [ d10b -> d20  ]
  !gapprompt@gap>A !gapinput@xp20 := CrossedPairing( XS20 );; A
  !gapprompt@gap>A !gapinput@xp20( p1^2, p2 );A
  (1,7,3,9,5)(2,8,4,10,6)
  
\end{Verbatim}
 

\subsection{\textcolor{Chapter }{ActorCrossedSquare}}
\logpage{[ 8, 2, 5 ]}\nobreak
\hyperdef{L}{X833362FE87ED3C48}{}
{\noindent\textcolor{FuncColor}{$\triangleright$\enspace\texttt{ActorCrossedSquare({\mdseries\slshape X0})\index{ActorCrossedSquare@\texttt{ActorCrossedSquare}}
\label{ActorCrossedSquare}
}\hfill{\scriptsize (attribute)}}\\
\noindent\textcolor{FuncColor}{$\triangleright$\enspace\texttt{CrossedPairingByDerivations({\mdseries\slshape X0})\index{CrossedPairingByDerivations@\texttt{CrossedPairingByDerivations}}
\label{CrossedPairingByDerivations}
}\hfill{\scriptsize (operation)}}\\


 The actor $\calA(\calX_0)$ of a crossed module $\calX_0$ has been described in Chapter 6 (see \texttt{ActorXMod} (\ref{ActorXMod})). The crossed pairing is given by 
\[ \boxtimes \;:\; R \times W \,\to\, S, \quad (r,\chi) \,\mapsto\, \chi r~. \]
 This is implemented as \texttt{ActorCrossedSquare(X0)}. 

 The example constructs \texttt{XSact}, the actor crossed square of the crossed module \texttt{X20}. This crossed square is converted to a cat$^2$\texttt{\symbol{45}}group \texttt{C2act} in section \ref{cat2-xsq}. }

 

 
\begin{Verbatim}[commandchars=!@|,fontsize=\small,frame=single,label=Example]
  
  !gapprompt@gap>| !gapinput@XSact := ActorCrossedSquare( X20 );|
  crossed square with crossed modules:
        up = Whitehead[d10a->d20]
      left = [d10a->d20]
     right = Actor[d10a->d20]
      down = Norrie[d10a->d20]
  !gapprompt@gap>| !gapinput@W := Range( Up2DimensionalGroup( XSact ) );|
  Group([ (2,5)(3,4), (2,3,5,4), (1,4,2,5,3) ])
  !gapprompt@gap>| !gapinput@StructureDescription( W );|
  "C5 : C4"
  !gapprompt@gap>| !gapinput@xpa := CrossedPairing( XSact );;|
  !gapprompt@gap>| !gapinput@xpa( p1, (2,3,5,4) );|
  (1,7,3,9,5)(2,8,4,10,6)
  
\end{Verbatim}
 

\subsection{\textcolor{Chapter }{CrossedSquareByAutomorphismGroup}}
\logpage{[ 8, 2, 6 ]}\nobreak
\hyperdef{L}{X82E4EA93824F7B26}{}
{\noindent\textcolor{FuncColor}{$\triangleright$\enspace\texttt{CrossedSquareByAutomorphismGroup({\mdseries\slshape G})\index{CrossedSquareByAutomorphismGroup@\texttt{CrossedSquareByAutomorphismGroup}}
\label{CrossedSquareByAutomorphismGroup}
}\hfill{\scriptsize (operation)}}\\
\noindent\textcolor{FuncColor}{$\triangleright$\enspace\texttt{CrossedPairingByConjugators({\mdseries\slshape G})\index{CrossedPairingByConjugators@\texttt{CrossedPairingByConjugators}}
\label{CrossedPairingByConjugators}
}\hfill{\scriptsize (operation)}}\\


 For $G$ a group let $\Inn(G)$ be its inner automorphism group and $\Aut(G)$ its full automorphism group. Then there is a crossed square with groups $[G,\Inn(G),\Inn(G),\Aut(G)]$ where the upper and left boundaries are the maps $g \mapsto \iota_g$, where $\iota_g$ is conjugation of $G$ by $g$, and the right and down boundaries are inclusions. The crossed pairing is
gived by $\iota_g \boxtimes \iota_h = [g,h]$. }

 

 
\begin{Verbatim}[commandchars=!@E,fontsize=\small,frame=single,label=Example]
  
  !gapprompt@gap>E !gapinput@AXS20 := CrossedSquareByAutomorphismGroup( d20 );E
  [      d20 -> Inn(d20) ]
  [     |          |     ]
  [ Inn(d20) -> Aut(d20) ]
  
  !gapprompt@gap>E !gapinput@StructureDescription( AXS20 );E
  [ "D20", "D10", "D10", "C2 x (C5 : C4)" ]
  !gapprompt@gap>E !gapinput@I20 := Range( Up2DimensionalGroup( AXS20 ) );;E
  !gapprompt@gap>E !gapinput@genI20 := GeneratorsOfGroup( I20 );           E
  [ ^(1,2,3,4,5,6,7,8,9,10), ^(2,10)(3,9)(4,8)(5,7) ]
  !gapprompt@gap>E !gapinput@xpi := CrossedPairing( AXS20 );;E
  !gapprompt@gap>E !gapinput@xpi( genI20[1], genI20[2] );E
  (1,9,7,5,3)(2,10,8,6,4)
  
\end{Verbatim}
 

\subsection{\textcolor{Chapter }{CrossedSquareByPullback}}
\logpage{[ 8, 2, 7 ]}\nobreak
\hyperdef{L}{X839E065783795CB8}{}
{\noindent\textcolor{FuncColor}{$\triangleright$\enspace\texttt{CrossedSquareByPullback({\mdseries\slshape X1, X2})\index{CrossedSquareByPullback@\texttt{CrossedSquareByPullback}}
\label{CrossedSquareByPullback}
}\hfill{\scriptsize (operation)}}\\


 If crossed modules $\calX_1 = (\nu : N \to P)$ and $\calX_2 = (\mu : M \to P)$ have a common range $P$, let $L$ be the pullback of $\{\nu,\mu\}$. Then $N$ acts on $L$ by $(n,m)^{n'} = (n^{n'},m^{\nu n'})$, and $M$ acts on $L$ by $(n,m)^{m'} = (n^{\mu m'}, m^{m'})$. So $(\pi_1 : L \to N)$ and $(\pi_2 : L \to M)$ are crossed modules, where $\pi_1,\pi_2$ are the two projections. The crossed pairing is given by: 
\[ \boxtimes \;:\; N \times M \to L, \quad (n,m) \mapsto (n^{-1}n^{\mu m},
(m^{-1})^{\nu n}m) . \]
 

 The second example below uses the central extension crossed module \texttt{X12=(D12\texttt{\symbol{45}}{\textgreater}S3)} which was constructed in subsection (\texttt{XModByCentralExtension} (\ref{XModByCentralExtension})), with pullback group \texttt{D12xC2}. }

 

 
\begin{Verbatim}[commandchars=!@A,fontsize=\small,frame=single,label=Example]
  
  !gapprompt@gap>A !gapinput@dn := Down2DimensionalGroup( XSconj );;A
  !gapprompt@gap>A !gapinput@rt := Right2DimensionalGroup( XSconj );;A
  !gapprompt@gap>A !gapinput@XSP := CrossedSquareByPullback( dn, rt ); A
  [ (d10b x_d20 d10a) -> d10a ]
  [         |             |   ]
  [              d10b -> d20  ]
  !gapprompt@gap>A !gapinput@StructureDescription( XSP );                  A
  [ "C5", "D10", "D10", "D20" ]
  !gapprompt@gap>A !gapinput@XS12 := CrossedSquareByPullback( X12, X12 );; A
  !gapprompt@gap>A !gapinput@StructureDescription( XS12 );                  A
  [ "C2 x C2 x S3", "D12", "D12", "S3" ]
  !gapprompt@gap>A !gapinput@xp12 := CrossedPairing( XS12 );; A
  !gapprompt@gap>A !gapinput@xp12( (1,2,3,4,5,6), (2,6)(3,5) );A
  (1,5,3)(2,6,4)(7,11,9)(8,12,10)
  
\end{Verbatim}
 

\subsection{\textcolor{Chapter }{CrossedSquareByXModSplitting}}
\logpage{[ 8, 2, 8 ]}\nobreak
\hyperdef{L}{X7F5554907AF73190}{}
{\noindent\textcolor{FuncColor}{$\triangleright$\enspace\texttt{CrossedSquareByXModSplitting({\mdseries\slshape X0})\index{CrossedSquareByXModSplitting@\texttt{CrossedSquareByXModSplitting}}
\label{CrossedSquareByXModSplitting}
}\hfill{\scriptsize (attribute)}}\\
\noindent\textcolor{FuncColor}{$\triangleright$\enspace\texttt{CrossedPairingByPreImages({\mdseries\slshape X1, X2})\index{CrossedPairingByPreImages@\texttt{CrossedPairingByPreImages}}
\label{CrossedPairingByPreImages}
}\hfill{\scriptsize (operation)}}\\


 For $\calX = (\partial : S \to R)$ let $Q$ be the image of $\partial$. Then $\partial = \partial' \circ \iota$ where $\partial' : S \to Q$ and $\iota$ is the inclusion of $Q$ in $R$. The diagonal of the square is then the initial $\calX$, and the crossed pairing is given by commutators of preimages. 

 A particular case is when $S$ is an $R$\texttt{\symbol{45}}module $A$ and $\partial$ is the zero map. 
\[ \vcenter{\xymatrix{ & & S \ar[rr]^{\partial'} \ar[dd]_{\partial'}
\ar[rrdd]^{\partial} && Q \ar[dd]^{\iota} && A \ar[rr]^0 \ar[dd]_0 \ar[rrdd]^0
&& 1 \ar[dd]^{\iota} & \\ & & && & && \\ & & Q \ar[rr]_{\iota} && R && 1
\ar[rr]_{\iota} && R }} \]
 }

 

 
\begin{Verbatim}[commandchars=!@|,fontsize=\small,frame=single,label=Example]
  
  !gapprompt@gap>| !gapinput@k4 := Group( (1,2), (3,4) );;|
  !gapprompt@gap>| !gapinput@AX4 := XModByAutomorphismGroup( k4 );;|
  !gapprompt@gap>| !gapinput@X4 := Image( IsomorphismPermObject( AX4 ) );;|
  !gapprompt@gap>| !gapinput@XSS4 := CrossedSquareByXModSplitting( X4 );;|
  !gapprompt@gap>| !gapinput@StructureDescription( XSS4 );|
  [ "C2 x C2", "1", "1", "S3" ]
  !gapprompt@gap>| !gapinput@XSS20 := CrossedSquareByXModSplitting( X20 );;|
  !gapprompt@gap>| !gapinput@up20 := Up2DimensionalGroup( XSS20 );; |
  !gapprompt@gap>| !gapinput@Range( up20 ) = d10a; |
  true
  !gapprompt@gap>| !gapinput@SetName( Range( up20 ), "d10a" ); |
  !gapprompt@gap>| !gapinput@Name( XSS20 );;|
  !gapprompt@gap>| !gapinput@XSS20;|
  [d10a->d10a,d10a->d20]
  !gapprompt@gap>| !gapinput@xps := CrossedPairing( XSS20 );;|
  !gapprompt@gap>| !gapinput@xps( p1^2, p2 );|
  (1,7,3,9,5)(2,8,4,10,6)
  
\end{Verbatim}
 

\subsection{\textcolor{Chapter }{CrossedSquare}}
\logpage{[ 8, 2, 9 ]}\nobreak
\hyperdef{L}{X87FBE3CE87DC8CD5}{}
{\noindent\textcolor{FuncColor}{$\triangleright$\enspace\texttt{CrossedSquare({\mdseries\slshape args})\index{CrossedSquare@\texttt{CrossedSquare}}
\label{CrossedSquare}
}\hfill{\scriptsize (function)}}\\


 The function \texttt{CrossedSquare} may be used to call some of the constructions described in the previous
subsections. 
\begin{itemize}
\item  \texttt{CrossedSquare(X0)} calls \texttt{CrossedSquareByXModSplitting}. 
\item  \texttt{CrossedSquare(C0)} calls \texttt{CrossedSquareOfCat2Group}. 
\item  \texttt{CrossedSquare(X0,X1)} calls \texttt{CrossedSquareByPullback} when there is a common range. 
\item  \texttt{CrossedSquare(X0,X1)} calls \texttt{CrossedSquareByNormalXMod} when \texttt{X1} is normal in \texttt{X0} . 
\item  \texttt{CrossedSquare(L,M,N,P)} calls \texttt{CrossedSquareByNormalSubgroups}. 
\end{itemize}
 }

 

 
\begin{Verbatim}[commandchars=!@|,fontsize=\small,frame=single,label=Example]
  
  !gapprompt@gap>| !gapinput@diag := Diagonal2DimensionalGroup( AXS20 );|
  [d20->Aut(d20)]
  !gapprompt@gap>| !gapinput@XSdiag := CrossedSquare( diag );;      |
  !gapprompt@gap>| !gapinput@StructureDescription( XSdiag );  |
  [ "D20", "D10", "D10", "C2 x (C5 : C4)" ]
  
\end{Verbatim}
 

\subsection{\textcolor{Chapter }{Transpose3DimensionalGroup (for crossed squares)}}
\logpage{[ 8, 2, 10 ]}\nobreak
\hyperdef{L}{X7F94830681EA19BE}{}
{\noindent\textcolor{FuncColor}{$\triangleright$\enspace\texttt{Transpose3DimensionalGroup({\mdseries\slshape S0})\index{Transpose3DimensionalGroup@\texttt{Transpose3DimensionalGroup}!for crossed squares}
\label{Transpose3DimensionalGroup:for crossed squares}
}\hfill{\scriptsize (attribute)}}\\


 The \emph{transpose} of a crossed square $\calS$ is the crossed square $\tilde{\calS}$ obtained by interchanging $M$ with $N$, $\kappa$ with $\lambda$, and $\nu$ with $\mu$. The crossed pairing is given by 
\[ \tilde{\boxtimes} \;:\; M \times N \to L, \quad (m,n) \;\mapsto\;
m\,\tilde{\boxtimes}\,n := (n \boxtimes m)^{-1}~. \]
 }

 

 
\begin{Verbatim}[commandchars=!@A,fontsize=\small,frame=single,label=Example]
  
  !gapprompt@gap>A !gapinput@XStrans := Transpose3DimensionalGroup( XSconj ); A
  [  c5d -> d10b ]
  [   |      |   ]
  [ d10a -> d20  ]
  
  
\end{Verbatim}
 

\subsection{\textcolor{Chapter }{CentralQuotient (for crossed modules)}}
\logpage{[ 8, 2, 11 ]}\nobreak
\hyperdef{L}{X7C2647CB82DDD065}{}
{\noindent\textcolor{FuncColor}{$\triangleright$\enspace\texttt{CentralQuotient({\mdseries\slshape X0})\index{CentralQuotient@\texttt{CentralQuotient}!for crossed modules}
\label{CentralQuotient:for crossed modules}
}\hfill{\scriptsize (attribute)}}\\


 The central quotient of a crossed module $\calX = (\partial : S \to R)$ is the crossed square where: 
\begin{itemize}
\item  the left crossed module is $\calX$; 
\item  the right crossed module is the quotient $\calX/Z(\calX)$ (see \texttt{CentreXMod} (\ref{CentreXMod})); 
\item  the up and down homomorphisms are the natural homomorphisms onto the quotient
groups; 
\item  the crossed pairing $\boxtimes : (R \times F) \to S$, where $F = \Fix(\calX,S,R)$, is the displacement element $\boxtimes(r,Fs) = \langle r,s \rangle = (s^{-1})^rs\quad$ (see \texttt{Displacement} (\ref{Displacement}) and section \ref{sect-isoclinic-xmods}). 
\end{itemize}
 This is the special case of an intended function \texttt{CrossedSquareByCentralExtension} which has not yet been implemented. In the example \texttt{Xn7} $\unlhd$ \texttt{X24}, constructed in section \ref{sect-more-xmod-ops}. }

 

 
\begin{Verbatim}[commandchars=!@|,fontsize=\small,frame=single,label=Example]
  
  !gapprompt@gap>| !gapinput@pos7 := Position( ids, [ [12,2], [24,5] ] );;|
  !gapprompt@gap>| !gapinput@Xn7 := nsx[pos7];; |
  !gapprompt@gap>| !gapinput@IdGroup( Xn7 );|
  [ [ 12, 2 ], [ 24, 5 ] ]
  !gapprompt@gap>| !gapinput@IdGroup( CentreXMod( Xn7 ) );  |
  [ [ 4, 1 ], [ 4, 1 ] ]
  !gapprompt@gap>| !gapinput@CQXn7 := CentralQuotient( Xn7 );;|
  !gapprompt@gap>| !gapinput@StructureDescription( CQXn7 );|
  [ "C12", "C3", "C4 x S3", "S3" ]
  
\end{Verbatim}
 

\subsection{\textcolor{Chapter }{IsCrossedSquare}}
\logpage{[ 8, 2, 12 ]}\nobreak
\hyperdef{L}{X8645AA3686F126D5}{}
{\noindent\textcolor{FuncColor}{$\triangleright$\enspace\texttt{IsCrossedSquare({\mdseries\slshape obj})\index{IsCrossedSquare@\texttt{IsCrossedSquare}}
\label{IsCrossedSquare}
}\hfill{\scriptsize (property)}}\\
\noindent\textcolor{FuncColor}{$\triangleright$\enspace\texttt{IsPreCrossedSquare({\mdseries\slshape obj})\index{IsPreCrossedSquare@\texttt{IsPreCrossedSquare}}
\label{IsPreCrossedSquare}
}\hfill{\scriptsize (property)}}\\
\noindent\textcolor{FuncColor}{$\triangleright$\enspace\texttt{Is3dObject({\mdseries\slshape obj})\index{Is3dObject@\texttt{Is3dObject}}
\label{Is3dObject}
}\hfill{\scriptsize (property)}}\\
\noindent\textcolor{FuncColor}{$\triangleright$\enspace\texttt{IsPerm3dObject({\mdseries\slshape obj})\index{IsPerm3dObject@\texttt{IsPerm3dObject}}
\label{IsPerm3dObject}
}\hfill{\scriptsize (property)}}\\
\noindent\textcolor{FuncColor}{$\triangleright$\enspace\texttt{IsPc3dObject({\mdseries\slshape obj})\index{IsPc3dObject@\texttt{IsPc3dObject}}
\label{IsPc3dObject}
}\hfill{\scriptsize (property)}}\\
\noindent\textcolor{FuncColor}{$\triangleright$\enspace\texttt{IsFp3dObject({\mdseries\slshape obj})\index{IsFp3dObject@\texttt{IsFp3dObject}}
\label{IsFp3dObject}
}\hfill{\scriptsize (property)}}\\


 These are the basic properties for 3d\texttt{\symbol{45}}groups, and crossed
squares in particular. }

 

\subsection{\textcolor{Chapter }{Up2DimensionalGroup (for crossed squares)}}
\logpage{[ 8, 2, 13 ]}\nobreak
\hyperdef{L}{X828CFC5A83097189}{}
{\noindent\textcolor{FuncColor}{$\triangleright$\enspace\texttt{Up2DimensionalGroup({\mdseries\slshape XS})\index{Up2DimensionalGroup@\texttt{Up2DimensionalGroup}!for crossed squares}
\label{Up2DimensionalGroup:for crossed squares}
}\hfill{\scriptsize (attribute)}}\\
\noindent\textcolor{FuncColor}{$\triangleright$\enspace\texttt{Left2DimensionalGroup({\mdseries\slshape XS})\index{Left2DimensionalGroup@\texttt{Left2DimensionalGroup}!for crossed squares}
\label{Left2DimensionalGroup:for crossed squares}
}\hfill{\scriptsize (attribute)}}\\
\noindent\textcolor{FuncColor}{$\triangleright$\enspace\texttt{Down2DimensionalGroup({\mdseries\slshape XS})\index{Down2DimensionalGroup@\texttt{Down2DimensionalGroup}!for crossed squares}
\label{Down2DimensionalGroup:for crossed squares}
}\hfill{\scriptsize (attribute)}}\\
\noindent\textcolor{FuncColor}{$\triangleright$\enspace\texttt{Right2DimensionalGroup({\mdseries\slshape XS})\index{Right2DimensionalGroup@\texttt{Right2DimensionalGroup}!for crossed squares}
\label{Right2DimensionalGroup:for crossed squares}
}\hfill{\scriptsize (attribute)}}\\
\noindent\textcolor{FuncColor}{$\triangleright$\enspace\texttt{CrossDiagonalActions({\mdseries\slshape XS})\index{CrossDiagonalActions@\texttt{CrossDiagonalActions}!for crossed squares}
\label{CrossDiagonalActions:for crossed squares}
}\hfill{\scriptsize (attribute)}}\\
\noindent\textcolor{FuncColor}{$\triangleright$\enspace\texttt{Diagonal2DimensionalGroup({\mdseries\slshape XS})\index{Diagonal2DimensionalGroup@\texttt{Diagonal2DimensionalGroup}!for crossed squares}
\label{Diagonal2DimensionalGroup:for crossed squares}
}\hfill{\scriptsize (attribute)}}\\
\noindent\textcolor{FuncColor}{$\triangleright$\enspace\texttt{Name({\mdseries\slshape S0})\index{Name@\texttt{Name}}
\label{Name}
}\hfill{\scriptsize (method)}}\\


 These are the basic attributes of a crossed square $\calS$. The six objects used in the construction of $\calS$ are the four crossed modules (2d\texttt{\symbol{45}}groups) on the sides of
the square (up; left; right and down); the diagonal action of $P$ on $L$; and the crossed pairing $\{M,N\} \to L$ (see the next subsection). The diagonal crossed module $(L \to P)$ is an additional attribute. 

 }

 

 
\begin{Verbatim}[commandchars=!@|,fontsize=\small,frame=single,label=Example]
  
  !gapprompt@gap>| !gapinput@Up2DimensionalGroup( XSconj );|
  [c5d->d10a]
  !gapprompt@gap>| !gapinput@Right2DimensionalGroup( XSact );|
  Actor[d10a->d20]
  !gapprompt@gap>| !gapinput@Name( XSconj ); |
  "[c5d->d10a,d10b->d20]"
  !gapprompt@gap>| !gapinput@cross1 := CrossDiagonalActions( XSconj )[1];; |
  !gapprompt@gap>| !gapinput@gensa := GeneratorsOfGroup( d10a );; |
  !gapprompt@gap>| !gapinput@gensb := GeneratorsOfGroup( d10a );; |
  !gapprompt@gap>| !gapinput@act1 := ImageElm( cross1, gensb[1] );;|
  !gapprompt@gap>| !gapinput@gensa[2]; ImageElm( act1, gensa[2] );|
  (2,10)(3,9)(4,8)(5,7)
  (1,5)(2,4)(6,10)(7,9)
  
\end{Verbatim}
 

\subsection{\textcolor{Chapter }{IsSymmetric3DimensionalGroup}}
\logpage{[ 8, 2, 14 ]}\nobreak
\hyperdef{L}{X7F81DA90820F4405}{}
{\noindent\textcolor{FuncColor}{$\triangleright$\enspace\texttt{IsSymmetric3DimensionalGroup({\mdseries\slshape obj})\index{IsSymmetric3DimensionalGroup@\texttt{IsSymmetric3DimensionalGroup}}
\label{IsSymmetric3DimensionalGroup}
}\hfill{\scriptsize (property)}}\\
\noindent\textcolor{FuncColor}{$\triangleright$\enspace\texttt{IsAbelian3DimensionalGroup({\mdseries\slshape obj})\index{IsAbelian3DimensionalGroup@\texttt{IsAbelian3DimensionalGroup}}
\label{IsAbelian3DimensionalGroup}
}\hfill{\scriptsize (property)}}\\
\noindent\textcolor{FuncColor}{$\triangleright$\enspace\texttt{IsTrivialAction3DimensionalGroup({\mdseries\slshape obj})\index{IsTrivialAction3DimensionalGroup@\texttt{IsTrivialAction3DimensionalGroup}}
\label{IsTrivialAction3DimensionalGroup}
}\hfill{\scriptsize (property)}}\\
\noindent\textcolor{FuncColor}{$\triangleright$\enspace\texttt{IsNormalSub3DimensionalGroup({\mdseries\slshape obj})\index{IsNormalSub3DimensionalGroup@\texttt{IsNormalSub3DimensionalGroup}}
\label{IsNormalSub3DimensionalGroup}
}\hfill{\scriptsize (property)}}\\
\noindent\textcolor{FuncColor}{$\triangleright$\enspace\texttt{IsCentralExtension3DimensionalGroup({\mdseries\slshape obj})\index{IsCentralExtension3DimensionalGroup@\texttt{IsCentralExtension3DimensionalGroup}}
\label{IsCentralExtension3DimensionalGroup}
}\hfill{\scriptsize (property)}}\\
\noindent\textcolor{FuncColor}{$\triangleright$\enspace\texttt{IsAutomorphismGroup3DimensionalGroup({\mdseries\slshape obj})\index{IsAutomorphismGroup3DimensionalGroup@\texttt{IsAutomorphism}\-\texttt{Group3}\-\texttt{Dimensional}\-\texttt{Group}}
\label{IsAutomorphismGroup3DimensionalGroup}
}\hfill{\scriptsize (property)}}\\


 These are further properties for 3d\texttt{\symbol{45}}groups, and crossed
squares in particular. A 3d\texttt{\symbol{45}}group is \emph{symmetric} if its \texttt{Up2DimensionalGroup} is equal to its \texttt{Left2DimensionalGroup}. }

 

\subsection{\textcolor{Chapter }{Crossed Pairing}}
\logpage{[ 8, 2, 15 ]}\nobreak
\hyperdef{L}{X7AE671C7798F99FD}{}
{\noindent\textcolor{FuncColor}{$\triangleright$\enspace\texttt{Crossed Pairing({\mdseries\slshape XS})\index{Crossed Pairing@\texttt{Crossed Pairing}}
\label{Crossed Pairing}
}\hfill{\scriptsize (attribute)}}\\


 Crossed pairings have been implemented as functions of two variables, $\{M,N\} \to L$. (Up until version 2.92 a more complicated two\texttt{\symbol{45}}variable
mapping was used.) }

 

 The first example shows the crossed pairing in the crossed square \texttt{XSconj}. 
\begin{Verbatim}[commandchars=!@|,fontsize=\small,frame=single,label=Example]
  
  !gapprompt@gap>| !gapinput@xp := CrossedPairing( XSconj );;|
  !gapprompt@gap>| !gapinput@xp( (1,6)(2,5)(3,4)(7,10)(8,9), (1,5)(2,4)(6,9)(7,8) );|
  (1,7,8,5,3)(2,9,10,6,4)
  
\end{Verbatim}
 The second example shows how to construct a crossed pairing. 
\begin{Verbatim}[commandchars=!@|,fontsize=\small,frame=single,label=Example]
  
  !gapprompt@gap>| !gapinput@F := FreeGroup(1);;|
  !gapprompt@gap>| !gapinput@x := GeneratorsOfGroup(F)[1];;|
  !gapprompt@gap>| !gapinput@z := GroupHomomorphismByImages( F, F, [x], [x^0] );;|
  !gapprompt@gap>| !gapinput@id := GroupHomomorphismByImages( F, F, [x], [x] );;|
  !gapprompt@gap>| !gapinput@h := function(n,m)|
  !gapprompt@>| !gapinput@            return x^(ExponentSumWord(n,x)*ExponentSumWord(m,x));|
  !gapprompt@>| !gapinput@        end;;|
  !gapprompt@gap>| !gapinput@h( x^3, x^4 );|
  f1^12
  !gapprompt@gap>| !gapinput@A := AutomorphismGroup( F );;|
  !gapprompt@gap>| !gapinput@a := GeneratorsOfGroup(A)[1];;|
  !gapprompt@gap>| !gapinput@act := GroupHomomorphismByImages( F, A, [x], [a^2] );;|
  !gapprompt@gap>| !gapinput@X0 := XModByBoundaryAndAction( z, act );;|
  !gapprompt@gap>| !gapinput@X1 := XModByBoundaryAndAction( id, act );;|
  !gapprompt@gap>| !gapinput@XSF := PreCrossedSquareByPreXMods( X0, X0, X1, X1, X0, h );; |
  !gapprompt@gap>| !gapinput@IsCrossedSquare( XSF ); |
  true
  
\end{Verbatim}
 }

 
\section{\textcolor{Chapter }{Substructures of Crossed Squares}}\logpage{[ 8, 3, 0 ]}
\hyperdef{L}{X79A59CED7C69BF18}{}
{
 To specify a sub\texttt{\symbol{45}}crossed square of $\calS$ we may specify four subgroups of the four groups in $\calS$. If these define five sub\texttt{\symbol{45}}crossed modules, and these comc
bine together satisfactorily, then a crossed square is produced. 

\subsection{\textcolor{Chapter }{SubCrossedSquare}}
\logpage{[ 8, 3, 1 ]}\nobreak
\hyperdef{L}{X83F16E94857407F3}{}
{\noindent\textcolor{FuncColor}{$\triangleright$\enspace\texttt{SubCrossedSquare({\mdseries\slshape XS, ul, ur, dl, dr})\index{SubCrossedSquare@\texttt{SubCrossedSquare}}
\label{SubCrossedSquare}
}\hfill{\scriptsize (operation)}}\\
\noindent\textcolor{FuncColor}{$\triangleright$\enspace\texttt{IsSubCrossedSquare({\mdseries\slshape XS, subXS})\index{IsSubCrossedSquare@\texttt{IsSubCrossedSquare}}
\label{IsSubCrossedSquare}
}\hfill{\scriptsize (operation)}}\\
\noindent\textcolor{FuncColor}{$\triangleright$\enspace\texttt{SubPreCrossedSquare({\mdseries\slshape XS, ul, ur, dl, dr})\index{SubPreCrossedSquare@\texttt{SubPreCrossedSquare}}
\label{SubPreCrossedSquare}
}\hfill{\scriptsize (operation)}}\\
\noindent\textcolor{FuncColor}{$\triangleright$\enspace\texttt{IsSubPreCrossedSquare({\mdseries\slshape XS, subXS})\index{IsSubPreCrossedSquare@\texttt{IsSubPreCrossedSquare}}
\label{IsSubPreCrossedSquare}
}\hfill{\scriptsize (operation)}}\\


 }

 

 
\begin{Verbatim}[commandchars=!@A,fontsize=\small,frame=single,label=Example]
  
  !gapprompt@gap>A !gapinput@Display( XS20 ); A
  [  c5d -> d10a ]
  [   |      |   ]
  [ d10b -> d20  ]
  !gapprompt@gap>A !gapinput@p5 := p1^5*p2;;  [p2,p12,p5];A
  [ (2,10)(3,9)(4,8)(5,7), (1,10)(2,9)(3,8)(4,7)(5,6), 
    (1,6)(2,5)(3,4)(7,10)(8,9) ]
  !gapprompt@gap>A !gapinput@L1 := TrivialSubgroup( c5d );;A
  !gapprompt@gap>A !gapinput@M1 := Subgroup( d10a, [ p2 ] );;A
  !gapprompt@gap>A !gapinput@N1 := Subgroup( d10b, [ p5 ] );;A
  !gapprompt@gap>A !gapinput@P1 := Subgroup( d20, [ p2, p5 ] );;A
  !gapprompt@gap>A !gapinput@sub20 := SubCrossedSquare( XS20, L1, M1, N1, P1 );A
  crossed square with crossed modules:
        up = [Group( () ) -> Group( [ ( 2,10)( 3, 9)( 4, 8)( 5, 7) ] )]
      left = [Group( () ) -> Group( [ ( 1, 6)( 2, 5)( 3, 4)( 7,10)( 8, 9) ] )]
     right = [Group( [ ( 2,10)( 3, 9)( 4, 8)( 5, 7) ] ) -> Group( 
  [ ( 2,10)( 3, 9)( 4, 8)( 5, 7), ( 1, 6)( 2, 5)( 3, 4)( 7,10)( 8, 9) ] )]
      down = [Group( [ ( 1, 6)( 2, 5)( 3, 4)( 7,10)( 8, 9) ] ) -> Group( 
  [ ( 2,10)( 3, 9)( 4, 8)( 5, 7), ( 1, 6)( 2, 5)( 3, 4)( 7,10)( 8, 9) ] )]
  
  !gapprompt@gap>A !gapinput@sxp := CrossedPairing( sub20 );;A
  !gapprompt@gap>A !gapinput@sxp( p5, p2 );A
  ()
  
\end{Verbatim}
 

\subsection{\textcolor{Chapter }{TrivialSub3DimensionalGroup}}
\logpage{[ 8, 3, 2 ]}\nobreak
\hyperdef{L}{X7D0D7F787D438D9A}{}
{\noindent\textcolor{FuncColor}{$\triangleright$\enspace\texttt{TrivialSub3DimensionalGroup({\mdseries\slshape 3dgp})\index{TrivialSub3DimensionalGroup@\texttt{TrivialSub3DimensionalGroup}}
\label{TrivialSub3DimensionalGroup}
}\hfill{\scriptsize (operation)}}\\
\noindent\textcolor{FuncColor}{$\triangleright$\enspace\texttt{TrivialSubCrossedSquare({\mdseries\slshape XS})\index{TrivialSubCrossedSquare@\texttt{TrivialSubCrossedSquare}}
\label{TrivialSubCrossedSquare}
}\hfill{\scriptsize (attribute)}}\\
\noindent\textcolor{FuncColor}{$\triangleright$\enspace\texttt{TrivialSubPreCrossedSquare({\mdseries\slshape XS})\index{TrivialSubPreCrossedSquare@\texttt{TrivialSubPreCrossedSquare}}
\label{TrivialSubPreCrossedSquare}
}\hfill{\scriptsize (attribute)}}\\


 A special case of the previous operation is when all the subgroups are
trivial. }

 

 
\begin{Verbatim}[commandchars=!@|,fontsize=\small,frame=single,label=Example]
  
  !gapprompt@gap>| !gapinput@TC2ab := TrivialSubCat2Group( C2ab );|
  (pre-)cat2-group with generating (pre-)cat1-groups:
  1 : [Group( () ) => Group( () )]
  2 : [Group( () ) => Group( () )]
  
\end{Verbatim}
 }

 
\section{\textcolor{Chapter }{Morphisms of crossed squares}}\logpage{[ 8, 4, 0 ]}
\hyperdef{L}{X78A79A7E85128C7B}{}
{
 \index{morphism of 3d\texttt{\symbol{45}}group} \index{crossed square morphism} \index{3d\texttt{\symbol{45}}mapping} This section describes an initial implementation of morphisms of
(pre\texttt{\symbol{45}})crossed squares. 

\subsection{\textcolor{Chapter }{CrossedSquareMorphism}}
\logpage{[ 8, 4, 1 ]}\nobreak
\hyperdef{L}{X83E733547A14FD61}{}
{\noindent\textcolor{FuncColor}{$\triangleright$\enspace\texttt{CrossedSquareMorphism({\mdseries\slshape args})\index{CrossedSquareMorphism@\texttt{CrossedSquareMorphism}}
\label{CrossedSquareMorphism}
}\hfill{\scriptsize (function)}}\\
\noindent\textcolor{FuncColor}{$\triangleright$\enspace\texttt{CrossedSquareMorphismByXModMorphisms({\mdseries\slshape src, rng, mors})\index{CrossedSquareMorphismByXModMorphisms@\texttt{Crossed}\-\texttt{Square}\-\texttt{Morphism}\-\texttt{By}\-\texttt{X}\-\texttt{Mod}\-\texttt{Morphisms}}
\label{CrossedSquareMorphismByXModMorphisms}
}\hfill{\scriptsize (operation)}}\\
\noindent\textcolor{FuncColor}{$\triangleright$\enspace\texttt{CrossedSquareMorphismByGroupHomomorphisms({\mdseries\slshape src, rng, homs})\index{CrossedSquareMorphismByGroupHomomorphisms@\texttt{Crossed}\-\texttt{Square}\-\texttt{Morphism}\-\texttt{By}\-\texttt{Group}\-\texttt{Homomorphisms}}
\label{CrossedSquareMorphismByGroupHomomorphisms}
}\hfill{\scriptsize (operation)}}\\
\noindent\textcolor{FuncColor}{$\triangleright$\enspace\texttt{PreCrossedSquareMorphismByPreXModMorphisms({\mdseries\slshape src, rng, mors})\index{PreCrossedSquareMorphismByPreXModMorphisms@\texttt{Pre}\-\texttt{Crossed}\-\texttt{Square}\-\texttt{Morphism}\-\texttt{By}\-\texttt{Pre}\-\texttt{X}\-\texttt{Mod}\-\texttt{Morphisms}}
\label{PreCrossedSquareMorphismByPreXModMorphisms}
}\hfill{\scriptsize (operation)}}\\
\noindent\textcolor{FuncColor}{$\triangleright$\enspace\texttt{PreCrossedSquareMorphismByGroupHomomorphisms({\mdseries\slshape src, rng, homs})\index{PreCrossedSquareMorphismByGroupHomomorphisms@\texttt{Pre}\-\texttt{Crossed}\-\texttt{Square}\-\texttt{Morphism}\-\texttt{By}\-\texttt{Group}\-\texttt{Homomorphisms}}
\label{PreCrossedSquareMorphismByGroupHomomorphisms}
}\hfill{\scriptsize (operation)}}\\


 }

 

\subsection{\textcolor{Chapter }{Source}}
\logpage{[ 8, 4, 2 ]}\nobreak
\hyperdef{L}{X7DE8173F80E07AB1}{}
{\noindent\textcolor{FuncColor}{$\triangleright$\enspace\texttt{Source({\mdseries\slshape map})\index{Source@\texttt{Source}}
\label{Source}
}\hfill{\scriptsize (attribute)}}\\
\noindent\textcolor{FuncColor}{$\triangleright$\enspace\texttt{Range({\mdseries\slshape map})\index{Range@\texttt{Range}}
\label{Range}
}\hfill{\scriptsize (attribute)}}\\
\noindent\textcolor{FuncColor}{$\triangleright$\enspace\texttt{Up2DimensionalMorphism({\mdseries\slshape map})\index{Up2DimensionalMorphism@\texttt{Up2DimensionalMorphism}}
\label{Up2DimensionalMorphism}
}\hfill{\scriptsize (attribute)}}\\
\noindent\textcolor{FuncColor}{$\triangleright$\enspace\texttt{Left2DimensionalMorphism({\mdseries\slshape map})\index{Left2DimensionalMorphism@\texttt{Left2DimensionalMorphism}}
\label{Left2DimensionalMorphism}
}\hfill{\scriptsize (attribute)}}\\
\noindent\textcolor{FuncColor}{$\triangleright$\enspace\texttt{Down2DimensionalMorphism({\mdseries\slshape map})\index{Down2DimensionalMorphism@\texttt{Down2DimensionalMorphism}}
\label{Down2DimensionalMorphism}
}\hfill{\scriptsize (attribute)}}\\
\noindent\textcolor{FuncColor}{$\triangleright$\enspace\texttt{Right2DimensionalMorphism({\mdseries\slshape map})\index{Right2DimensionalMorphism@\texttt{Right2DimensionalMorphism}}
\label{Right2DimensionalMorphism}
}\hfill{\scriptsize (attribute)}}\\


 Morphisms of \texttt{3dObjects} are implemented as \texttt{3dMappings}. These have a pair of 3d\texttt{\symbol{45}}groups as source and range,
together with four 2d\texttt{\symbol{45}}morphisms mapping between the four
pairs of crossed modules on the four sides of the squares. These functions
return \texttt{fail} when invalid data is supplied. }

 

\subsection{\textcolor{Chapter }{IsCrossedSquareMorphism}}
\logpage{[ 8, 4, 3 ]}\nobreak
\hyperdef{L}{X8284240C7B9BB783}{}
{\noindent\textcolor{FuncColor}{$\triangleright$\enspace\texttt{IsCrossedSquareMorphism({\mdseries\slshape map})\index{IsCrossedSquareMorphism@\texttt{IsCrossedSquareMorphism}}
\label{IsCrossedSquareMorphism}
}\hfill{\scriptsize (property)}}\\
\noindent\textcolor{FuncColor}{$\triangleright$\enspace\texttt{IsPreCrossedSquareMorphism({\mdseries\slshape map})\index{IsPreCrossedSquareMorphism@\texttt{IsPreCrossedSquareMorphism}}
\label{IsPreCrossedSquareMorphism}
}\hfill{\scriptsize (property)}}\\
\noindent\textcolor{FuncColor}{$\triangleright$\enspace\texttt{IsBijective({\mdseries\slshape mor})\index{IsBijective@\texttt{IsBijective}}
\label{IsBijective}
}\hfill{\scriptsize (method)}}\\
\noindent\textcolor{FuncColor}{$\triangleright$\enspace\texttt{IsEndomorphism3dObject({\mdseries\slshape mor})\index{IsEndomorphism3dObject@\texttt{IsEndomorphism3dObject}}
\label{IsEndomorphism3dObject}
}\hfill{\scriptsize (property)}}\\
\noindent\textcolor{FuncColor}{$\triangleright$\enspace\texttt{IsAutomorphism3dObject({\mdseries\slshape mor})\index{IsAutomorphism3dObject@\texttt{IsAutomorphism3dObject}}
\label{IsAutomorphism3dObject}
}\hfill{\scriptsize (property)}}\\


 A morphism \texttt{mor} between two pre\texttt{\symbol{45}}crossed squares $\calS_{1}$ and $\calS_{2}$ consists of four crossed module morphisms \texttt{Up2DimensionalMorphism(mor)}, mapping the \texttt{Up2DimensionalGroup} of $\calS_1$ to that of $\calS_2$, \texttt{Left2DimensionalMorphism(mor)}, \texttt{Right2DimensionalMorphism(mor)} and \texttt{Down2DimensionalMorphism(mor)}. These four morphisms are required to commute with the four boundary maps and
to preserve the rest of the structure. The current version of \texttt{IsCrossedSquareMorphism} does not perform all the required checks. }

 

 
\begin{Verbatim}[commandchars=!@|,fontsize=\small,frame=single,label=Example]
  
  !gapprompt@gap>| !gapinput@ad20 := GroupHomomorphismByImages( d20, d20, [p1,p2], [p1,p2^p1] );;|
  !gapprompt@gap>| !gapinput@ad10a := GroupHomomorphismByImages( d10a, d10a, [p1^2,p2], [p1^2,p2^p1] );;|
  !gapprompt@gap>| !gapinput@ad10b := GroupHomomorphismByImages( d10b, d10b, [p1^2,p12], [p1^2,p12^p1] );;|
  !gapprompt@gap>| !gapinput@idc5d := IdentityMapping( c5d );;|
  !gapprompt@gap>| !gapinput@up := Up2DimensionalGroup( XSconj );;|
  !gapprompt@gap>| !gapinput@lt := Left2DimensionalGroup( XSconj );; |
  !gapprompt@gap>| !gapinput@rt := Right2DimensionalGroup( XSconj );; |
  !gapprompt@gap>| !gapinput@dn := Down2DimensionalGroup( XSconj );; |
  !gapprompt@gap>| !gapinput@mup := XModMorphismByGroupHomomorphisms( up, up, idc5d, ad10a );|
  [[c5d->d10a] => [c5d->d10a]]
  !gapprompt@gap>| !gapinput@mlt := XModMorphismByGroupHomomorphisms( lt, lt, idc5d, ad10b );|
  [[c5d->d10b] => [c5d->d10b]]
  !gapprompt@gap>| !gapinput@mrt := XModMorphismByGroupHomomorphisms( rt, rt, ad10a, ad20 );|
  [[d10a->d20] => [d10a->d20]]
  !gapprompt@gap>| !gapinput@mdn := XModMorphismByGroupHomomorphisms( dn, dn, ad10b, ad20 );|
  [[d10b->d20] => [d10b->d20]]
  !gapprompt@gap>| !gapinput@autoconj := CrossedSquareMorphism( XSconj, XSconj, [mup,mlt,mrt,mdn] );; |
  !gapprompt@gap>| !gapinput@ord := Order( autoconj );;|
  !gapprompt@gap>| !gapinput@Display( autoconj );|
  Morphism of crossed squares :- 
  : Source = [c5d->d10a,d10b->d20]
  : Range = [c5d->d10a,d10b->d20]
  :     order = 5
  :    up-left: [ [ ( 1, 3, 5, 7, 9)( 2, 4, 6, 8,10) ], 
    [ ( 1, 3, 5, 7, 9)( 2, 4, 6, 8,10) ] ]
  :   up-right: 
  [ [ ( 1, 3, 5, 7, 9)( 2, 4, 6, 8,10), ( 2,10)( 3, 9)( 4, 8)( 5, 7) ], 
    [ ( 1, 3, 5, 7, 9)( 2, 4, 6, 8,10), ( 1, 3)( 4,10)( 5, 9)( 6, 8) ] ]
  :  down-left: 
  [ [ ( 1, 3, 5, 7, 9)( 2, 4, 6, 8,10), ( 1,10)( 2, 9)( 3, 8)( 4, 7)( 5, 6) ], 
    [ ( 1, 3, 5, 7, 9)( 2, 4, 6, 8,10), ( 1, 2)( 3,10)( 4, 9)( 5, 8)( 6, 7) ] ]
  : down-right: 
  [ [ ( 1, 2, 3, 4, 5, 6, 7, 8, 9,10), ( 2,10)( 3, 9)( 4, 8)( 5, 7) ], 
    [ ( 1, 2, 3, 4, 5, 6, 7, 8, 9,10), ( 1, 3)( 4,10)( 5, 9)( 6, 8) ] ]
  !gapprompt@gap>| !gapinput@IsAutomorphismHigherDimensionalDomain( autoconj );|
  true
  !gapprompt@gap>| !gapinput@kpo := KnownPropertiesOfObject( autoconj );;|
  !gapprompt@gap>| !gapinput@Set( kpo );|
  [ "CanEasilyCompareElements", "CanEasilySortElements", 
    "IsAutomorphismHigherDimensionalDomain", "IsCrossedSquareMorphism", 
    "IsEndomorphismHigherDimensionalDomain", "IsInjective", 
    "IsPreCrossedSquareMorphism", "IsSingleValued", "IsSurjective", "IsTotal" ]
  
\end{Verbatim}
 

\subsection{\textcolor{Chapter }{InclusionMorphismHigherDimensionalDomains}}
\logpage{[ 8, 4, 4 ]}\nobreak
\hyperdef{L}{X8614E38E7A67690E}{}
{\noindent\textcolor{FuncColor}{$\triangleright$\enspace\texttt{InclusionMorphismHigherDimensionalDomains({\mdseries\slshape obj, sub})\index{InclusionMorphismHigherDimensionalDomains@\texttt{Inclusion}\-\texttt{Morphism}\-\texttt{Higher}\-\texttt{Dimensional}\-\texttt{Domains}}
\label{InclusionMorphismHigherDimensionalDomains}
}\hfill{\scriptsize (operation)}}\\


 }

 }

         
\section{\textcolor{Chapter }{Definitions and constructions for cat$^2$\texttt{\symbol{45}}groups and their morphisms }}\label{sect-cat2-definitions}
\logpage{[ 8, 5, 0 ]}
\hyperdef{L}{X86D5AA247B64ED51}{}
{
  \index{cat$^2$\texttt{\symbol{45}}group} We give here three equivalent definitions of cat$^2$\texttt{\symbol{45}}groups. When we come to define cat$^n$\texttt{\symbol{45}}groups we shall give a similar set of definitions. 

 Firstly, we take the definition of a cat$^2$\texttt{\symbol{45}}group from Section 5 of Brown and Loday \cite{brow:lod}, suitably modified. A cat$^2$\texttt{\symbol{45}}group $\calC = (C_{[2]},C_{\{2\}},C_{\{1\}},C_{\emptyset})$ comprises four groups (one for each of the subsets of $[2]$) and $15$ homomorphisms, as shown in the following diagram: 
\[ \vcenter{\xymatrix{ & C_{[2]} \ar[ddd] <-1.2ex> \ar[ddd]
<-2.0ex>_{\ddot{t}_2,\ddot{h}_2} \ar[rrr] <+1.2ex> \ar[rrr]
<+2.0ex>^{\ddot{t}_1,\ddot{h}_1} \ar[dddrrr] <-0.2ex> \ar[dddrrr]
<-1.0ex>_(0.55){t_{[2]},h_{[2]}} &&& C_{\{2\}} \ar[lll]^{\ddot{e}_1}
\ar[ddd]<+1.2ex> \ar[ddd] <+2.0ex>^{\dot{t}_2,\dot{h}_2} \\ \calC \quad =
\quad & &&& \\ & &&& \\ & C_{\{1\}} \ar[uuu]_{\ddot{e}_2} \ar[rrr] <-1.2ex>
\ar[rrr] <-2.0ex>_{\dot{t}_1,\dot{h}_1} &&& C_{\emptyset} \ar[uuu]^{\dot{e}_2}
\ar[lll]_{\dot{e}_1} \ar[uuulll] <-1.0ex>_{e_{[2]}} \\ }} \]
 The following axioms are satisfied by these homomorphisms: 
\begin{itemize}
\item  the four sides of the square (up, left, right, down) are cat$^1$\texttt{\symbol{45}}groups, denoted $\ddot{\calC}_1, \ddot{\calC}_2, \dot{\calC}_1, \dot{\calC}_2$; 
\item  $ \dot{t}_1\circ\ddot{h}_2 = \dot{h}_2\circ\ddot{t}_1, ~
\dot{t}_2\circ\ddot{h}_1 = \dot{h}_1\circ\ddot{t}_2, ~ \dot{e}_1\circ\dot{t}_2
= \ddot{t}_2\circ\ddot{e}_1, ~ \dot{e}_2\circ\dot{t}_1 =
\ddot{t}_1\circ\ddot{e}_2, ~ \dot{e}_1\circ\dot{h}_2 =
\ddot{h}_2\circ\ddot{e}_1, ~ \dot{e}_2\circ\dot{h}_1 =
\ddot{h}_1\circ\ddot{e}_2; $ 
\item  $ \dot{t}_1\circ\ddot{t}_2 = \dot{t}_2\circ\ddot{t}_1 = t_{[2]}, ~
\dot{h}_1\circ\ddot{h}_2 = \dot{h}_2\circ\ddot{h}_1 = h_{[2]}, ~
\dot{e}_1\circ\ddot{e}_2 = \dot{e}_2\circ\ddot{e}_1 = e_{[2]}, $ making the diagonal a pre\texttt{\symbol{45}}cat$^1$\texttt{\symbol{45}}group $(e_{[2]}; t_{[2]}, h_{[2]} : C_{[2]} \to C_{\emptyset})$. 
\end{itemize}
 It follows from these identities that $(\ddot{t}_1,\dot{t}_1),\,(\ddot{h}_1,\dot{h}_1)$ and $(\ddot{e}_1,\dot{e}_1)$ are morphisms of cat$^1$\texttt{\symbol{45}}groups, and similarly in the vertical direction. 

 Secondly, we give the simplest of the three definitions, adapted from
Ellis\texttt{\symbol{45}}Steiner \cite{ell:st}. A cat$^2$\texttt{\symbol{45}}group $\calC$ consists of groups $G, R_1,R_2$ and six homomorphisms $t_1,h_1 : G \to R_2,~ e_1 : R_2 \to G,~ t_2,h_2 : G \to R_1,~ e_2 : R_1 \to G$, satisfying the following axioms for all $1 \leqslant i \leqslant 2$, 
\begin{itemize}
\item  $ (t_i \circ e_i)r = r,~ (h_i \circ e_i)r = r,~ \forall r \in R_{[2] \setminus
\{i\}}, \quad [\ker t_i, \ker h_i] = 1, $ 
\item  $ (e_1 \circ t_1) \circ (e_2 \circ t_2) = (e_2 \circ t_2) \circ (e_1 \circ t_1),
\quad (e_1 \circ h_1) \circ (e_2 \circ h_2) = (e_2 \circ h_2) \circ (e_1 \circ
h_1), $ 
\item  $ (e_1 \circ t_1) \circ (e_2 \circ h_2) = (e_2 \circ h_2) \circ (e_1 \circ t_1),
\quad (e_2 \circ t_2) \circ (e_1 \circ h_1) = (e_1 \circ h_1) \circ (e_2 \circ
t_2). $ 
\end{itemize}
 

 Our third definition defines a cat$^2$\texttt{\symbol{45}}group as a "cat$^1$\texttt{\symbol{45}}group of cat$^1$\texttt{\symbol{45}}groups". A cat$^2$\texttt{\symbol{45}}group $\calC$ consists of two cat$^1$\texttt{\symbol{45}}groups $\calC_1 = (e_1;t_1,h_1 : G_1 \to R_1)$ and $\calC_2 = (e_2;t_2,h_2 : G_2 \to R_2)$ and cat$^1$\texttt{\symbol{45}}morphisms $t = (\ddot{t},\dot{t}),\; h = (\ddot{h},\dot{h}) : \calC_1 \to \calC_2,\; e =
(\ddot{e},\dot{e}) : \calC_2 \to \calC_1$, subject to the following conditions: 
\[ (t \circ e) ~\mbox{and}~ (h \circ e) ~\mbox{are the identity mapping on}~
\calC_2, \qquad [\ker t, \ker h] = \{ 1_{\calC_1} \}, \]
 where $\ker t = (\ker \ddot{t},\ \ker \dot{t})$, and similarly for $\ker h$. 

 In a recent paper \emph{Computing 3\texttt{\symbol{45}}Dimensional Groups : Crossed Squares and
Cat2\texttt{\symbol{45}}Groups}, by Arvasi, Odabas and Wensley \cite{AOW}, there are tables listing the numbers of isomorphism classes of cat$^2$\texttt{\symbol{45}}groups on groups of order at most $30$ {\textendash} a total of $1007$ cat$^2$\texttt{\symbol{45}}groups. 

 As with cat$^1$\texttt{\symbol{45}}groups, there is no requirement for the various tail and
head maps to be endomorphisms, but most of the cat$^2$\texttt{\symbol{45}}groups constructed by this package are of this type, so
that groups $C_{\{2\}},C_{\{1\}}$ and $C_{\emptyset}$ are all subgroups of $C_{[2]}$. 

\subsection{\textcolor{Chapter }{Cat2Group}}
\logpage{[ 8, 5, 1 ]}\nobreak
\label{cat2-group}
\hyperdef{L}{X849D845586F92444}{}
{\noindent\textcolor{FuncColor}{$\triangleright$\enspace\texttt{Cat2Group({\mdseries\slshape args})\index{Cat2Group@\texttt{Cat2Group}}
\label{Cat2Group}
}\hfill{\scriptsize (function)}}\\
\noindent\textcolor{FuncColor}{$\triangleright$\enspace\texttt{PreCat2Group({\mdseries\slshape args})\index{PreCat2Group@\texttt{PreCat2Group}}
\label{PreCat2Group}
}\hfill{\scriptsize (function)}}\\
\noindent\textcolor{FuncColor}{$\triangleright$\enspace\texttt{IsCat2Group({\mdseries\slshape C})\index{IsCat2Group@\texttt{IsCat2Group}}
\label{IsCat2Group}
}\hfill{\scriptsize (property)}}\\
\noindent\textcolor{FuncColor}{$\triangleright$\enspace\texttt{PreCat2GroupByPreCat1Groups({\mdseries\slshape L})\index{PreCat2GroupByPreCat1Groups@\texttt{PreCat2GroupByPreCat1Groups}}
\label{PreCat2GroupByPreCat1Groups}
}\hfill{\scriptsize (operation)}}\\


 Following the second definition above, the global functions \texttt{Cat2Group} and \texttt{PreCat2Group} are normally called with two arguments \texttt{\symbol{45}} the generating up
and left cat$^1$\texttt{\symbol{45}}groups $\calC_1$ and $\calC_2$. In this case the groups $C_{[2]}, C_{\{2\}}, C_{\{1\}}$ are the common source and the ranges of $\calC_1,\calC_2$, while $C_{\emptyset}$ is the range of $(e_1 \circ t_1) \circ (e_2 \circ t_2) = (e_2 \circ t_2) \circ (e_1 \circ t_1)$. 

 Alternatively, they may be called with a single argument which is a crossed
square. 

 The operation \texttt{PreCat2GroupByPreCat1Groups} has five arguments \texttt{\symbol{45}} the up, left, right, down and diagonal
cat$^1$\texttt{\symbol{45}}groups. 

 The two cat$^2$\texttt{\symbol{45}}groups \texttt{C2a, C2b} constructed in the following example are isomorphic. They differ in the
down\texttt{\symbol{45}}right group \texttt{P}. }

 

 
\begin{Verbatim}[commandchars=!@|,fontsize=\small,frame=single,label=Example]
  
  !gapprompt@gap>| !gapinput@a := (1,2,3,4,5,6);;  b := (2,6)(3,5);; |
  !gapprompt@gap>| !gapinput@G := Group( a, b );;  SetName( G, "d12" );|
  !gapprompt@gap>| !gapinput@t1 := GroupHomomorphismByImages( G, G, [a,b], [a^3,b] );; |
  !gapprompt@gap>| !gapinput@up := PreCat1GroupByEndomorphisms( t1, t1 );;|
  !gapprompt@gap>| !gapinput@t2 := GroupHomomorphismByImages( G, G, [a,b], [a^4,b] );; |
  !gapprompt@gap>| !gapinput@left := PreCat1GroupByEndomorphisms( t2, t2 );;|
  !gapprompt@gap>| !gapinput@C2a := Cat2Group( up, left );|
  (pre-)cat2-group with generating (pre-)cat1-groups:
  1 : [d12 => Group( [ (1,4)(2,5)(3,6), (2,6)(3,5) ] )]
  2 : [d12 => Group( [ (1,5,3)(2,6,4), (2,6)(3,5) ] )]
  !gapprompt@gap>| !gapinput@IsCat2Group( C2a );|
  true
  !gapprompt@gap>| !gapinput@genR := [ (1,4)(2,5)(3,6), (2,6)(3,5) ];;|
  !gapprompt@gap>| !gapinput@R := Subgroup( G, genR );; |
  !gapprompt@gap>| !gapinput@genQ := [ (1,3,5)(2,4,6), (2,6)(3,5) ];; |
  !gapprompt@gap>| !gapinput@Q := Subgroup( G, genQ );; |
  !gapprompt@gap>| !gapinput@Pa := Group( b );;  SetName( Pa, "c2a" ); |
  !gapprompt@gap>| !gapinput@Pb := Group( (7,8) );;  SetName( Pb, "c2b" ); |
  !gapprompt@gap>| !gapinput@t3 := GroupHomomorphismByImages( R, P, genR, [(),(7,8)] );; |
  !gapprompt@gap>| !gapinput@e3 := GroupHomomorphismByImages( P, R, [(7,8)], [(2,6)(3,5)] );; |
  !gapprompt@gap>| !gapinput@right := PreCat1GroupByTailHeadEmbedding( t3, t3, e3 );;|
  !gapprompt@gap>| !gapinput@t4 := GroupHomomorphismByImages( Q, P, genQ, [(),(7,8)] );; |
  !gapprompt@gap>| !gapinput@e4 := GroupHomomorphismByImages( P, Q, [(7,8)], [(2,6)(3,5)] );; |
  !gapprompt@gap>| !gapinput@down := PreCat1GroupByTailHeadEmbedding( t4, t4, e4 );;|
  !gapprompt@gap>| !gapinput@t0 := t1 * t3;; |
  !gapprompt@gap>| !gapinput@e0 := GroupHomomorphismByImages( P, G, [(7,8)], [(2,6)(3,5)] );; |
  !gapprompt@gap>| !gapinput@diag := PreCat1GroupByTailHeadEmbedding( t0, t0, e0 );;|
  !gapprompt@gap>| !gapinput@C2b := PreCat2GroupByPreCat1Groups( up, left, right, down, diag ); |
  (pre-)cat2-group with generating (pre-)cat1-groups:
  1 : [d12 => Group( [ (1,4)(2,5)(3,6), (2,6)(3,5) ] )]
  2 : [d12 => Group( [ (1,5,3)(2,6,4), (2,6)(3,5) ] )]
  !gapprompt@gap>| !gapinput@C2a = C2b;|
  false
  !gapprompt@gap>| !gapinput@GroupsOfHigherDimensionalGroup( C2a )[4];|
  Group([ (), (2,6)(3,5) ])
  !gapprompt@gap>| !gapinput@GroupsOfHigherDimensionalGroup( C2b )[4];|
  Group([ (7,8) ])
  
\end{Verbatim}
 

\subsection{\textcolor{Chapter }{Up2DimensionalGroup (for cat2-groups)}}
\logpage{[ 8, 5, 2 ]}\nobreak
\hyperdef{L}{X81BD4011837BCC2E}{}
{\noindent\textcolor{FuncColor}{$\triangleright$\enspace\texttt{Up2DimensionalGroup({\mdseries\slshape C2G})\index{Up2DimensionalGroup@\texttt{Up2DimensionalGroup}!for cat2-groups}
\label{Up2DimensionalGroup:for cat2-groups}
}\hfill{\scriptsize (attribute)}}\\
\noindent\textcolor{FuncColor}{$\triangleright$\enspace\texttt{Left2DimensionalGroup({\mdseries\slshape C2G})\index{Left2DimensionalGroup@\texttt{Left2DimensionalGroup}!for cat2-groups}
\label{Left2DimensionalGroup:for cat2-groups}
}\hfill{\scriptsize (attribute)}}\\
\noindent\textcolor{FuncColor}{$\triangleright$\enspace\texttt{Down2DimensionalGroup({\mdseries\slshape C2G})\index{Down2DimensionalGroup@\texttt{Down2DimensionalGroup}!for cat2-groups}
\label{Down2DimensionalGroup:for cat2-groups}
}\hfill{\scriptsize (attribute)}}\\
\noindent\textcolor{FuncColor}{$\triangleright$\enspace\texttt{Right2DimensionalGroup({\mdseries\slshape C2G})\index{Right2DimensionalGroup@\texttt{Right2DimensionalGroup}!for cat2-groups}
\label{Right2DimensionalGroup:for cat2-groups}
}\hfill{\scriptsize (attribute)}}\\
\noindent\textcolor{FuncColor}{$\triangleright$\enspace\texttt{Diagonal2DimensionalGroup({\mdseries\slshape C2G})\index{Diagonal2DimensionalGroup@\texttt{Diagonal2DimensionalGroup}!for cat2-groups}
\label{Diagonal2DimensionalGroup:for cat2-groups}
}\hfill{\scriptsize (attribute)}}\\


 These six attributes were defined for crossed squares but apply equally to cat$^2$\texttt{\symbol{45}}groups. }

 

 
\begin{Verbatim}[commandchars=!@|,fontsize=\small,frame=single,label=Example]
  
  !gapprompt@gap>| !gapinput@Diagonal2DimensionalGroup( C2a );|
  [d12 => Group( [ (), (2,6)(3,5) ] )]
  !gapprompt@gap>| !gapinput@Diagonal2DimensionalGroup( C2b );|
  [d12 => Group( [ (7,8) ] )]
  
\end{Verbatim}
 

\subsection{\textcolor{Chapter }{DirectProduct}}
\logpage{[ 8, 5, 3 ]}\nobreak
\hyperdef{L}{X861BA02C7902A4F4}{}
{\noindent\textcolor{FuncColor}{$\triangleright$\enspace\texttt{DirectProduct({\mdseries\slshape C2a, C2b})\index{DirectProduct@\texttt{DirectProduct}}
\label{DirectProduct}
}\hfill{\scriptsize (operation)}}\\


 The direct product $\calC_{1} \times \calC_{2}$ has as its four up, left, right and down cat$^1$\texttt{\symbol{45}}groups the direct products of those in $\calC_{1}$ and $\calC_{2}$. The embeddings and projections are constructed automatically, and placed in
the \texttt{DirectProductInfo} attribute, together with the two \emph{objects} $\calC_{1}$ and $\calC_{2}$. }

 

 
\begin{Verbatim}[commandchars=!@|,fontsize=\small,frame=single,label=Example]
  
  !gapprompt@gap>| !gapinput@C2ab := DirectProductOp( [ C2a, C2b ], C2a ); |
  (pre-)cat2-group with generating (pre-)cat1-groups:
  1 : [Group( [ (1,2,3,4,5,6), (2,6)(3,5), ( 7, 8, 9,10,11,12), ( 8,12)( 9,11) 
   ] ) => Group( [ (1,4)(2,5)(3,6), (2,6)(3,5), ( 7,10)( 8,11)( 9,12), 
    ( 8,12)( 9,11) ] )]
  2 : [Group( [ (1,2,3,4,5,6), (2,6)(3,5), ( 7, 8, 9,10,11,12), ( 8,12)( 9,11) 
   ] ) => Group( [ (1,5,3)(2,6,4), (2,6)(3,5), ( 7, 9,11)( 8,10,12), 
    ( 8,12)( 9,11) ] )]
  !gapprompt@gap>| !gapinput@StructureDescription( C2ab );  |
  [ "C2 x C2 x S3 x S3", "C2 x C2 x C2 x C2", "S3 x S3", "C2 x C2" ]
  !gapprompt@gap>| !gapinput@SetName( C2ab, "C2ab" );|
  !gapprompt@gap>| !gapinput@Embedding( C2ab, 1 );   |
  <mapping: (pre-)cat2-group with generating (pre-)cat1-groups:
  1 : [d12 => Group( [ (1,4)(2,5)(3,6), (2,6)(3,5) ] )]
  2 : [d12 => Group( [ (1,5,3)(2,6,4), (2,6)(3,5) ] )] -> C2ab >
  !gapprompt@gap>| !gapinput@Projection( C2ab, 2 );|
  <mapping: C2ab -> (pre-)cat2-group with generating (pre-)cat1-groups:
  1 : [d12 => Group( [ (1,4)(2,5)(3,6), (2,6)(3,5) ] )]
  2 : [d12 => Group( [ (1,5,3)(2,6,4), (2,6)(3,5) ] )] >
  
\end{Verbatim}
 

\subsection{\textcolor{Chapter }{DisplayLeadMaps}}
\logpage{[ 8, 5, 4 ]}\nobreak
\hyperdef{L}{X85713B1C7C34324E}{}
{\noindent\textcolor{FuncColor}{$\triangleright$\enspace\texttt{DisplayLeadMaps({\mdseries\slshape C0})\index{DisplayLeadMaps@\texttt{DisplayLeadMaps}}
\label{DisplayLeadMaps}
}\hfill{\scriptsize (operation)}}\\


 This operation provides an alternative to \texttt{Display} giving a shorter output. Generators of the up\texttt{\symbol{45}}left group
are output, together with their images under the up and left tail and head
maps. }

 

 
\begin{Verbatim}[commandchars=!@|,fontsize=\small,frame=single,label=Example]
  
  !gapprompt@gap>| !gapinput@DisplayLeadMaps( C2b );|
  (pre-)cat2-group with up-left group: [ (1,2,3,4,5,6), (2,6)(3,5) ]
     up tail=head images: [ (1,4)(2,5)(3,6), (2,6)(3,5) ]
   left tail=head images: [ (1,5,3)(2,6,4), (2,6)(3,5) ]
  
\end{Verbatim}
 

\subsection{\textcolor{Chapter }{Transpose3DimensionalGroup (for cat2-groups)}}
\logpage{[ 8, 5, 5 ]}\nobreak
\hyperdef{L}{X864B346686AAA522}{}
{\noindent\textcolor{FuncColor}{$\triangleright$\enspace\texttt{Transpose3DimensionalGroup({\mdseries\slshape S0})\index{Transpose3DimensionalGroup@\texttt{Transpose3DimensionalGroup}!for cat2-groups}
\label{Transpose3DimensionalGroup:for cat2-groups}
}\hfill{\scriptsize (attribute)}}\\


 The \emph{transpose} of a cat$^2$\texttt{\symbol{45}}group $\calC$ with groups $[G,R,Q,P]$ is the cat$^2$\texttt{\symbol{45}}group $\tilde{\calC}$ with groups $[G,Q,R,P]$. }

 

 
\begin{Verbatim}[commandchars=!@|,fontsize=\small,frame=single,label=Example]
  
  !gapprompt@gap>| !gapinput@TC2a := Transpose3DimensionalGroup( C2a );|
  (pre-)cat2-group with generating (pre-)cat1-groups:
  1 : [d12 => Group( [ (1,5,3)(2,6,4), (2,6)(3,5) ] )]
  2 : [d12 => Group( [ (1,4)(2,5)(3,6), (2,6)(3,5) ] )]
  
\end{Verbatim}
 

\subsection{\textcolor{Chapter }{Cat2GroupMorphism}}
\logpage{[ 8, 5, 6 ]}\nobreak
\label{cat2-mor}
\hyperdef{L}{X83616C237F4745FB}{}
{\noindent\textcolor{FuncColor}{$\triangleright$\enspace\texttt{Cat2GroupMorphism({\mdseries\slshape args})\index{Cat2GroupMorphism@\texttt{Cat2GroupMorphism}}
\label{Cat2GroupMorphism}
}\hfill{\scriptsize (function)}}\\
\noindent\textcolor{FuncColor}{$\triangleright$\enspace\texttt{Cat2GroupMorphismByCat1GroupMorphisms({\mdseries\slshape src, rng, upmor, ltmor})\index{Cat2GroupMorphismByCat1GroupMorphisms@\texttt{Cat2}\-\texttt{Group}\-\texttt{Morphism}\-\texttt{By}\-\texttt{Cat1}\-\texttt{Group}\-\texttt{Morphisms}}
\label{Cat2GroupMorphismByCat1GroupMorphisms}
}\hfill{\scriptsize (operation)}}\\
\noindent\textcolor{FuncColor}{$\triangleright$\enspace\texttt{Cat2GroupMorphismByGroupHomomorphisms({\mdseries\slshape src, rng, homs})\index{Cat2GroupMorphismByGroupHomomorphisms@\texttt{Cat2}\-\texttt{Group}\-\texttt{Morphism}\-\texttt{By}\-\texttt{Group}\-\texttt{Homomorphisms}}
\label{Cat2GroupMorphismByGroupHomomorphisms}
}\hfill{\scriptsize (operation)}}\\
\noindent\textcolor{FuncColor}{$\triangleright$\enspace\texttt{PreCat2GroupMorphism({\mdseries\slshape args})\index{PreCat2GroupMorphism@\texttt{PreCat2GroupMorphism}}
\label{PreCat2GroupMorphism}
}\hfill{\scriptsize (function)}}\\
\noindent\textcolor{FuncColor}{$\triangleright$\enspace\texttt{PreCat2GroupMorphismByPreCat1GroupMorphisms({\mdseries\slshape src, rng, upmor, ltmor})\index{PreCat2GroupMorphismByPreCat1GroupMorphisms@\texttt{Pre}\-\texttt{Cat2}\-\texttt{Group}\-\texttt{Morphism}\-\texttt{By}\-\texttt{Pre}\-\texttt{Cat1}\-\texttt{Group}\-\texttt{Morphisms}}
\label{PreCat2GroupMorphismByPreCat1GroupMorphisms}
}\hfill{\scriptsize (operation)}}\\
\noindent\textcolor{FuncColor}{$\triangleright$\enspace\texttt{PreCat2GroupMorphismByGroupHomomorphisms({\mdseries\slshape src, rng, homs})\index{PreCat2GroupMorphismByGroupHomomorphisms@\texttt{Pre}\-\texttt{Cat2}\-\texttt{Group}\-\texttt{Morphism}\-\texttt{By}\-\texttt{Group}\-\texttt{Homomorphisms}}
\label{PreCat2GroupMorphismByGroupHomomorphisms}
}\hfill{\scriptsize (operation)}}\\


 A (pre\texttt{\symbol{45}})cat$^2$\texttt{\symbol{45}}group morphism $\mu : \calC = (G,R,Q,P) \to \calC' = (G',R',Q',P')$ is a list of four group homomorphisms $\gamma : G \to G',~ \rho : R \to R',~ \xi : Q \to Q'$ and $\pi : P \to P'$ which commute with all the tail, head and embedding maps so that $(\gamma,\rho), (\gamma,\xi), (\rho,\pi)$ and $(\xi,\pi)$ are all (pre\texttt{\symbol{45}})cat$^1$\texttt{\symbol{45}}group morphisms. 

 For the operations \texttt{(Pre)Cat2GroupMorphismByPreCat1GroupMorphisms} the third and fourth parameters \texttt{upmor, ltmor} are two cat$^1$\texttt{\symbol{45}}group morphisms with source the up and left cat$^1$\texttt{\symbol{45}}groups in $\calC$. 

 For the operations \texttt{(Pre)Cat2GroupMorphismByGroupHomomorphisms} the third parameter \texttt{mors} is the list $[\gamma,\rho,\xi,\pi]$. 

 The example constructs an automorphism of \texttt{c2a} is two ways, using the two methods described above, an d verifies that the
result is the same in each case. }

 
\begin{Verbatim}[commandchars=!@|,fontsize=\small,frame=single,label=Example]
  
  !gapprompt@gap>| !gapinput@gamma := GroupHomomorphismByImages( G, G, [a,b], [a^-1,b] );;|
  !gapprompt@gap>| !gapinput@rho := IdentityMapping( R );;|
  !gapprompt@gap>| !gapinput@xi := GroupHomomorphismByImages( Q, Q, [a^2,b], [a^-2,b] );;|
  !gapprompt@gap>| !gapinput@pi := IdentityMapping( Pa );;|
  !gapprompt@gap>| !gapinput@homs := [ gamma, rho, xi, pi ];;|
  !gapprompt@gap>| !gapinput@mor1 := Cat2GroupMorphismByGroupHomomorphisms( C2a, C2a, homs );   |
  <mapping: (pre-)cat2-group with generating (pre-)cat1-groups:
  1 : [d12 => Group( [ (1,4)(2,5)(3,6), (2,6)(3,5) ] )]
  2 : [d12 => Group( [ (1,5,3)(2,6,4), (2,6)(3,5) ] )] -> (pre-)cat
  2-group with generating (pre-)cat1-groups:
  1 : [d12 => Group( [ (1,4)(2,5)(3,6), (2,6)(3,5) ] )]
  2 : [d12 => Group( [ (1,5,3)(2,6,4), (2,6)(3,5) ] )] >
  !gapprompt@gap>| !gapinput@upmor := Cat1GroupMorphism( up, up, gamma, rho );; |
  !gapprompt@gap>| !gapinput@ltmor := Cat1GroupMorphism( left, left, gamma, xi );; |
  !gapprompt@gap>| !gapinput@mor2 := Cat2GroupMorphismByCat1GroupMorphisms( C2a, C2a, upmor, ltmor );; |
  !gapprompt@gap>| !gapinput@mor1 = mor2; |
  true
  
\end{Verbatim}
 

\subsection{\textcolor{Chapter }{Cat2GroupOfCrossedSquare}}
\logpage{[ 8, 5, 7 ]}\nobreak
\label{cat2-xsq}
\hyperdef{L}{X7D46CB517FCAA83B}{}
{\noindent\textcolor{FuncColor}{$\triangleright$\enspace\texttt{Cat2GroupOfCrossedSquare({\mdseries\slshape xsq})\index{Cat2GroupOfCrossedSquare@\texttt{Cat2GroupOfCrossedSquare}}
\label{Cat2GroupOfCrossedSquare}
}\hfill{\scriptsize (attribute)}}\\
\noindent\textcolor{FuncColor}{$\triangleright$\enspace\texttt{CrossedSquareOfCat2Group({\mdseries\slshape CC})\index{CrossedSquareOfCat2Group@\texttt{CrossedSquareOfCat2Group}}
\label{CrossedSquareOfCat2Group}
}\hfill{\scriptsize (attribute)}}\\


 These functions provide for conversion between crossed squares and cat$^2$\texttt{\symbol{45}}groups. (They are the 3\texttt{\symbol{45}}dimensional
equivalents of \texttt{Cat1GroupOfXMod} (\ref{Cat1GroupOfXMod}) and \texttt{XModOfCat1Group} (\ref{XModOfCat1Group}).) The actor crossed square \texttt{XSact} was constructed in section \texttt{ActorCrossedSquare} (\ref{ActorCrossedSquare}). }

 

 
\begin{Verbatim}[commandchars=!@A,fontsize=\small,frame=single,label=Example]
  
  !gapprompt@gap>A !gapinput@xsC2a := CrossedSquareOfCat2Group( C2a );A
  crossed square with crossed modules:
        up = [Group( () ) -> Group( [ (1,4)(2,5)(3,6) ] )]
      left = [Group( () ) -> Group( [ (1,3,5)(2,4,6) ] )]
     right = [Group( [ (1,4)(2,5)(3,6) ] ) -> Group( [ (2,6)(3,5) ] )]
      down = [Group( [ (1,3,5)(2,4,6) ] ) -> Group( [ (2,6)(3,5) ] )]
  
  !gapprompt@gap>A !gapinput@IdGroup( xsC2a );A
  [ [ 1, 1 ], [ 2, 1 ], [ 3, 1 ], [ 2, 1 ] ]
  
  !gapprompt@gap>A !gapinput@SetName( Source( Right2DimensionalGroup( XSact ) ), "c5:c4" );A
  !gapprompt@gap>A !gapinput@SetName( Range( Right2DimensionalGroup( XSact ) ), "c5:c4" );A
  !gapprompt@gap>A !gapinput@Name( XSact );A
  "[d10a->c5:c4,d20->c5:c4]"
  
  !gapprompt@gap>A !gapinput@C2act := Cat2GroupOfCrossedSquare( XSact );             A
  (pre-)cat2-group with generating (pre-)cat1-groups:
  1 : [((c5:c4 |X c5:c4) |X (d20 |X d10a))=>(c5:c4 |X c5:c4)]
  2 : [((c5:c4 |X c5:c4) |X (d20 |X d10a))=>(c5:c4 |X d20)]
  !gapprompt@gap>A !gapinput@Size3d( C2act );A
  [ 80000, 400, 400, 20 ]
  
\end{Verbatim}
 

\subsection{\textcolor{Chapter }{Subdiagonal2DimensionalGroup}}
\logpage{[ 8, 5, 8 ]}\nobreak
\label{subdiag-cat1}
\hyperdef{L}{X82F896C7851294B7}{}
{\noindent\textcolor{FuncColor}{$\triangleright$\enspace\texttt{Subdiagonal2DimensionalGroup({\mdseries\slshape obj})\index{Subdiagonal2DimensionalGroup@\texttt{Subdiagonal2DimensionalGroup}}
\label{Subdiagonal2DimensionalGroup}
}\hfill{\scriptsize (attribute)}}\\


 The diagonal of a crossed square is always a crossed module, but the diagonal
of a cat$^2$\texttt{\symbol{45}}group need only be a pre\texttt{\symbol{45}}cat$^1$\texttt{\symbol{45}}group. There is, however, a sub\texttt{\symbol{45}}cat$^1$\texttt{\symbol{45}}group of this diagonal which, in the case of a cat$^2$\texttt{\symbol{45}}group constructed from a crossed square, is $(P \ltimes L => P)$. (The name of this operation is very provisional.) 

 }

 

 
\begin{Verbatim}[commandchars=!@|,fontsize=\small,frame=single,label=Example]
  
  !gapprompt@gap>| !gapinput@G24 := SmallGroup( 24, 10 );; |
  !gapprompt@gap>| !gapinput@w := G24.1;; x := G24.2;; y := G24.3;; z := G24.4;; o := One(G24);; |
  !gapprompt@gap>| !gapinput@R := Subgroup( G24, [x,y] );; |
  !gapprompt@gap>| !gapinput@txy := GroupHomomorphismByImages( G24, R, [w,x,y,z], [o,x,y,o] );; |
  !gapprompt@gap>| !gapinput@exy := GroupHomomorphismByImages( R, G24, [x,y], [x,y] );; |
  !gapprompt@gap>| !gapinput@C1xy := PreCat1GroupByTailHeadEmbedding( txy, txy, exy );; |
  !gapprompt@gap>| !gapinput@Q := Subgroup( G24, [w,y] );; |
  !gapprompt@gap>| !gapinput@twy := GroupHomomorphismByImages( G24, Q, [w,x,y,z], [w,o,y,o] );; |
  !gapprompt@gap>| !gapinput@ewy := GroupHomomorphismByImages( Q, G24, [w,y], [w,y] );; |
  !gapprompt@gap>| !gapinput@C1wy := PreCat1GroupByTailHeadEmbedding( twy, twy, ewy );; |
  !gapprompt@gap>| !gapinput@C2wxy := PreCat2Group( C1xy, C1xy );; |
  !gapprompt@gap>| !gapinput@dg := Diagonal2DimensionalGroup( C2wxy );;|
  !gapprompt@gap>| !gapinput@C1sub := Subdiagonal2DimensionalGroup( C2wxy );; |
  !gapprompt@gap>| !gapinput@[ IsCat1Group(dg), IsCat1Group(C1sub), IsSub2DimensionalGroup(dg,C1sub) ];|
  [ false, true, true ]
  
\end{Verbatim}
 

\subsection{\textcolor{Chapter }{SubCat2Group}}
\logpage{[ 8, 5, 9 ]}\nobreak
\hyperdef{L}{X7DB082D282C52855}{}
{\noindent\textcolor{FuncColor}{$\triangleright$\enspace\texttt{SubCat2Group({\mdseries\slshape C2G, ul, ur, dl})\index{SubCat2Group@\texttt{SubCat2Group}}
\label{SubCat2Group}
}\hfill{\scriptsize (operation)}}\\
\noindent\textcolor{FuncColor}{$\triangleright$\enspace\texttt{IsSubCat2Group({\mdseries\slshape C2G, subC2G})\index{IsSubCat2Group@\texttt{IsSubCat2Group}}
\label{IsSubCat2Group}
}\hfill{\scriptsize (operation)}}\\
\noindent\textcolor{FuncColor}{$\triangleright$\enspace\texttt{SubPreCat2Group({\mdseries\slshape C2G, ul, ur, dl})\index{SubPreCat2Group@\texttt{SubPreCat2Group}}
\label{SubPreCat2Group}
}\hfill{\scriptsize (operation)}}\\
\noindent\textcolor{FuncColor}{$\triangleright$\enspace\texttt{IsSubPreCat2Group({\mdseries\slshape C2G, subC2G})\index{IsSubPreCat2Group@\texttt{IsSubPreCat2Group}}
\label{IsSubPreCat2Group}
}\hfill{\scriptsize (operation)}}\\


 As with crossed squares, we may construct sub\texttt{\symbol{45}}structures of
cat$^2$\texttt{\symbol{45}}groups. In this case we need only prescribe three
subgroups, rather than four, so that the up and left
sub\texttt{\symbol{45}}cat$^1$\texttt{\symbol{45}}groups can be specified. The remaining
sub\texttt{\symbol{45}}cat$^1$\texttt{\symbol{45}}groups are then calculated automatically. }

 

 
\begin{Verbatim}[commandchars=!@|,fontsize=\small,frame=single,label=Example]
  
  !gapprompt@gap>| !gapinput@gps := GroupsOfHigherDimensionalGroup( C2ab );;|
  !gapprompt@gap>| !gapinput@c6c2 := Subgroup( gps[1], [ (1,2,3,4,5,6), (8,12)(9,11) ] );;|
  !gapprompt@gap>| !gapinput@c2c2 := Subgroup( gps[2], [ (1,4)(2,5)(3,6), (8,12)(9,11) ] );;|
  !gapprompt@gap>| !gapinput@c3c2 := Subgroup( gps[3], [ (1,5,3)(2,6,4), (8,12)(9,11) ] );;|
  !gapprompt@gap>| !gapinput@SC2ab := SubCat2Group( C2ab, c6c2, c2c2, c3c2 );;|
  !gapprompt@gap>| !gapinput@Display( SC2ab );              |
  (pre-)cat2-group with groups: [ Group( [ (1,2,3,4,5,6), ( 8,12)( 9,11) ] ), 
    Group( [ (1,4)(2,5)(3,6), ( 8,12)( 9,11) ] ), 
    Group( [ (1,5,3)(2,6,4), ( 8,12)( 9,11) ] ), 
    Group( [ (), ( 8,12)( 9,11) ] ) ]
     up tail=head: [ [ (1,2,3,4,5,6), ( 8,12)( 9,11) ], 
    [ (1,4)(2,5)(3,6), ( 8,12)( 9,11) ] ]
   left tail=head: [ [ (1,2,3,4,5,6), ( 8,12)( 9,11) ], 
    [ (1,5,3)(2,6,4), ( 8,12)( 9,11) ] ]
  right tail=head: [ [ (1,4)(2,5)(3,6), ( 8,12)( 9,11) ], 
    [ (), ( 8,12)( 9,11) ] ]
   down tail=head: [ [ (1,5,3)(2,6,4), ( 8,12)( 9,11) ], [ (), ( 8,12)( 9,11) ] 
   ]
  
\end{Verbatim}
 

\subsection{\textcolor{Chapter }{TrivialSubCat2Group}}
\logpage{[ 8, 5, 10 ]}\nobreak
\hyperdef{L}{X838423B47FEE09B2}{}
{\noindent\textcolor{FuncColor}{$\triangleright$\enspace\texttt{TrivialSubCat2Group({\mdseries\slshape C2G})\index{TrivialSubCat2Group@\texttt{TrivialSubCat2Group}}
\label{TrivialSubCat2Group}
}\hfill{\scriptsize (attribute)}}\\
\noindent\textcolor{FuncColor}{$\triangleright$\enspace\texttt{TrivialSubPreCat2Group({\mdseries\slshape C2G})\index{TrivialSubPreCat2Group@\texttt{TrivialSubPreCat2Group}}
\label{TrivialSubPreCat2Group}
}\hfill{\scriptsize (attribute)}}\\


 A special case of the previous operation is when all the subgroups are
trivial. }

 

 
\begin{Verbatim}[commandchars=!@|,fontsize=\small,frame=single,label=Example]
  
  !gapprompt@gap>| !gapinput@TC2ab := TrivialSubCrossedSquare( C2ab );|
  crossed square with crossed modules:
        up = [Group( () ) -> Group( () )]
      left = [Group( () ) -> Group( () )]
     right = [Group( () ) -> Group( () )]
      down = [Group( () ) -> Group( () )]
  
\end{Verbatim}
 }

 
\section{\textcolor{Chapter }{Enumerating cat$^2$\texttt{\symbol{45}}groups with a given source}}\label{sect-allcat2}
\logpage{[ 8, 6, 0 ]}
\hyperdef{L}{X80FB2B328578DE42}{}
{
  This section mirrors that for cat$^1$\texttt{\symbol{45}}groups (\ref{sect-allcat1}). As the size of a group $G$ increases, the number of cat$^2$\texttt{\symbol{45}}groups with source $G$ increases rapidly. However, one is usually only interested in the isomorphism
classes of cat$^2$\texttt{\symbol{45}}groups with source $G$. An iterator \texttt{AllCat2GroupsIterator} is provided, which runs through the various cat$^2$\texttt{\symbol{45}}groups. This iterator finds, for each unordered pair of
subgroups $R,Q$ of $G$, the cat$^2$\texttt{\symbol{45}}groups whose \texttt{Up2DimensionalGroup} has range $R$, and whose \texttt{Left2DimensionalGroup} has range $Q$. It does this by running through \texttt{UnoderedPairsIterator(AllSubgroupsIterator(G))} provided by the \textsf{Utils} package, and then using the iterator \texttt{AllCat2GroupsWithImagesIterator(G,R,Q)}. 

\subsection{\textcolor{Chapter }{AllCat2GroupsWithImagesIterator}}
\logpage{[ 8, 6, 1 ]}\nobreak
\hyperdef{L}{X7D421DE57B44F37A}{}
{\noindent\textcolor{FuncColor}{$\triangleright$\enspace\texttt{AllCat2GroupsWithImagesIterator({\mdseries\slshape G, R, Q})\index{AllCat2GroupsWithImagesIterator@\texttt{AllCat2GroupsWithImagesIterator}}
\label{AllCat2GroupsWithImagesIterator}
}\hfill{\scriptsize (operation)}}\\
\noindent\textcolor{FuncColor}{$\triangleright$\enspace\texttt{AllCat2GroupsWithImagesNumber({\mdseries\slshape G, R, Q})\index{AllCat2GroupsWithImagesNumber@\texttt{AllCat2GroupsWithImagesNumber}}
\label{AllCat2GroupsWithImagesNumber}
}\hfill{\scriptsize (attribute)}}\\
\noindent\textcolor{FuncColor}{$\triangleright$\enspace\texttt{AllCat2GroupsWithImages({\mdseries\slshape G, R, Q})\index{AllCat2GroupsWithImages@\texttt{AllCat2GroupsWithImages}}
\label{AllCat2GroupsWithImages}
}\hfill{\scriptsize (operation)}}\\
\noindent\textcolor{FuncColor}{$\triangleright$\enspace\texttt{AllCat2GroupsWithImagesUpToIsomorphism({\mdseries\slshape G, R, Q})\index{AllCat2GroupsWithImagesUpToIsomorphism@\texttt{All}\-\texttt{Cat2}\-\texttt{Groups}\-\texttt{With}\-\texttt{Images}\-\texttt{Up}\-\texttt{To}\-\texttt{Isomorphism}}
\label{AllCat2GroupsWithImagesUpToIsomorphism}
}\hfill{\scriptsize (operation)}}\\


 The iterator \texttt{AllCat2GroupsWithImagesIterator(G)} iterates through all the cat$^2$\texttt{\symbol{45}}groups with source \texttt{G} and generating cat$^1$\texttt{\symbol{45}}groups \texttt{(G={\textgreater}R)} and \texttt{(G={\textgreater}Q)}. The attribute \texttt{AllCat2GroupsWithImagesNumber(G)} runs through this iterator to determine the number $n$ of these cat$^2$\texttt{\symbol{45}}groups. The operation \texttt{AllCat2GroupsWithImages(G)} returns a list containing these $n$ cat$^2$\texttt{\symbol{45}}groups. Since these lists can get very long, this
operation should only be used for simple cases. The operation \texttt{AllCat2GroupsWithImagesUpToIsomorphism(G)} returns representatives of the isomorphism classes of these cat$^2$\texttt{\symbol{45}}groups. }

 

 
\begin{Verbatim}[commandchars=!@|,fontsize=\small,frame=single,label=Example]
  
  !gapprompt@gap>| !gapinput@G8 := Group( (1,2), (3,4), (5,6) );;|
  !gapprompt@gap>| !gapinput@A := Subgroup( G8, [ (1,2) ] );; |
  !gapprompt@gap>| !gapinput@B := Subgroup( G8, [ (3,4) ] );;|
  !gapprompt@gap>| !gapinput@AllCat2GroupsWithImagesNumber( G8, A, A );|
  4
  !gapprompt@gap>| !gapinput@all := AllCat2GroupsWithImages( G8, A, A );;     |
  !gapprompt@gap>| !gapinput@for C2 in all do DisplayLeadMaps( C2 ); od;|
  (pre-)cat2-group with up-left group: [ (1,2), (3,4), (5,6) ]
     up tail=head images: [ (1,2), (1,2), () ]
   left tail=head images: [ (1,2), (1,2), () ]
  (pre-)cat2-group with up-left group: [ (1,2), (3,4), (5,6) ]
     up tail=head images: [ (1,2), (), () ]
   left tail=head images: [ (1,2), (), () ]
  (pre-)cat2-group with up-left group: [ (1,2), (3,4), (5,6) ]
     up tail=head images: [ (1,2), (), (1,2) ]
   left tail=head images: [ (1,2), (), (1,2) ]
  (pre-)cat2-group with up-left group: [ (1,2), (3,4), (5,6) ]
     up tail=head images: [ (1,2), (1,2), (1,2) ]
   left tail=head images: [ (1,2), (1,2), (1,2) ]
  !gapprompt@gap>| !gapinput@AllCat2GroupsWithImagesNumber( G8, A, B );|
  16
  !gapprompt@gap>| !gapinput@iso := AllCat2GroupsWithImagesUpToIsomorphism( G8, A, B );;|
  !gapprompt@gap>| !gapinput@for C2 in iso do DisplayLeadMaps( C2 ); od;|
  (pre-)cat2-group with up-left group: [ (1,2), (3,4), (5,6) ]
     up tail=head images: [ (1,2), (), () ]
   left tail=head images: [ (), (3,4), () ]
  (pre-)cat2-group with up-left group: [ (1,2), (3,4), (5,6) ]
     up tail=head images: [ (1,2), (), () ]
   left tail/head images: [ (), (3,4), () ], [ (), (3,4), (3,4) ]
  (pre-)cat2-group with up-left group: [ (1,2), (3,4), (5,6) ]
     up tail/head images: [ (1,2), (), () ], [ (1,2), (), (1,2) ]
   left tail/head images: [ (), (3,4), () ], [ (), (3,4), (3,4) ]
  
\end{Verbatim}
 

\subsection{\textcolor{Chapter }{AllCat2GroupsWithFixedUp}}
\logpage{[ 8, 6, 2 ]}\nobreak
\hyperdef{L}{X783F7FEB7B79B062}{}
{\noindent\textcolor{FuncColor}{$\triangleright$\enspace\texttt{AllCat2GroupsWithFixedUp({\mdseries\slshape C})\index{AllCat2GroupsWithFixedUp@\texttt{AllCat2GroupsWithFixedUp}}
\label{AllCat2GroupsWithFixedUp}
}\hfill{\scriptsize (operation)}}\\
\noindent\textcolor{FuncColor}{$\triangleright$\enspace\texttt{AllCat2GroupsWithFixedUpAndLeftRange({\mdseries\slshape C, R})\index{AllCat2GroupsWithFixedUpAndLeftRange@\texttt{All}\-\texttt{Cat2}\-\texttt{Groups}\-\texttt{With}\-\texttt{Fixed}\-\texttt{Up}\-\texttt{And}\-\texttt{Left}\-\texttt{Range}}
\label{AllCat2GroupsWithFixedUpAndLeftRange}
}\hfill{\scriptsize (operation)}}\\


 The operation \texttt{AllCat2GroupsWithFixedUp(C)} constructs all the cat$^2$\texttt{\symbol{45}}groups with a fixed \texttt{Up2DimensionalGroup} $C$. In the second operation the user may also specify the range of the \texttt{Left2DimensionalGroup}. }

 

 
\begin{Verbatim}[commandchars=!@|,fontsize=\small,frame=single,label=Example]
  
  !gapprompt@gap>| !gapinput@up := Up2DimensionalGroup( iso[1] );                |
  [Group( [ (1,2), (3,4), (5,6) ] )=>Group( [ (1,2), (), () ] )]
  !gapprompt@gap>| !gapinput@AllCat2GroupsWithFixedUp( up );;                    |
  !gapprompt@gap>| !gapinput@Length(last);                                       |
  28
  !gapprompt@gap>| !gapinput@L := AllCat2GroupsWithFixedUpAndLeftRange( up, B );;|
  !gapprompt@gap>| !gapinput@for C in L do DisplayLeadMaps( C ); od;             |
  (pre-)cat2-group with up-left group: [ (1,2), (3,4), (5,6) ]
     up tail=head images: [ (1,2), (), () ]
   left tail=head images: [ (), (3,4), () ]
  (pre-)cat2-group with up-left group: [ (1,2), (3,4), (5,6) ]
     up tail=head images: [ (1,2), (), () ]
   left tail/head images: [ (), (3,4), () ], [ (), (3,4), (3,4) ]
  (pre-)cat2-group with up-left group: [ (1,2), (3,4), (5,6) ]
     up tail=head images: [ (1,2), (), () ]
   left tail/head images: [ (), (3,4), (3,4) ], [ (), (3,4), () ]
  (pre-)cat2-group with up-left group: [ (1,2), (3,4), (5,6) ]
     up tail=head images: [ (1,2), (), () ]
   left tail=head images: [ (), (3,4), (3,4) ]
  
\end{Verbatim}
 

\subsection{\textcolor{Chapter }{AllCat2GroupsMatrix}}
\logpage{[ 8, 6, 3 ]}\nobreak
\hyperdef{L}{X84636E047CDF2DF3}{}
{\noindent\textcolor{FuncColor}{$\triangleright$\enspace\texttt{AllCat2GroupsMatrix({\mdseries\slshape G})\index{AllCat2GroupsMatrix@\texttt{AllCat2GroupsMatrix}}
\label{AllCat2GroupsMatrix}
}\hfill{\scriptsize (attribute)}}\\


 The operation \texttt{AllCat2GroupsMatrix(G)} constructs a symmetric matrix $M$ with rows and columns labelled by the cat$^1$\texttt{\symbol{45}}groups $C_i$ on $G$, where $M_{ij}$ is $1$ if $C_i,C_j$ combine to form a cat$^2$\texttt{\symbol{45}}group, and $0$ otherwise. The matrix is automatically printed out with dots in place of
zeroes. 

 In the example we see that the dihedral group $D_{12}$ has $12$ cat$^1$\texttt{\symbol{45}}groups and $41$ cat$^2$\texttt{\symbol{45}}groups, $12$ of which are symmetric. This operation is intended to be used to illustrate
how cat$^2$\texttt{\symbol{45}}groups are formed, and should only be used with groups of
low order. 

 The attribute \texttt{AllCat2GroupsNumber(G)} returns the number $n$ of these cat$^2$\texttt{\symbol{45}}groups. }

 

 
\begin{Verbatim}[commandchars=!@|,fontsize=\small,frame=single,label=Example]
  
  !gapprompt@gap>| !gapinput@AllCat2GroupsMatrix(d12);;                |
  number of cat2-groups found = 41
  1.....1..1.1
  .1.....1.1.1
  ..1.....11.1
  ...1....1.11
  ....1.1...11
  .....1.1..11
  1...1.1..111
  .1...1.1.111
  ..11....1111
  111...1111.1
  ...111111.11
  111111111111
  !gapprompt@gap>| !gapinput@AllCat2GroupsNumber(d12); |
  41
  
\end{Verbatim}
 

\subsection{\textcolor{Chapter }{AllCat2GroupsIterator}}
\logpage{[ 8, 6, 4 ]}\nobreak
\hyperdef{L}{X7EFCF9697E845B2C}{}
{\noindent\textcolor{FuncColor}{$\triangleright$\enspace\texttt{AllCat2GroupsIterator({\mdseries\slshape G})\index{AllCat2GroupsIterator@\texttt{AllCat2GroupsIterator}}
\label{AllCat2GroupsIterator}
}\hfill{\scriptsize (operation)}}\\
\noindent\textcolor{FuncColor}{$\triangleright$\enspace\texttt{AllCat2Groups({\mdseries\slshape G})\index{AllCat2Groups@\texttt{AllCat2Groups}}
\label{AllCat2Groups}
}\hfill{\scriptsize (operation)}}\\
\noindent\textcolor{FuncColor}{$\triangleright$\enspace\texttt{AllCat2GroupsUpToIsomorphism({\mdseries\slshape G})\index{AllCat2GroupsUpToIsomorphism@\texttt{AllCat2GroupsUpToIsomorphism}}
\label{AllCat2GroupsUpToIsomorphism}
}\hfill{\scriptsize (operation)}}\\
\noindent\textcolor{FuncColor}{$\triangleright$\enspace\texttt{AllCat2GroupFamilies({\mdseries\slshape G})\index{AllCat2GroupFamilies@\texttt{AllCat2GroupFamilies}}
\label{AllCat2GroupFamilies}
}\hfill{\scriptsize (operation)}}\\
\noindent\textcolor{FuncColor}{$\triangleright$\enspace\texttt{CatnGroupNumbers({\mdseries\slshape G})\index{CatnGroupNumbers@\texttt{CatnGroupNumbers}!for cat2-groups}
\label{CatnGroupNumbers:for cat2-groups}
}\hfill{\scriptsize (attribute)}}\\
\noindent\textcolor{FuncColor}{$\triangleright$\enspace\texttt{CatnGroupLists({\mdseries\slshape G})\index{CatnGroupLists@\texttt{CatnGroupLists}!for cat2-groups}
\label{CatnGroupLists:for cat2-groups}
}\hfill{\scriptsize (attribute)}}\\


 The iterator \texttt{AllCat2GroupsIterator(G)} iterates through all the cat$^2$\texttt{\symbol{45}}groups with source \texttt{G}. The operation \texttt{AllCat2Groups(G)} returns a list containing these $n$ cat$^2$\texttt{\symbol{45}}groups. Since these lists can get very long, this
operation should only be used for simple cases. The operation \texttt{AllCat2GroupsUpToIsomorphism(G)} returns representatives of the isomorphism classes of these subgroups. The
operation \texttt{AllCat2GroupFamilies(G)} returns a list of lists. The $k$\texttt{\symbol{45}}th list contains the positions of the cat$^2$\texttt{\symbol{45}}groups in \texttt{AllCat2Groups(G)} which are isomorphic to the $k$\texttt{\symbol{45}}th representative. So, for \texttt{d12}, the $41$ cat$^2$\texttt{\symbol{45}}groups form $10$ classes, and the sizes of these classes are \texttt{[6,6,6,6,3,6,3,2,2,1]}. Four of these classes contain symmetric cat$^2$\texttt{\symbol{45}}groups. 

 The field \texttt{CatnGroupNumbers(G).cat2} is the number of cat$^2$\texttt{\symbol{45}}groups on $G$, while \texttt{CatnGroupNumbers(G).iso2} is the number of isomorphism classes of these cat$^2$\texttt{\symbol{45}}groups. Also \texttt{CatnGroupNumbers(G).symm} is the number of cat$^2$\texttt{\symbol{45}}groups whose \texttt{Up2DimensionalGroup} is the same as the \texttt{Left2DimensionalGroup}, while \texttt{CatnGroupNumbers(G).siso} is the number of isomorphism classes of these symmetric cat$^2$\texttt{\symbol{45}}groups. 

 Provided that \texttt{CatnGroupLists(G).omit} is not set to \texttt{true}, \emph{sorted} lists of generating pairs, and of the classes they belong to, are added to the
record \texttt{CatnGroupLists}. For example \texttt{[5,7]} in these lists for \texttt{d12} indicates that there is a cat$^2$\texttt{\symbol{45}}group generated by the fifth and seventh cat$^1$\texttt{\symbol{45}}groups and that this is in the second class whose
representative is \texttt{[1,7]}. Classes \texttt{[1,5,8,10]} contain symmetric cat$^2$\texttt{\symbol{45}}groups. }

 

 
\begin{Verbatim}[commandchars=!@|,fontsize=\small,frame=single,label=Example]
  
  !gapprompt@gap>| !gapinput@AllCat2GroupsNumber( d12 );|
  41
  !gapprompt@gap>| !gapinput@reps2 := AllCat2GroupsUpToIsomorphism( d12 );;|
  !gapprompt@gap>| !gapinput@Length( reps2 );|
  10
  !gapprompt@gap>| !gapinput@List( reps2, C -> StructureDescription( C ) );|
  [ [ "D12", "C2", "C2", "C2" ], [ "D12", "C2", "C2 x C2", "C2" ], 
    [ "D12", "C2", "S3", "C2" ], [ "D12", "C2", "D12", "C2" ], 
    [ "D12", "C2 x C2", "C2 x C2", "C2 x C2" ], [ "D12", "C2 x C2", "S3", "C2" ]
      , [ "D12", "C2 x C2", "D12", "C2 x C2" ], [ "D12", "S3", "S3", "S3" ], 
    [ "D12", "S3", "D12", "S3" ], [ "D12", "D12", "D12", "D12" ] ]
  !gapprompt@gap>| !gapinput@fams := AllCat2GroupFamilies( d12 );|
  [ [ 1, 2, 3, 4, 5, 6 ], [ 7, 8, 10, 11, 13, 14 ], [ 16, 17, 18, 23, 24, 25 ], 
    [ 30, 31, 32, 33, 34, 35 ], [ 9, 12, 15 ], [ 19, 20, 21, 26, 27, 28 ], 
    [ 36, 37, 38 ], [ 22, 29 ], [ 39, 40 ], [ 41 ] ]
  !gapprompt@gap>| !gapinput@CatnGroupNumbers( d12 );|
  rec( cat1 := 12, cat2 := 41, idem := 21, iso1 := 4, iso2 := 10, 
    isopredg := 0, predg := 0, siso := 4, symm := 12 )
  !gapprompt@gap>| !gapinput@CatnGroupLists( d12 );|
  rec( allcat2pos := [ 1, 7, 9, 16, 19, 22, 30, 36, 39, 41 ],
    cat2classes := 
      [ [ [ 1, 1 ], [ 2, 2 ], [ 3, 3 ], [ 4, 4 ], [ 5, 5 ], [ 6, 6 ] ], 
        [ [ 1, 7 ], [ 5, 7 ], [ 2, 8 ], [ 6, 8 ], [ 3, 9 ], [ 4, 9 ] ], 
        [ [ 1, 10 ], [ 2, 10 ], [ 3, 10 ], [ 4, 11 ], [ 5, 11 ], [ 6, 11 ] ], 
        [ [ 1, 12 ], [ 2, 12 ], [ 3, 12 ], [ 4, 12 ], [ 5, 12 ], [ 6, 12 ] ], 
        [ [ 7, 7 ], [ 8, 8 ], [ 9, 9 ] ], 
        [ [ 7, 10 ], [ 8, 10 ], [ 9, 10 ], [ 7, 11 ], [ 8, 11 ], [ 9, 11 ] ], 
        [ [ 7, 12 ], [ 8, 12 ], [ 9, 12 ] ], [ [ 10, 10 ], [ 11, 11 ] ], 
        [ [ 10, 12 ], [ 11, 12 ] ], [ [ 12, 12 ] ] ], 
    cat2pairs := [ [ 1, 1 ], [ 1, 7 ], [ 1, 10 ], [ 1, 12 ], [ 2, 2 ], 
        [ 2, 8 ], [ 2, 10 ], [ 2, 12 ], [ 3, 3 ], [ 3, 9 ], [ 3, 10 ], 
        [ 3, 12 ], [ 4, 4 ], [ 4, 9 ], [ 4, 11 ], [ 4, 12 ], [ 5, 5 ], 
        [ 5, 7 ], [ 5, 11 ], [ 5, 12 ], [ 6, 6 ], [ 6, 8 ], [ 6, 11 ], 
        [ 6, 12 ], [ 7, 7 ], [ 7, 10 ], [ 7, 11 ], [ 7, 12 ], [ 8, 8 ], 
        [ 8, 10 ], [ 8, 11 ], [ 8, 12 ], [ 9, 9 ], [ 9, 10 ], [ 9, 11 ], 
        [ 9, 12 ], [ 10, 10 ], [ 10, 12 ], [ 11, 11 ], [ 11, 12 ], [ 12, 12 ] ],
    omit := false, pisopos := [  ], sisopos := [ 1, 5, 8, 10 ] )
  
\end{Verbatim}
 }

 }

         
\chapter{\textcolor{Chapter }{Cat$^3$\texttt{\symbol{45}}groups and Crossed cubes}}\label{chap-obj4}
\logpage{[ 9, 0, 0 ]}
\hyperdef{L}{X7DBA3A7E81C71A64}{}
{
  \index{4d\texttt{\symbol{45}}group} \index{4d\texttt{\symbol{45}}domain} \index{cat$^3$\texttt{\symbol{45}}group} The term \emph{4d\texttt{\symbol{45}}group} refers to a set of equivalent categories of which the ones we are most
interested are the categories of \emph{crossed cubes} and \emph{cat$^3$\texttt{\symbol{45}}groups}. A \emph{4d\texttt{\symbol{45}}mapping} is a function between two 4d\texttt{\symbol{45}}groups which preserves all the
structure. 

 The material in this chapter should be considered very experimental. As yet
there are no functions for crossed cubes. 
\section{\textcolor{Chapter }{Functions for (pre\texttt{\symbol{45}})cat$^3$\texttt{\symbol{45}}groups}}\label{sect-cat3-functions}
\logpage{[ 9, 1, 0 ]}
\hyperdef{L}{X7CC52AF4840F478E}{}
{
  We shall use the following standard orientation of a cat$^3$\texttt{\symbol{45}}group $\calE$ on a group $G$. $\calE$ contains $8$ groups; $12$ cat$^1$\texttt{\symbol{45}}groups and $6$ cat$^2$\texttt{\symbol{45}}groups forming the vertices; edges and faces of a cube, as
shown in the following diagram. 
\[ \vcenter{\xymatrix{ & H \ar[dddd] <+0.5ex> \ar[dddd] <+0.0ex> \ar[rrrr]
<+0.5ex> \ar[rrrr] <+0.0ex> \ar[dr] <+0.6ex> & & & & N \ar[llll] <+0.6ex>
\ar[dddd] <+0.6ex> \ar[dddd] <+0.0ex> \ar[dr] <+0.6ex> & \\ & & G \ar[dddd]
<+0.5ex> \ar[dddd] <+0.0ex> \ar[rrrr] <+0.5ex> \ar[rrrr] <+0.0ex> \ar[ul]
<+0.5ex> \ar[ul] <+0.0ex> & & & & R \ar[llll] <+0.6ex> \ar[dddd] <+0.5ex>
\ar[dddd] <+0.0ex> \ar[ul] <+0.5ex> \ar[ul] <+0.0ex> \\ & & & & & & \\ & & & &
& & \\ & M \ar[uuuu] <+0.6ex> \ar[rrrr] <+0.5ex> \ar[rrrr] <+0.0ex> \ar[dr]
<+0.6ex> & & & & L \ar[uuuu] <+0.6ex> \ar[llll] <+0.6ex> \ar[dr] <+0.6ex> & \\
& & Q \ar[uuuu] <+0.6ex> \ar[rrrr] <+0.5ex> \ar[rrrr] <+0.0ex> \ar[ul]
<+0.5ex> \ar[ul] <+0.0ex> & & & & P \ar[uuuu] <+0.6ex> \ar[llll] <+0.6ex>
\ar[ul] <+0.5ex> \ar[ul] <+0.0ex> \\ }} \]
 

 By definition, $\calE$ is generated by three commuting cat$^1$\texttt{\symbol{45}}groups $(G \Rightarrow R), (G \Rightarrow Q)$ and $(G \Rightarrow H)$, but it is more convenient to think of $\calE$ as generated by two cat$^2$\texttt{\symbol{45}}groups 
\begin{itemize}
\item  \emph{front}$(\calE)$, generated by $(G \Rightarrow R)$ and $(G \Rightarrow Q)$; 
\item  \emph{left}$(\calE)$, generated by $(G \Rightarrow Q)$ and $(G \Rightarrow H)$. 
\end{itemize}
 Because the tail, head and embedding maps all commute, it follows that \emph{up}$(\calE)$, generated by $(G \Rightarrow H)$ and $(G \Rightarrow R)$, is a third cat$^2$\texttt{\symbol{45}}group. The three remaining faces (cat$^2$\texttt{\symbol{45}}groups) \emph{right}$(\calE)$, \emph{down}$(\calE)$ and \emph{back}$(\calE)$ are then easily constructed. We shall always use the order [\emph{front, left, up, right, down, back}] for the six faces. 

 

\subsection{\textcolor{Chapter }{Cat3Group}}
\logpage{[ 9, 1, 1 ]}\nobreak
\label{cat3-group}
\hyperdef{L}{X7828496F7D72E232}{}
{\noindent\textcolor{FuncColor}{$\triangleright$\enspace\texttt{Cat3Group({\mdseries\slshape args})\index{Cat3Group@\texttt{Cat3Group}}
\label{Cat3Group}
}\hfill{\scriptsize (function)}}\\
\noindent\textcolor{FuncColor}{$\triangleright$\enspace\texttt{PreCat3Group({\mdseries\slshape args})\index{PreCat3Group@\texttt{PreCat3Group}}
\label{PreCat3Group}
}\hfill{\scriptsize (function)}}\\
\noindent\textcolor{FuncColor}{$\triangleright$\enspace\texttt{IsCat3Group({\mdseries\slshape C})\index{IsCat3Group@\texttt{IsCat3Group}}
\label{IsCat3Group}
}\hfill{\scriptsize (property)}}\\
\noindent\textcolor{FuncColor}{$\triangleright$\enspace\texttt{PreCat3GroupByPreCat2Groups({\mdseries\slshape L})\index{PreCat3GroupByPreCat2Groups@\texttt{PreCat3GroupByPreCat2Groups}}
\label{PreCat3GroupByPreCat2Groups}
}\hfill{\scriptsize (operation)}}\\


 The global functions \texttt{Cat3Group} and \texttt{PreCat3Group} normally take as arguments a pair of cat$^2$\texttt{\symbol{45}}groups [\emph{front, left}], or a trio of cat$^1$\texttt{\symbol{45}}groups [\emph{front\texttt{\symbol{45}}up, front\texttt{\symbol{45}}left =
left\texttt{\symbol{45}}up, left\texttt{\symbol{45}}left}]. 

 In subsection \texttt{AllCat2GroupsIterator} (\ref{AllCat2GroupsIterator}) the list of pairs \texttt{CatnGroupLists(d12).pairs} contains the three entries \texttt{[6,8],[8,11]} and \texttt{[6,11]}. It follows that the sixth, eighth and eleventh cat$^1$\texttt{\symbol{45}}groups for \texttt{d12} generate a cat$^3$\texttt{\symbol{45}}group. }

 

 
\begin{Verbatim}[commandchars=!@|,fontsize=\small,frame=single,label=Example]
  
  !gapprompt@gap>| !gapinput@alld12 := AllCat1Groups( d12 );; |
  !gapprompt@gap>| !gapinput@C68 := Cat2Group( alld12[6], alld12[8] );; |
  !gapprompt@gap>| !gapinput@C811 := Cat2Group( alld12[8], alld12[11] );;|
  !gapprompt@gap>| !gapinput@C3Ga := Cat3Group( C68, C811 );|
  cat3-group with generating (pre-)cat1-groups:
  1 : [d12 => Group( [ (), (1,6)(2,5)(3,4) ] )]
  2 : [d12 => Group( [ (1,4)(2,5)(3,6), (1,3)(4,6) ] )]
  3 : [d12 => Group( [ (1,5,3)(2,6,4), (1,4)(2,3)(5,6) ] )]
  !gapprompt@gap>| !gapinput@C3Gb := Cat3Group( alld12[6], alld12[8], alld12[11] );;|
  !gapprompt@gap>| !gapinput@C3Ga = C3Gb;|
  true
  
\end{Verbatim}
 

\subsection{\textcolor{Chapter }{Front3DimensionalGroup (for cat3-groups)}}
\logpage{[ 9, 1, 2 ]}\nobreak
\hyperdef{L}{X85A5A6967D942463}{}
{\noindent\textcolor{FuncColor}{$\triangleright$\enspace\texttt{Front3DimensionalGroup({\mdseries\slshape C3})\index{Front3DimensionalGroup@\texttt{Front3DimensionalGroup}!for cat3-groups}
\label{Front3DimensionalGroup:for cat3-groups}
}\hfill{\scriptsize (attribute)}}\\
\noindent\textcolor{FuncColor}{$\triangleright$\enspace\texttt{Left3DimensionalGroup({\mdseries\slshape C3})\index{Left3DimensionalGroup@\texttt{Left3DimensionalGroup}!for cat3-groups}
\label{Left3DimensionalGroup:for cat3-groups}
}\hfill{\scriptsize (attribute)}}\\
\noindent\textcolor{FuncColor}{$\triangleright$\enspace\texttt{Up3DimensionalGroup({\mdseries\slshape C3})\index{Up3DimensionalGroup@\texttt{Up3DimensionalGroup}!for cat3-groups}
\label{Up3DimensionalGroup:for cat3-groups}
}\hfill{\scriptsize (attribute)}}\\
\noindent\textcolor{FuncColor}{$\triangleright$\enspace\texttt{Right3DimensionalGroup({\mdseries\slshape C3})\index{Right3DimensionalGroup@\texttt{Right3DimensionalGroup}!for cat3-groups}
\label{Right3DimensionalGroup:for cat3-groups}
}\hfill{\scriptsize (attribute)}}\\
\noindent\textcolor{FuncColor}{$\triangleright$\enspace\texttt{Down3DimensionalGroup({\mdseries\slshape C3})\index{Down3DimensionalGroup@\texttt{Down3DimensionalGroup}!for cat3-groups}
\label{Down3DimensionalGroup:for cat3-groups}
}\hfill{\scriptsize (attribute)}}\\
\noindent\textcolor{FuncColor}{$\triangleright$\enspace\texttt{Back3DimensionalGroup({\mdseries\slshape C3})\index{Back3DimensionalGroup@\texttt{Back3DimensionalGroup}!for cat3-groups}
\label{Back3DimensionalGroup:for cat3-groups}
}\hfill{\scriptsize (attribute)}}\\


 The six faces of a cat$^3$\texttt{\symbol{45}}group are stored as these attributes. }

 

 
\begin{Verbatim}[commandchars=!@|,fontsize=\small,frame=single,label=Example]
  
  !gapprompt@gap>| !gapinput@C116 := Cat2Group( alld12[11], alld12[6] );;|
  !gapprompt@gap>| !gapinput@Up3DimensionalGroup( C3Ga ) = C116;|
  true
  
\end{Verbatim}
 }

 
\section{\textcolor{Chapter }{Enumerating cat$^3$\texttt{\symbol{45}}groups with a given source}}\label{sect-allcat3}
\logpage{[ 9, 2, 0 ]}
\hyperdef{L}{X80E074E37D02B2F6}{}
{
  Once the list \texttt{CatnGroupLists(G).pairs} has been obtained we may seek all triples $[i,j],[j,k]$ and $[k,i]$ or $[i,k]$ of pairs in this list and then, for each such triple, construct a cat$^3$\texttt{\symbol{45}}group generated by the $i$\texttt{\symbol{45}}th, $j$\texttt{\symbol{45}}th and $k$\texttt{\symbol{45}}th cat$^1$\texttt{\symbol{45}}group on $G$. 

 

\subsection{\textcolor{Chapter }{AllCat3GroupTriples}}
\logpage{[ 9, 2, 1 ]}\nobreak
\hyperdef{L}{X83E4A7367DD13598}{}
{\noindent\textcolor{FuncColor}{$\triangleright$\enspace\texttt{AllCat3GroupTriples({\mdseries\slshape G})\index{AllCat3GroupTriples@\texttt{AllCat3GroupTriples}}
\label{AllCat3GroupTriples}
}\hfill{\scriptsize (operation)}}\\
\noindent\textcolor{FuncColor}{$\triangleright$\enspace\texttt{AllCat3GroupsNumber({\mdseries\slshape G})\index{AllCat3GroupsNumber@\texttt{AllCat3GroupsNumber}}
\label{AllCat3GroupsNumber}
}\hfill{\scriptsize (attribute)}}\\
\noindent\textcolor{FuncColor}{$\triangleright$\enspace\texttt{AllCat3Groups({\mdseries\slshape G})\index{AllCat3Groups@\texttt{AllCat3Groups}}
\label{AllCat3Groups}
}\hfill{\scriptsize (operation)}}\\


 The list of triples returned by the operation \texttt{AllCat3GroupTriples} is saved as \texttt{CatnGroupLists(G).cat3triples}. The length of this list is the number of cat$^3$\texttt{\symbol{45}}groups on $G$, and is saved as \texttt{CatnGroupNumbers(G).cat3}. 

 As yet there is no operation \texttt{AllCat3GroupsUpToIsomorphism(G)}. }

 

 
\begin{Verbatim}[commandchars=!@|,fontsize=\small,frame=single,label=Example]
  
  !gapprompt@gap>| !gapinput@triples := AllCat3GroupTriples( d12 );;|
  !gapprompt@gap>| !gapinput@CatnGroupNumbers( d12 ).cat3; |
  94
  !gapprompt@gap>| !gapinput@triples[46];|
  [ 5, 7, 11 ]
  !gapprompt@gap>| !gapinput@alld12 := AllCat1Groups( d12 );; |
  !gapprompt@gap>| !gapinput@Cat3Group( alld12[5], alld12[7], alld12[11] );|
  cat3-group with generating (pre-)cat1-groups:
  1 : [d12 => Group( [ (), (1,4)(2,3)(5,6) ] )]
  2 : [d12 => Group( [ (1,4)(2,5)(3,6), (2,6)(3,5) ] )]
  3 : [d12 => Group( [ (1,5,3)(2,6,4), (1,4)(2,3)(5,6) ] )]
  
\end{Verbatim}
 }

 
\section{\textcolor{Chapter }{ Definition and constructions for cat$^n$\texttt{\symbol{45}}groups and their morphisms }}\label{sect-catn-definitions}
\logpage{[ 9, 3, 0 ]}
\hyperdef{L}{X7F8538B580847268}{}
{
  \index{cat$^n$\texttt{\symbol{45}}group} In this chapter and the previous one we are interested in cat$^2$\texttt{\symbol{45}}groups and cat$^3$\texttt{\symbol{45}}groups, and it is convenient in this section to give the
more general definition. There are three equivalent descriptions of a cat$^n$\texttt{\symbol{45}}group. 

 A \emph{cat$^n$\texttt{\symbol{45}}group} consists of the following. 
\begin{itemize}
\item  $2^n$ groups $G_A$, one for each subset $A$ of $[n]$, the \emph{vertices} of an $n$\texttt{\symbol{45}}cube. 
\item  Group homomorphisms forming $n2^{n-1}$ commuting cat$^1$\texttt{\symbol{45}}groups, 
\[ \calC_{A,i} ~=~ (e_{A,i};\; t_{A,i},\; h_{A,i} \ :\ G_A \to G_{A \setminus
\{i\}}), \quad\mbox{for all} \quad A \subseteq [n],~ i \in A, \]
 the \emph{edges} of the cube. 
\item  These cat$^1$\texttt{\symbol{45}}groups combine (in sets of $4$) to form $n(n-1)2^{n-3}$ cat$^2$\texttt{\symbol{45}}groups $\calC_{A,\{i,j\}}$ for all $\{i,j\} \subseteq A \subseteq [n],~ i \neq j$, the \emph{faces} of the cube. 
\end{itemize}
 Note that, since the $t_{A,i}, h_{A,i}$ and $e_{A,i}$ commute, composite homomorphisms $t_{A,B}, h_{A,B} : G_A \to G_{A \setminus B}$ and $e_{A,B} : G_{A \setminus B} \to G_A$ are well defined for all $B \subseteq A \subseteq [n]$. 

 Secondly, we give the simplest of the three descriptions, again adapted from
Ellis\texttt{\symbol{45}}Steiner \cite{ell:st}. 

 A cat$^n$\texttt{\symbol{45}}group $\calC$ consists of $2^n$ groups $G_A$, one for each subset $A$ of $[n]$, and $3n$ homomorphisms 
\[ t_{[n],i}, h_{[n],i} : G_{[n]} \to G_{[n] \setminus \{i\}},~ e_{[n],i} :
G_{[n] \setminus \{i\}} \to G_{[n]}, \]
 satisfying the following axioms for all $1 \leqslant i \leqslant n$,\texttt{\symbol{125}} 
\begin{itemize}
\item  the $\calC_{[n],i} ~=~ (e_{[n],i};\; t_{[n],i},\; h_{[n],i} \ :\ G_{[n]} \to G_{[n]
\setminus \{i\}})~$ are \emph{commuting} cat$^1$\texttt{\symbol{45}}groups, so that: 
\item  $ (e_1 \circ t_1) \circ (e_2 \circ t_2) = (e_2 \circ t_2) \circ (e_1 \circ t_1),
\quad (e_1 \circ h_1) \circ (e_2 \circ h_2) = (e_2 \circ h_2) \circ (e_1 \circ
h_1), $ 
\item  $ (e_1 \circ t_1) \circ (e_2 \circ h_2) = (e_2 \circ h_2) \circ (e_1 \circ t_1),
\quad (e_2 \circ t_2) \circ (e_1 \circ h_1) = (e_1 \circ h_1) \circ (e_2 \circ
t_2). $ 
\end{itemize}
 Our third description defines a cat$^n$\texttt{\symbol{45}}group as a "cat$^1$\texttt{\symbol{45}}group of cat$^{(n-1)}$\texttt{\symbol{45}}groups". 

 A \emph{cat$^n$\texttt{\symbol{45}}group} $\calC$ consists of two cat$^{(n-1)}$\texttt{\symbol{45}}groups: 
\begin{itemize}
\item  $\calA$ with groups $G_A,\; A \subseteq [n-1]$, and homomorphisms $\ddot{t}_{A,i}, \ddot{h}_{A,i}, \ddot{e}_{A,i}$, 
\item  $\calB$ with groups $H_B,\; B \subseteq [n-1]$, and homomorphisms $\dot{t}_{B,i}, \dot{h}_{B,i}, \dot{e}_{B,i}$, and 
\item  cat$^{(n-1)}$\texttt{\symbol{45}}morphisms $t,h : \calA \to \calB$ and $e : \calB \to \calA$ subject to the following conditions: 
\[ (t \circ e) ~\mbox{and}~ (h \circ e) ~\mbox{are the identity mapping on}~
\calB, \qquad [\ker t, \ker h] = \{ 1_{\calA} \}. \]
 
\end{itemize}
 

\subsection{\textcolor{Chapter }{PreCatnGroup}}
\logpage{[ 9, 3, 1 ]}\nobreak
\hyperdef{L}{X81C3D39A81B20D76}{}
{\noindent\textcolor{FuncColor}{$\triangleright$\enspace\texttt{PreCatnGroup({\mdseries\slshape L})\index{PreCatnGroup@\texttt{PreCatnGroup}}
\label{PreCatnGroup}
}\hfill{\scriptsize (operation)}}\\
\noindent\textcolor{FuncColor}{$\triangleright$\enspace\texttt{CatnGroup({\mdseries\slshape L})\index{CatnGroup@\texttt{CatnGroup}}
\label{CatnGroup}
}\hfill{\scriptsize (operation)}}\\


 The operation \texttt{(Pre)CatnGroup} expects as input a list of cat$^1$\texttt{\symbol{45}}groups. For our group \texttt{d12} we may construct various cat$^4$\texttt{\symbol{45}}groups, and here is one of them. 

 }

 

 
\begin{Verbatim}[commandchars=!@|,fontsize=\small,frame=single,label=Example]
  
  !gapprompt@gap>| !gapinput@PC4 := PreCatnGroup( [ alld12[5], alld12[7], alld12[11], alld12[12] ] );|
  (pre-)cat4-group with generating (pre-)cat1-groups:
  1 : [d12 => Group( [ (), (1,4)(2,3)(5,6) ] )]
  2 : [d12 => Group( [ (1,4)(2,5)(3,6), (2,6)(3,5) ] )]
  3 : [d12 => Group( [ (1,5,3)(2,6,4), (1,4)(2,3)(5,6) ] )]
  4 : [d12 => Group( [ (1,2,3,4,5,6), (2,6)(3,5) ] )]
  !gapprompt@gap>| !gapinput@IsCatnGroup( PC4 );                                             |
  true
  !gapprompt@gap>| !gapinput@HigherDimension( PC4 );|
  5
  
\end{Verbatim}
 For a cat$^5$\texttt{\symbol{45}}group we may start with the cyclic group whose order is
the product of the first five primes. With this group we may form $32$ cat$^1$\texttt{\symbol{45}}groups and $528$ cat$^2$\texttt{\symbol{45}}groups. 

 
\begin{Verbatim}[commandchars=!@|,fontsize=\small,frame=single,label=Example]
  
  !gapprompt@gap>| !gapinput@G := Group( (1,2), (3,4,5), (6,7,8,9,10), (11,12,13,14,15,16,17),|
  !gapprompt@>| !gapinput@               (20,21,22,23,24,25,26,27,28,29,30) );;|
  !gapprompt@gap>| !gapinput@SetName( G, "C2310" );|
  !gapprompt@gap>| !gapinput@all1 := AllCat1Groups( G );;|
  !gapprompt@gap>| !gapinput@Print( "G has ", CatnGroupNumbers( G ).cat1, " cat1-groups\n" );|
  G has 32 cat1-groups
  !gapprompt@gap>| !gapinput@PC5 := PreCatnGroup( [ all1[2], all1[5], all1[13], all1[25], all1[32] ] );|
  (pre-)cat5-group with generating (pre-)cat1-groups:
  1 : [C2310 => Group( [ (), (), (), (), (1,2) ] )]
  2 : [C2310 => Group( [ (), (), (), (3,4,5), (1,2) ] )]
  3 : [C2310 => Group( [ (), (), ( 6, 7, 8, 9,10), (3,4,5), (1,2) ] )]
  4 : [C2310 => Group( [ (), (11,12,13,14,15,16,17), ( 6, 7, 8, 9,10), (3,4,5), 
    (1,2) ] )]
  5 : [C2310 => Group( [ (20,21,22,23,24,25,26,27,28,29,30), 
    (11,12,13,14,15,16,17), ( 6, 7, 8, 9,10), (3,4,5), (1,2) ] )]
  !gapprompt@gap>| !gapinput@HigherDimension( PC5 );|
  6
  
\end{Verbatim}
 }

 }

         
\chapter{\textcolor{Chapter }{Crossed modules of groupoids}}\label{chap-gpd2o}
\logpage{[ 10, 0, 0 ]}
\hyperdef{L}{X80B3A81B7E5CA3A9}{}
{
  \index{crossed module of groupoids} The material documented in this chapter is experimental, and is likely to be
changed in due course. 
\section{\textcolor{Chapter }{Constructions for crossed modules of groupoids}}\logpage{[ 10, 1, 0 ]}
\hyperdef{L}{X847F4ED77F50528C}{}
{
  \index{crossed module over a groupoid} \index{2dimensional\texttt{\symbol{45}}domain with objects} A typical example of a crossed module $\calX$ over a groupoid has for its range a connected groupoid. This is a direct
product of a group with a complete graph, and we call the vertices of the
graph the \emph{objects} of the crossed module. The source of $\calX$ is a groupoid, with the same objects, which is either discrete or connected.
The boundary morphism is constant on objects. For details and other references
see \cite{AW2}. 

\subsection{\textcolor{Chapter }{PreXModWithObjectsByBoundaryAndAction}}
\logpage{[ 10, 1, 1 ]}\nobreak
\hyperdef{L}{X78F89CAB7A281B8F}{}
{\noindent\textcolor{FuncColor}{$\triangleright$\enspace\texttt{PreXModWithObjectsByBoundaryAndAction({\mdseries\slshape bdy, act})\index{PreXModWithObjectsByBoundaryAndAction@\texttt{Pre}\-\texttt{X}\-\texttt{Mod}\-\texttt{With}\-\texttt{Objects}\-\texttt{By}\-\texttt{Boundary}\-\texttt{And}\-\texttt{Action}}
\label{PreXModWithObjectsByBoundaryAndAction}
}\hfill{\scriptsize (operation)}}\\


 This is the groupoid generalisation of the operation \texttt{PreXModByBoundaryAndAction}. }

 

\subsection{\textcolor{Chapter }{SinglePiecePreXModWithObjects}}
\logpage{[ 10, 1, 2 ]}\nobreak
\hyperdef{L}{X86CD034F82F5F029}{}
{\noindent\textcolor{FuncColor}{$\triangleright$\enspace\texttt{SinglePiecePreXModWithObjects({\mdseries\slshape pxmod, obs, isdisc})\index{SinglePiecePreXModWithObjects@\texttt{SinglePiecePreXModWithObjects}}
\label{SinglePiecePreXModWithObjects}
}\hfill{\scriptsize (operation)}}\\


 At present the experimental operation \texttt{SinglePiecePreXModWithObjects} accepts a precrossed module \texttt{pxmod}, a set of objects \texttt{obs}, and a boolean \texttt{isdisc} which is \texttt{true} when the source groupoid is homogeneous and discrete and \texttt{false} when the source groupoid is connected. Other operations will be added as time
permits. 

 In the example the crossed module \texttt{DX4} has discrete source, while the crossed module \texttt{CX4} has connected source. These are groupoid generalisations of \texttt{XModByNormalSubgroup} (\ref{XModByNormalSubgroup}) and the example \texttt{X4} in \texttt{NormalSubXMods} (\ref{NormalSubXMods}).      }

 

 
\begin{Verbatim}[commandchars=!@|,fontsize=\small,frame=single,label=Example]
  
  !gapprompt@gap>| !gapinput@CX4 := SinglePiecePreXModWithObjects( X4, [-6,-5,-4], false );|
  single piece precrossed module with objects
    source groupoid:
      single piece groupoid: < a4, [ -6, -5, -4 ] >
    and range groupoid:
      single piece groupoid: < s4, [ -6, -5, -4 ] >
  !gapprompt@gap>| !gapinput@SetName( CX4, "CX4" );|
  !gapprompt@gap>| !gapinput@Ca4 := Source( CX4 );;  SetName( Ca4, "Ca4" );|
  !gapprompt@gap>| !gapinput@Cs4 := Range( CX4 );;  SetName( Cs4, "Cs4" );|
  !gapprompt@gap>| !gapinput@DX4 := SinglePiecePreXModWithObjects( X4, [-9,-8,-7], true );|
  single piece precrossed module with objects
    source groupoid:
      homogeneous, discrete groupoid: < a4, [ -9, -8, -7 ] >
    and range groupoid:
      single piece groupoid: < s4, [ -9, -8, -7 ] >
  !gapprompt@gap>| !gapinput@SetName( DX4, "DX4" );|
  !gapprompt@gap>| !gapinput@Da4 := Source( DX4 );;  SetName( Da4, "Da4" );|
  !gapprompt@gap>| !gapinput@Ds4 := Range( DX4 );;  SetName( Ds4, "Ds4" );|
  
\end{Verbatim}
 

\subsection{\textcolor{Chapter }{IsXModWithObjects}}
\logpage{[ 10, 1, 3 ]}\nobreak
\hyperdef{L}{X7B76F2BF82E075FF}{}
{\noindent\textcolor{FuncColor}{$\triangleright$\enspace\texttt{IsXModWithObjects({\mdseries\slshape pxmod})\index{IsXModWithObjects@\texttt{IsXModWithObjects}}
\label{IsXModWithObjects}
}\hfill{\scriptsize (property)}}\\
\noindent\textcolor{FuncColor}{$\triangleright$\enspace\texttt{IsPreXModWithObjects({\mdseries\slshape pxmod})\index{IsPreXModWithObjects@\texttt{IsPreXModWithObjects}}
\label{IsPreXModWithObjects}
}\hfill{\scriptsize (property)}}\\


 \index{Is2DimensionalGroupWithObjects} The precrossed module \texttt{CX4} belongs to the category \texttt{Is2DimensionalGroupWithObjects} and is, of course, a crossed module. 

 }

 

 
\begin{Verbatim}[commandchars=!@|,fontsize=\small,frame=single,label=Example]
  
  !gapprompt@gap>| !gapinput@Set( KnownPropertiesOfObject( CX4 ) );|
  [ "CanEasilyCompareElements", "CanEasilySortElements", "IsAssociative", 
    "IsDuplicateFree", "IsGeneratorsOfSemigroup", "IsPreXModWithObjects", 
    "IsSinglePieceDomain", "IsXModWithObjects" ]
  
\end{Verbatim}
 

\subsection{\textcolor{Chapter }{IsPermPreXModWithObjects}}
\logpage{[ 10, 1, 4 ]}\nobreak
\hyperdef{L}{X858EB4F97D04D012}{}
{\noindent\textcolor{FuncColor}{$\triangleright$\enspace\texttt{IsPermPreXModWithObjects({\mdseries\slshape pxmod})\index{IsPermPreXModWithObjects@\texttt{IsPermPreXModWithObjects}}
\label{IsPermPreXModWithObjects}
}\hfill{\scriptsize (property)}}\\
\noindent\textcolor{FuncColor}{$\triangleright$\enspace\texttt{IsPcPreXModWithObjects({\mdseries\slshape pxmod})\index{IsPcPreXModWithObjects@\texttt{IsPcPreXModWithObjects}}
\label{IsPcPreXModWithObjects}
}\hfill{\scriptsize (property)}}\\
\noindent\textcolor{FuncColor}{$\triangleright$\enspace\texttt{IsFpPreXModWithObjects({\mdseries\slshape pxmod})\index{IsFpPreXModWithObjects@\texttt{IsFpPreXModWithObjects}}
\label{IsFpPreXModWithObjects}
}\hfill{\scriptsize (property)}}\\


 To test these properties we test the precrossed modules from which they were
constructed. 

 }

 

 
\begin{Verbatim}[commandchars=!@|,fontsize=\small,frame=single,label=Example]
  
  !gapprompt@gap>| !gapinput@IsPermPreXModWithObjects( CX4 );|
  true
  !gapprompt@gap>| !gapinput@IsPcPreXModWithObjects( CX4 );|
  false
  !gapprompt@gap>| !gapinput@IsFpPreXModWithObjects( CX4 );|
  false
  
\end{Verbatim}
 

\subsection{\textcolor{Chapter }{Root2dGroup}}
\logpage{[ 10, 1, 5 ]}\nobreak
\hyperdef{L}{X797B1CD07C3682EE}{}
{\noindent\textcolor{FuncColor}{$\triangleright$\enspace\texttt{Root2dGroup({\mdseries\slshape pxmod})\index{Root2dGroup@\texttt{Root2dGroup}}
\label{Root2dGroup}
}\hfill{\scriptsize (attribute)}}\\
\noindent\textcolor{FuncColor}{$\triangleright$\enspace\texttt{XModAction({\mdseries\slshape pxmod})\index{XModAction@\texttt{XModAction}!for crossed modules of groupoids}
\label{XModAction:for crossed modules of groupoids}
}\hfill{\scriptsize (attribute)}}\\


 \index{ObjectList} The attributes of a precrossed module with objects include the standard \texttt{Source}; \texttt{Range}; \texttt{Boundary} (\ref{Boundary:for crossed modules}); and \texttt{XModAction} (\ref{XModAction:for crossed modules of groups}) as with precrossed modules of groups. There is also \texttt{ObjectList}, as in the \textsf{groupoids} package. Additionally there is \texttt{Root2dGroup} which is the underlying precrossed module used in the construction. 

 Note that \texttt{XModAction} is now a groupoid homomorphism from the source groupoid to a
one\texttt{\symbol{45}}object groupoid (with object \texttt{0}) where the group is the automorphism group of the range groupoid. 

 }

 

 
\begin{Verbatim}[commandchars=!@|,fontsize=\small,frame=single,label=Example]
  
  !gapprompt@gap>| !gapinput@Root2dGroup( CX4 );|
  [a4->s4]
  !gapprompt@gap>| !gapinput@actC := XModAction( CX4 );;|
  !gapprompt@gap>| !gapinput@Size( Range( actC ) );|
  20736
  !gapprompt@gap>| !gapinput@r1 := Arrow( Cs4, (1,2,3,4), -4, -5 );;|
  !gapprompt@gap>| !gapinput@ir1 := ImageElm( actC, r1 );|
  [groupoid homomorphism : Ca4 -> Ca4
  [ [ [(1,2,3) : -6 -> -6], [(2,3,4) : -6 -> -6], [() : -6 -> -5], 
        [() : -6 -> -4] ], 
    [ [(2,3,4) : -4 -> -4], [(1,3,4) : -4 -> -4], [() : -4 -> -6], 
        [() : -4 -> -5] ] ] : 0 -> 0]
  !gapprompt@gap>| !gapinput@s1 := Arrow( Ca4, (1,2,4), -5, -5 );;|
  !gapprompt@gap>| !gapinput@##  calculate s1^r1 |
  !gapprompt@gap>| !gapinput@ims1 := ImageElmXModAction( CX4, s1, r1 );|
  [(1,2,3) : -6 -> -6]
  !gapprompt@gap>| !gapinput@## repeat using DX4|
  !gapprompt@gap>| !gapinput@actD := XModAction( DX4 );;|
  !gapprompt@gap>| !gapinput@SetName( Range( actD ), "autD" );|
  !gapprompt@gap>| !gapinput@Size( Range( actD ) );|
  82944
  !gapprompt@gap>| !gapinput@r2 := Arrow( Ds4, (1,2,3,4), -7, -8 );;|
  !gapprompt@gap>| !gapinput@ImageElm( actD, r2 );|
  [groupoid homomorphism : Da4 -> Da4
   : 0 -> 0]
  !gapprompt@gap>| !gapinput@s2 := Arrow( Da4, (1,2,4), -8, -8 );;|
  !gapprompt@gap>| !gapinput@##  calculate s2^r2 |
  !gapprompt@gap>| !gapinput@ims2 := ImageElmXModAction( DX4, s2, r2 );|
  [(1,2,3) : -9 -> -9]
  
\end{Verbatim}
 There is much more to be done with these constructions. }

 }

         
\chapter{\textcolor{Chapter }{Double Groupoids}}\label{chap-double}
\logpage{[ 11, 0, 0 ]}
\hyperdef{L}{X83B7E8A287C9284A}{}
{
  A \emph{double groupoid} is a \emph{double category} in which all the category structures are groupoids. There is also a
pre\texttt{\symbol{45}}crossed module associated to the double groupoid. In a
double groupoid, as well as objects and arrows we need a set of \emph{squares}. A square is bounded by four arrows, two horizontal and two vertical, and
there is a \emph{horizontal} groupoid structure and a \emph{vertical} groupoid structure on these squares. A square also contains an element of the
source group of the pre\texttt{\symbol{45}}crossed module, and its image under
the boundary map is equal to the boundary of the square. When printing a
square, this element is located at the centre, 

 The double groupoids constructed here are special in that all four arrows come
from the same groupoid. We call these \emph{edge\texttt{\symbol{45}}symmetric} double groupoids. 

 This material in this chapter is experimental. It was started in 2023, in
version 2.91, and extensively revised in 2025, for version 2.95. Further
extensions are likely in due course. 
\section{\textcolor{Chapter }{Constructions for Double Groupoids}}\label{sect-basic-double-groupoids}
\logpage{[ 11, 1, 0 ]}
\hyperdef{L}{X87AC8EF586C35CD4}{}
{
  It is assumed in this chapter that the reader is familiar with constructions
for groupoids given in the \textsf{Groupoids} package, such as \texttt{SinglePieceBasicDoubleGroupoid}. Such groupoids are \emph{basic}, in that there is no pre\texttt{\symbol{45}}crossed module involvement. In \textsf{XMod} the operation \texttt{SinglePieceDoubleGroupoid} requires a groupoid and a crossed module as input parameters. 

\subsection{\textcolor{Chapter }{SinglePieceDoubleGroupoid}}
\logpage{[ 11, 1, 1 ]}\nobreak
\hyperdef{L}{X78936F448231692E}{}
{\noindent\textcolor{FuncColor}{$\triangleright$\enspace\texttt{SinglePieceDoubleGroupoid({\mdseries\slshape gpd, pxmod})\index{SinglePieceDoubleGroupoid@\texttt{SinglePieceDoubleGroupoid}}
\label{SinglePieceDoubleGroupoid}
}\hfill{\scriptsize (operation)}}\\


 For an example we take for our groupoid the product of the group $S_3 = \langle (7,8,9), (8,9) \rangle$ with the complete graph on $[-6 \ldots -1]$ and, for our pre\texttt{\symbol{45}}crossed module, the \texttt{X12}, isomorphic to $(D_{12} \to S_3)$, constructed using \texttt{XModByCentralExtension} (\ref{XModByCentralExtension}). The source of \texttt{X12} has generating set $\left\{ g = (11,12,13,14,15,16),~ h = (12,16)(13,15) \right\}$. }

 
\begin{Verbatim}[commandchars=@|F,fontsize=\small,frame=single,label=Example]
  
  @gapprompt|gap>F @gapinput|gens3 := [ (7,8,9), (8,9) ];;F
  @gapprompt|gap>F @gapinput|s3 := Group( gens3 );;F
  @gapprompt|gap>F @gapinput|SetName( s3, "s3" ); F
  @gapprompt|gap>F @gapinput|Gs3 := Groupoid( s3, [-6..-1] );;F
  @gapprompt|gap>F @gapinput|SetName( Gs3, "Gs3" ); F
  @gapprompt|gap>F @gapinput|g := (11,12,13,14,15,16);;  h := (12,16)(13,15);;F
  @gapprompt|gap>F @gapinput|gend12 := [ g, h ];;F
  @gapprompt|gap>F @gapinput|d12 := Group( gend12 );;F
  @gapprompt|gap>F @gapinput|SetName( d12, "d12" ); F
  @gapprompt|gap>F @gapinput|pr12 := GroupHomomorphismByImages( d12, s3, gend12, gens3 );;F
  @gapprompt|gap>F @gapinput|X12 := XModByCentralExtension( pr12 );; F
  @gapprompt|gap>F @gapinput|SetName( X12, "X12" ); F
  @gapprompt|gap>F @gapinput|Display( X12 ); F
  
  Crossed module X12 :- 
  : Source group d12 has generators:
    [ (11,12,13,14,15,16), (12,16)(13,15) ]
  : Range group s3 has generators:
    [ (7,8,9), (8,9) ]
  : Boundary homomorphism maps source generators to:
    [ (7,8,9), (8,9) ]
  : Action homomorphism maps range generators to automorphisms:
    (7,8,9) --> { source gens --> [ (11,12,13,14,15,16), (11,13)(14,16) ] }
    (8,9) --> { source gens --> [ (11,16,15,14,13,12), (12,16)(13,15) ] }
    These 2 automorphisms generate the group of automorphisms.
  
  @gapprompt|gap>F @gapinput|D1 := SinglePieceDoubleGroupoid( Gs3, X12 );; F
  @gapprompt|gap>F @gapinput|D1!.groupoid;F
  Gs3
  @gapprompt|gap>F @gapinput|D1!.prexmod;F
  X12
  
\end{Verbatim}
 

\subsection{\textcolor{Chapter }{SquareOfArrows}}
\logpage{[ 11, 1, 2 ]}\nobreak
\hyperdef{L}{X823A3A7481B90EB7}{}
{\noindent\textcolor{FuncColor}{$\triangleright$\enspace\texttt{SquareOfArrows({\mdseries\slshape bdy, act})\index{SquareOfArrows@\texttt{SquareOfArrows}}
\label{SquareOfArrows}
}\hfill{\scriptsize (operation)}}\\


 Let $G$ be a groupoid with object set $\Omega$ and object group $g$. Let $\Box$ be the set of squares with objects from $\Omega$ at each corner; plus two vertical arrows and two horizontal arrows from Arr$(G)$. Further, let $\calP = (\partial : S \to R)$ be a pre\texttt{\symbol{45}}crossed module, and let $m_1 \in S$ be placed at the centre of the square. The following picture illustrates the
situation where $u_1,v_1,w_1,x_1 \in \Omega$ and $a_1,b_1,c_1,d_1 \in g$. 
\[  \vcenter{\xymatrix @=4pc{ u_1 \ar[r]^{a_1} \ar[d]_{b_1}\ar@{}[dr] |{m_1} & v_1
\ar[d]^{c_1} \\ w_1 \ar[r]_{d_1} & x_1 }} \]
 We think of the square being \emph{based} at the bottom, right\texttt{\symbol{45}}hand corner, $x_1$. The \emph{boundary} of the square is the loop $(x_1, d_1^{-1}b_1^{-1}a_1c_1,x_1) = (x_1,p_1,x_1)$. The \emph{boundary condition} which $m_1$ has to satisfy is that $\partial m_1 = p_1$. (Beware a possible source of confusion. In the example code it is convenient
of define arrow \texttt{a1 := Arrow( Gs3, (7,8), \texttt{\symbol{45}}6, \texttt{\symbol{45}}5 );} whereas, in the description, $a_1 = (7,8)$. Similarly for all the other arrows.) }

 
\begin{Verbatim}[commandchars=@CF,fontsize=\small,frame=single,label=Example]
  
  @gappromptCgap>F @gapinputCa1 := Arrow( Gs3, (7,8), -6, -5 );;F
  @gappromptCgap>F @gapinputCb1 := Arrow( Gs3, (7,9), -6, -1 );;F
  @gappromptCgap>F @gapinputCc1 := Arrow( Gs3, (8,9), -5, -3 );;F
  @gappromptCgap>F @gapinputCd1 := Arrow( Gs3, (7,8,9), -1, -3 );;F
  @gappromptCgap>F @gapinputCsq1 := SquareOfArrows( D1, g*h, a1, b1, c1, d1 ); F
  [-6] ------------- (7,8) ------------> [-5]
    |                                     |
  (7,9)        (11,16)(12,15)(13,14)        (8,9)
    V                                     V
  [-1] ------------ (7,8,9) -----------> [-3]
  @gappromptCgap>F @gapinputCm1 := ElementOfSquare( sq1 ); F
  (11,16)(12,15)(13,14)
  @gappromptCgap>F @gapinputCUpArrow( sq1 );F
  [(7,8) : -6 -> -5]
  @gappromptCgap>F @gapinputCLeftArrow( sq1 );F
  [(7,9) : -6 -> -1]
  @gappromptCgap>F @gapinputCRightArrow( sq1 );F
  [(8,9) : -5 -> -3]
  @gappromptCgap>F @gapinputCDownArrow( sq1 );F
  [(7,8,9) : -1 -> -3]
  @gappromptCgap>F @gapinputC## check the boundary condition:F
  @gappromptCgap>F @gapinputCbdy1 := Boundary( D1!.prexmod );;F
  @gappromptCgap>F @gapinputCp1:= BoundaryOfSquare( sq1 );F
  [(7,9) : -3 -> -3]
  @gappromptCgap>F @gapinputCImageElm( bdy1, m1 );F
  (7,9)
  
\end{Verbatim}
 

\subsection{\textcolor{Chapter }{HorizontalProduct}}
\logpage{[ 11, 1, 3 ]}\nobreak
\hyperdef{L}{X7D3737FA7E9E3ECA}{}
{\noindent\textcolor{FuncColor}{$\triangleright$\enspace\texttt{HorizontalProduct({\mdseries\slshape sq1, sq2})\index{HorizontalProduct@\texttt{HorizontalProduct}}
\label{HorizontalProduct}
}\hfill{\scriptsize (operation)}}\\


 When defining a \emph{horizontal composition}, as illustrated by 
\[ \vcenter{\xymatrix @=4pc{ u_1 \ar[r]^{a_1} \ar[d]_{b_1} \ar@{}[dr]|{m_1} & v_1
\ar[r]^{a_2} \ar[d]^{c_1} \ar@{}[dr]|{m_2} & v_2 \ar[d]^{c_2} \ar@{}[dr]|= &
u_1 \ar[r]^{a_1a_2} \ar[d]_{b_1} \ar@{}[dr]|{m_1^{d_2}m_2} & v_2 \ar[d]^{c_2}
\\ w_1 \ar[r]_{d_1} & x_1 \ar[r]_{d_2} & x_2 & w_1 \ar[r]_{d_1d_2} & x_2 }} \]
 we have to move $m_1$, based at $x_1$, to the new base $x_2$, and we do this by using the action of the pre\texttt{\symbol{45}}crossed
module of $d_2$ on $m_1$. Notice that the boundary condition for the product is satisfied, since the
first pre\texttt{\symbol{45}}crossed module axiom applies: 
\[ \partial(m_1^{d_2}m_2) ~=~ \partial(m_1^{d_2}) (\partial m_2) ~=~
d_2^{-1}(d_1^{-1}b_1^{-1}a_1c_1)d_2(d_2^{-1}c_1^{-1}a_2c_2) ~=~
(d_1d_2)^{-1}b_1^{-1}(a_1a_2)c_2. \]
 

 }

 
\begin{Verbatim}[commandchars=!@B,fontsize=\small,frame=single,label=Example]
  
  !gapprompt@gap>B !gapinput@a2 := Arrow( Gs3, (8,9), -5, -4 );; B
  !gapprompt@gap>B !gapinput@c2 := Arrow( Gs3, (7,8), -4, -4 );;B
  !gapprompt@gap>B !gapinput@d2 := Arrow( Gs3, (7,9), -3, -4 );;B
  !gapprompt@gap>B !gapinput@sq2 := SquareOfArrows( D1, g^2, a2, c1, c2, d2 ); B
  [-5] ------------ (8,9) -----------> [-4]
    |                                    |
  (8,9)        (11,13,15)(12,14,16)        (7,8)
    V                                    V
  [-3] ------------ (7,9) -----------> [-4]
  !gapprompt@gap>B !gapinput@m2 := ElementOfSquare( sq2 ); B
  (11,13,15)(12,14,16)
  !gapprompt@gap>B !gapinput@LeftArrow( sq2 ) = RightArrow( sq1 ); B
  true
  !gapprompt@gap>B !gapinput@sq12 := HorizontalProduct( sq1, sq2 );B
  [-6] ------------ (7,9,8) -----------> [-4]
    |                                     |
  (7,9)        (11,12)(13,16)(14,15)        (7,8)
    V                                     V
  [-1] ------------- (7,8) ------------> [-4]
  !gapprompt@gap>B !gapinput@## check the boundary condition:B
  !gapprompt@gap>B !gapinput@ed2 := ElementOfArrow( d2 );;B
  !gapprompt@gap>B !gapinput@im1d2 := ImageElmXModAction( X12, m1, ed2 );B
  (11,16)(12,15)(13,14)
  !gapprompt@gap>B !gapinput@m12 := ElementOfSquare( sq12 );B
  (11,12)(13,16)(14,15)
  !gapprompt@gap>B !gapinput@im1d2 * m2 = m12;B
  true
  
\end{Verbatim}
 

\subsection{\textcolor{Chapter }{VerticalProduct}}
\logpage{[ 11, 1, 4 ]}\nobreak
\hyperdef{L}{X873F01287A2DC41F}{}
{\noindent\textcolor{FuncColor}{$\triangleright$\enspace\texttt{VerticalProduct({\mdseries\slshape sq1, sq2})\index{VerticalProduct@\texttt{VerticalProduct}}
\label{VerticalProduct}
}\hfill{\scriptsize (operation)}}\\


 Similarly, vertical composition is illustrated by 
\[ \vcenter{\xymatrix @=2pc{ u_1 \ar[rr]^{a_1} \ar[dd]_{b_1} \ar@{}[ddrr]|{m_1}
&& v_1 \ar[dd]^{c_1} & & && \\ && & & u_1 \ar[rr]^{a_1} \ar[dd]_{b_1b_3}
\ar@{}[ddrr]|{m_3m_1^{c_3}} && v_1 \ar[dd]^{c_1c_3} \\ w_1 \ar[rr]_{d_1}
\ar[dd]_{b_3} \ar@{}[ddrr]|{m_3} && x_1 \ar[dd]^{c_3} &=& && \\ && & & w_3
\ar[rr]_{d_3} && x_3 \\ w_3 \ar[rr]_{d_3} && x_3 }} \]
 

 Again the boundary condition is satisfied: 
\[ \partial(m_3m_1^{c_3}) ~=~ (\partial m_3) \partial(m_1^{c_3}) ~=~
(d_3^{-1}b_3^{-1}d_1c_3)c_3^{-1}(d_1^{-1}b_1^{-1}a_1c_1)c_3 ~=~
d_3^{-1}(b_1b_3)^{-1}a_1(c_1c_3). \]
 

 }

 
\begin{Verbatim}[commandchars=!@B,fontsize=\small,frame=single,label=Example]
  
  !gapprompt@gap>B !gapinput@b3 := Arrow( Gs3, (8,9), -1, -2 );;B
  !gapprompt@gap>B !gapinput@c3 := Arrow( Gs3, (7,9,8), -3, -2 );;B
  !gapprompt@gap>B !gapinput@d3 := Arrow( Gs3, (7,9), -2, -2 );;B
  !gapprompt@gap>B !gapinput@sq3 := SquareOfArrows( D1, g, d1, b3, c3, d3 );B
  [-1] ------------ (7,8,9) -----------> [-3]
    |                                     |
  (8,9)        (11,12,13,14,15,16)        (7,9,8)
    V                                     V
  [-2] ------------- (7,9) ------------> [-2]
  !gapprompt@gap>B !gapinput@m3 := ElementOfSquare( sq3 ); B
  (11,12,13,14,15,16)
  !gapprompt@gap>B !gapinput@UpArrow( sq3 ) = DownArrow( sq1 ); B
  true
  !gapprompt@gap>B !gapinput@sq13 := VerticalProduct( sq1, sq3 ); B
  [-6] ---------- (7,8) ---------> [-5]
    |                                |
  (7,8,9)        (11,13)(14,16)        (7,9)
    V                                V
  [-2] ---------- (7,9) ---------> [-2]
  !gapprompt@gap>B !gapinput@## check the boundary condition:B
  !gapprompt@gap>B !gapinput@ec3 := ElementOfArrow( c3 );;B
  !gapprompt@gap>B !gapinput@im1c3 := ImageElmXModAction( X12, m1, ec3 );B
  (11,14)(12,13)(15,16)
  !gapprompt@gap>B !gapinput@m13 := ElementOfSquare( sq13 );B
  (11,13)(14,16)
  !gapprompt@gap>B !gapinput@m3 * im1c3 = m13;B
  true
  
\end{Verbatim}
 The \texttt{HorizontalProduct} and \texttt{VerticalProduct} commute, so we may construct products such as: 
\[ \vcenter{\xymatrix @=2pc{ u_1 \ar[rr]^{a_1} \ar[dd]_{b_1} \ar@{}[ddrr]|{m_1}
&& v_1 \ar[rr]^{a_2} \ar[dd]|{c_1} \ar@{}[ddrr]|{m_2} && v_2 \ar[dd]^{c_2} & &
&&& \\ && && & & u_1 \ar[rrr]^{a_1a_2} \ar[dd]_{b_1b_3}
\ar@{}[ddrrr]|{m_3^{d_4}m_4\left(m_1^{d_2}m_2\right)^{c_4}} &&& v_2
\ar[dd]^{c_1c_4} \\ w_1 \ar[rr]|{d_1} \ar[dd]_{b_3} \ar@{}[ddrr]|{m_3} && x_1
\ar[rr]|{d_2} \ar[dd]|{c_3} \ar@{}[ddrr]|{m_4} && x_2 \ar[dd]^{c_4} &=& &&& \\
&& && & & w_3 \ar[rrr]_{d_3d_4} &&& x_4 \\ w_3 \ar[rr]_{d_3} && x_3
\ar[rr]_{d_4} && x_4 }} \]
 where 
\[ m_3^{d_4}m_4 (m_1^{d_2}m_2)^{c_4} ~=~ (m_3m_1^{c_3})^{d_4} m_4m_2^{c_4} ~=~
(d_3d_4)^{-1}(b_1b_3)^{-1}(a_1a_2)(c_3c_4). \]
 

 Continuing with our example, we check that the two ways of computing the
product of four squares below agree. 
\[ \vcenter{\xymatrix @=2pc{ -6 \ar[rr]^{(7,8)} \ar[dd]_{(7,9)} \ar@{}[ddrr]|{gh}
&& -5 \ar[rr]^{(8,9)} \ar[dd]|{(8,9)} \ar@{}[ddrr]|{g^2} && -4 \ar[dd]^{(7,8)}
& & &&& \\ && && & & -6 \ar[rrr]^{(7,9,8)} \ar[dd]_{(7,8,9)}
\ar@{}[ddrrr]|{(11,15,13)(12,16,14)} &&& -4 \ar[dd]^{(7,9,8)} \\ -1
\ar[rr]|{(7,8,9)} \ar[dd]_{(8,9)} \ar@{}[ddrr]|{g} && -3 \ar[rr]|{(7,9)}
\ar[dd]|{(7,9,8)} \ar@{}[ddrr]|{h} && -4 \ar[dd]^{(8,9)} &=& &&& \\ && && & &
-2 \ar[rrr]_{(7,9,8)} &&& -3 \\ -2 \ar[rr]_{(7,9)} && -2 \ar[rr]_{(7,8)} && -3
}} \]
 
\begin{Verbatim}[commandchars=!@B,fontsize=\small,frame=single,label=Example]
  
  !gapprompt@gap>B !gapinput@c4 := Arrow( Gs3, (8,9), -4, -3 );; B
  !gapprompt@gap>B !gapinput@d4 := Arrow( Gs3, (7,8), -2, -3 );; B
  !gapprompt@gap>B !gapinput@sq4 := SquareOfArrows( D1, h, d2, c3, c4, d4 );;B
  !gapprompt@gap>B !gapinput@UpArrow(sq4)=DownArrow(sq2) and LeftArrow(sq4)=RightArrow(sq3); B
  true
  !gapprompt@gap>B !gapinput@sq24 := VerticalProduct( sq2, sq4 ); B
  [-5] ---------- (8,9) ---------> [-4]
    |                                |
  (7,9)        (11,15)(12,14)        (7,9,8)
    V                                V
  [-2] ---------- (7,8) ---------> [-3]
  !gapprompt@gap>B !gapinput@sq34 := HorizontalProduct( sq3, sq4 );B
  [-1] ------------- (7,8) ------------> [-4]
    |                                     |
  (8,9)        (11,12)(13,16)(14,15)        (8,9)
    V                                     V
  [-2] ------------ (7,9,8) -----------> [-3]
  !gapprompt@gap>B !gapinput@sq1324 := HorizontalProduct( sq13, sq24 );B
  [-6] ------------- (7,9,8) ------------> [-4]
    |                                        |
  (7,8,9)        (11,15,13)(12,16,14)        (7,9,8)
    V                                        V
  [-2] ------------- (7,9,8) ------------> [-3]
  !gapprompt@gap>B !gapinput@sq1234 := VerticalProduct( sq12, sq34 );;B
  !gapprompt@gap>B !gapinput@sq1324 = sq1234;B
  true
  
\end{Verbatim}
 }

 
\section{\textcolor{Chapter }{Conversion of Basic Double Groupoids}}\label{sec-basic-dgpds}
\logpage{[ 11, 2, 0 ]}
\hyperdef{L}{X7D69B5A680FE4C81}{}
{
  As mentioned earlier, double groupoids were introduced in the \textsf{Groupoids} package, but these were \emph{basic double groupoids}, without any pre\texttt{\symbol{45}}crossed module. The element of a square
was simply its boundary. Here we introduce an operation which converts such a
basic double groupoid into the more general case considered in this package. 

\subsection{\textcolor{Chapter }{EnhancedBasicDoubleGroupoid}}
\logpage{[ 11, 2, 1 ]}\nobreak
\hyperdef{L}{X7CB177EF78B559DB}{}
{\noindent\textcolor{FuncColor}{$\triangleright$\enspace\texttt{EnhancedBasicDoubleGroupoid({\mdseries\slshape bdg})\index{EnhancedBasicDoubleGroupoid@\texttt{EnhancedBasicDoubleGroupoid}}
\label{EnhancedBasicDoubleGroupoid}
}\hfill{\scriptsize (operation)}}\\


 We need to add a pre\texttt{\symbol{45}}crossed module to the definition of a
basic double groupoid. We choose to add $(G \to G)$ where $G$ is the root group of the underlying groupoid. (This is only valid for
groupoids which are the direct product with a complete graph.) The example is
taken from section 7.1 of the \textsf{Groupoids} package, converting basic \texttt{B0} to \texttt{D0}, and we check that the same square is produced in each case. }

 
\begin{Verbatim}[commandchars=@CF,fontsize=\small,frame=single,label=Example]
  
  @gappromptCgap>F @gapinputCg2 := (1,2,3,4);;  h2 := (1,3);;F
  @gappromptCgap>F @gapinputCgend8 := [ g2, h2 ];;F
  @gappromptCgap>F @gapinputCd8 := Group( gend8 );;F
  @gappromptCgap>F @gapinputCSetName( d8, "d8" ); F
  @gappromptCgap>F @gapinputCGd8 := Groupoid( d8, [-9..-7] );;F
  @gappromptCgap>F @gapinputCSetName( Gd8, "Gd8" ); F
  @gappromptCgap>F @gapinputCB0 := SinglePieceBasicDoubleGroupoid( Gd8 );; F
  @gappromptCgap>F @gapinputCB0!.groupoid;F
  Gd8
  @gappromptCgap>F @gapinputCB0!.objects;F
  [ -9 .. -7 ]
  @gappromptCgap>F @gapinputCa0 := Arrow(Gd8,(),-9,-7);;    b0 := Arrow(Gd8,g2,-9,-9);;  F
  @gappromptCgap>F @gapinputCc0 := Arrow(Gd8,h2,-7,-8);;    d0 := Arrow(Gd8,(2,4),-9,-8);;      F
  @gappromptCgap>F @gapinputCbdy0 := d0![2]^-1 * b0![2]^-1 * a0![2] * c0![2];; F
  @gappromptCgap>F @gapinputCbsq0 := SquareOfArrows( B0, bdy0, a0, b0, c0, d0 ); F
  [-9] ---------- () ---------> [-7]
    |                             |
  (1,2,3,4)        (1,4,3,2)        (1,3)
    V                             V
  [-9] --------- (2,4) --------> [-8]
  
  @gappromptCgap>F @gapinputCD0 := EnhancedBasicDoubleGroupoid( B0 );;F
  @gappromptCgap>F @gapinputCD0!.prexmod;F
  [d8->d8]
  @gappromptCgap>F @gapinputCbsq0 = SquareOfArrows( D0, bdy0, a0, b0, c0, d0 ); F
  true
  
\end{Verbatim}
 }

 
\section{\textcolor{Chapter }{Commutative double groupoids}}\label{sec-commutative-dgpds}
\logpage{[ 11, 3, 0 ]}
\hyperdef{L}{X853B15F483477D5C}{}
{
  A double groupoid square 
\[  \vcenter{\xymatrix @=4pc{ u_1 \ar[r]^{a_1} \ar[d]_{b_1}\ar@{}[dr] |{m} & v_1
\ar[d]^{c_1} \\ w_1 \ar[r]_{d_1} & x_1 }} \]
 is \emph{commutative} if $a_1c_1 = b_1d_1$, which means that its boundary is the identity. So a double groupoid which
consists only of commutative squares must have a
pre\texttt{\symbol{45}}crossed module with zero boundary. Commutative squares
compose horizontally and vertically provided only that they have the correct
common arrow. 

\subsection{\textcolor{Chapter }{DoubleGroupoidWithZeroBoundary}}
\logpage{[ 11, 3, 1 ]}\nobreak
\hyperdef{L}{X7DC35C557E498880}{}
{\noindent\textcolor{FuncColor}{$\triangleright$\enspace\texttt{DoubleGroupoidWithZeroBoundary({\mdseries\slshape gpd, src})\index{DoubleGroupoidWithZeroBoundary@\texttt{DoubleGroupoidWithZeroBoundary}}
\label{DoubleGroupoidWithZeroBoundary}
}\hfill{\scriptsize (operation)}}\\


 The data for a double groupoid of commutative squares therefore consists of a
groupoid and a source group. We may use the operation \texttt{PreXModWithTrivialRange} (\ref{PreXModWithTrivialRange}) to provide a pre\texttt{\symbol{45}}crossed module. We take for our example
the groupoid \texttt{Gd8} and the pre\texttt{\symbol{45}}crossed module \texttt{Q16} of section \ref{sect-precrossed-modules}. We introduce a new right arrow to construct a square which commutes. }

 
\begin{Verbatim}[commandchars=@CE,fontsize=\small,frame=single,label=Example]
  
  @gappromptCgap>E @gapinputCD16 := DoubleGroupoidWithZeroBoundary( Gs3, d16 );;E
  @gappromptCgap>E @gapinputCD16!.prexmod;E
  [d16->Group( [ () ] )]
  @gappromptCgap>E @gapinputCe16 := Arrow( Gs3, (7,9,8), -5, -3 );;E
  @gappromptCgap>E @gapinputCsq16 := SquareOfArrows( D16, (12,18)(13,17)(14,16), a1, b1, e16, d1 );E
  [-6] -------------- (7,8) -------------> [-5]
    |                                       |
  (7,9)        (12,18)(13,17)(14,16)        (7,9,8)
    V                                       V
  [-1] ------------- (7,8,9) ------------> [-3]
  
\end{Verbatim}
 }

 }

         
\chapter{\textcolor{Chapter }{Applications}}\label{chap-apps}
\logpage{[ 12, 0, 0 ]}
\hyperdef{L}{X7DD7F0847FF2B96C}{}
{
  This chapter was added in April 2018 for version 2.66 of \textsf{XMod}. Initially it describes crossed modules for free loop spaces. Further
applications may arise in due course. 
\section{\textcolor{Chapter }{Free Loop Spaces}}\logpage{[ 12, 1, 0 ]}
\hyperdef{L}{X8575260A80F735BD}{}
{
 \index{loop space} \index{free loop space} These functions have been used to produce examples for Ronald Brown's paper \emph{Crossed modules, and the homotopy $2$\texttt{\symbol{45}}type of a free loop space} \cite{B18}. The relevant theorem in that paper is as follows. 

 \textsc{Theorem 2.1} \emph{ Let $\calM = (\partial : M \to P)$ be a crossed module of groups and let $X = B\calM$ be the classifying space of $\calM$. Then the components of $LX$, the free loop space on $X$, are determined by equivalence classes of elements $a \in P$ where $a,a'$ are equivalent if and only if there are elements $m \in M,\, p \in P$ such that $a'= p + a - \partial m - p$. } 

 \emph{ Further the homotopy $2$\texttt{\symbol{45}}type of a component of $LX$ given by $a \in P$ is determined by the crossed module of groups $L\calM[a] = (\partial_a : M \to P(a))$ where: } 
\begin{itemize}
\item \emph{ $P(a)$ is the subgroup of the cat$^1$\texttt{\symbol{45}}group $G = P \ltimes M$ such that $\partial m = [p,a] = -p-a+p+a$; }
\item \emph{ $\partial_a(m) = (\partial m, m^{-1}m^a)$ for $m \in M$; }
\item \emph{ the action of $P(a)$ on $M$ is given by $n^{(p,m)} = n^p$ for $n \in M,\, (p,m) \in P(a)$. }
\end{itemize}
 

 \emph{ In particular $\pi_1(LX,a)$ is isomorphic to $\mathrm{cokernel}(\partial_a)$, and $\pi_2(LX,a) \cong \pi_2(X,*)^{\bar{a}}$, the elements of $\pi_2(X,*)$ fixed under the action of $\bar{a}$, the class of $a$ in $\pi_1(X,*)$. } 

 \emph{ There is an exact sequence $ \pi \stackrel{\,\phi\,}{\to} \pi \to \pi_1(LX,a) \to C_{\bar{a}}(\pi_1(X,*))
\to 1$, in which $\pi = \pi_2(X,*)$, and $\phi$ is the morphism $m \mapsto m^{-1}m^a$. } 

 

\subsection{\textcolor{Chapter }{LoopClasses}}
\logpage{[ 12, 1, 1 ]}\nobreak
\hyperdef{L}{X87781C76804E783E}{}
{\noindent\textcolor{FuncColor}{$\triangleright$\enspace\texttt{LoopClasses({\mdseries\slshape M})\index{LoopClasses@\texttt{LoopClasses}}
\label{LoopClasses}
}\hfill{\scriptsize (operation)}}\\
\noindent\textcolor{FuncColor}{$\triangleright$\enspace\texttt{LoopsXMod({\mdseries\slshape M, a})\index{LoopsXMod@\texttt{LoopsXMod}}
\label{LoopsXMod}
}\hfill{\scriptsize (operation)}}\\
\noindent\textcolor{FuncColor}{$\triangleright$\enspace\texttt{AllLoopsXMod({\mdseries\slshape M})\index{AllLoopsXMod@\texttt{AllLoopsXMod}}
\label{AllLoopsXMod}
}\hfill{\scriptsize (operation)}}\\


 The operation \texttt{LoopClasses} computes the equivalence classes $[a]$ described above. These are all unions of conjugacy classes. 

 The operation \texttt{LoopsXMod(M,a)} calculates the crossed module $L\calM[a]$ described in the theorem. 

 The operation \texttt{AllLoopsXMod(M)} returns a list of crossed modules, one for each equivalence class of elements $[a] \subseteq P$. 

 In the example below the automorphism crossed module \texttt{X8} has $M \cong C_2^3$ and $P = PSL(3,2)$ is the automorphism group of $M$. There are $6$ equivalence classes which, in this case, are identical with the conjugacy
classes. For each $LX$ calculated, the \texttt{IdGroup} (\ref{IdGroup:for 2d-groups}) is printed out. }

 

 
\begin{Verbatim}[commandchars=!@|,fontsize=\small,frame=single,label=Example]
  
  !gapprompt@gap>| !gapinput@SetName( k8, "k8" ); |
  !gapprompt@gap>| !gapinput@Y8 := XModByAutomorphismGroup( k8 );; |
  !gapprompt@gap>| !gapinput@X8 := Image( IsomorphismPerm2DimensionalGroup( Y8 ) );;|
  !gapprompt@gap>| !gapinput@SetName( X8, "X8" );|
  !gapprompt@gap>| !gapinput@Print( "X8: ", Size( X8 ), " : ", StructureDescription( X8 ), "\n" );  |
  X8: [ 8, 168 ] : [ "C2 x C2 x C2", "PSL(3,2)" ]
  !gapprompt@gap>| !gapinput@classes := LoopClasses( X8 );;|
  !gapprompt@gap>| !gapinput@List( classes, c -> Length(c) );|
  [ 1, 21, 56, 42, 24, 24 ]
  !gapprompt@gap>| !gapinput@LX := LoopsXMod( X8, (1,2)(5,6) );;|
  !gapprompt@gap>| !gapinput@Size2d( LX ); |
  [ 8, 64 ]
  !gapprompt@gap>| !gapinput@IdGroup( LX );|
  [ [ 8, 5 ], [ 64, 138 ] ]
  !gapprompt@gap>| !gapinput@SetInfoLevel( InfoXMod, 1 );|
  !gapprompt@gap>| !gapinput@LX8 := AllLoopsXMod( X8 );;|
  #I  LoopsXMod with a = (),  IdGroup = [ [ 8, 5 ], [ 1344, 11686 ] ]
  #I  LoopsXMod with a = (4,5)(6,7),  IdGroup = [ [ 8, 5 ], [ 64, 138 ] ]
  #I  LoopsXMod with a = (2,3)(4,6,5,7),  IdGroup = [ [ 8, 5 ], [ 32, 6 ] ]
  #I  LoopsXMod with a = (2,4,6)(3,5,7),  IdGroup = [ [ 8, 5 ], [ 24, 13 ] ]
  #I  LoopsXMod with a = (1,2,4,3,6,7,5),  IdGroup = [ [ 8, 5 ], [ 56, 11 ] ]
  #I  LoopsXMod with a = (1,2,4,5,7,3,6),  IdGroup = [ [ 8, 5 ], [ 56, 11 ] ]
  !gapprompt@gap>| !gapinput@iso := IsomorphismGroups( Range( LX ), Range( LX8[2] ) );;|
  !gapprompt@gap>| !gapinput@iso = fail;|
  false
  
\end{Verbatim}
 }

 }

         
\chapter{\textcolor{Chapter }{Interaction with HAP }}\label{chap-hap}
\logpage{[ 13, 0, 0 ]}
\hyperdef{L}{X81EC8C8A82C15298}{}
{
  This chapter describes functions which allow functions in the package \textsf{HAP} to be called from \textsf{XMod}. 
\section{\textcolor{Chapter }{Calling HAP functions}}\label{sect-hap}
\logpage{[ 13, 1, 0 ]}
\hyperdef{L}{X865CE53A827FBE6F}{}
{
  In \textsf{HAP} a cat$^1$\texttt{\symbol{45}}group is called a \texttt{CatOneGroup} and the traditional terms \emph{source} and \emph{target} are used for the \texttt{TailMap} and \texttt{HeadMap}. A \texttt{CatOneGroup} is a record \texttt{C} with fields \texttt{C!.sourceMap} and \texttt{C!.targetMap}. 

\subsection{\textcolor{Chapter }{SmallCat1Group}}
\logpage{[ 13, 1, 1 ]}\nobreak
\hyperdef{L}{X8699357D7DC6279E}{}
{\noindent\textcolor{FuncColor}{$\triangleright$\enspace\texttt{SmallCat1Group({\mdseries\slshape n, i, j})\index{SmallCat1Group@\texttt{SmallCat1Group}}
\label{SmallCat1Group}
}\hfill{\scriptsize (operation)}}\\


 This operation calls the \textsf{HAP} function \texttt{SmallCatOneGroup(n,i,j)} which returns a \texttt{CatOneGroup} from the \textsf{HAP} database. This is then converted into an \textsf{XMod} cat$^1$\texttt{\symbol{45}}group. Note that the numbering is not the same as that
used by the \textsf{XMod} operation \texttt{Cat1Select}. In the example \texttt{C12} is the converted form of \texttt{H12}. }

 

 
\begin{Verbatim}[commandchars=!@|,fontsize=\small,frame=single,label=Example]
  
  !gapprompt@gap>| !gapinput@H12 := SmallCatOneGroup( 12, 4, 3 );|
  Cat-1-group with underlying group Group( [ f1, f2, f3 ] ) . 
  !gapprompt@gap>| !gapinput@C12 := SmallCat1Group( 12, 4, 3 );|
  [Group( [ f1, f2, f3 ] )=>Group( [ f1, f2, <identity> of ... ] )]
  
\end{Verbatim}
 

\subsection{\textcolor{Chapter }{CatOneGroupToXMod}}
\logpage{[ 13, 1, 2 ]}\nobreak
\hyperdef{L}{X7B00E3FB82DC305D}{}
{\noindent\textcolor{FuncColor}{$\triangleright$\enspace\texttt{CatOneGroupToXMod({\mdseries\slshape C})\index{CatOneGroupToXMod@\texttt{CatOneGroupToXMod}}
\label{CatOneGroupToXMod}
}\hfill{\scriptsize (operation)}}\\
\noindent\textcolor{FuncColor}{$\triangleright$\enspace\texttt{Cat1GroupToHAP({\mdseries\slshape C})\index{Cat1GroupToHAP@\texttt{Cat1GroupToHAP}}
\label{Cat1GroupToHAP}
}\hfill{\scriptsize (operation)}}\\


 These two functions convert between the two alternative implementations. }

 

 
\begin{Verbatim}[commandchars=!@|,fontsize=\small,frame=single,label=Example]
  
  !gapprompt@gap>| !gapinput@C12 := CatOneGroupToXMod( H12 );    |
  [Group( [ f1, f2, f3 ] )=>Group( [ f1, f2, <identity> of ... ] )]
  !gapprompt@gap>| !gapinput@C18 := Cat1Select( 18, 4, 3 );|
  [(C3 x C3) : C2=>Group( [ f1, <identity> of ..., f3 ] )]
  !gapprompt@gap>| !gapinput@H18 := Cat1GroupToHAP( C18 ); |
  Cat-1-group with underlying group (C3 x C3) : C2 . 
  
\end{Verbatim}
 

\subsection{\textcolor{Chapter }{IdCat1Group}}
\logpage{[ 13, 1, 3 ]}\nobreak
\hyperdef{L}{X84B7160284FD454A}{}
{\noindent\textcolor{FuncColor}{$\triangleright$\enspace\texttt{IdCat1Group({\mdseries\slshape C})\index{IdCat1Group@\texttt{IdCat1Group}}
\label{IdCat1Group}
}\hfill{\scriptsize (operation)}}\\


 This function calls the \textsf{HAP} function \texttt{IdCatOneGroup} on a cat$^1$\texttt{\symbol{45}}group $C$. This returns $[n,i,j]$ if the cat$^1$\texttt{\symbol{45}}group is the $j$\texttt{\symbol{45}}th structure on the \texttt{SmallGroup(n,i)}. }

 

 
\begin{Verbatim}[commandchars=!@|,fontsize=\small,frame=single,label=Example]
  
  !gapprompt@gap>| !gapinput@IdCatOneGroup( H18 ); |
  [ 18, 4, 4 ]
  !gapprompt@gap>| !gapinput@IdCat1Group( C18 ); |
  [ 18, 4, 4 ]
  
\end{Verbatim}
 }

 }

         
\chapter{\textcolor{Chapter }{Utility functions}}\label{chap-util}
\logpage{[ 14, 0, 0 ]}
\hyperdef{L}{X810FFB1C8035C8BE}{}
{
  By a utility function we mean a \textsf{GAP} function which is 
\begin{itemize}
\item  needed by other functions in this package, 
\item  not (as far as we know) provided by the standard \textsf{GAP} library, 
\item  more suitable for inclusion in the main library than in this package. 
\end{itemize}
 Sections on \emph{Printing Lists} and \emph{Distinct and Common Representatives} were moved to the \textsf{Utils} package with version 2.56. 
\section{\textcolor{Chapter }{Mappings}}\logpage{[ 14, 1, 0 ]}
\hyperdef{L}{X7C9734B880042C73}{}
{
 \index{inclusion mapping} \index{restriction mapping} The following two functions have been moved to the \textsf{gpd} package, but are still documented here. 

\subsection{\textcolor{Chapter }{InclusionMappingGroups}}
\logpage{[ 14, 1, 1 ]}\nobreak
\hyperdef{L}{X7F8E297F7C84DE51}{}
{\noindent\textcolor{FuncColor}{$\triangleright$\enspace\texttt{InclusionMappingGroups({\mdseries\slshape G, H})\index{InclusionMappingGroups@\texttt{InclusionMappingGroups}}
\label{InclusionMappingGroups}
}\hfill{\scriptsize (operation)}}\\
\noindent\textcolor{FuncColor}{$\triangleright$\enspace\texttt{MappingToOne({\mdseries\slshape G, H})\index{MappingToOne@\texttt{MappingToOne}}
\label{MappingToOne}
}\hfill{\scriptsize (operation)}}\\


 This set of utilities concerns mappings. The map \texttt{incd8} is the inclusion of \texttt{d8} in \texttt{d16} used in Section \ref{sect-oper-mor}. \texttt{MappingToOne(G,H)} maps the whole of $G$ to the identity element in $H$. }

 

 
\begin{Verbatim}[commandchars=!@|,fontsize=\small,frame=single,label=Example]
  
  !gapprompt@gap>| !gapinput@Print( incd8, "\n" );|
  [ (11,13,15,17)(12,14,16,18), (11,18)(12,17)(13,16)(14,15) ] ->
  [ (11,13,15,17)(12,14,16,18), (11,18)(12,17)(13,16)(14,15) ]
  !gapprompt@gap>| !gapinput@imd8 := Image( incd8 );;|
  !gapprompt@gap>| !gapinput@MappingToOne( c4, imd8 );|
  [ (11,13,15,17)(12,14,16,18) ] -> [ () ]
  
\end{Verbatim}
 

\subsection{\textcolor{Chapter }{InnerAutomorphismsByNormalSubgroup}}
\logpage{[ 14, 1, 2 ]}\nobreak
\hyperdef{L}{X81D29E737F3D4878}{}
{\noindent\textcolor{FuncColor}{$\triangleright$\enspace\texttt{InnerAutomorphismsByNormalSubgroup({\mdseries\slshape G, N})\index{InnerAutomorphismsByNormalSubgroup@\texttt{InnerAutomorphismsByNormalSubgroup}}
\label{InnerAutomorphismsByNormalSubgroup}
}\hfill{\scriptsize (operation)}}\\


 Inner automorphisms of a group \texttt{G} by the elements of a normal subgroup \texttt{N} are calculated, often with \texttt{G} = \texttt{N}. }

 

 
\begin{Verbatim}[commandchars=!@|,fontsize=\small,frame=single,label=Example]
  
  !gapprompt@gap>| !gapinput@autd8 := AutomorphismGroup( d8 );;|
  !gapprompt@gap>| !gapinput@innd8 := InnerAutomorphismsByNormalSubgroup( d8, d8 );;|
  !gapprompt@gap>| !gapinput@GeneratorsOfGroup( innd8 );|
  [ ^(1,2,3,4), ^(1,3) ]
  
\end{Verbatim}
 

\subsection{\textcolor{Chapter }{IsGroupOfAutomorphisms}}
\logpage{[ 14, 1, 3 ]}\nobreak
\hyperdef{L}{X7FC631B786C1DC8B}{}
{\noindent\textcolor{FuncColor}{$\triangleright$\enspace\texttt{IsGroupOfAutomorphisms({\mdseries\slshape A})\index{IsGroupOfAutomorphisms@\texttt{IsGroupOfAutomorphisms}}
\label{IsGroupOfAutomorphisms}
}\hfill{\scriptsize (property)}}\\


 Tests whether the elements of a group are automorphisms. }

 

 
\begin{Verbatim}[commandchars=!@|,fontsize=\small,frame=single,label=Example]
  
  !gapprompt@gap>| !gapinput@IsGroupOfAutomorphisms( innd8 );|
  true
  
\end{Verbatim}
 }

 
\section{\textcolor{Chapter }{Abelian Modules}}\label{sect-abelian-modules}
\logpage{[ 14, 2, 0 ]}
\hyperdef{L}{X852BD9CA84C2AFF0}{}
{
  \index{abelian module} 

\subsection{\textcolor{Chapter }{AbelianModuleObject}}
\logpage{[ 14, 2, 1 ]}\nobreak
\hyperdef{L}{X806DEFCC859BB4F1}{}
{\noindent\textcolor{FuncColor}{$\triangleright$\enspace\texttt{AbelianModuleObject({\mdseries\slshape grp, act})\index{AbelianModuleObject@\texttt{AbelianModuleObject}}
\label{AbelianModuleObject}
}\hfill{\scriptsize (operation)}}\\
\noindent\textcolor{FuncColor}{$\triangleright$\enspace\texttt{IsAbelianModule({\mdseries\slshape obj})\index{IsAbelianModule@\texttt{IsAbelianModule}}
\label{IsAbelianModule}
}\hfill{\scriptsize (property)}}\\
\noindent\textcolor{FuncColor}{$\triangleright$\enspace\texttt{AbelianModuleGroup({\mdseries\slshape obj})\index{AbelianModuleGroup@\texttt{AbelianModuleGroup}}
\label{AbelianModuleGroup}
}\hfill{\scriptsize (attribute)}}\\
\noindent\textcolor{FuncColor}{$\triangleright$\enspace\texttt{AbelianModuleAction({\mdseries\slshape obj})\index{AbelianModuleAction@\texttt{AbelianModuleAction}}
\label{AbelianModuleAction}
}\hfill{\scriptsize (attribute)}}\\


 An abelian module is an abelian group together with a group action. These are
used by the crossed module constructor \texttt{XModByAbelianModule} (\ref{XModByAbelianModule}). 

 The resulting \texttt{Xabmod} is isomorphic to the output from \texttt{XModByAutomorphismGroup( k4 );}. }

 

 
\begin{Verbatim}[commandchars=!@|,fontsize=\small,frame=single,label=Example]
  
  !gapprompt@gap>| !gapinput@x := (6,7)(8,9);;  y := (6,8)(7,9);;  z := (6,9)(7,8);;|
  !gapprompt@gap>| !gapinput@k4a := Group( x, y );;  SetName( k4a, "k4a" );|
  !gapprompt@gap>| !gapinput@gens3a := [ (1,2), (2,3) ];;|
  !gapprompt@gap>| !gapinput@s3a := Group( gens3a );;  SetName( s3a, "s3a" );|
  !gapprompt@gap>| !gapinput@alpha := GroupHomomorphismByImages( k4a, k4a, [x,y], [y,x] );;|
  !gapprompt@gap>| !gapinput@beta := GroupHomomorphismByImages( k4a, k4a, [x,y], [x,z] );;|
  !gapprompt@gap>| !gapinput@auta := Group( alpha, beta );;|
  !gapprompt@gap>| !gapinput@acta := GroupHomomorphismByImages( s3a, auta, gens3a, [alpha,beta] );;|
  !gapprompt@gap>| !gapinput@abmod := AbelianModuleObject( k4a, acta );;|
  !gapprompt@gap>| !gapinput@Xabmod := XModByAbelianModule( abmod );|
  [k4a->s3a]
  !gapprompt@gap>| !gapinput@Display( Xabmod );|
  
  Crossed module [k4a->s3a] :- 
  : Source group k4a has generators:
    [ (6,7)(8,9), (6,8)(7,9) ]
  : Range group s3a has generators:
    [ (1,2), (2,3) ]
  : Boundary homomorphism maps source generators to:
    [ (), () ]
  : Action homomorphism maps range generators to automorphisms:
    (1,2) --> { source gens --> [ (6,8)(7,9), (6,7)(8,9) ] }
    (2,3) --> { source gens --> [ (6,7)(8,9), (6,9)(7,8) ] }
    These 2 automorphisms generate the group of automorphisms.
  
  
\end{Verbatim}
 }

 }

         
\chapter{\textcolor{Chapter }{Development history}}\label{chap-history}
\logpage{[ 15, 0, 0 ]}
\hyperdef{L}{X810C43BC7F63C4B4}{}
{
  This chapter, which contains details of the major changes to the package as it
develops, was first created in April 2002. Details of the changes from \textsf{XMod} 1 to \textsf{XMod} 2.001 are far from complete. Starting with version 2.009 the file \texttt{CHANGES} lists the minor changes as well as the more fundamental ones. 

 The inspiration for this package was the need, in the
mid\texttt{\symbol{45}}1990's, to calculate induced crossed modules (see \cite{BW1}, \cite{BW2}, \cite{BW3}). \textsf{GAP} was chosen over other computational group theory systems because the code was
freely available, and it was possible to modify the Tietze transformation code
so as to record the images of the original generators of a presentation as
words in the simplified presentation. (These modifications are now a standard
part of the Tietze transformation package in \textsf{GAP}.) 
\section{\textcolor{Chapter }{Changes from version to version}}\logpage{[ 15, 1, 0 ]}
\hyperdef{L}{X7ACE7E8384B73156}{}
{
 
\subsection{\textcolor{Chapter }{Version 1 for \textsf{GAP} 3}}\logpage{[ 15, 1, 1 ]}
\hyperdef{L}{X848198AA862249C4}{}
{
 \index{version 1 for \textsf{GAP} 3} The first version of \textsf{XMod} became an accepted package for \textsf{GAP} 3.4.3 in December 1996. }

 
\subsection{\textcolor{Chapter }{Version 2}}\logpage{[ 15, 1, 2 ]}
\hyperdef{L}{X7CF8E72D80AAB54F}{}
{
 Conversion of \textsf{XMod} 1 from \textsf{GAP} 3.4.3 to the new \textsf{GAP} syntax began soon after \textsf{GAP} 4 was released, and had a lengthy gestation. The new \textsf{GAP} syntax encouraged a re\texttt{\symbol{45}}naming of many of the function
names. An early decision was to introduce generic categories \texttt{2dDomain} for (pre\texttt{\symbol{45}})crossed modules and
(pre\texttt{\symbol{45}})cat1\texttt{\symbol{45}}groups, and \texttt{2dMapping} for the various types of morphism. In 2.009 \texttt{3dDomain} was used for crossed squares and cat2\texttt{\symbol{45}}groups, and \texttt{3dMapping} for their morphisms. A generic name for derivations and sections is also
required, and \texttt{Up2dMapping} is currently used. }

 
\subsection{\textcolor{Chapter }{Version 2.001 for \textsf{GAP} 4}}\logpage{[ 15, 1, 3 ]}
\hyperdef{L}{X7F9CE0487BB6F660}{}
{
 \index{version 2.001 for \textsf{GAP} 4} This was the first version of \textsf{XMod} for \textsf{GAP} 4, completed in April 2002 in time for the release of \textsf{GAP} 4.3. Functions for actors and induced crossed modules were not included, nor
many of the functions for derivations and sections, for example \texttt{InnerDerivation}. }

 
\subsection{\textcolor{Chapter }{Induced crossed modules}}\logpage{[ 15, 1, 4 ]}
\hyperdef{L}{X7966FF497C36C465}{}
{
 During May 2002 converted the code for induced crossed modules. (Induced
cat1\texttt{\symbol{45}}groups may be converted one day.) }

 
\subsection{\textcolor{Chapter }{Versions 2.002 \texttt{\symbol{45}}\texttt{\symbol{45}} 2.006}}\logpage{[ 15, 1, 5 ]}
\hyperdef{L}{X7E0B70FD82DC5BA8}{}
{
 Version 2.004 of April 14th 2004 added the \texttt{Cat1Select} (\ref{Cat1Select}) functionality of version 1 to the \texttt{Cat1Group} (\ref{Cat1Group}) function. 

 A significant addition in Version 2.005 was the conversion of the actor
crossed module functions from the \texttt{3.4.4} version. This included \texttt{AutomorphismPermGroup} (\ref{AutomorphismPermGroup}) for a crossed module; \texttt{WhiteheadXMod} (\ref{WhiteheadXMod}); \texttt{NorrieXMod} (\ref{NorrieXMod}); \texttt{LueXMod} (\ref{LueXMod}); \texttt{ActorXMod} (\ref{ActorXMod}); \texttt{CentreXMod} (\ref{CentreXMod}) of a crossed module; \texttt{InnerMorphism} (\ref{InnerMorphism}); and \texttt{InnerActorXMod} (\ref{InnerActorXMod}). }

 
\subsection{\textcolor{Chapter }{Versions 2.007 \texttt{\symbol{45}}\texttt{\symbol{45}} 2.010}}\logpage{[ 15, 1, 6 ]}
\hyperdef{L}{X7F6E650E85384C25}{}
{
 These versions contain changes made between September 2004 and October 2007. 
\begin{itemize}
\item  Added basic functions for crossed squares, considered as \texttt{3dObjects} with crossed pairings, and their morphisms. Groups with two normal subgroups,
and the actor of a crossed module, provide standard examples of crossed
squares. (Cat2\texttt{\symbol{45}}groups are not yet implemented.) 
\item  Converted the documentation to the format of the \textsf{GAPDoc} package. 
\item  Improved \texttt{AutomorphismPermGroup} (\ref{AutomorphismPermGroup}) for crossed modules, and introduced a special method for conjugation crossed
modules. 
\item  Substantial revisons made to \texttt{XModByCentralExtension} (\ref{XModByCentralExtension}); \texttt{NorrieXMod} (\ref{NorrieXMod}); \texttt{LueXMod} (\ref{LueXMod}); \texttt{ActorXMod} (\ref{ActorXMod}); and \texttt{InducedXModByCopower} (\ref{InducedXModByCopower}). 
\item  Version 2.010, of October 2007, was timed to coincide with the release of \textsf{GAP} 4.4.10, and included a change of filenames; and correct file protection codes. 
\end{itemize}
 }

 }

 
\section{\textcolor{Chapter }{Versions for \textsf{GAP} [4.5 .. 4.12]}}\logpage{[ 15, 2, 0 ]}
\hyperdef{L}{X80A8A3FB82048ADD}{}
{
 Version 2.19, released on 9th June 2012, included the following changes: 
\begin{itemize}
\item  The file \texttt{makedocrel.g} was copied, with appropriate changes, from \textsf{GAPDoc}, and now provides the correct way to update the documentation. 
\item  The first functions for crossed modules of groupoids were introduced. 
\item  A GNU General Public License declaration was added. 
\end{itemize}
 
\subsection{\textcolor{Chapter }{AllCat1s}}\logpage{[ 15, 2, 1 ]}
\hyperdef{L}{X794BBE42839F2E18}{}
{
 Version \textsc{2.21} contained major changes to the \texttt{Cat1Select} (\ref{Cat1Select}) operation: the list \texttt{CAT1{\textunderscore}LIST} of cat1\texttt{\symbol{45}}structures in the data file \texttt{cat1data.g} was changed from permutation groups to pc\texttt{\symbol{45}}groups, with the
generators of subgroups; images of the tail map; and images of the head map
being given as \texttt{ExtRepOfObj} of words in the generators. 

 The \texttt{AllCat1s} function was reintroduced from the \textsf{GAP}3 version, and renamed as the operation \texttt{AllCat1DataGroupsBasic}. 

 In version \textsc{2.25} the data in \texttt{cat1data.g} was extended from groups of size up to $48$ to groups of size up to $70$. In particular, the $267$ groups of size 64 give rise to a total of $1275$ cat1\texttt{\symbol{45}}groups. The authors are indebted to Van Luyen Le in
Galway for pointing out a number of errors in the version of this list
distributed with version \textsc{2.24} of this package. }

 
\subsection{\textcolor{Chapter }{Versions 2.43 \texttt{\symbol{45}} 2.56}}\logpage{[ 15, 2, 2 ]}
\hyperdef{L}{X78C26CC27D48B1A8}{}
{
 Version \textsc{2.43}, released on 11th November 2015, included the following changes: 
\begin{itemize}
\item  The material on isoclinism in Chapter 4 was added. 
\item  The package webpage has moved to \href{https://github.com/cdwensley} {\texttt{https://github.com/cdwensley}}. 
\item  A GitHub repository was started at: \href{https://github.com/gap-packages/xmod} {\texttt{https://github.com/gap\texttt{\symbol{45}}packages/xmod}}. 
\item  The section on \emph{Distinct and Common Representatives} was moved to the \textsf{Utils} package. 
\end{itemize}
 }

 
\subsection{\textcolor{Chapter }{Version 2.61}}\logpage{[ 15, 2, 3 ]}
\hyperdef{L}{X8715310378F0F8D2}{}
{
 Major changes in names took place, with \texttt{2dDomain, 2dGroup, 2dMapping}, etc. becoming \texttt{2DimensionalDomain, 2DimensionalGroup, 2DimensionalMapping}, etc., and similarly for 3\texttt{\symbol{45}}dimensional versions. Also \texttt{HigherDimensionalDomain} and related categories, domains, properties, attributes and operations were
introduced. At the same time, functions for cat2\texttt{\symbol{45}}groups
were introduced by Alper Odabas. }

 
\subsection{\textcolor{Chapter }{Versions 2.63 \texttt{\symbol{45}} 2.74}}\logpage{[ 15, 2, 4 ]}
\hyperdef{L}{X87EE8E70786CAF46}{}
{
 
\begin{itemize}
\item  Added an implementation of crossed modules of groupoids. 
\item  Lots more work on crossed squares and cat2\texttt{\symbol{45}}groups. 
\item  Added an implementation of group groupoids. 
\end{itemize}
 }

 
\subsection{\textcolor{Chapter }{Versions 2.75 \texttt{\symbol{45}} 2.85}}\logpage{[ 15, 2, 5 ]}
\hyperdef{L}{X85F63D6979B72CA5}{}
{
 
\begin{itemize}
\item  Added conversion functions between \textsf{XMod} and \textsf{Hap} and a new chapter in the manual about these functions. 
\item  Added functions for quasi\texttt{\symbol{45}}isomorphisms. 
\end{itemize}
 }

 
\subsection{\textcolor{Chapter }{Versions 2.86 \texttt{\symbol{45}} 2.91}}\logpage{[ 15, 2, 6 ]}
\hyperdef{L}{X81023EBF7CD2352E}{}
{
 
\begin{itemize}
\item  Added attributes \texttt{Size2d} for 2d\texttt{\symbol{45}}objects and \texttt{Size3d} for 3d\texttt{\symbol{45}}objects since lists are inappropriate values for the
standard function \texttt{Size}. 
\item  Added \texttt{PreXModWithTrivialRange} and started work on double groupoids. 
\end{itemize}
 }

 
\subsection{\textcolor{Chapter }{Versions from 2.92}}\logpage{[ 15, 2, 7 ]}
\hyperdef{L}{X87062F217BAC0B6E}{}
{
 
\begin{itemize}
\item  Renamed \texttt{PermAutomorphismAsXModMorphism} as \texttt{PermAutomorphismAs2dGroupMorphism} and added a cat$^1$\texttt{\symbol{45}}group method for this and for \texttt{AutomorphismPermGroup}. 
\end{itemize}
 }

 
\subsection{\textcolor{Chapter }{Versions from 2.98}}\logpage{[ 15, 2, 8 ]}
\hyperdef{L}{X790B71B07A845736}{}
{
 
\begin{itemize}
\item  Documented \texttt{InnerAutomorphismXMod} and \texttt{InnerAutomorphismCat1Group}, and added \texttt{InnerAutomorphismPermGroup} for a 2d\texttt{\symbol{45}}group. 
\item  Included the verification that \texttt{CentreXMod(X)} and \texttt{XModCentre(X)} are equivalent. 
\item  Added the operation \texttt{WhiteheadGroupElement} for a derivation. 
\end{itemize}
 }

 }

 
\section{\textcolor{Chapter }{What needs doing next?}}\logpage{[ 15, 3, 0 ]}
\hyperdef{L}{X83D1530487593182}{}
{
 
\begin{itemize}
\item  Speed up the calculation of Whitehead groups. 
\item  Add more functions for \texttt{3dObjects} and implement \texttt{cat2\texttt{\symbol{45}}groups}. 
\item  Improve interaction with the package \textsf{groupoids} implementing the group groupoid version of a crossed module, and adding more
functions for crossed modules of groupoids. 
\item  Add interaction with \textsf{IdRel} (and possibly \textsf{XRes} and \textsf{natp}) . 
\item  Need \texttt{InverseGeneralMapping} for morphisms and more features for \texttt{FpXMods}, \texttt{PcXMods}, etc. 
\item  Implement actions of a crossed module. 
\item  Implement \texttt{FreeXMods} and an operation \texttt{Isomorphism2dDomains}. 
\item  Allow the construction of a group of morphisms of crossed modules. 
\item  Complete the conversion from Version \textsc{1} of the calculation of sections using \texttt{EndoClasses}. 
\item  Another crossed square construction. If $M, N$ are ordinary $P$\texttt{\symbol{45}}modules and $A$ is an arbitrary abelian group on which $P$ acts trivially, then there is a crossed square with sides 
\[ 0 : A \to N,\quad 0 : A \to M,\quad 0 : M \to P,\quad 0 : N \to P. \]
 
\item  Improve the interaction with the \textsf{HAP} package. 
\item  Implement cat1\texttt{\symbol{45}}groups with objects. 
\item  Lots more work on double groupoids. 
\end{itemize}
 }

 }

 \def\bibname{References\logpage{[ "Bib", 0, 0 ]}
\hyperdef{L}{X7A6F98FD85F02BFE}{}
}

\bibliographystyle{alpha}
\bibliography{bib.xml}

\addcontentsline{toc}{chapter}{References}

\def\indexname{Index\logpage{[ "Ind", 0, 0 ]}
\hyperdef{L}{X83A0356F839C696F}{}
}

\cleardoublepage
\phantomsection
\addcontentsline{toc}{chapter}{Index}


\printindex

\immediate\write\pagenrlog{["Ind", 0, 0], \arabic{page},}
\newpage
\immediate\write\pagenrlog{["End"], \arabic{page}];}
\immediate\closeout\pagenrlog
\end{document}
